% !TeX program = xelatex
% Standalone English volume. Use XeLaTeX until the references stabilize.
\documentclass[a4paper,11pt]{article}
\usepackage[T1]{fontenc}
\usepackage{lmodern}
\usepackage[english]{babel}
\usepackage[left=26mm,right=26mm,top=25mm,bottom=25mm,headheight=15pt,headsep=7mm,footskip=11mm]{geometry}
\usepackage{amsmath,amssymb,mathtools,amsthm}
\newtheorem{proposition}{Proposition}[section]
\newtheorem{lemma}[proposition]{Lemma}
\newtheorem{theorem}[proposition]{Theorem}
\newtheorem{corollary}[proposition]{Corollary}
\newtheorem{maintheorem}[proposition]{Main Theorem}
\theoremstyle{definition}
\newtheorem{definition}[proposition]{Definition}
\theoremstyle{remark}
\newtheorem{remark}[proposition]{Remark}
\makeatletter
\renewcommand*\l@section{\@dottedtocline{1}{0em}{2.4em}}
\renewcommand*\l@subsection{\@dottedtocline{2}{1.45em}{3.45em}}
\renewcommand*\l@subsubsection{\@dottedtocline{3}{3.8em}{4.55em}}
\makeatother
\usepackage{xcolor,fancyhdr,titlesec,needspace,xurl}
\usepackage{longtable,multicol}
\usepackage{hyperref}
\definecolor{ink}{HTML}{1A3542}
\definecolor{muted}{HTML}{4F606A}
\hypersetup{unicode=true,colorlinks=true,linkcolor=ink,filecolor=ink,urlcolor=ink,citecolor=ink,bookmarksnumbered=true,
 pdfdisplaydoctitle=true,pdflang={en},
 pdftitle={The generalized Riemann hypothesis for Dirichlet L-functions. Edition 930, Volume XXVII},
 pdfsubject={A natural-language reconstruction of the proposed proof},
 pdfauthor={A mathematical reconstruction}}
\titleformat{\section}{\Large\bfseries\color{ink}}{\thesection}{.65em}{}
\titlespacing*{\section}{0pt}{.8em}{.65em}
\setlength{\parindent}{1.2em}
\linespread{1.045}
\setlength{\parskip}{.22em plus .05em minus .03em}
\setlength{\emergencystretch}{2em}
\setcounter{topnumber}{2}
\setcounter{tocdepth}{3}\setcounter{secnumdepth}{3}
\widowpenalty=4500
\clubpenalty=4500
\pagestyle{fancy}\fancyhf{}
\fancyhead[L]{\small\color{muted}Dirichlet GRH \textperiodcentered\ Volume I}
\fancyhead[R]{\small\color{muted}Dirichlet zeros and power intervals}
\fancyfoot[C]{\small\thepage}
\renewcommand{\headrulewidth}{.3pt}
\newcommand{\lead}[1]{\par\smallskip\noindent{\bfseries\color{ink}#1.}\enspace}
\newcommand{\sour}[2]{\unskip\penalty150\hspace{.35em}{\small\hyperlink{ref#1}{[#2]}}}
\titleformat{\subsection}{\large\bfseries}{\thesubsection}{.6em}{}
\titlespacing*{\subsection}{0pt}{1.15em plus .15em}{.4em}
\setlength{\abovedisplayskip}{9pt plus 2pt minus 3pt}
\setlength{\belowdisplayskip}{9pt plus 2pt minus 3pt}
\newcommand{\RR}{\mathbb R}\newcommand{\CC}{\mathbb C}\newcommand{\QQ}{\mathbb Q}
\newcommand{\ZZ}{\mathbb Z}\newcommand{\ord}{\operatorname{ord}}
\newcommand{\Frob}{\operatorname{Frob}}\newcommand{\Gal}{\operatorname{Gal}}
\newcommand{\len}{\operatorname{length}}\newcommand{\Res}{\operatorname*{Res}}
\newcommand{\wmean}{\operatorname{av}}\newcommand{\Ker}{\operatorname{ker}}
\newcommand{\Lh}{\mathcal L}\newcommand{\Gg}{\mathcal G}
\newcommand{\real}{\operatorname{Re}}\newcommand{\imag}{\operatorname{Im}}
\newcommand{\abs}[1]{\left|#1\right|}\newcommand{\eps}{\varepsilon}
\newcommand{\anchoritem}[2]{\par\Needspace{2.8\baselineskip}\hypertarget{ref#1}{}\noindent{\bfseries\color{ink}#2.}\enspace}
\newcommand{\srcloc}[2]{{\small\nolinkurl{#1.lean}}, line~#2}
\numberwithin{equation}{section}
\raggedbottom

\newcommand{\Sel}{\mathcal S}
\newcommand{\Cl}{\operatorname{Cl}}
\newcommand{\Aut}{\operatorname{Aut}}
\newcommand{\GL}{\operatorname{GL}}
\newcommand{\OO}{\mathcal O}
\hypersetup{bookmarksdepth=3}


\addto\captionsenglish{\renewcommand{\partname}{Volume}}


\makeatletter
\expandafter\def\csname r@cor:class66\endcsname{{65.28}{45}{}{proposition.65.28}{Dirichlet-GRH-v930-Volume-XIII.pdf}}
\expandafter\def\csname r@cor:count111\endcsname{{9.15}{119}{}{proposition.9.15}{Dirichlet-GRH-v930-Volume-I.pdf}}
\expandafter\def\csname r@cor:fullperiodmoments400\endcsname{{128.26}{77}{}{proposition.128.26}{Dirichlet-GRH-v930-Volume-XXIV.pdf}}
\expandafter\def\csname r@cor:inertcorrection461\endcsname{{68.3}{89}{}{proposition.68.3}{Dirichlet-GRH-v930-Volume-XIII.pdf}}
\expandafter\def\csname r@cor:interval56\endcsname{{65.23}{41}{}{proposition.65.23}{Dirichlet-GRH-v930-Volume-XIII.pdf}}
\expandafter\def\csname r@cor:normaldegree347\endcsname{{127.2}{27}{}{proposition.127.2}{Dirichlet-GRH-v930-Volume-XXIV.pdf}}
\expandafter\def\csname r@cor:rank111\endcsname{{9.14}{119}{}{proposition.9.14}{Dirichlet-GRH-v930-Volume-I.pdf}}
\expandafter\def\csname r@cor:sameprimes517\endcsname{{9.34}{134}{}{proposition.9.34}{Dirichlet-GRH-v930-Volume-I.pdf}}
\expandafter\def\csname r@cor:shortcenter460\endcsname{{9.27}{130}{}{proposition.9.27}{Dirichlet-GRH-v930-Volume-I.pdf}}
\expandafter\def\csname r@cor:singular349\endcsname{{104.33}{29}{}{proposition.104.33}{Dirichlet-GRH-v930-Volume-XX.pdf}}
\expandafter\def\csname r@cor:smallresidual515\endcsname{{22.23}{81}{}{proposition.22.23}{Dirichlet-GRH-v930-Volume-IV.pdf}}
\expandafter\def\csname r@cor:splitcost456\endcsname{{120.18}{29}{}{proposition.120.18}{Dirichlet-GRH-v930-Volume-XXIII.pdf}}
\expandafter\def\csname r@eq:Gchi\endcsname{{1.1}{3}{}{equation.1.1}{Dirichlet-GRH-v930-Volume-I.pdf}}
\expandafter\def\csname r@eq:cp73\endcsname{{66.1}{46}{}{equation.66.1}{Dirichlet-GRH-v930-Volume-XIII.pdf}}
\expandafter\def\csname r@eq:gap\endcsname{{6.3}{26}{}{equation.6.3}{Dirichlet-GRH-v930-Volume-I.pdf}}
\expandafter\def\csname r@eq:letters111\endcsname{{9.1}{117}{}{equation.9.1}{Dirichlet-GRH-v930-Volume-I.pdf}}
\expandafter\def\csname r@eq:sub81\endcsname{{66.22}{67}{}{equation.66.22}{Dirichlet-GRH-v930-Volume-XIII.pdf}}
\expandafter\def\csname r@eq:uncond12\endcsname{{3.1}{13}{}{equation.3.1}{Dirichlet-GRH-v930-Volume-I.pdf}}
\expandafter\def\csname r@eq:window\endcsname{{{W}}{2}{}{AMS.2}{Dirichlet-GRH-v930-Volume-I.pdf}}
\expandafter\def\csname r@lem:addext492\endcsname{{68.32}{108}{}{proposition.68.32}{Dirichlet-GRH-v930-Volume-XIII.pdf}}
\expandafter\def\csname r@lem:basis59\endcsname{{65.8}{27}{}{proposition.65.8}{Dirichlet-GRH-v930-Volume-XIII.pdf}}
\expandafter\def\csname r@lem:basis63\endcsname{{65.14}{34}{}{proposition.65.14}{Dirichlet-GRH-v930-Volume-XIII.pdf}}
\expandafter\def\csname r@lem:cmorbit487\endcsname{{68.26}{103}{}{proposition.68.26}{Dirichlet-GRH-v930-Volume-XIII.pdf}}
\expandafter\def\csname r@lem:coefficientdecay697\endcsname{{132.8}{57}{}{proposition.132.8}{Dirichlet-GRH-v930-Volume-XXV.pdf}}
\expandafter\def\csname r@lem:convexcount666\endcsname{{6.37}{85}{}{proposition.6.37}{Dirichlet-GRH-v930-Volume-I.pdf}}
\expandafter\def\csname r@lem:cosetpacking699\endcsname{{132.13}{60}{}{proposition.132.13}{Dirichlet-GRH-v930-Volume-XXV.pdf}}
\expandafter\def\csname r@lem:cyclicneighbor495\endcsname{{68.41}{114}{}{proposition.68.41}{Dirichlet-GRH-v930-Volume-XIII.pdf}}
\expandafter\def\csname r@lem:far14\endcsname{{7.1}{88}{}{proposition.7.1}{Dirichlet-GRH-v930-Volume-I.pdf}}
\expandafter\def\csname r@lem:green486\endcsname{{68.21}{100}{}{proposition.68.21}{Dirichlet-GRH-v930-Volume-XIII.pdf}}
\expandafter\def\csname r@lem:grid85\endcsname{{66.33}{72}{}{proposition.66.33}{Dirichlet-GRH-v930-Volume-XIII.pdf}}
\expandafter\def\csname r@lem:hook86\endcsname{{92.4}{5}{}{proposition.92.4}{Dirichlet-GRH-v930-Volume-XVIII.pdf}}
\expandafter\def\csname r@lem:kernel\endcsname{{2.4}{11}{}{proposition.2.4}{Dirichlet-GRH-v930-Volume-I.pdf}}
\expandafter\def\csname r@lem:mass32\endcsname{{121.17}{48}{}{proposition.121.17}{Dirichlet-GRH-v930-Volume-XXIII.pdf}}
\expandafter\def\csname r@lem:mean74\endcsname{{66.5}{49}{}{proposition.66.5}{Dirichlet-GRH-v930-Volume-XIII.pdf}}
\expandafter\def\csname r@lem:mellin\endcsname{{2.3}{11}{}{proposition.2.3}{Dirichlet-GRH-v930-Volume-I.pdf}}
\expandafter\def\csname r@lem:residuelifts701\endcsname{{133.1}{65}{}{proposition.133.1}{Dirichlet-GRH-v930-Volume-XXV.pdf}}
\expandafter\def\csname r@lem:sphere60\endcsname{{66.8}{51}{}{proposition.66.8}{Dirichlet-GRH-v930-Volume-XIII.pdf}}
\expandafter\def\csname r@lem:unipotentvolume703\endcsname{{133.9}{70}{}{proposition.133.9}{Dirichlet-GRH-v930-Volume-XXV.pdf}}
\expandafter\def\csname r@lem:unitframe455\endcsname{{120.13}{25}{}{proposition.120.13}{Dirichlet-GRH-v930-Volume-XXIII.pdf}}
\expandafter\def\csname r@lem:weyldual351\endcsname{{120.4}{16}{}{proposition.120.4}{Dirichlet-GRH-v930-Volume-XXIII.pdf}}
\expandafter\def\csname r@nav:I-28\endcsname{{6.6}{32}{}{subsection.6.6}{Dirichlet-GRH-v930-Volume-I.pdf}}
\expandafter\def\csname r@prop:allzeros\endcsname{{3.2}{13}{}{proposition.3.2}{Dirichlet-GRH-v930-Volume-I.pdf}}
\expandafter\def\csname r@prop:averagedschur684\endcsname{{131.6}{19}{}{proposition.131.6}{Dirichlet-GRH-v930-Volume-XXV.pdf}}
\expandafter\def\csname r@prop:column110\endcsname{{86.29}{23}{}{proposition.86.29}{Dirichlet-GRH-v930-Volume-XVII.pdf}}
\expandafter\def\csname r@prop:field12\endcsname{{6.4}{27}{}{proposition.6.4}{Dirichlet-GRH-v930-Volume-I.pdf}}
\expandafter\def\csname r@prop:fock187\endcsname{{104.20}{17}{}{proposition.104.20}{Dirichlet-GRH-v930-Volume-XX.pdf}}
\expandafter\def\csname r@prop:normalframe349\endcsname{{104.34}{29}{}{proposition.104.34}{Dirichlet-GRH-v930-Volume-XX.pdf}}
\expandafter\def\csname r@prop:outsidewordcost451\endcsname{{9.23}{127}{}{proposition.9.23}{Dirichlet-GRH-v930-Volume-I.pdf}}
\expandafter\def\csname r@prop:polarcoefficient697\endcsname{{132.7}{56}{}{proposition.132.7}{Dirichlet-GRH-v930-Volume-XXV.pdf}}
\expandafter\def\csname r@prop:ray50\endcsname{{6.30}{63}{}{proposition.6.30}{Dirichlet-GRH-v930-Volume-I.pdf}}
\expandafter\def\csname r@prop:res23\endcsname{{6.11}{32}{}{proposition.6.11}{Dirichlet-GRH-v930-Volume-I.pdf}}
\expandafter\def\csname r@prop:rowmean690\endcsname{{131.25}{35}{}{proposition.131.25}{Dirichlet-GRH-v930-Volume-XXV.pdf}}
\expandafter\def\csname r@sec:chain81\endcsname{{66.10}{67}{}{subsection.66.10}{Dirichlet-GRH-v930-Volume-XIII.pdf}}
\expandafter\def\csname r@sec:markedline350\endcsname{{119.4}{7}{}{subsection.119.4}{Dirichlet-GRH-v930-Volume-XXIII.pdf}}
\expandafter\def\csname r@sec:missing\endcsname{{6.6}{32}{}{subsection.6.6}{Dirichlet-GRH-v930-Volume-I.pdf}}
\expandafter\def\csname r@sec:period63\endcsname{{65.5}{34}{}{subsection.65.5}{Dirichlet-GRH-v930-Volume-XIII.pdf}}
\expandafter\def\csname r@sec:weylperiod351\endcsname{{119.6}{9}{}{subsection.119.6}{Dirichlet-GRH-v930-Volume-XXIII.pdf}}
\expandafter\def\csname r@seriesguide\endcsname{{1.1}{4}{}{subsection.1.1}{Dirichlet-GRH-v930-Volume-I.pdf}}
\expandafter\def\csname r@thm:adapted314\endcsname{{122.14}{72}{}{proposition.122.14}{Dirichlet-GRH-v930-Volume-XXIII.pdf}}
\expandafter\def\csname r@thm:besselmatrix701\endcsname{{133.2}{65}{}{proposition.133.2}{Dirichlet-GRH-v930-Volume-XXV.pdf}}
\expandafter\def\csname r@thm:centerwords451\endcsname{{9.22}{126}{}{proposition.9.22}{Dirichlet-GRH-v930-Volume-I.pdf}}
\expandafter\def\csname r@thm:class61\endcsname{{65.11}{31}{}{proposition.65.11}{Dirichlet-GRH-v930-Volume-XIII.pdf}}
\expandafter\def\csname r@thm:classden312\endcsname{{122.13}{71}{}{proposition.122.13}{Dirichlet-GRH-v930-Volume-XXIII.pdf}}
\expandafter\def\csname r@thm:common111\endcsname{{9.13}{118}{}{proposition.9.13}{Dirichlet-GRH-v930-Volume-I.pdf}}
\expandafter\def\csname r@thm:conductor232\endcsname{{19.15}{12}{}{proposition.19.15}{Dirichlet-GRH-v930-Volume-IV.pdf}}
\expandafter\def\csname r@thm:covectorpaths352\endcsname{{73.1}{73}{}{proposition.73.1}{Dirichlet-GRH-v930-Volume-XIV.pdf}}
\expandafter\def\csname r@thm:criticalmarks339\endcsname{{22.12}{71}{}{proposition.22.12}{Dirichlet-GRH-v930-Volume-IV.pdf}}
\expandafter\def\csname r@thm:cycleweyl461\endcsname{{68.2}{88}{}{proposition.68.2}{Dirichlet-GRH-v930-Volume-XIII.pdf}}
\expandafter\def\csname r@thm:cyclicunit495\endcsname{{68.42}{115}{}{proposition.68.42}{Dirichlet-GRH-v930-Volume-XIII.pdf}}
\expandafter\def\csname r@thm:definitelevels678\endcsname{{130.2}{4}{}{proposition.130.2}{Dirichlet-GRH-v930-Volume-XXV.pdf}}
\expandafter\def\csname r@thm:direct76\endcsname{{66.24}{63}{}{proposition.66.24}{Dirichlet-GRH-v930-Volume-XIII.pdf}}
\expandafter\def\csname r@thm:divisionphase493\endcsname{{68.35}{110}{}{proposition.68.35}{Dirichlet-GRH-v930-Volume-XIII.pdf}}
\expandafter\def\csname r@thm:elliptic312\endcsname{{67.12}{86}{}{proposition.67.12}{Dirichlet-GRH-v930-Volume-XIII.pdf}}
\expandafter\def\csname r@thm:factor83\endcsname{{66.30}{70}{}{proposition.66.30}{Dirichlet-GRH-v930-Volume-XIII.pdf}}
\expandafter\def\csname r@thm:fouriersmith338\endcsname{{126.12}{22}{}{proposition.126.12}{Dirichlet-GRH-v930-Volume-XXIV.pdf}}
\expandafter\def\csname r@thm:framechange314\endcsname{{122.3}{58}{}{proposition.122.3}{Dirichlet-GRH-v930-Volume-XXIII.pdf}}
\expandafter\def\csname r@thm:fullrealperiod698\endcsname{{132.12}{60}{}{proposition.132.12}{Dirichlet-GRH-v930-Volume-XXV.pdf}}
\expandafter\def\csname r@thm:greenenergy488\endcsname{{68.29}{106}{}{proposition.68.29}{Dirichlet-GRH-v930-Volume-XIII.pdf}}
\expandafter\def\csname r@thm:lattice55\endcsname{{65.3}{23}{}{proposition.65.3}{Dirichlet-GRH-v930-Volume-XIII.pdf}}
\expandafter\def\csname r@thm:level85\endcsname{{66.34}{73}{}{proposition.66.34}{Dirichlet-GRH-v930-Volume-XIII.pdf}}
\expandafter\def\csname r@thm:levels88\endcsname{{7.11}{97}{}{proposition.7.11}{Dirichlet-GRH-v930-Volume-I.pdf}}
\expandafter\def\csname r@thm:long232\endcsname{{5.6}{23}{}{proposition.5.6}{Dirichlet-GRH-v930-Volume-I.pdf}}
\expandafter\def\csname r@thm:lowerlattice347\endcsname{{83.28}{74}{}{proposition.83.28}{Dirichlet-GRH-v930-Volume-XVI.pdf}}
\expandafter\def\csname r@thm:mean56\endcsname{{65.22}{40}{}{proposition.65.22}{Dirichlet-GRH-v930-Volume-XIII.pdf}}
\expandafter\def\csname r@thm:neighbormove356\endcsname{{73.3}{75}{}{proposition.73.3}{Dirichlet-GRH-v930-Volume-XIV.pdf}}
\expandafter\def\csname r@thm:neutral315\endcsname{{122.4}{59}{}{proposition.122.4}{Dirichlet-GRH-v930-Volume-XXIII.pdf}}
\expandafter\def\csname r@thm:one61\endcsname{{65.13}{33}{}{proposition.65.13}{Dirichlet-GRH-v930-Volume-XIII.pdf}}
\expandafter\def\csname r@thm:orientation46\endcsname{{6.23}{57}{}{proposition.6.23}{Dirichlet-GRH-v930-Volume-I.pdf}}
\expandafter\def\csname r@thm:poincaregram704\endcsname{{133.12}{72}{}{proposition.133.12}{Dirichlet-GRH-v930-Volume-XXV.pdf}}
\expandafter\def\csname r@thm:pointlattice353\endcsname{{99.20}{60}{}{proposition.99.20}{Dirichlet-GRH-v930-Volume-XIX.pdf}}
\expandafter\def\csname r@thm:pol49\endcsname{{64.10}{12}{}{proposition.64.10}{Dirichlet-GRH-v930-Volume-XIII.pdf}}
\expandafter\def\csname r@thm:polycenter460\endcsname{{9.26}{129}{}{proposition.9.26}{Dirichlet-GRH-v930-Volume-I.pdf}}
\expandafter\def\csname r@thm:prime209\endcsname{{93.9}{28}{}{proposition.93.9}{Dirichlet-GRH-v930-Volume-XVIII.pdf}}
\expandafter\def\csname r@thm:prime210\endcsname{{93.17}{34}{}{proposition.93.17}{Dirichlet-GRH-v930-Volume-XVIII.pdf}}
\expandafter\def\csname r@thm:projectedrank705\endcsname{{133.18}{77}{}{proposition.133.18}{Dirichlet-GRH-v930-Volume-XXV.pdf}}
\expandafter\def\csname r@thm:rangenew50\endcsname{{6.29}{62}{}{proposition.6.29}{Dirichlet-GRH-v930-Volume-I.pdf}}
\expandafter\def\csname r@thm:rankonecost455\endcsname{{120.14}{25}{}{proposition.120.14}{Dirichlet-GRH-v930-Volume-XXIII.pdf}}
\expandafter\def\csname r@thm:residuejet492\endcsname{{68.33}{109}{}{proposition.68.33}{Dirichlet-GRH-v930-Volume-XIII.pdf}}
\expandafter\def\csname r@thm:terminalformula347\endcsname{{127.1}{26}{}{proposition.127.1}{Dirichlet-GRH-v930-Volume-XXIV.pdf}}
\expandafter\def\csname r@thm:terminaltypes327\endcsname{{123.15}{96}{}{proposition.123.15}{Dirichlet-GRH-v930-Volume-XXIII.pdf}}
\expandafter\def\csname r@thm:threeaverages667\endcsname{{6.40}{86}{}{proposition.6.40}{Dirichlet-GRH-v930-Volume-I.pdf}}
\expandafter\def\csname r@thm:torsion312\endcsname{{66.35}{74}{}{proposition.66.35}{Dirichlet-GRH-v930-Volume-XIII.pdf}}
\expandafter\def\csname r@thm:trans57\endcsname{{65.6}{25}{}{proposition.65.6}{Dirichlet-GRH-v930-Volume-XIII.pdf}}
\expandafter\def\csname r@thm:translates706\endcsname{{133.20}{78}{}{proposition.133.20}{Dirichlet-GRH-v930-Volume-XXV.pdf}}
\expandafter\def\csname r@thm:weyllower351\endcsname{{120.6}{18}{}{proposition.120.6}{Dirichlet-GRH-v930-Volume-XXIII.pdf}}
\expandafter\def\csname r@thm:weylperiod351\endcsname{{119.10}{10}{}{proposition.119.10}{Dirichlet-GRH-v930-Volume-XXIII.pdf}}
\expandafter\def\csname r@thm:weylrowjet449\endcsname{{120.12}{23}{}{proposition.120.12}{Dirichlet-GRH-v930-Volume-XXIII.pdf}}
\makeatother
\begin{document}
\setcounter{part}{26}
\fancyhead[L]{\small\color{muted}Dirichlet GRH \textperiodcentered\ Volume XXVII}
\fancyhead[R]{\small\color{muted}Heights of projected theta periods}
\thispagestyle{plain}\begin{center}
{\LARGE\bfseries The generalized Riemann hypothesis\\[2pt]for Dirichlet \(L\)-functions}\\[6pt]
{\small 27 November 2026 \quad\textperiodcentered\quad Edition 930}\end{center}
\part{Heights, Gram operators and divisors of projected theta periods}\label{part:XXVII}
\begin{abstract}\noindent We compare finite period ideals with complex norms of the signed columns. Gram operators of the first lifts, marked determinants and terminal Fourier coefficients enter the height and common-divisor estimates.\end{abstract}
\setcounter{section}{140}
\noindent The \hyperref[seriesguide]{guide in Volume~I} lists all twenty-seven volumes.
\begingroup\small\setlength{\parskip}{0pt}\tableofcontents\endgroup
\clearpage

\section{Complex norms and finite period ideals}\label{nav:chapter-norms-and-period-ideals}\label{nav:volume-XXVII}

\subsection{Testing a signed period on rational unitary frames}
\label{sec:frameperiods446}

Fix \(x\) on one of the intervals of Section~\ref{sec:missing}, and put
\(E=E_x\), \(F=E^+\). Use the matched small endpoint, finite characters
and finite class measure of Section~\ref{sec:weylperiod351}.
A small class \(a\) has an underlying unimodular Hermitian lattice
\(M_a\), stabilizer \(\Gamma_a\), and a polynomial coefficient space at
each member \(\tau\) of the CM type. Choose an \(E\)-basis of the
rational space giving a basis of \(M_a\) at every place above
\(g_1,g_2,T\). Such a basis exists. In any starting basis the condition
at a finite place is an open condition on an invertible matrix, and
weak approximation on the entries meets these finitely many conditions.
Use the dual basis on the opposite coefficient line. These choices write
an algebraic output at \(a\) as a flag polynomial \(Q_a\) times the
fixed basis of its determinant line. We evaluate its polynomial part.
Restoring a determinant character on a unitary frame multiplies the
value by a unit at every place where the frame and its inverse are
integral.

\begin{lemma}[Simultaneous rational frames for the small classes]\leavevmode
\label{lem:rationalframes446}
At each \(a\) there is a finite set \(\mathcal F_a\) in
\(U(M_a\otimes E)(F)\), integral on \(M_a\) at all three marked
rational primes, with the following property. At each of these primes
\(p\), and at every coefficient place \(w\mid p\), its reductions in
the CM coordinates range over all of
\(\prod_{\tau}\mathrm{GL}_n(\mathbb F_p)\). One set works for all
polynomial weights on this Hermitian lattice.
\end{lemma}
\begin{proof}
The primes split completely in \(E\). For each \(v\mid p\) of \(F\)
the two factors of \(E\otimes_F F_v\) are paired by the Hermitian
form. An isometry is determined by its matrix on the first factor.
On the other it acts by the inverse transpose in paired bases.
Unimodularity makes the second lattice the full dual of the first,
so every invertible integral matrix on the first preserves both.
Choose one of these two factors at each \(v\), and lift any list
of invertible residue matrices to integral isometries there.

To meet a prescribed list of residue classes by a single rational
isometry, let \(H\) be the Hermitian matrix over \(E\) in our basis
and let \(X^\ast=H^{-1}\bar X^{\mathsf t}H\).
The transformation \(X\mapsto(I-X)(I+X)^{-1}\) identifies open
subsets of the \(F\)-vector space \(X^\ast=-X\) and the unitary
group. Its inverse is given by the same formula, wherever \(I+g\)
is invertible. Any of the chosen local isometries can be moved into
this open subset within its residue class. Multiply it on the
two paired factors by \(t,t^{-1}\) with \(t\in 1+p\mathbb Z_p\).
Only finitely many \(t\) give eigenvalue \(-1\). The inverse
transforms of the resulting isometries have open neighborhoods
which the transform carries into the prescribed classes. Weak
approximation in an \(F\)-basis of the skew-Hermitian space meets
all these neighborhoods at once. Its rational transform is defined,
since its denominator is nonzero at the chosen completions.
Taking one such lift of each list gives a finite set. Changing \(w\)
permutes the factors and may change a chosen factor to its paired
one. In integral bases inverse transpose and changes of basis induce
bijections on the invertible residue matrices. This proves the
assertion for every \(w\).
\end{proof}

For the small weight \(\mu_\tau\) put
\(m_\tau=\mu_\tau-\mu_{\tau n}\mathbf 1\) and
\(h=\max_\tau m_{\tau1}\). Use the bilinear highest-vector
normalization \(B=\bigotimes_\tau B_{m_\tau}\). For \(g\in\mathcal F_a\)
let \(l_{a,g}\) be the vector in the opposite coefficient fiber such that
\[
             B(Q,l_{a,g})=Q\big((\tau g)_\tau\big)
\]
in the chosen determinant frames. Concretely its polynomial at each
embedding is \(p_{m_\tau}((\tau g)^{\mathsf t}X_\tau)\), where
\(p_m\) is the product of the leading minors. This is the polynomial
identity of Lemma~\ref{lem:weyldual351}. Form a test \(t_{a,g}\)
supported on \(a\) by taking the sum of the \(\Gamma_a\)-translates
of \(l_{a,g}\), with the dual coefficient action and the inverse
finite character. At the other representatives set the value to zero
and extend by equivariance. Only the coefficient argument of the test
varies. The finite representative, its residue flags and its level
character remain those of \(a\).

\begin{proposition}[Integral generation of the terminal tests]\leavevmode
\label{prop:framegeneration446}
Let \(w\mid p\) be a place of the coefficient field above one of the
three marked primes. Suppose \(p>h+n\) and \(p\nmid|\Gamma_a|\) for
each small class. Then over the valuation ring \(O_w\) the tests
\(t_{a,g}\) generate the entire integral small test lattice \(S_w\).
For every invariant output table \(Y\), the period with \(t_{a,g}\)
is exactly \(Y(a)((\tau g)_\tau)\) in these frames.
\end{proposition}
\begin{proof}
In a single factor the dual Weyl lattice \(\nabla_m\) consists
of the flag polynomials with integral coefficients. The Weyl
lattice \(\Delta_m\) lies in it, by its construction from integral
row substitutions. The form \(B_m\) on this polynomial lattice
has monomial weights \(\alpha!/H_m\), with \(H_m\) the hook
product. Each hook is at most \(m_1+n-1\). Thus \(B_m\) is
integral on \(\nabla_m\times\nabla_m\) at \(w\).
Lemma~\ref{lem:weyldual351} identifies the dual of \(\nabla_m\)
with \(\Delta_m\), giving equality of the two lattices and a
perfect pairing. The same holds on their tensor product with the
paired determinant lines. In its polynomial realization this
tensor product is saturated in the ambient coefficient lattice.
Indeed each individual integral flag space is the intersection of
a vector subspace with the monomial lattice, hence a direct summand
over a discrete valuation ring. Tensoring the splittings proves
the assertion.

Each \(l_{a,g}\) is integral by its displayed polynomial and our
choice of bases. Suppose its evaluations failed to detect some
nonzero reduction of a vector in the paired lattice. By saturation
the reduced polynomial \(Q\) of that vector is nonzero. Each matrix
entry at \(\tau\) has degree at most \(m_{\tau1}\), since its entire
matrix column has degree at most \(m_{\tau1}\). The polynomial
\[
                    Q(Y)\prod_\tau \det Y_\tau
\]
is nonzero over the residue field of \(O_w\), of degree at most
\(h+1<p\) in each variable. A polynomial of degrees strictly less
than \(p\) in the separate variables and zero on the full grid
\(\mathbb F_p^d\) is zero, even with coefficients in an extension
field. Successive use of the one-variable root bound proves this.
Thus \(Q\) is nonzero on a tuple of invertible residue matrices.
Lemma~\ref{lem:rationalframes446} supplies a frame giving that tuple.
By perfect duality the reductions of the \(l_{a,g}\) therefore span
the opposite fiber, and Nakayama's lemma lifts this spanning
statement to \(O_w\).

The stabilizer trace maps that fiber onto its invariant submodule,
since on an invariant vector it is multiplication by the unit
\(|\Gamma_a|\). Tracing the generators, then summing over the
disjoint classes, proves generation of \(S_w\). Contracting an
invariant output with each term in the trace gives the same
value \(B(Y(a),l_{a,g})\) by duality of the character actions.
There are \(|\Gamma_a|\) terms and the class has measure
\(1/|\Gamma_a|\), proving the period identity.
\end{proof}

\begin{theorem}[Determinantal ideals of the actual signed columns at the marked primes]\leavevmode
\label{thm:frameperiods446}
At every sufficiently large \(x\) on the stated intervals, take any
finite list of the prescribed signed theta columns with a common
small testing space. For each column \(j\) let \(Y_j=R_1P_j\) be
its actual lowered small table, including all earlier operations.
Fix the finite central projector \(e\) of
Theorem~\ref{thm:terminaltypes327}, allowing also \(e=1\).
Write \(Z_j=e^{\mathsf t}Y_j\), where the transpose is for the
finite class pairing. At every \(w\mid g_1g_2T\), and for every
choice of \(s\) columns, the fractional ideal of the \(s\)-rowed
period determinants tested on \(eS_w\) is generated by
\[
   \det\big[Z_j(a_i)((\tau g_i)_\tau)\big]_{i,j=1}^{s},
        \qquad a_i\ \text{small classes},\quad g_i\in\mathcal F_{a_i}.
\]
It is also generated by the coefficient minors
\(\det[(Z_j(a_i))_{\alpha_i}]_{i,j=1}^{s}\), in the same polynomial
frames, with one joint monomial \(\alpha_i\) for each row. Both
identities include zero ideals. There is no loss in valuation at
any exterior degree in passing between these tests and coefficients.
\end{theorem}
\begin{proof}
The cutoff and stabilizer assertion in the proof of
Theorem~\ref{thm:terminalformula347} give the hypotheses of
Proposition~\ref{prop:framegeneration446} for all three primes
on these fields and intervals. Indeed the small weight width is
bounded by the large weight width by interlacing. The character
values and the original column entries lie in the common
coefficient field. The projector can have a denominator there.
Its transpose exists over the field because the pairings on the
opposite finite table spaces are perfect in characteristic zero,
by the classwise nondegenerate pairing and stabilizer averaging.
Since \(eS_w\) is generated by \(e t_{a,g}\), multilinearity of
the determinant proves the first equality of ideals from
\(\langle Y_j,e t_{a,g}\rangle=\langle Z_j,t_{a,g}\rangle\).

For fixed classes use independent matrix variables \(X_i\) in
each of the \(s\) rows of \(\det[Z_j(a_i)(X_i)]\). Its coefficients
are exactly the asserted determinants, by expansion in those
independent variables. Each evaluation is an integral combination
of them. For a nonzero determinant polynomial divide it by a
coefficient of least valuation. Its reduction is nonzero and its
degree in any matrix entry is at most \(h\), since any determinant
term uses just one polynomial in each row. Multiplication by one
determinant for every row and embedding allows the grid argument
in the proposition to find a tuple of invertible residues giving a
nonzero value. Choose the rational lifts separately in each row.
The resulting determinant attains the least coefficient valuation.
If the polynomial is zero, every evaluation and coefficient is zero.
Taking all class lists proves the second equality.
\end{proof}

For the saturated projected tests \(S_{0,w}=S_w\cap e(S_w\otimes
\operatorname{Frac}O_w)\) the distinction between the two integral
domains can likewise be kept exact. Let \(L\) send the standard
generators of a free \(O_w\)-module to the tests \(t_{a,g}\).
Compute \(K=\ker((1-e)L)\) over \(O_w\), clearing any scalar
denominator of \(1-e\) before computing the kernel.
Proposition~\ref{prop:framegeneration446} gives \(L(K)=S_{0,w}\):
lift any element of \(S_{0,w}\) through \(L\), and the lift
belongs to \(K\). Consequently a basis of \(K\), applied as row
combinations to the unprojected frame evaluations of \(Y_j\),
gives the period minors for \(S_{0,w}\) by multilinearity.
This computes both choices of integral domain on the given columns.
Their minimum valuations still depend on the coefficients of the
signed lifts. An estimate for those valuations and an identification
with the analytic incoming row are required in Section~\ref{sec:missing}.


\subsection{A common set of tests for compact norms and marked period ideals}
\label{sec:compactquadrature447}

We work on the small definite space at a class \(a\) of
Section~\ref{sec:frameperiods446}. Put \(\mathcal S=\{g_1,g_2,T\}\)
and \(G=\operatorname{Res}_{F/\mathbb Q}U(M_a\otimes E)\). Its real
points form a compact group. Give this group Haar measure of total
mass one. Each rational point of \(G\) gives a frame in the same
Hermitian space over \(F\). We shall use one finite set of such
frames both for the integral periods and for integration on the
compact group. For positive finite integration formulas over the
real numbers, see Bayer and Teichmann~\cite{XXVI-XX-bayer447}.

Write \(V\) for the underlying rational vector space of \(M_a\),
and \(\rho\) for the action of \(G\) on \(V\). In a rational basis
let \(W_d\) be the space of restrictions to \(G\) of polynomials
of total degree at most \(d\) in the entries of \(\rho\). It is a
finite-dimensional rational space of regular functions, unchanged
by a rational change of basis. In this section \(d\ge1\).

\begin{lemma}[The local denominator of compact integration]\leavevmode
\label{lem:integralhaar447}
Normalized integration on \(W_d\) is a rational linear functional
\(I\). At a prime \(p\) split completely in \(E\), let
\(K_p=\prod_{v\mid p}\mathrm{GL}_n(\mathbb Z_p)\) be the stabilizer
of the underlying unimodular lattice. If \(p>\max(2d+1,d+n)\),
the extension of \(I\) to \(\mathbb Q_p\) is integral on every
function in \(W_d\otimes\mathbb Q_p\) taking integral values on
\(K_p\).
\end{lemma}
\begin{proof}
The rational points of \(G\) are dense in \(G(\mathbb R)\).
Indeed the Cayley argument of Lemma~\ref{lem:rationalframes446}
also works at every archimedean place of \(F\). Multiplication by a
central unit complex number sufficiently near one avoids eigenvalue
\(-1\) there. Weak approximation on the skew-Hermitian coordinates
can include these neighborhoods and any finite list of the open
conditions in that lemma. The same Cayley chart shows that
\(G(\mathbb R)\) is Zariski dense, since a polynomial vanishing on
an open real coordinate box is identically zero.

The translates \(R_g f(x)=f(xg)\) preserve the real space of
functions \(W_d\otimes\mathbb R\). With its Haar \(L^2\) inner
product, the orthogonal complement of the span of
\(R_g f-f\), for rational \(g\), consists of the invariant
functions. By continuity and density these are the constants.
Thus \(I\) is the unique functional annihilating the translation
differences with \(I(1)=1\). A rational basis and the rational
translations give rational linear equations for it. Finitely many
of their rows have the full rank, so the solution is rational.
Uniqueness for this finite system persists over \(\mathbb Q_p\).

In paired local bases write \(Y_v\) for the matrix on one of the
two factors above \(v\). The action on the other factor is its
inverse transpose. After any fixed change of basis each entry of
\(\rho\) is a linear combination of entries of these matrices and
their inverse transposes. Write \(\Delta(Y)=\prod_{v\mid p}\det Y_v\).
Every \(f\in W_d\otimes\mathbb Q_p\) has a unique polynomial
numerator \(P_f\) in \(f=\Delta^{-d}P_f\). Its degree in any entry
is at most \(2d\). Indeed an inverse entry contributes a cofactor
divided by a determinant. In one factor, \(b\) inverse entries
and \(c\) forward entries give at most \(b+c+(d-b)=d+c\)
in each variable after clearing the denominator. A function
integral on \(K_p\) has \(P_f\) integral. Otherwise scale the
numerator to a primitive integral polynomial by a positive power
of \(p\). Its reduction times \(\Delta\) is nonzero, with every
individual degree less than \(p\). The grid argument of
Proposition~\ref{prop:framegeneration446} yields an invertible
residue tuple on which it is nonzero, contradicting integrality
of the original function. The converse follows since \(\Delta\)
is a unit on \(K_p\). The argument uses just one representative
of each invertible residue tuple.

For a square variable \(Y\) put \(\partial_Y=(\partial/\partial
Y_{ij})_{ij}\) and \(H=\prod_{j=0}^{n-1}(d+j)!/j!\). Then
\[
 I(f)=
 \left.\frac{\displaystyle\prod_{v\mid p}
              \det(\partial_{Y_v})^d P_f(Y)}
             {H^{[F:\mathbb Q]}}\right|_{Y=0}.
\]
To verify the formula, first take \(f=1\). The Fischer squared
norm of \((\det Y)^d\) is the hook product of the rectangle
\((d^n)\), namely \(H\), by Lemma~\ref{lem:hook86}.
For a right translation \(Y_v\mapsto Y_v A_v\) the numerator
polynomial is replaced by \(P_f((Y_v A_v)_v)/\Delta(A)^d\).
The chain rule multiplies its indicated determinant of derivatives
by \(\prod_v\det(A_v)^d\) at zero, cancelling the denominator.
The formula therefore satisfies the same rational translation
equations and normalization as \(I\). Uniqueness proves it.
The factorial ratios defining \(H\) are products of integers
at most \(d+n-1\). The differential operators have integral
coefficients, so an integral numerator gives an integral answer.
\end{proof}

\begin{theorem}[Positive rational integration weights at the three marked primes]\leavevmode
\label{thm:compactquadrature447}
Suppose each \(p\in\mathcal S\) exceeds \(\max(2d+1,d+n)\).
There is a finite set \(A_a\) of rational unitary frames integral
at every \(p\in\mathcal S\), containing the complete residue tests
\(\mathcal F_a\) of Lemma~\ref{lem:rationalframes446}, and positive
rational numbers \(\lambda_{a,g}\) such that
\[
       \sum_{g\in A_a}\lambda_{a,g}=1,\qquad
       I(f)=\sum_{g\in A_a}\lambda_{a,g} f(g)\quad(f\in W_d).
\]
Every \(\lambda_{a,g}\) is a unit at every prime in \(\mathcal S\).
The same formula integrates \(W_d\) after complex scalar extension.
\end{theorem}
\begin{proof}
In \(W_d^\vee\otimes\mathbb R\) write \(u_g\) for evaluation at
\(g\). The compact convex hull of the evaluations on \(G(\mathbb R)\)
has affine span \(u(1)=1\). Indeed any further affine relation
would give a nonconstant element of \(W_d\) equal to a constant
on the group. Its barycenter \(I\) is in the relative interior.
Otherwise a supporting affine functional nonconstant on this hull
would attain its maximum at its integral. Continuity and positive
Haar mass on every nonempty open set exclude this.

There is a finite set of rational frames integral at the chosen primes
whose evaluations have \(I\) in the interior of their convex hull
relative to the same hyperplane. To see this, take points in the hull
on both sides of \(I\) along each coordinate axis of that hyperplane,
all at a common small distance. Each is a finite convex combination of
evaluations. The resulting polytope contains a ball about \(I\).
Perturb the finitely many frames to rational frames integral at the
chosen primes by the simultaneous approximation in the preceding
proof. If all the evaluations move by less than half the radius,
every supporting linear functional of norm one still has maximum
strictly above its value at \(I\). Separation proves the assertion.
The zero-dimensional case requires just one frame.
Adjoin \(\mathcal F_a\) to this set and call the result \(B\).
Adjoin also a distinct frame \(g'\) for each \(g\in B\) congruent
to \(g\) at all the marked primes. Such frames exist by the
same approximation with varying archimedean neighborhoods, since
the unitary center has positive real dimension. Call the enlarged
set \(A_a\). There are strictly positive real coefficients on all
of \(A_a\) representing \(I\). In fact assign a sufficiently small
positive mass at every point. The barycenter required of the
remaining mass is still in the original interior neighborhood.

At each marked \(p\), the functions integral on \(K_p\) form a
full \(\mathbb Z_p\)-lattice \(\Lambda_p\) in \(W_d\otimes\mathbb Q_p\).
This follows by intersecting its numerator space with the finite
monomial lattice. The evaluations from \(\mathcal F_a\) span
\(\Lambda_p^\vee\) over \(\mathbb Z_p\). Indeed their dual lattice
is exactly \(\Lambda_p\) by the grid argument in
Lemma~\ref{lem:integralhaar447}, and double duality over a
discrete valuation ring gives the assertion. That lemma puts
\(I\) in their integral span. Write it as \(\sum_{g\in B}c_g u_g\)
with \(c_g\in\mathbb Z_p\). Since \(p>2\), write each residue
\(\bar c_g\) as a sum of two nonzero residues: choose the first
different from both zero and \(\bar c_g\). Lift the two to units
at \(g,g'\). Their weighted evaluations reproduce \(I\) modulo
\(p\Lambda_p^\vee\), because integral numerators divided by
\(\Delta^d\) take the same residue at \(g\) and \(g'\).
The old evaluations span \(\Lambda_p^\vee\), so changing their
coefficients by elements of \(p\mathbb Z_p\) removes the error.
This gives an exact local solution with all coefficients units.

The equations on the coefficients are rational linear equations.
Gaussian elimination makes their solution space a rational affine
space, since they admit a real solution. Approximate the strictly
positive solution at infinity and the constructed unit solutions
at the three primes in its free rational coordinates. Positivity
and being a local unit are open conditions. The resulting single
rational solution gives the asserted weights.
\end{proof}

We now apply the integration formula to the signed columns
\(Y_j=R_1P_j\) and \(Z_j=e^{\mathsf t}Y_j\) of
Theorem~\ref{thm:frameperiods446}. For a small weight set
\[
                  d=\max\bigl(1,2\sum_\tau |m_\tau|\bigr),
           \qquad |m_\tau|=\sum_{i=1}^n m_{\tau i}.
\]
One may take the largest such \(d\) for all the common small types
in use at \(x\). This is \(O(x^C b)\) for a fixed \(C\).
Indeed the field degrees and ranks have polynomial bounds and
each width is \(O(xb+x^{C'})\) by Theorem~\ref{thm:terminalformula347}.
The bounds \(g_i>(kb)^8\) and \(\log T\gg m x^\eps\), with
\(m\longrightarrow\infty\), put these same three primes above
\(\max(2d+1,d+n)\) for all sufficiently distant \(x\).

For a complex embedding \(\sigma\) use the transported positive
Hermitian form on the small fiber in the highest-vector Fischer
normalization. Write \(r_\mu\) for its complex dimension. We may
and do choose the rational basis at each class to be orthonormal
and integral at all three marks. Indeed the small definite group
in Theorem~\ref{thm:terminaltypes327} is the isometry group of
the standard Hermitian space \(E^n\), as in
\srcloc{Formalism/LocalToGlobal/Columns/StdLatUnimodular}{29}.
Start in the standard rational orthonormal basis. At each split
\(v\) choose a basis of the first factor of \(M_a\) and its paired
dual basis in the second. The resulting basis in the product
algebra is orthonormal, by unimodularity. A local unitary matrix
takes the starting basis to this one. Generating that same local
lattice is an open condition on the matrix. The rational Cayley
approximation used in Lemma~\ref{lem:rationalframes446} meets all
these conditions at once. Its rational matrix is exactly unitary,
so the resulting basis is orthonormal at every infinite place.
In these frames the representing vector of evaluation at the
identity is the tensor of the \(p_{m_\tau}\), of norm one by
Lemma~\ref{lem:weyldual351}, with the unit determinant vector for
the separate line. All the tests in what follows use these frames.

\begin{theorem}[The definite Gram determinant on the same rational tests]\leavevmode
\label{thm:commoncompactgram447}
Let \(G_\sigma\) be the Hermitian Gram matrix of \(s\) such projected
signed columns in the small class metric with masses
\(1/|\Gamma_a|\). Put
\[
  z_{a,g;j}=\langle Y_j,e t_{a,g}\rangle,\qquad
  q_{a,g}=\frac{r_\mu\,\lambda_{a,g}}{|\Gamma_a|}.
\]
For \(g\in A_a\) use the stabilizer trace defining \(t_{a,g}\)
in Section~\ref{sec:frameperiods446}. Then
\[
    \det G_\sigma=
       \sum_{\substack{B\subseteq\{(a,g):g\in A_a\}\\ |B|=s}}
        \left(\prod_{(a,g)\in B}q_{a,g}\right)
        \left|\det[\sigma z_{a,g;j}]_{(a,g)\in B,j}\right|^2 .
\]
All the weights are positive rational numbers, independent of
\(\sigma\) and units at the three marks. At every marked coefficient
place the minors in the sum generate the exact determinantal ideal
on \(eS_w\).
\end{theorem}
\begin{proof}
Fix a class and an embedding. The product of the small flag
modules, meaning their external tensor product over the real
places, is an irreducible unitary module of the product of the
compact unitary groups. This follows from the irreducibility of
each highest-weight representation and the tensor product rule for
a direct product. Let \(l\) be its functional evaluating at the
identity. The average over the group of the positive operators
\((l\circ g)^*(l\circ g)\) commutes with the group. Its trace is
\(\|l\|^2=1\), so Schur's lemma makes it \(1/r_\mu\) times the
identity. Row substitution translates the functional at the
identity into matrix evaluation. In the orthonormal bases the
transpose and inverse conventions both preserve normalized Haar
measure on \(U(n)\). Removing the
determinant character multiplies each translated functional by
a scalar of absolute value one, which cancels in these products.

For vectors \(P,Q\) in this fiber, the product of the polynomial
value \(P(g)\) and the conjugate of \(Q(g)\) belongs to
\(W_d\otimes\mathbb C\). Each entry of the complex matrix at a
CM embedding is a complex linear combination of entries of
\(\rho(g)\), by extension of scalars in the restriction of
scalars representation. Those real matrix entries are real on
\(G(\mathbb R)\), so the same assertion holds for their complex
conjugates. The product degree is at most \(2\sum_\tau |m_\tau|\).
This also applies at every \(\sigma\), which permutes the CM
embeddings or selects their conjugates. Theorem~\ref{thm:compactquadrature447}
integrates each such product exactly with the same rational
weights. Applying the preceding operator identity to the small
values of the \(Z_j\), and inserting the class masses, gives the
Gram matrix of the array \((\sqrt{q_{a,g}}\,\sigma
z_{a,g;j})\). Here the trace identity in
Proposition~\ref{prop:framegeneration446} identifies each entry
before taking absolute values. Cauchy--Binet gives the formula.
Finally \(\mathcal F_a\subseteq A_a\), and every additional frame
is integral on the same lattices at the marked primes. Its trace
test belongs to \(S_w\), so applying \(e\) gives a test in \(eS_w\).
The two inclusions of the determinantal ideals follow from
Theorem~\ref{thm:frameperiods446}. The stabilizer orders are units
at the marks by Theorem~\ref{thm:terminalformula347}. The factors
in the Weyl dimension formula for \(r_\mu\) are
\((m_{\tau i}-m_{\tau j}+j-i)/(j-i)\), \(1\le i<j\le n\).
Both integers in each ratio lie strictly between zero and \(p\)
at all three primes. The weights \(\lambda_{a,g}\) are also units
there by construction, proving the assertion about \(q_{a,g}\).
\end{proof}

\begin{corollary}[A simultaneous minor at three prescribed coefficient places]\leavevmode
\label{cor:threeminor447}
For each \(p\in\mathcal S\) choose one coefficient place \(w_p\mid p\).
For fixed \(s\) columns of full rank on the projected tests,
there are \(s\) semilocally integral tests before projection for which
the same period determinant \(z\), after projection, attains the
least valuation at each of the three \(w_p\). Set
\(P_{\mathcal S}=g_1g_2T\), and let \(M\) bound all \(|\Gamma_a|\)
as in Theorem~\ref{thm:weylperiod351}. For every complex embedding
one has
\[
                |\sigma z|^2\le(9 M P_{\mathcal S}^2)^s\det G_\sigma .
\]
The logarithm of this comparison factor per exterior degree is
\(o(b/x^A)\) for every fixed \(A>0\) on the same intervals.
\end{corollary}
\begin{proof}
For \(w_p\), Theorem~\ref{thm:frameperiods446} gives rows
\(t_{i,p}=t_{a,g}\) achieving its minimum. Choose rational integers
\(0\le c_p<P_{\mathcal S}\) equal to one modulo \(p\) and zero
modulo the other two primes, by the Chinese remainder theorem.
Use \(t_i=\sum_p c_p t_{i,p}\) before projection. At \(w_p\) the
determinant expansion has one term of minimum valuation multiplied
by \(c_p^s\). Every other term has one coefficient divisible by
\(p\), and its period determinant has valuation at least the same
minimum since its rows too are permissible. This proves simultaneous
attainment.

At infinity the norm of the linear functional on a small class
table given by \(t_{a,g}\) is at most
\(\sqrt{|\Gamma_a|}\). In fact the trace and mass
cancel on invariant values, evaluation on that fiber has norm
one by unitarity, and taking its class value
costs at most \(\sqrt{|\Gamma_a|}\) in the mass norm. Thus each
combined row functional has norm at most
\(3P_{\mathcal S}\sqrt{M}\). Apply their exterior product to
\(Z_1\wedge\cdots\wedge Z_s\). Its norm is at most the product
of their norms by the Gram determinant inequality, and the squared
norm of the column wedge is \(\det G_\sigma\). The upper bounds
on the prime sizes in Theorems~\ref{thm:levels88}
and~\ref{thm:common111} make \(\log P_{\mathcal S}\) polynomial
in \(x\), and Theorem~\ref{thm:weylperiod351} gives the same for
\(\log M\). The asserted comparison follows from \(b=\exp(x^\eps)\).
\end{proof}

The primes in the corollary can be those of
Theorem~\ref{thm:criticalmarks339}. Thus on these same \(E_x\)
the critical-sign summands \(e_{r+1}\Sel_{p,r}(E_x)\) vanish
through that theorem's weight range, while the comparison above
uses the same three primes for the same signed determinant.
The rational frames can introduce denominators away from these
primes. The formula determines the Gram matrix of the definite
class tables. Comparison with the indefinite cuspidal period and
a uniform estimate of the signed coefficients are still needed
for Section~\ref{sec:missing}.


\subsection{All finite period ideals and the definite Gram determinant}
\label{sec:globalperiodcontent448}

Fix the fields, weights and class measures of
Section~\ref{sec:compactquadrature447}. Work in a common finite
Galois coefficient field \(J\subset\mathbb C\) and write \(c\) for
complex conjugation in this embedding. At \(w\mid p\) use the
valuation \(v_w(p)=1\). Let \(Y_1,\ldots,Y_s\) be the actual lowered
signed columns for a common small type, and set \(Z_i=e^{\mathsf t}Y_i\).
The finite central projector \(e\), which may be the identity, acts
on the test side. On each class the bilinear period pairs opposite
coefficient spaces by the highest-vector Fischer form with weight
\(1/|\Gamma_a|\). Denote this sum by \([Y,F]\).

For \(D=eS\) or \(D=S_0=S\cap e(S\otimes J)\), define
\(\mathfrak a(D)\) to be the fractional ideal generated by
\(\det([Y_i,F_j])_{i,j=1}^s\) for \(F_1,\ldots,F_s\in D\).
Use the abbreviations \(\mathfrak a_{\rm im}=\mathfrak a(eS)\) and
\(\mathfrak a_0=\mathfrak a(S_0)\). We assume the columns are
independent on the projected tests, so these ideals are nonzero.
They can be computed after localization on local generators of the
corresponding lattices, by multilinearity and flatness. Write
\(b_{{\rm im},w}=v_w(\mathfrak a_{\rm im})\) and
\(b_{0,w}=v_w(\mathfrak a_0)\).

There is a conjugate-linear identification \(C\) from the algebraic
column space to the opposite test space. In standard rational
orthonormal coordinates conjugate the coefficients of the flag
polynomials and use the dual unit vector of each determinant line.
The bilinear Fischer contraction of \(P\) with \(CQ\) is the
Hermitian form of the small fiber, linear in \(P\). In particular
\[
             \mathcal D=\det\big([Z_i,CZ_j]\big)_{i,j=1}^s
\]
is an element of \(J\), strictly positive at the chosen complex
embedding by the independence hypothesis. Here all finite
characters have modulus one, so conjugation changes a column's
stabilizer character to the inverse character of a test.

\begin{lemma}[Conjugating a primitive exterior column into the test lattice]\leavevmode
\label{lem:exteriorcontent448}
Fix \(w\). In the direct sum of the class coefficient Weyl
lattices at \(w\), every exterior coordinate of
\(\xi=Z_1\wedge\cdots\wedge Z_s\) has valuation at least
\(b_{{\rm im},w}\). At the conjugate place \(cw\), conjugation
carries the corresponding local Weyl lattice integrally into
the opposite dual Weyl lattice at \(w\). Consequently, if
\(t\in J_w^\times\) has \(v_w(t)=b_{{\rm im},cw}\), then
\[
                    t^{-1}C\xi\ \in\ \bigwedge\nolimits^s S_{0,w}.
\]
Here \(C\) on an exterior power acts on every factor.
\end{lemma}
\begin{proof}
Choose a basis at \(w\) of the direct sum of the full coefficient
lattices. Their integral duals for the classwise contraction
are the opposite dual Weyl lattices by
Lemma~\ref{lem:weyldual351}. They include the dual basis
functionals. For any one such functional, take the stabilizer
trace on each class, including its inverse level character.
This gives an element of \(S_w\). The class-mass pairing of this
test with \(Z_i\) equals the value of the original functional
on \(Z_i\). Indeed each term of the trace has the same pairing
on an invariant output, and the number of terms cancels the
class mass. This calculation holds also when \(p\mid|\Gamma_a|\),
since only the unnormalized trace was used in the integral test.
Pairing \(Z_i\) with a test in \(S_w\) equals pairing \(Y_i\)
with its image under \(e\). Thus every coordinate minor for
\(\xi\) belongs to \(\mathfrak a_{\rm im}O_w\).

For the integral effect of \(C\), start with the Hermitian
lattice \(M_a\) itself. In standard orthonormal coordinates its
conjugate is its bilinear dual: integrality of
\(\sum_i x_i\bar y_i\) for all \(x\in M_a\) is equivalent to
\(y\in M_a\), by unimodularity. After extension of scalars the
same equality pairs the factor at \(cw\) with the opposite one
at \(w\). Thus conjugating an integral basis of a factor at
\(cw\) gives an integral basis of the opposite lattice at \(w\).
The flag construction respects changes of basis in
\(\mathrm{GL}_n(O_w)\). In particular the leading-minor polynomial
and its iterated integral row coefficients defining \(\Delta_m\)
have integral coefficients in the paired frames. Hence conjugation
carries \(\Delta_m\) into the opposite integral polynomial
lattice \(\nabla_m\) there. Exterior determinant lines of any
integer power on the two sides pair integrally by the same
unimodularity, using the dual line for a negative power.
Tensoring and taking the direct sum gives the required integral
map in every residue characteristic.

The central idempotent is self-adjoint for the Hermitian class
metric. On a common full residue marking the group acts
unitarily, and its central coefficients
\(\sum_\chi\chi(1)\chi(g^{-1})\) at \(g^{-1}\) are the conjugates
of those at \(g\). Restriction to the original level gives
Theorem~\ref{thm:terminaltypes327}. In particular
\(C e^{\mathsf t}=e C\), by the defining bilinear and Hermitian
adjoint identities. It follows that \(C\xi\) lies rationally in
the exterior power of the projected test space. By the first
part at \(cw\) and the integral map just proved, \(t^{-1}C\xi\)
belongs to the exterior power of the full ambient test lattice.
Inside this lattice the submodule \(S_{0,w}\) is the intersection
with the rational projected test space. Indeed the invariant table
is obtained by the rational stabilizer equations and intersection
with the ambient lattice, and \(S_0\) imposes the additional
rational equation \(eF=F\). This intersection is saturated and
hence a direct summand over the valuation ring. A basis extending
a basis of the summand shows that the same intersection equality
holds in exterior degree \(s\). This proves the assertion.
\end{proof}

\begin{theorem}[An all-prime inequality for the actual projected definite periods]\leavevmode
\label{thm:globalperiodcontent448}
For the independent list above one has the fractional ideal
memberships
\[
   \mathcal D\ \in\
     \mathfrak a_0\,c(\mathfrak a_{\rm im})
                   \ \cap\ \mathfrak a_{\rm im}\,c(\mathfrak a_0).
\]
For every embedding \(\sigma:J\hookrightarrow\mathbb C\) let
\(G_\sigma\) be the positive Gram matrix of \(\sigma Z_1,\ldots,
\sigma Z_s\) in the class metric as in
Theorem~\ref{thm:commoncompactgram447}. With the absolute norm on
fractional ideals, this gives
\[
   \log N\mathfrak a_0+\log N\mathfrak a_{\rm im}
        \ \le\ \sum_{\sigma:J\hookrightarrow\mathbb C}\log\det G_\sigma .
\]
In particular this comparison between finite valuations and complex
Gram determinants includes every coefficient prime of the original
definite columns, with the specified signs, level and class masses.
\end{theorem}
\begin{proof}
At \(w\) consider the alternating functional on \(s\) tests given
by the determinant of the periods with \(Y_1,\ldots,Y_s\).
On \(\bigwedge^s S_{0,w}\) its image is
\(\mathfrak a_0O_w\) by definition. The action of \(e\) on
these tests is the identity, so using \(Z_i\) gives the same
functional. At \(C\xi\) its value is \(\mathcal D\).
Lemma~\ref{lem:exteriorcontent448} therefore gives
\[
                  v_w(\mathcal D)\ge b_{0,w}+b_{{\rm im},cw}.
\]
Since \(c(\mathcal D)=\mathcal D\), applying this at the conjugate
place gives the second membership. No choice of a single
simultaneous generator of the local ideals is needed.

For an embedding \(\sigma\) write
\(\sigma'=\operatorname{conj}_{\mathbb C}\circ\sigma\circ c\).
In standard rational orthonormal polynomial coordinates use the
direct sum of the complex coefficient spaces over the class
representatives, with the rational positive masses and normalized
Fischer forms. In this ambient Hilbert space \(\sigma\mathcal D\)
is the Hermitian product of the wedges with coordinate arrays
\(\sigma\xi\) and \(\sigma'\xi\). Indeed conjugating the latter
coordinates gives exactly \(\sigma(c\xi)\), while the monomial
weights in the Fischer pairing and the masses are rational.
In particular this computation does not require \(c\) to
commute with each embedding. Cauchy--Schwarz gives
\[
          |\sigma\mathcal D|\le
                 (\det G_\sigma\,\det G_{\sigma'})^{1/2}.
\]
Both determinants are positive, since linear independence persists
under the embeddings. The map \(\sigma\mapsto\sigma'\) permutes
the embeddings. Consequently the sum of their log determinant
bounds is \(\sum_\sigma\log\det G_\sigma\).
At the finite places the first membership gives
\(\log N\mathfrak a_0+\log N\mathfrak a_{\rm im}
\le \log|N_{J/\mathbb Q}\mathcal D|\) by factorization of
fractional ideals. Summing the archimedean inequality gives
the result.
\end{proof}

\begin{corollary}[Uniform cost of using saturated projected tests]\leavevmode
\label{cor:saturatedcontent448}
Define \(B(D)=(s[J:\mathbb Q])^{-1}\log N\mathfrak a(D)\) and
\[
       H=\frac{1}{s[J:\mathbb Q]}
             \sum_{\sigma:J\hookrightarrow\mathbb C}\log\det G_\sigma .
\]
On the original fields and intervals,
\[
      2B(eS)\le H,\qquad
      2B(S_0)\le H+\log m_G ,
\]
where \(m_G\) is the order of the two marked finite block groups
taken together. For \(e=1\) the latter cost can also be taken
to be zero. On the prescribed intervals \(\log m_G=o(b/x^A)\)
for every fixed \(A>0\).
\end{corollary}
\begin{proof}
The inclusions \(m_G eS\subseteq S_0\subseteq eS\) from
Theorem~\ref{thm:terminaltypes327} give
\[
       m_G^s\mathfrak a_{\rm im}\subseteq\mathfrak a_0
                     \subseteq\mathfrak a_{\rm im}
\]
by evaluating the determinant. Thus
\(0\le B(S_0)-B(eS)\le\log m_G\), while the theorem gives
\(B(S_0)+B(eS)\le H\). These prove both inequalities.
For \(e=1\) the two lattices agree. The logarithm of the finite
group order is polynomially bounded in \(x\) by
Theorem~\ref{thm:terminaltypes327}, which proves the final
estimate. In the local normalization above,
\(\log N\mathfrak a(D)=\sum_p\log p\sum_{w\mid p}
[J_w:\mathbb Q_p]v_w(\mathfrak a(D))\).
All three normalized quantities are unchanged on extending
the coefficient field: localization and flat base change preserve
the period ideals, the rational invariant subspaces and their
integral intersections, while each complex embedding contributes
with the same multiplicity. The period field is still \(E_x\).
\end{proof}

The right side is the actual definite Gram quantity computed in
Theorem~\ref{thm:commoncompactgram447}. To use this result in
Section~\ref{sec:missing} one still needs a lower valuation
estimate encoding the prescribed arithmetic counts, and an
upper estimate for these Gram determinants in the specified
analytic row, on a family of sufficient rank. The complementary
strict cohomology summands in Proposition~\ref{prop:res23}
also enter the interval estimate.


\subsection{The signal prime count on the signed definite columns}
\label{sec:signedsignal450}

Fix a translate and the ordinary words of Theorem~\ref{thm:common111}
on its field \(E_x\). We work with the algebraic class pairing of
Theorem~\ref{thm:terminalformula347}, in the marked elliptic
differential. The high set is \(T_{\rm hi}\subset\Phi_0\).
Write \(n=|\Phi_1|\), \(a=\lceil xb\rceil\),
\(b=\exp(x^\varepsilon)\), \(X'=\exp(x^{\varepsilon/3})\),
\(b'=xb/X'\). For \(1\le i\le m\), the count uses the weights
\(e_{i,\tau}\) following \eqref{eq:letters111}.
At an unramified prime inert in \(K_d\) and split in \(K_q\) put
\[
 U_p=\sum_{i=1}^m\sum_{t\in T_{\rm hi}}
          \sum_{\substack{j\ge1\\j\ {\rm odd}}}
      [a/p^j-e_{i,\pi^jt}]_+,\qquad \pi=\sigma_pc_E .
\]
Its contribution to the uncentered count in this sign is
\(2U_p\log p\) by \eqref{eq:cp73}.

\begin{lemma}[The first row at the actual word weights]\leavevmode
\label{lem:wordfirstrow450}
On every sufficiently distant common interval, uniformly in the
word, translate and edge, the partition \(\nu_i\) of the \(i\)-th
definite kernel satisfies
\[
     G_p(\nu_i)\ \ge\
          \sum_{t\in T_{\rm hi}}\sum_{j\ {\rm odd}}
                 [a/p^j-e_{i,\pi^jt}]_+ .
\]
Here \(G_p\) is the integer in Theorem~\ref{thm:weylrowjet449}.
\end{lemma}
\begin{proof}
Write \(d=t_{\rm r}+1\), \(r=k+(i-1)d\), \(s=k+id\).
At stage \(i-1\) the first highest weight is
\(u_1-(k-1)/2-(i-1)d/2\). Indeed the initial algebraic weight
subtracts the half-sum of positive roots, with first coordinate
\((k-1)/2\), from \(u_1\). At each step the input determinant
shift is \(e_{i,\tau}-s/2\) and the output subtracts
\(e_{i,\tau}-r/2\) from the zero-padded partition.
Thus previous entries decrease by \(d/2\) per step, and a
new entry equals \(-e_{i,\tau}+r/2\). These are the shifts
in the prescribed polynomial substitution, given in
\srcloc{Formalism/LocalRingCardinality/Params}{71}.
In particular the partition obtained by adding the input shift
has first row
\[
               \nu_{i,\tau,1}=u_1(\tau)+e_{i,\tau}-(r+s-1)/2.
\]
In the polynomial model one has first subtracted the last
incoming weight \(c\) to trivialize its determinant line.
The determinant power in \eqref{eq:sub81} is consequently
\(c+e_{i,\tau}-s/2\). At an edge after the first this is
\(2H_0-d\), by the formula for a new entry. At the first edge it
is nonnegative by the choices of \(B_\tau,H_0\) in
\eqref{eq:letters111}, also at a high place with \(c=-\lceil b'\rceil\).
This identifies the partition in the operator.

At \(t\in T_{\rm hi}\), \(u_1=a+(k-1)/2\) and
\(e_{i,t}\ge b'\). At \(\mu\in\Phi_1\), \(u_1\) is bounded by
a fixed power of \(x\). The ranks and \(n\) also have polynomial
bounds. Thus \(\nu_{i,t,1}\ge a+b'/2\) and
\(\nu_{i,\mu,1}\le e_{i,\mu}+C x^B\) for fixed \(B,C\).
All \(e_{i,\mu}\) on \(\Phi_1\) are at least \(b/(nX')\).
If \(M=p^j\) contributes positively at \(t\), then
\(M<a/e_{i,\pi^jt}\le2nxX'\). Hence
\[
    \frac{b'}{2M}\ \ge\ \frac{b}{4n(X')^2}>s+Cx^B
\]
eventually. The last inequality follows on taking logarithms,
since \(x^\varepsilon-2x^{\varepsilon/3}\) dominates every
constant multiple of \(\log x\). As
\(d_M(t',s)\ge t'/M-s\), it follows for each active \(M\) that
\[
       d_M(\nu_{i,t,1},s)-\nu_{i,\pi^jt,1}
                  \ \ge\ a/M-e_{i,\pi^jt}.
\]
The positive parts also give the claim for inactive \(M\).
Sum the \(l=1\) terms of \(G_p\). Its other terms are nonnegative.
\end{proof}

Let \(Y_1,\ldots,Y_q\) be columns with this translate and a
common small testing space, arising from the full integral
initial Weyl tables, with the prescribed signs and coupled last
edges. Use \(S_0,eS,m_G\) as in Theorem~\ref{thm:weylperiod351};
here \(S_0\) denotes the saturated projected test lattice.
Let \(\mathcal S\) be any finite set of the unramified primes of
the indicated sign, \(p\ge5\), at which the finite class cosets
in the path formula have prime-to-\(p\) shifts and the good
local vacuum of Theorem~\ref{thm:direct76}. These requirements
concern the supplied tests before applying any derivatives.
For instance, in the presentation of Theorem~\ref{thm:neutral315}
the primes prime to all its original shift orders \(Q_i^0\)
meet the shift condition. One can use representatives equal to
the identity at these primes. The proof of that theorem chooses
representatives supported at completely split primes outside any
prescribed finite set. Thus at \(p\) the standard frames are
integral frames on the common lattices of consecutive class maps.

We record separately the scalar coefficients that may occur if
the finite tests have denominators. In a given column expand the
finite coset sections in translates of the good common vacuum,
keeping the matrix coset restrictions and the shared coupling
label on the last two edges. For a path \(\gamma\), including
one summand of each such expansion, let \(a_\gamma\) be the
product of its external scalar weights, its coefficients of
translation, and its determinant-character scalars. Include
any algebraic change of the marked scalar normalization in this
product. Use integral determinant frames at \(p\). Exclude the
stabilizer masses and the terminal test projector from \(a_\gamma\).
Zero terms can be omitted. Define
\[
 L_{i,w}=\max\big(\{0\}\cup
    \{-v_w(a_\gamma):\gamma\text{ in column }i\}\big),\qquad w\mid p .
\]
These numbers use only the coefficients before applying the
theta derivatives. In particular they vanish for integral
coefficients. The expansion by normalized finite lattice
intertwiners, including the reciprocal indices and the splitting
characters, is given in the proof of Theorem~\ref{thm:classden312}.
Scaling an additive differential by \(c_\tau\) scales a kernel
of block degree \(|\nu_\tau|\) by \(c_\tau^{|\nu_\tau|}\).
This explains how to include an algebraic differential rescaling
if a different normalization of that line is used.

For each of the \(m\) source vertices of an edge, and once for
the small endpoint, take \(M_v=|\mathrm{GL}_{n_v[E_x:\mathbb Q]}
(\mathbb F_3)|\), where \(n_v\) is its rank. Set \(D=\prod_v M_v\).
In the definitions of \(\mathfrak a_0,\mathfrak a_{\rm im}\)
in Section~\ref{sec:globalperiodcontent448} use these \(q\) columns.
Their valuations are denoted by \(b_{0,w},b_{{\rm im},w}\),
allowing \(+\infty\) for a zero ideal.

\begin{theorem}[The good signal primes on both projected test lattices]\leavevmode
\label{thm:signedsignal450}
At each \(p\in\mathcal S\) and each \(w\mid p\), one has
\[
 \begin{aligned}
 b_{0,w}&\ge q U_p-\sum_{i=1}^q L_{i,w}-qv_p(D),\\
 b_{{\rm im},w}&\ge q U_p-\sum_{i=1}^q L_{i,w}-qv_p(Dm_G).
 \end{aligned}
\]
For the identity projection one takes \(m_G=1\).
The logarithm of \(D m_G\) is \(o(b/x^A)\) for every fixed
\(A>0\) on the same intervals.
\end{theorem}
\begin{proof}
At an unrefined initial class, any nonzero terminal covector
tuple on a path gives an allowed refined input on its Weyl
lattice, as in Theorem~\ref{thm:neighbormove356}. The lattice
modification made at \(T\) localizes to the original lattice at
\(p\) since \(T\) splits in \(K_d\). Starting at this input,
apply the coefficient map on each matrix coset. In the marked
differential the rule is to twist the source polynomial by the
determinant power, substitute \(D^{\mathsf t}Y_{[r]}\), and
replace \(D^\alpha\) by the undivided derivative of the
corresponding CM section. For a finite coset indicator this is
the same rule on its finite sum of translated vacua, by the
linearity of differentiation. Its scalar factors are exactly
the external factors recorded above. This describes
\eqref{eq:sub81} at consecutive vertices and the coefficient
maps in Theorem~\ref{thm:covectorpaths352}. In the archive the
substitution and evaluation are written in
\srcloc{Formalism/LocalToGlobal/KernelLattices/ThDistC}{27}.

At \(w\), each prime-to-\(p\) translated product with its
coefficient removed meets the local germ hypothesis. Its
ratio to the reference vacuum is an integral elliptic germ
by Corollary~\ref{cor:class66} and its proof; this also follows
directly by multiplying the individual translation ratios.
Theorem~\ref{thm:weylrowjet449} puts the output in the Weyl
lattice times \(p^{G_p(\nu_j)}\) at edge \(j\). Changes between
two integral frames preserve this lattice. Indeed the Fischer
adjoint is transpose substitution on the integral polynomial
dual, which is preserved by every integral invertible change
of variables. Determinant lines are changed by units.
Induction on the edges and Lemma~\ref{lem:wordfirstrow450} give
a valuation of at least \(U_p\) in every terminal Weyl coordinate
before inserting the scalar weights and masses. The inequality
involving a real \(U_p\) means that each coordinate has at
least that valuation. It follows from the integer exponent
\(\sum_jG_p(\nu_j)\).

Now use the path identity of Theorem~\ref{thm:covectorpaths352}.
It sums only paths with the specified initial covector, and keeps
each coupling label common to both incident maps. On each path
the original scalar product costs at most \(L_{i,w}\).
Its stabilizer order at each vertex divides \(M_v\), by the
faithful integral action and torsion-free level-three kernel
used in Theorem~\ref{thm:weylperiod351}. Clearing \(D\) thus
clears every path mass, including the endpoint. Summing with
the original signs preserves the common bound. The terminal
restriction acts integrally on Weyl lattices in matched frames
by Theorem~\ref{thm:weyllower351} and
Proposition~\ref{prop:normalframe349}. The contraction with an
integral small test is integral by Fischer duality.
This proves the bound for a single entry with tests in \(S_0\).
For tests in \(eS\), multiplication by \(m_G\) puts every test
in the integral small lattice, by the finite character sum for
the projector. In each term of a \(q\)-row determinant every
column occurs exactly once. Adding the entrywise bounds and
then taking the ideal of all such determinants proves both
inequalities. Finally \(\log|\mathrm{GL}_d(\mathbb F_3)|
\le d^2\log3\); all vertex ranks, their number and the field
degree are polynomial in \(x\). The same polynomial bound for
\(\log m_G\) is given by Theorem~\ref{thm:weylperiod351}.
\end{proof}

\begin{corollary}[The good-prime count and the Gram determinant of the same columns]\leavevmode
\label{cor:gramsignal450}
Suppose the chosen columns have rank \(q\) on the projected
tests. Write \(\omega_w=[J_w:\mathbb Q_p]/[J:\mathbb Q]\) and set
\[
 \begin{split}
 L_{\mathcal S}&=\sum_{p\in\mathcal S}\log p
                  \sum_{w\mid p}\omega_w\sum_{i=1}^q L_{i,w},\\
 B_{\rm out}&=\sum_{p\notin\mathcal S}\log p
     \sum_{w\mid p}\omega_w
           \big([-b_{0,w}]_++[-b_{{\rm im},w}]_+\big).
 \end{split}
\]
Then
\[
 2\sum_{p\in\mathcal S}U_p\log p
 \ \le\ \frac{\sum_\sigma\log\det G_\sigma}{q[J:\mathbb Q]}
      +\frac{2L_{\mathcal S}+B_{\rm out}}q+2\log D+\log m_G .
\]
\end{corollary}
\begin{proof}
For a nonzero fractional ideal \(\mathfrak a\), its normalized
log norm is \(\sum_p\log p\sum_{w\mid p}\omega_w v_w(\mathfrak a)\).
This follows from the residue degrees in the ideal norm and
\(v_w(p)=1\). Sum both inequalities of
Theorem~\ref{thm:signedsignal450} with these weights on
\(\mathcal S\); their sum on the other primes is bounded below
by \(-B_{\rm out}\). Theorem~\ref{thm:globalperiodcontent448}
bounds the sum of the two normalized ideal log norms above by
the mean of the log Gram determinants. Divide by \(q\) and
use \(\sum_{w\mid p}\omega_w=1\).
\end{proof}

This compares the signal-prime contribution and the complex
Gram determinant for one set of actual definite columns. The
rank assumption, the costs \(L_{\mathcal S},B_{\rm out}\),
and the size of that Gram determinant enter any application to
\eqref{eq:gap}.


\subsection{All active signal primes for the prescribed finite tests}
\label{sec:signaltestsupport452}

Work with one \(K_d\)-oriented field and translate of
Theorem~\ref{thm:centerwords451}. Let \(\mathcal P\) consist of the
unramified \(p\ge5\) split in \(K_q\), inert in \(K_d\), for which
\(U_p>0\). This is a finite set since \(p\le 2nxX'\).
Use the marked elliptic differential and the elementary finite
tests of the definite column in Theorem~\ref{thm:covectorpaths352}.
These tests include the row and flag conditions, the terminal
covector condition, and the common vector coset on the last two
edges. Stabilizer masses and the small test projector have the
normalizations of Theorem~\ref{thm:signedsignal450}.

\begin{lemma}[Support and coefficients of the finite tests]\leavevmode
\label{lem:signalsupport452}
At each \(p\in\mathcal P\) the elementary finite tests have the
standard local lattice vector. On every edge there is a single
choice of product basis and class representatives for which their
expansions into elliptic products have prime-to-\(p\) shifts and
coefficients integral above \(p\), simultaneously for
\(p\in\mathcal P\). This holds for each allowed auxiliary
character, each matrix coset, and each value of the common vector
coset. The coefficients here precede insertion of the stabilizer
masses, the small test projector, and any additional scalar
linear combination of elementary columns.
\end{lemma}
\begin{proof}
Write \(F=E^+\). At a place \(v\) of \(F\) an allowed auxiliary
character is trivial on all local units unless \(v\) splits in
\(E/F\), its norm is a rational prime \(\ell\), and
\(\ell\notin v^2\). These are the defining conductor restrictions
in \srcloc{Formalism/LocalToGlobal/Divisibility/Tuples}{59--66}
and \srcloc{Formalism/LocalToGlobal/Divisibility/Tuples}{110--118}.
The norm condition gives residue degree one and the valuation
condition gives ramification index one. At either place of \(E\)
over a split \(v\) the completion is therefore \(\mathbb Q_\ell\).
Since \(E/\mathbb Q\) is Galois, this forces \(\ell\) to split
completely in \(E\). The two flag primes \(g_1,g_2\) and the
word's terminal prime \(T\) also split completely by construction.
Let \(S_{\rm row}\) be the selected places at the conductor of
the auxiliary character and at \(g_1,g_2\). Every prime below
\(S_{\rm row}\) or \(T\) is split in \(K_d\).

To check that this list includes every modification of the finite
vector, consider the matrix test in its position coordinates.
At the row places it is a function on lattice cosets. The two
integer bounds on denominators and moduli are
\[
       c_A=\Big(\prod_{P\in S_{\rm row}} N P\Big)^A,\qquad
       c_B=T\Big(\prod_{P\in S_{\rm row}} N P\Big)^B .
\]
Away from their support the position lattice and its congruence
sublattice agree. At \(T\) the extra test is the indicator of the
linear congruence on the terminal covectors. On the last two
edges the shared variable \(Y\) is a vector coset with conditions
only at \(S_{\rm row}\). In its Fourier transform the normal
covector is multiplied by \(\pi_T^{-1}\), where
\(\pi_T\overline{\pi_T}=T\). Hence it too is integral at every
\(p\in\mathcal P\). These are the tests written in
\srcloc{Formalism/LocalToGlobal/Columns/FiniteSumNonzeroLargeDepth}{1110--1139}
and \srcloc{Formalism/LocalToGlobal/Columns/FiniteSumNonzeroLargeDepth}{1150--1170}.
Keeping \(Y\) fixed throughout this calculation preserves the
coupling. At \(p\in\mathcal P\) the matrix lattice is consequently
the original unramified Hom lattice, with the reference vacuum.

We check the scalar weights as well. At a modified place of
residue characteristic \(\ell\), the flag and block tests take the
value zero or a product of residue-field characters, their inverses,
and an additive character of a finite quotient. Every nonzero
such character value is a root of unity. The free entries
in the unipotent integral are summed over lattice cosets, each
with measure a power of \(\ell\). The Fourier transform of \(Y\)
has the same form, with a power of \(\ell\) from the coset volume.
The determinant changes of variables contribute powers of
\((NP)^{1/2}\) and values of the splitting character at
ideles supported over \(\ell\). For the latter, raise the idele
to a ray class power to write it as a principal idele times a
unit idele on which the finite character is trivial. This
can be done since the quotient of the ideal class group with
the finite unit conditions is finite. The principal generator
is a unit at every prime over \(p\). The algebraic infinity type,
after squaring to clear the half weights, is a monomial in its
conjugates with integral exponents. Every factor is a \(p\)-unit
because \(E/\mathbb Q\) permutes the primes above \(p\). This
proves the same unit assertion for the original character value.

There is also a ratio of two compact open volumes on the
unipotent radical in the row integral. In a basis adapted to its
flag this group acts by upper unitriangular matrices. Conjugation
by a diagonal matrix of sufficiently separated uniformizer
powers puts both compact subgroups in the group of integral
upper unitriangular matrices. Every congruence quotient of this
ambient compact group has order a power of \(\ell\), by counting
its strict upper entries. A sufficiently deep congruence kernel
intersects each subgroup inside their common intersection, by
openness in the radical. The two indices over that intersection
are therefore powers of \(\ell\). This proves the assertion also
for the volume ratio and its inverse. The change of variables,
finite additive sum, reciprocal radical volume, and terminal
Fourier transform occur respectively in
\srcloc{Formalism/LocalToGlobal/Columns/FiniteSumNonzeroLargeDepth}{552--586},
\srcloc{Formalism/LocalToGlobal/Columns/FiniteSumNonzeroLargeDepth}{941--950}
and \srcloc{Formalism/LocalToGlobal/Columns/FiniteSumNonzeroLargeDepth}{1150--1155}.
All these factors are integral above \(p\).

It remains to pass from cosets to the vacuum expansion. At a
split place let \(L\) be the position lattice of the vacuum and
refine a coset to \(z+L'\) with \(L'\subset L\) an open sublattice.
On \(z+L\) its indicator is
\[
      [L:L']^{-1}\sum_{\eta\in(L/L')^\vee}\eta(y-z).
\]
This follows from the character sum on the finite additive
group. The index is a power of \(\ell\). Hence every term is a
position and frequency translate of the vacuum with
\(p\)-integral coefficient. Its shifts are supported over
\(\ell\). Clearing a common denominator and applying the Chinese
remainder theorem identifies these finite local shift classes
with a rational class in the period space modulo its lattice.
Its order is prime to every member of \(\mathcal P\).

Apply Theorem~\ref{thm:class61} at each vertex with the whole
finite set \(\mathcal P\) excluded. The representatives move
only at other completely split primes. Their lattice
intertwiners add shifts there, with root-of-unity factors and
inverse square roots of indices supported there. Their character
factors are units by the ray class argument above. At \(p\) the
class lattices at a shared vertex still agree on both incident
edges, as in Theorem~\ref{thm:neutral315}. Finally \(p\) is
unramified and inert in the elliptic CM field \(K_d\).
Lemma~\ref{lem:basis63} gives one orthonormal product basis per
edge generating the trace lattice at all \(p\in\mathcal P\).
The intertwiner to this product lattice again has index prime to
every \(p\) in the set. In the discriminant-one elliptic model
of Section~\ref{sec:period63} prime-to-\(p\) translations
have integral germs relative to the marked vacuum, including
when the translated section vanishes at the mark, by
Corollary~\ref{cor:class66}. This proves the assertions.
\end{proof}

\begin{theorem}[The complete signal-prime contribution to the period ideals]\leavevmode
\label{thm:fullsignal452}
Let column \(Y_i\) use an integral initial Weyl table and the
prescribed elementary tests. One may use a scalar linear
combination of such columns on the same directed tower, with
the same weights, for \(Y_i\). Let \(C_i\) be a positive integer
clearing these additional scalar coefficients. For an elementary
column or an integral combination take \(C_i=1\). With
\(\omega_w=[J_w:\mathbb Q_p]/[J:\mathbb Q]\), one has for
every \(p\in\mathcal P\)
\[
 \begin{aligned}
 \sum_{w\mid p}\omega_w b_{0,w}
      &\ge q U_p-\sum_{i=1}^{q}v_p(C_i)-qv_p(D),\\
 \sum_{w\mid p}\omega_w b_{{\rm im},w}
      &\ge q U_p-\sum_{i=1}^{q}v_p(C_i)-qv_p(Dm_G).
 \end{aligned}
\]
The bounds use the same period ideals in any marked product
bases, with the differential and physical coordinates fixed.
\end{theorem}
\begin{proof}
In the bases and representatives of Lemma~\ref{lem:signalsupport452},
clearing \(C_i\) makes every premass path coefficient integral
at \(w\mid p\), including the character factors. At a common
vertex use the integral frames on its Weyl lattice; a unit
change of determinant frame preserves it. Theorem~\ref{thm:signedsignal450}
thus bounds the valuations of the minors in these bases at each
individual \(w\) by the right sides. This argument applies
termwise to the integral combination \(C_iY_i\) and then divides
by \(C_i\) once per column in the minor.

Changing to any other fixed marked product basis multiplies all
entries of column \(i\) by the same scalar \(u_i\), by
Theorem~\ref{thm:framechange314}. Every term of a combination
on that tower contains one copy of each directed edge, so it
has that same scalar. Every minor is multiplied by
\(\prod_i u_i\). Enlarge \(J\) to contain these scalars if
necessary. Since \(|\sigma u_i|=1\) for all embeddings,
\(|N_{J/\mathbb Q}u_i|=1\), and therefore
\[
 \sum_{w\mid p}\omega_w v_w(u_i)
       =[J:\mathbb Q]^{-1}v_p(N_{J/\mathbb Q}u_i)=0
\]
at each rational \(p\). Both weighted ideal valuations are
therefore preserved. The local degree identity in a tower
preserves the averages after extending coefficients.
\end{proof}

\begin{corollary}[The centered bound for the prescribed column tests]\leavevmode
\label{cor:fullcenter452}
If these \(q\) columns are independent on the projected small
tests, the centered sum on the same field and interval satisfies
\[
 \begin{split}
 \Delta_u\le{}&
   \frac{\sum_\sigma\log\det G_\sigma}{q[J:\mathbb Q]}
        -A_{\rm sig}(x)+\frac{B_{\rm data}}q\\
   &+\frac{2}{q}\sum_{i=1}^{q}\log C_i
        +2\log D+\log m_G+o_{\psi,\varepsilon,A}(b/x^A)
 \end{split}
\]
for every fixed \(A>0\). In \(B_{\rm data}\) use
\(\mathcal S=\mathcal P\) and the fixed normalization used to
evaluate the Gram determinants.
\end{corollary}
\begin{proof}
Theorem~\ref{thm:fullsignal452} supplies both lower bounds at
every active prime in Theorem~\ref{thm:centerwords451}.
At the other primes use Proposition~\ref{prop:outsidewordcost451}
in the original bases. Sum and apply the product formula
inequality of Theorem~\ref{thm:globalperiodcontent448} in
those bases. The restricted sum of \(v_p(C_i)\log p\) is at
most \(\log C_i\), and the identical assertion applies to
\(D,m_G\). Substitution of the centered identity proves the
claim. Basis change preserves the determinant of the Gram matrix
at each embedding, since it multiplies the columns by scalars
of absolute value one.
\end{proof}


\subsection{The different and the split tests in the centered bound}
\label{sec:outsideprimeperiods453}

Use the \(q\) columns of Theorem~\ref{thm:fullsignal452}. Denote
the total derivative degree at embedding \(\tau\) on edge \(j\)
of column \(i\) by \(d_{ij,\tau}=|\nu_{ij,\tau}|\), and let
\(d_{ij}=\sum_\tau d_{ij,\tau}\). In a linear combination take
the common directed tower and weights of that theorem.
Let \(\mathcal M\) be the rational primes below the row places,
the two flag primes, and all terminal primes of the selected
columns. Write \(\mathcal T\subset\mathcal M\) for those terminal
primes. All of \(\mathcal M\) splits completely in \(E\) by
Lemma~\ref{lem:signalsupport452}. We use the prescribed integral
Weyl tables at the incoming endpoint.

\begin{lemma}[The loss of an undivided derivative on the Weyl lattices]\leavevmode
\label{lem:jetweyl453}
Work over a coefficient valuation ring, using integral coefficient
frames at the two ends of an edge. Suppose each undivided
derivative of its class section of the given block degrees has
valuation at least \(t\), for a real \(t\). Every coordinate of
the resulting edge operator between the source and target Weyl
lattices has valuation at least \(t\).
\end{lemma}
\begin{proof}
Take a generator of the source as a coefficient of the
highest polynomial on a product of elementary row matrices
with independent parameters. After the determinant twist,
its polynomial is the corresponding coefficient of
\(p_\nu(AX)\), since \(\det A=1\).
Pair with an integral dual polynomial \(P_\tau\) at each target
embedding. As in the proof of Theorem~\ref{thm:weylrowjet449},
the differential symbol in that block is
\(P_\tau([D A_\tau^{\mathsf t},0])\). Its coefficients in
the indeterminates and in the entries of \(D\) are integral.
The product over embeddings is homogeneous of degrees
\(d_{ij,\tau}\). After parameter extraction its evaluation
is an integral linear combination of precisely the derivatives
in the hypothesis. Such generators span the source, and the
tensor products of these tests span the integral dual of the
target. Testing on a dual basis proves the assertion. A
determinant line trivialized in an integral frame contributes
a unit under an integral change of frame.
\end{proof}

Here is finite data for a prime \(p\in\mathcal M\setminus\mathcal T\).
At each \(v\mid p\) in \(F\), let \(L_v\) be the standard
split position lattice of the given edge. Its
\(\mathbb Z_p\)-rank is \(D_{jv}=r_js_j\) for edge ranks
\(r_j,s_j\), since \(F_v=\mathbb Q_p\). For all elementary
functions of this edge in this column, divided by the amplitude
of the reference position vacuum, choose nonnegative integers
\(a_v,b_v\) so that their support lies in \(p^{-a_v}L_v\)
and they are constant on cosets of \(p^{b_v}L_v\). Use the
matrix-coset indicators appearing in the path sum, with the
value of the coupled coset held fixed. Such bounds exist by
the denominator and period lattices in Lemma~\ref{lem:signalsupport452}.
For a spherical factor take \(a_v=b_v=0\). Let \(c_{ij,p}\)
be the maximum of zero and the negatives of the valuations of
the nonzero products of these original local values at \(v\mid p\),
over all the cosets in the finite list and all coefficient
places above \(p\). The sum of the corresponding nonnegative
maxima place by place also gives an admissible upper bound. This includes the row
integral scalars, prior to inserting the class masses.

Label the two places of \(K=K_d\) above \(p\) by \(+,-\).
Let \(\epsilon_v\) be the label of the position factor at \(v\).
For \(\epsilon\in\{+,-\}\), put
\[
 h_\epsilon=
 \max\big(\{0\}\cup\{a_v:\epsilon_v=\epsilon\}
                 \cup\{b_v:\epsilon_v\ne\epsilon\}\big),
 \qquad
 J_{ij,p}^{\rm row}
  =c_{ij,p}+\sum_{v\mid p}D_{jv}b_v
                 +\frac{d_{ij}}2(h_++h_-).
\]
The lattices, bounds and exponents here refer to the specified
column and edge; they can be read off before any differentiation.

At every rational prime put
\(\rho_p=v_w(\mathfrak D_{E/K})\), with \(v_w(p)=1\),
independent of \(w\mid p\) by normality. Let \(\kappa_p\) be
the bound for the marked elliptic differential in
Theorem~\ref{thm:elliptic312}. For \(p\notin\mathcal M\) set
\[
      J_{ij,p}^{\rm vac}=
        \tfrac12 d_{ij}\rho_p+
                      \kappa_p\lfloor d_{ij}/2\rfloor .
\]

\begin{theorem}[Outside-prime valuations of the two period ideals]\leavevmode
\label{thm:outsidedifferent453}
For \(p\notin\mathcal T\cup\mathcal P\) use
\(J_{ij,p}=J_{ij,p}^{\rm vac}\) off \(\mathcal M\)
and \(J_{ij,p}=J_{ij,p}^{\rm row}\) on \(\mathcal M\).
In the original normalization, writing \(b_{S_0,w}=b_{0,w}\)
and \(b_{eS,w}=b_{{\rm im},w}\), take \(m'=1\) in the saturated
case and \(m'=m_G\) in the image case. One has
\[
 \sum_{w\mid p}\omega_w b_{D',w}
 \ \ge\ -\sum_{i,j}J_{ij,p}-\sum_i v_p(C_i)-qv_p(D)
      -qv_p(m')
\]
for the saturated and image test lattices, respectively. In
particular an unmarked, unramified prime with \(\kappa_p=0\)
costs only the explicitly displayed scalar, mass, and projector
valuations.
\end{theorem}
\begin{proof}
Fix \(p\) and choose class representatives equal to the identity
over \(p\), by Theorem~\ref{thm:neutral315}. All class movements
occur at other split primes and contribute \(p\)-unit multipliers
and shifts with zero \(p\)-primary component. Also the finite
tests and changes of variables at the other marked primes have
integral coefficients above \(p\), by the local character,
volume and finite Fourier calculations in
Lemma~\ref{lem:signalsupport452}. Choose on each edge the product
basis of Theorem~\ref{thm:adapted314} generating the local
Hermitian trace lattice. Its vacuum intertwiner adds shifts of
order prime to \(p\) and has unit normalization.

Suppose first \(p\notin\mathcal M\). Every translated elliptic
factor now has order prime to \(p\) in its translation parameter.
Theorem~\ref{thm:torsion312} bounds the loss at derivative order
\(s\) by \(\kappa_p\lfloor s/2\rfloor\). In an undivided
derivative of degree \(d_{ij}\) the orders sum to \(d_{ij}\),
so these losses sum to at most
\(\kappa_p\lfloor d_{ij}/2\rfloor\). In physical block \(\tau\)
the change of coordinates has entries of valuation at least
\(-v_w(\bar\tau\mathfrak b)\), since product basis vectors
belong locally to \(\mathfrak b^{-1}\operatorname{Mat}(\mathcal O_E)\)
and the inverse embedding matrix is the conjugate transpose.
Here \(\mathfrak b\overline{\mathfrak b}=\mathfrak D_{E/K}\).
There are exactly \(d_{ij,\tau}\) differentiations in this block.
All multinomial coefficients in the chain rule are integral.
Thus a bound for the loss above \(p\) is
\[
 H_{ij,w}=\kappa_p\lfloor d_{ij}/2\rfloor+
                \sum_\tau d_{ij,\tau}v_w(\bar\tau\mathfrak b).
\]
Use \(\tau\) in place of \(\bar\tau\) on a conjugated edge.
In a common normal coefficient field, conjugation permutes
the places above \(p\) and preserves their local degrees. For
each individual \(\tau\) the two conjugate ideals consequently
have the same weighted average valuation. Their product has
valuation \(\rho_p\) at every place above \(p\).
It follows that \(\sum_{w\mid p}\omega_w H_{ij,w}=J_{ij,p}^{\rm vac}\).

Suppose next \(p\in\mathcal M\setminus\mathcal T\). For one
position coset \(z+p^{b_v}L_v\), Fourier expansion on \(L_v\)
costs exactly the reciprocal index \(p^{-D_{jv}b_v}\) before
cancellation. Its position translate has order dividing
\(p^{a_v}\), and its frequency character has order dividing
\(p^{b_v}\) modulo the dual lattice. This is the coset character
identity in Lemma~\ref{lem:signalsupport452}. On taking products,
the \(p\)-parts of the two idempotent components of any resulting
shift have orders at most \(p^{h_+}\), \(p^{h_-}\), respectively.
This follows by taking the maximum of the exponents on the
factors of the direct product. The local basis change is
\(\mathcal O_K\otimes\mathbb Z_p\)-linear, so these exponents
also work for each elliptic coordinate. The conjugate
annihilator ideal in Theorem~\ref{thm:torsion312} exchanges
the two components. Since \(\kappa_p=0\), at a place over
\(\epsilon\) the loss of all \(d_{ij}\) differentiations is
at most \(d_{ij}h_{-\epsilon}\). The physical coordinate
conversion costs zero here. Indeed \(\mathfrak b\) in
Theorem~\ref{thm:pol49} is integral, and its product with its
conjugate is a unit at every prime over this unramified \(p\).
The two places of \(K\) each have weight \(1/2\) on averaging
over \(w\mid p\), by the local degree formula in a tower.
Including the original scalar and the Fourier indices gives
the bound \(J_{ij,p}^{\rm row}\) on the weighted average.

Apply Lemma~\ref{lem:jetweyl453} to every edge at each \(w\),
before averaging. Since \(p\notin\mathcal T\), the terminal
refinements localize to the unrefined lattice at each vertex.
Thus the identity representatives give matching integral
frames at both ends of every composition. The terminal
restriction is integral by Theorem~\ref{thm:weyllower351}.
Each term of a column contains exactly one derivative operator
per edge, hence its coordinate loss is at most the sum of the
individual bounds plus \(v_p(C_i)\). The finite signed sums
preserve this lower valuation bound. Clearing \(D\) clears
the path masses and \(m_G eS\subset S_0\) clears the projector.
Expand the minors, add over the columns, and take the minimum
over the tests. Changing from these product bases to the original
ones preserves each weighted ideal valuation by the norm
calculation in Theorem~\ref{thm:fullsignal452}. This proves
both inequalities.
\end{proof}

\begin{corollary}[A finite cost supported on the terminal and row primes]\leavevmode
\label{cor:outsidecenter453}
For independent projected columns let \(B_\mathcal T\) be the
sum in Proposition~\ref{prop:outsidewordcost451} restricted
to \(p\in\mathcal T\), in the original normalization. On the
same intervals let \(d_0=|\operatorname{disc}K|\). Then
\[
\begin{split}
 \Delta_u\le{}&
  \frac{\sum_\sigma\log\det G_\sigma}{q[J:\mathbb Q]}-A_{\rm sig}(x)\\
 &+\frac1q\sum_{i,j}\!
  \left(d_{ij}\log\operatorname{rd}(E)
       -\tfrac12{\bf1}_{d_{ij}\ {\rm odd}}\log d_0\right)\\
 &+\frac2q\sum_{p\in\mathcal M\setminus\mathcal T}
                \sum_{i,j}J_{ij,p}^{\rm row}\log p
     +\frac{B_\mathcal T}{q}\\
 &+\frac2q\sum_i\log C_i+2\log D+\log m_G+
                            o_{\psi,\varepsilon,A}(b/x^A).
\end{split}
\]
\end{corollary}
\begin{proof}
The different tower formula gives
\(\sum_p\rho_p\log p=\log\operatorname{rd}(E)-\tfrac12\log d_0\).
Here \(3\) splits in \(K\) on the selected intervals.
Consequently the only nonzero \(\kappa_p\) is \(\kappa_{d_0}=1/2\).
Every prime of \(\mathcal M\cup\mathcal P\) is unramified
in \(E\) and has \(\kappa_p=0\). Twice the sum of the vacuum
costs of edge \(ij\) is therefore
\[
 d_{ij}\big(\log\operatorname{rd}(E)-\tfrac12\log d_0\big)
                      +\lfloor d_{ij}/2\rfloor\log d_0 .
\]
This is the displayed discriminant term. At the terminal primes
use the finite symbolic bound, which includes the actual input
and terminal frames there. At active signal primes use
Theorem~\ref{thm:fullsignal452}; at all others use
Theorem~\ref{thm:outsidedifferent453}. Sum the two normalized
ideal log norms before taking the archimedean bound of
Theorem~\ref{thm:globalperiodcontent448}. For each
\(C_i,D,m_G\) the sum off the terminal primes is at most its
full logarithm. Inserting the centered identity of
Theorem~\ref{thm:centerwords451} proves the inequality.
\end{proof}

The loss at unmarked primes in this estimate depends on the
field and the degrees alone. The quantitative comparison still
involves the size and rank of the signed projected periods and
the finite tests at their specified levels.


\subsection{The terminal covector congruence in the centered bound}
\label{sec:terminalramps454}

Fix a translate and the common terminal generator \(\pi\) of
Corollary~\ref{cor:count111}. Thus \(T=\pi\bar\pi\) splits
completely in \(E\), \(S_T=\{P:P\mid\pi\}\) contains one place
over each place \(v\mid T\) of \(F\), and \(f=|S_T|=[F:\mathbb Q]\).
The auxiliary characters are unramified there and \(T\) is
distinct from the flag primes. We use the elementary column
tests of Theorem~\ref{thm:fullsignal452}. Let \(r_j,s_j\) be the
ranks at the two ends of edge \(j\). At a fixed pair of unrefined
classes write \(L_{j,P}\) for the lattice of integral position
matrices at \(P\); its reduction consists of all \(s_j\)-by-\(r_j\)
matrices over \(\mathbb F_T\).

\begin{lemma}[Fourier expansion with the intermediate covectors fixed]\leavevmode
\label{lem:terminalfourier454}
In the signed path sum one can sum over nonzero terminal
covectors at the vertices and use on each edge the indicator
of the single affine linear condition
\[
                       \xi' \bar X=\xi
\]
on \(L_{j,P}\), for each \(P\in S_T\). Relative to the unit
vacuum, this indicator is a sum of pure frequency translates
with common denominator \(T^{r_j}\) at \(P\). Their \(T\)-primary
shift classes are killed by \(T\) and lie in the component
at the conjugate place \(\bar P\).
\end{lemma}
\begin{proof}
Keep the original classes, their masses, all nonterminal local
data and the common vector coset on the last two edges.
For an integral matrix path and the specified final covector,
successive multiplication backwards gives a unique tuple of
intermediate covectors. If any intermediate covector is zero,
the initial one is zero and the path is discarded. Conversely
every surviving path occurs once on summing over the nonzero
tuples and imposing the displayed equations. This is the path
identity of Theorem~\ref{thm:covectorpaths352}. With the tuples
fixed the conditions on the separate edges are independent.
The matrix-coset terms can be summed into their union because
their indicator functions add, the theta section is linear in
the finite test, and the iterated map is linear in the section
on each edge separately. Chinese remaindering makes this
regrouping independent of the conditions at the row places.
The finite test at \(T\) in
\srcloc{Formalism/LocalToGlobal/Columns/FiniteSumNonzeroLargeDepth}{646}
is exactly this covector indicator.

Fix a nontrivial additive character \(\psi_T\) of \(\mathbb F_T\).
On \(L_{j,P}\) the indicator is
\[
  T^{-r_j}\sum_{u\in\mathbb F_T^{r_j}}
                 \psi_T\big(\xi'\bar X u-\xi u\big).
\]
The sum is one after division if \(\xi'\bar X=\xi\), and zero
otherwise by character orthogonality. Multiplication by
\({\bf1}_{L_{j,P}}\) supplies the support condition. The
coefficient phase is a root of unity. The linear functional
\(X\mapsto\xi'\bar X u\) is a character of \(L_{j,P}/T L_{j,P}\),
so its frequency is a class in \(T^{-1}L_{j,P}^\vee/L_{j,P}^\vee\).
Under the split Hermitian pairing the dual of the position
component at \(P\) is the component at \(\bar P\). This proves
all the assertions, including the support of each individual
shift.
\end{proof}

For column \(i\), let \(\lambda^{\rm in}_{i,\tau,k}\) be the
smallest entry of its initial highest weight at embedding
\(\tau\), and let \(\delta_{i,\tau}\) be its terminal normal
degree in Corollary~\ref{cor:normaldegree347}. Define
\[
 z_i=\prod_\tau\tau(\bar\pi)^{\lambda^{\rm in}_{i,\tau,k}},
 \qquad Z_i^{\rm norm}=\prod_\tau\tau(\bar\pi)^{-\delta_{i,\tau}},
 \qquad
       \widetilde Y_i=(z_i Z_i^{\rm norm})^{-1}Y_i .
\]
The two scalars are units away from \(T\). Use tildes on the
valuation of an ideal or on a Gram matrix to indicate that
these rescaled columns are used, with the original small
test lattices.

\begin{theorem}[The terminal valuation on both projected period ideals]\leavevmode
\label{thm:terminaldegrees454}
For \(D'=S_0,eS\), take \(m'=1,m_G\), respectively. One has
\[
 \sum_{w\mid T}\omega_w\widetilde b_{D',w}\ \ge\
 -\sum_{i,j}\left(fr_j+\tfrac12d_{ij}\right)
 -\sum_i v_T(C_i)-qv_T(Dm').
\]
\end{theorem}
\begin{proof}
Choose unrefined class representatives equal to the identity
above \(T\), by Theorem~\ref{thm:neutral315}. The small class
and its lift then have the standard unrefined frames at \(T\).
For each initial residue covector \(\xi\) its refined Weyl
lattice satisfies \(D_{a,\xi}\subset z_iD_a\), where \(D_a\)
is the unrefined lattice. Indeed the prescribed kernel at \(P\)
changes the column lattice at \(\bar P\) by replacing one basis
vector by \(T\) times that vector, by the two dualizations
in the proof of Theorem~\ref{thm:pointlattice353}.
In the polynomial representation with its last weight
subtracted, all degrees of that vector are nonnegative.
Reinstating the determinant power gives the factor
\(T^{\lambda^{\rm in}_{i,\tau,k}}\) at the affected embedding.
Their product differs from \(z_i\) by a local unit.
Thus \(z_i^{-1}Y_i\) may be computed with initial values integral
in the unrefined frames. At the last class, in those frames
the actual restriction is \(Z_i^{\rm norm}\) times the integral
standard restriction, by Theorem~\ref{thm:terminalformula347}
and Theorem~\ref{thm:weyllower351}. Removing this second scalar
allows us to use unrefined large frames all the way through
the path sum. There is no choice of a test or a new projection
at an intermediate vertex.

We make a single choice of product coordinates at \(T\) for
each directed edge. The split trace lattice is the orthogonal
sum over \(v\mid T\) of the position lattice and its dual,
paired with each other. For a position basis and its exact
dual, the sums of corresponding vectors give a local
orthonormal \(K\otimes\mathbb Q_T\)-basis, as in
Theorem~\ref{thm:adapted314}. Label its coordinates by \(v\)
and a matrix-entry index. Approximate this local basis by a
rational orthonormal basis agreeing modulo \(T\) at both
places of \(K\), using the Cayley approximation in that theorem.
The congruence is an open condition. Both the change of basis
and its inverse then preserve the individual \(v\)-blocks
modulo \(T\). In particular an order-\(T\) shift has in each
new local coordinate the same component modulo the period
lattice as under the exact local basis, after identification
by these bases. The vacuum intertwiner for this rational basis
has index prime to \(T\). Its shifts and the shifts from the
other finite primes add no \(T\)-primary components.

Fix \(w\mid T\) and an embedding \(\tau\) in the CM type.
Let \(P_{\tau,w}\) be the prime of \(E\) obtained by pulling
\(w\) back under \(\tau\), and put \(v=P_{\tau,w}\cap F\).
For a coordinate in the \(v\)-block, the conjugate annihilator
bound of Theorem~\ref{thm:torsion312} uses the part of the
shift at the place of \(K\) opposite the one induced by \(\tau,w\).
If \(P_{\tau,w}\notin S_T\), this opposite component of the
shift at \(v\) is zero modulo the lattice. In fact the position
prime there is \(\overline{P_{\tau,w}}\), and the lemma puts
the entire frequency shift at \(P_{\tau,w}\). Hence the
undivided derivatives on that coordinate are integral in
every order. If \(P_{\tau,w}\in S_T\), the same argument permits
a loss of one per differentiation, since \(T\) kills the
primary shift. Here \(\kappa_T=0\).

The physical derivative in a \(\tau\)-matrix entry is a sum of
product derivatives with integral coefficients. For coordinates
belonging to \(v'\ne v\) these coefficients lie in \(T\mathcal O_w\).
Indeed the coefficient is the conjugate embedded basis entry,
which vanishes in those blocks for the exact local basis and
agrees with it modulo \(T\) for the chosen basis. Conjugation
exchanges \(P_{\tau,w}\) and its partner over the same \(v\).
Each such off-block factor therefore pays the loss of at most
one from a differentiation there. To see this in all orders,
expand the product of the linear differential operators. In a
product coordinate whose primary shift can cost one per
derivative, its derivative order counts exactly the individual
differentiations landing there. Each arriving from an uncharged
block has a factor of \(T\). Terms from charged blocks have
integral coefficients, and integral multinomial coefficients
preserve the resulting lower bound.

Consequently the loss in a derivative of the edge's block
degrees, apart from the Fourier indices, is at most
\[
       \sum_\tau d_{ij,\tau}\,
                       {\bf1}_{P_{\tau,w}\in S_T}.
\]
For a conjugated edge evaluate this formula at the conjugate
place. There are \(f\) Fourier indices from the lemma, costing
\(fr_j\) in either orientation. Other finite inputs have
coefficients integral at \(T\) by the split-prime character
and volume calculation in Lemma~\ref{lem:signalsupport452}.
At the fixed \(w\), use Lemma~\ref{lem:jetweyl453} on the
unrefined Weyl lattices, then compose the maps at the common
vertices. Taking determinants adds the costs for the columns.
The path masses, projector, and coefficient denominators
contribute precisely the allowances displayed in the statement.

For each fixed \(\tau\), the total local degree weight of
\(P_{\tau,w}\in S_T\) is \(f/[E:\mathbb Q]=1/2\), because
\(T\) splits completely and local degrees add in the tower
\(J/\tau(E)/\mathbb Q\). The same holds at the conjugate places.
Average the preceding bound. A return to the original product
bases multiplies the minors by a common product of column
phases of rational norm of absolute value one, as proved in
Theorem~\ref{thm:fullsignal452}. This preserves the average
and proves the theorem.
\end{proof}

\begin{corollary}[The centered inequality with explicit finite terms]\leavevmode
\label{cor:allcost454}
At the common prime \(T\) of the original family, the \(q\)
independent projected columns satisfy, for every fixed \(A>0\),
\[
\begin{split}
 \Delta_u\le{}&
 \frac{\sum_\sigma\log\det\widetilde G_\sigma}{q[J:\mathbb Q]}
          -A_{\rm sig}(x)\\
 &+\frac1q\sum_{i,j}
  \left(d_{ij}\log\operatorname{rd}(E)
          -\tfrac12{\bf1}_{d_{ij}\ {\rm odd}}\log d_0\right)\\
 &+\frac2q\sum_{p\in\mathcal M\setminus\{T\}}
             \sum_{i,j}J^{\rm row}_{ij,p}\log p
   +\frac1q\sum_{i,j}(2fr_j+d_{ij})\log T\\
 &+\frac2q\sum_i\log C_i+2\log D+\log m_G
                         +o_{\psi,\varepsilon,A}(b/x^A).
\end{split}
\]
In the first term one may equivalently use the original
\(\log\det G_\sigma\) corrected by
\(-\sum_{i,\tau}(\lambda^{\rm in}_{i,\tau,k}-\delta_{i,\tau})\log T\)
at every embedding.
\end{corollary}
\begin{proof}
Use Theorem~\ref{thm:terminaldegrees454} at \(T\) in the proof
of Corollary~\ref{cor:outsidecenter453} and use the rescaled
columns throughout. The two scalars removed from each column
are units at every other prime, so the previous estimates for
both ideals hold there with the same valuations. The \(T\)-terms
from the two test lattices sum to the displayed terminal cost.
The original and rescaled Gram determinants differ by
\(\prod_i|\sigma(z_iZ_i^{\rm norm})|^{-2}\), since the
class metric and projector are fixed. Now
\(|\sigma\tau(\bar\pi)|=\sqrt T\) for every \(\sigma,\tau\).
This gives the asserted correction.
\end{proof}

All terms of this inequality are evaluated on the field and
translate of the prime count. Its use in \eqref{eq:gap} requires
a uniform bound for their sum and a supply of nonzero projected
columns on those intervals.


\subsection{The contribution of primes inert in both quadratic fields}
\label{sec:inertperiodmass462}

Work on the polynomial-conductor fields and intervals of
Theorem~\ref{thm:polycenter460}. Put \(Y=X'=\exp(x^{\varepsilon/3})\).
Write \(T_{\rm hi}\) for the high set of size \(h\).
The extra divisor \(\mathcal D_u\) of
Corollary~\ref{cor:inertcorrection461} measures divisibility of
the signed word periods beyond the primes used in \(U_p\).
We estimate a part of it for which the dyadic Frobenius moves
every high place to a place of small weight.

Let \(G_j=\operatorname{Gal}(C_j/\QQ)\) with \(j=j_x\).
For a power of two \(M\mid 2^j\) use the cyclic quotient
\(G_j/MG_j\) of order \(M\). For a prime unramified in \(E_x\)
let \(\bar p\) denote its Frobenius in this quotient. It is also
the image of \(\pi_p=\sigma_pc_E\), since \(C_j\) is totally
real. For a function on a finite abelian group normalize its
Fourier coefficients so that it equals their sum times the
corresponding characters, and write \(\widehat w(1)\) for the
average of \(w\).

\begin{lemma}[A weighted prime sum on the dyadic quotient]\leavevmode
\label{lem:farprimedensity462}
Let \(w\) be real valued on \(G_j/MG_j\) with
\(\sum_\chi|\widehat w(\chi)|\le K\). Uniformly in \(M,w,x\),
\[
 \sum_{\substack{p\le Y,\ p\ge5,\ p\ {\rm unramified\ in}\ E_x\\
                       p\ {\rm inert\ in}\ K_q,\ K_d}}
                 w(\bar p)\frac{\log p}{p}
   =\frac{\widehat w(1)}4\log Y
                  +O_{\psi,\varepsilon}\bigl(K\log^2(2x)\bigr).
\]
\end{lemma}
\begin{proof}
On unramified primes multiply \(w\) by
\((1-\xi_q)(1-\xi_d)/4\). The two quadratic fields and the
dyadic quotient have pairwise coprime conductors, so their
compositum has the direct product Galois group. This follows
also by independent variation of the units at the three
moduli. The product has constant Fourier coefficient
\(\widehat w(1)/4\), and its sum of absolute Fourier
coefficients is at most \(K\). Every character in the
expansion belongs to \(E_x\).

For the primitive representative of each character, use
\(R_\chi(Y)/Y\) from \eqref{eq:Gchi}. Replacing the sum over
unramified primes by these expressions costs \(O_\psi(K)\).
Indeed removing prime powers of exponent at least two costs
at most \(\sum_p(\log p)/(p(p-1))<\infty\) per character.
The linear smoothing term costs at most
\(Y^{-1}\sum_{p\le Y}\log p=O(1)\) by Chebyshev's bound.
At \(2,3\) and the primes dividing \(N_x\), the remaining
difference costs at most a constant times
\(\sum_{p\mid 6Fqd_0}\log(p)/(p-1)=O_\psi(1)\).
This accounts also for primitive representatives with smaller
conductors.

The large-argument estimate in the proof of
Theorem~\ref{thm:polycenter460} gives \(G_\chi(Y)/Y=o(1)\)
uniformly. For the principal character the term remaining in
\(R_\chi(Y)/Y\) is \(\log Y-1-\gamma\). For a nonprincipal
character it is \(-L'(1,\chi)/L(1,\chi)\), bounded in absolute
value by \(O_\psi(\log^2(2N_x))\) by the zero-free estimate
and the logarithmic derivative calculation of
Lemma~\ref{lem:far14}. Since \(N_x\le x^{D_0}\), Fourier
summation gives the assertion.
\end{proof}

\begin{theorem}[A lower bound for the extra divisor of the projected periods]\leavevmode
\label{thm:inertmass462}
For any of the high coset choices in Theorem~\ref{thm:rangenew50},
the additional term in Corollary~\ref{cor:inertcorrection461}
satisfies, uniformly over translates and the admissible lists
of independent signed columns,
\[
          \mathcal D_u\ \ge\
          \bigl(\tfrac14-o_{\psi,\varepsilon}(1)\bigr)\,
                              hma\log X'.
\]
\end{theorem}
\begin{proof}
Let \(D\) be the dyadic index used to choose the high set,
so \(2x^{\varepsilon/8}\le D<4x^{\varepsilon/8}\).
In \(G_j/DG_j\) the image of \(T_{\rm hi}\) contains at most
\(3o+1\) points, \(o=[K_\psi:\QQ]\), since \(T_{\rm hi}\)
is a union of that many cosets of \(DA_0\). Write \(S\)
for their difference set. It contains zero and has size at
most \((3o+1)^2\). Take \(w=1-\mathbf1_S\). Each point
indicator has Fourier coefficient sum in absolute value one,
so the lemma applies with \(K\le1+|S|\) and
\(\widehat w(1)=1-|S|/D\).

For a prime counted with this weight, \(\pi_p\) has nonzero
image in the dyadic quotient and hence even order. Its cycles
on the embedding labels all have that order, since the action
is translation in the abelian group. The two negative signs
ensure that it preserves the \(K_q\)-sign on the \(K_d\)-type.
Thus \(\pi_p T_{\rm hi}\subset \Phi_0\setminus T_{\rm hi}\).
Choose in each cycle the half containing at least half of its
high places. At least \(h/2\) high places are then in the
chosen halves.

For each such \(t\), each column and each edge the high
partition has \(\nu_{t,1}\ge a\) by the word calculation in
Lemma~\ref{lem:wordfirstrow450}. At \(\mu=\pi_p t\) the
partition has \(\nu_{\mu,1}\le b/(nY)+O(x^B)\) for a fixed
\(B\): this place is neither high nor distinguished, so its
weight is \(\lceil b/(nY)\rceil+O(mH_0+1)\) and its starting
first coordinate has polynomial size. The rank \(s\) is
polynomially bounded as well. Keep only the first high row
and the first Frobenius power in the row contribution of
Theorem~\ref{thm:cycleweyl461}. Since
\(d_p(\nu_{t,1},s)\ge a/p-s\), the contribution of this place
is at least \(a/p-b/(nY)-O(x^B)\), increasing \(B\) if
necessary. This is positive throughout \(p\le Y\), since
\((b/(nY)+x^B)Y/a=o(1)\). The constants are independent
of the word, edge and translate.

Here \(U_p=0\). The factor two in \(\mathcal D_u\) comes from
the two integral test lattices. After division by \(hma\), the
lower bound from these primes is their sum of \((\log p)/p\)
up to an additive loss
\[
  O\!\left(\frac{(b/(nY)+x^B)Y}{a}\right)=o(1),
\]
by the Chebyshev bound again. The lemma gives the more precise
lower bound \((1-|S|/D)\log Y/4-O_{\psi,\varepsilon}(\log^2(2x))\)
for this normalized expression. Since \(D\) tends to infinity
and \(\log^2(2x)=o(\log Y)\), the theorem follows.
\end{proof}

\begin{corollary}[Aligning the high cosets in the dyadic quotient]\leavevmode
\label{cor:dyadicchoice462}
Let \(e_2=[K_\psi\cap C_{j_x}:\QQ]\). One may choose the high
cosets in the same fields so that the preceding inequality
improves to
\[
          \mathcal D_u\ \ge\
     \bigl(\tfrac14(1+e_2^{-1})-o_{\psi,\varepsilon}(1)\bigr)
                                  hma\log X'.
\]
In particular the coefficient is \(1/2-o(1)\) for the
conductor-one character.
\end{corollary}
\begin{proof}
Use a generator to identify \(G_j\) with \(\ZZ/2^j\).
Restriction identifies \(A_0=\operatorname{Gal}(K_\psi C_j/K_\psi)\)
with \(e_2 G_j\), and \(D\mid |A_0|\) on these intervals.
In each coset of \(A_0\), denote its common residue modulo
\(e_2\) by \(v\), represented in \([0,e_2)\). Inside the
quotient modulo \(e_2D\), choose the high residues
\[
          v+t\,e_2 D/4,\qquad t=0,1,2,
\]
and also \(t=3\) in the identity coset. For large \(x\)
the integer \(D\) is divisible by four. These are three or
four distinct \(DA_0\)-cosets in the required coset, by the
restriction identification. This is an allowed selection in
Theorem~\ref{thm:rangenew50}. The character cancellations use
only that \(T_{\rm hi}\) is a union of \(DA_0\)-cosets, with
the prescribed number in each \(A_0\)-coset, and that \(T_2\)
is balanced modulo \(D\). The sizes and word weights per type
are therefore those of Corollary~\ref{cor:count111} on this
field.

Now take \(M=e_2D/4\) and form \(S\) in \(G_j/MG_j\).
Its elements are differences \(v-v'\), so \(|S|\) is bounded
in terms of \(\psi\) and \(S\cap e_2(G_j/MG_j)=\{0\}\).
For \(\bar p\notin S\), the theorem's half-count argument still
works. In addition, if \(\bar p\) is a nonzero multiple of
\(e_2\), all high places of any one cycle lie in the same
alternating half. Indeed if \(t'\) and \(t\) are high and
\(t'=\pi_p^n t\), their projection difference is \(n\bar p\).
It is a multiple of \(e_2\) in \(S\), so it is zero modulo
\(M\). The nonzero element \(\bar p\) has even order in that
quotient, forcing \(n\) even. Choose the half containing all
these high places in each cycle.

The additional \(h/2\) high places, beyond the original lower
bound \(h/2\), can thus be counted with the second weight
\(\mathbf1_{e_2(G_j/MG_j)}-\mathbf1_{\{0\}}\). Its mean is
\(1/e_2-1/M\), and its Fourier coefficient sum in absolute
value is at most two: a subgroup indicator is the average of
the characters annihilating it. Summing the lemma for the
first and second weights gives the normalized lower bound
\[
  \frac14\left(1+\frac1{e_2}-\frac{|S|+1}{M}\right)\log Y
                     -O_{\psi,\varepsilon}(\log^2(2x)).
\]
Since \(M\) grows and \(e_2\le o\), this proves the result.
\end{proof}

These lower bounds enter the centered inequality by subtraction
on its right hand side. The archimedean Gram determinants there
are calculated on the same projected signed columns. Their
required upper bound, the existence of enough independent
columns, and the remaining strict cohomology lengths are open
steps in the argument for Dirichlet GRH.


\subsection{Rational divisors and the height of the signed period columns}
\label{sec:normalizedperiodheight463}

Use the common terminal normalization of
Corollary~\ref{cor:allcost454} on the polynomial-conductor
family of Corollary~\ref{cor:shortcenter460}. Let
\(\widetilde Y_1,\ldots,\widetilde Y_{q_c}\) be the signed
columns in this normalization. Write \(\widetilde G_\sigma\)
for the Gram matrix of their projections in the positive class
metric at an embedding \(\sigma:J\hookrightarrow\mathbb C\).
Assume the projections are independent when a logarithm of a
Gram determinant occurs.

At each unramified \(p\ge5\) inert in \(K_d\) let
\(g_{i,p}=\sum_{j=1}^m V_p(\nu_{ij})\). Extend \(U_p\) by
zero as in Corollary~\ref{cor:inertcorrection461}. Define
\[
    b_i^{\rm per}=\prod_p p^{\,g_{i,p}-\lceil U_p\rceil},
       \qquad \rho_x=2\sum_p(\lceil U_p\rceil-U_p)\log p ,
\]
with both operations restricted to these primes. The data here
are calculated before any values of period integrals are known.

\begin{lemma}[An integer divisor for each individual column]\leavevmode
\label{lem:periodinteger463}
Each \(b_i^{\rm per}\) is a positive integer. Division by it
leaves a gain of at least \(\lceil U_p\rceil\) at each indicated
prime on the individual column in adapted frames, before its
specified scalar, stabilizer and projector denominators.
For every fixed \(A>0\),
\[
       \frac2{q_c}\sum_i\log b_i^{\rm per}
                    =\mathcal D_u-\rho_x,\qquad
       0\le\rho_x=o_{\psi,\varepsilon,A}(b/x^A).
\]
\end{lemma}
\begin{proof}
On each word the integral gains from the edges give
\(g_{i,p}\ge U_p\), by
Theorem~\ref{thm:cycleweyl461} and
Lemma~\ref{lem:wordfirstrow450}. At a prime inert in both
quadratic fields this uses only \(g_{i,p}\ge0\). Since
\(g_{i,p}\) is an integer, all the exponents in the product
are nonnegative. They have finite support by
Corollary~\ref{cor:inertcorrection461}. The divisibility in
adapted frames follows from the edge inclusions on each path
in the proof of that corollary. Removing the rational integer
subtracts the same valuation at every coefficient place over
\(p\). The terminal normalization is a unit at these primes.

Taking logarithms gives the identity from the definition of
\(\mathcal D_u\). A rounding term can occur only when \(U_p>0\),
which forces \(p\le2nxX'\) by Theorem~\ref{thm:centerwords451}
and the actual word cutoff. Each term before multiplication by
\(2\log p\) is less than one. Chebyshev's bound gives
\(\rho_x=O(nxX')\), which is \(o(b/x^A)\) since \(n\) is
polynomial, \(X'=\exp(x^{\varepsilon/3})\), and
\(b=\exp(x^\varepsilon)\).
\end{proof}

Let \(G^\circ_\sigma\) be the Gram matrix for the projections
of \(\widetilde Y_i/b_i^{\rm per}\). To specify all other terms,
put
\[
\begin{aligned}
 K_u={}&\frac1{q_c}\sum_{i,j}
  \left(d_{ij}\log\operatorname{rd}(E_x)
       -\tfrac12{\bf1}_{d_{ij}\ {\rm odd}}\log d_0\right)\\
 &+\frac2{q_c}\sum_{p\in\mathcal M\setminus\{T\}}
                   \sum_{i,j} J^{\rm mat}_{ij,p}\log p
  +\frac1{q_c}\sum_{i,j}
       \left(2fr_j+\sum_\tau\nu_{ij,\tau,1}\right)\log T\\
 &+\frac2{q_c}\sum_i\log C_i+2\log D+\log m_G .
\end{aligned}
\]
Here \(f=[E_x^+:\QQ]\), \(r_j\) is the source rank on the
edge and \(d_{ij}\) its total differential degree.
The row-prime terms are the finite Fourier costs of
Corollary~\ref{cor:splitcost456}; \(C_i,D,m_G\) clear the given
scalar, class mass and projector denominators. For any
such family of Gram matrices \(G\) use the real number
\[
            H(G)=\frac{1}{q_c[J:\QQ]}\sum_\sigma\log\det G_\sigma
                      +K_u .
\]

\begin{theorem}[A centered bound for the divided columns]\leavevmode
\label{thm:primitivegram463}
With the specified tests on each \(E_x\), for every fixed \(A>0\),
\[
       \Delta_u\le H(G^\circ)-A_{\rm sig}(x)+W_{d_0}(u)
                     -\rho_x+o_{\psi,\varepsilon,A}(b/x^A).
\]
The two heights satisfy the exact identity
\[
           H(\widetilde G)-H(G^\circ)
                  =\frac{2}{q_c}\sum_i\log b_i^{\rm per}.
\]
\end{theorem}
\begin{proof}
Projection commutes with multiplication of a column by a
rational scalar. At every embedding the Gram determinant is
therefore multiplied by \(\prod_i(b_i^{\rm per})^{-2}\).
This proves the second identity. At all row and terminal
primes, and at primes ramified in \(E_x\), the integers being
removed are units. At the other finite places the precise
subtraction is accounted for in Lemma~\ref{lem:periodinteger463}.
In particular there is no need to add the integer once more
to the scalar denominators in \(K_u\). Corollaries
~\ref{cor:shortcenter460} and~\ref{cor:inertcorrection461}
give \(\Delta_u\le H(\widetilde G)-A_{\rm sig}(x)
 +W_{d_0}(u)-\mathcal D_u+o(b/x^A)\).
Insert the two exact identities to obtain the result.
\end{proof}

\begin{corollary}[A necessary height for independent projected columns]\leavevmode
\label{cor:heightfloor463}
For the word weights and sparse sets on these intervals one has
\[
 \begin{split}
        H(G^\circ)&\ge(\tfrac12-o(1))hma\log X',\\
        H(\widetilde G)&\ge(\tfrac34-o(1))hma\log X'.
 \end{split}
\]
For the allowed cosets in Corollary~\ref{cor:dyadicchoice462}
the second coefficient is at least
\(\tfrac34+1/(4e_2)-o(1)\). The lower bounds concern the full
quantity \(H\), including \(K_u\).
\end{corollary}
\begin{proof}
By Corollary~\ref{cor:count111},
\(2\Delta_u/(hmb\delta)=\Gg_\psi(x)
+O(x^{1/2+\varepsilon/2})\), where the notation on the right
denotes the fixed weighted sum of the \(G_\chi(x)\) in
\eqref{eq:window}. Each fixed \(G_\chi(x)=o_\chi(x)\) by
Proposition~\ref{prop:allzeros}. Since \(0<\delta\le1\),
\(\Delta_u=o(hma)\) uniformly over \(u\). No estimate for
zeros inside the critical strip is required beyond that
proposition's zero-free boundary argument.

In \(A_{\rm sig}(x)\) the three quadratic logarithmic
derivatives are \(O_{\psi,\varepsilon}(\log x)\) by
Theorem~\ref{thm:long232}. Also \(\log n=O_\psi(\log x)\)
and \(\delta=|T_2|/n=o(1)\) by the representative bound in
the proof of Theorem~\ref{thm:polycenter460}. Its displayed
formula for the complementary main term therefore gives
\[
                A_{\rm sig}(x)
                  =(\tfrac12+o(1))hma\log X'.
\]
The bound of Corollary~\ref{cor:shortcenter460} makes
\(W_{d_0}(u)=o(hma)\). Rearranging the centered inequality
proves the first claim. The exact height difference, the
rounding bound, and Theorem~\ref{thm:inertmass462} now give
the second. Apply Corollary~\ref{cor:dyadicchoice462} for the
last coefficient.
\end{proof}

For comparison with the precision actually needed, the size
conditions preceding \eqref{eq:gap} give
\(\mathcal B/(hmb)\ll_{\psi,\varepsilon}x^{1/2+2\varepsilon}\).
Indeed the three contributions after division are bounded by
\((2x/m)(1+m^2/k)\), \(2(n/h)mx^\varepsilon\), and
\((h_2/h)(k/n^2)x^{2\varepsilon}\). Use
\(h\ge n x^{-\varepsilon/8}/2\), \(h_2\le n\), the stated
bounds for \(m,k\), and \(m^2/k\le m/x\).
Thus \(x^{5\varepsilon}\mathcal B/(hma)
\ll x^{-1/2+7\varepsilon}\). Proving a sufficiently precise
upper bound for the divided Gram determinants at this scale
still requires information on the actual signed periods.





\subsection{First prime powers in the actual projected columns}
\label{sec:firstpowers668}

Take the polynomial-conductor intervals of
Theorem~\ref{thm:polycenter460}. Fix a node \(x\) and one of its
translates \(u\). Work with the ordinary word, its integral Weyl
table and its projected small test lattice \(eS\) as in
Theorem~\ref{thm:fullsignal452}. Write \(L=2nx\exp(x^{\eps/3})\)
and let \(\mathcal P\) consist of the unramified rational primes
\(5\le p\le L\) which split in \(K_q\) and are inert in \(K_d\).
In the sum defining \(U_p\) in Section~\ref{sec:signedsignal450},
denote the summands of a given odd exponent \(j\), summed over
all edges and high embeddings, by \(U_{p,j}\). For \(Z=x^{3/4}\)
put
\[
 H_p=\sum_{\substack{j\ge3\ \mathrm{odd}\\p^j\le Z}}U_{p,j},
 \qquad g_p=\lfloor H_p\rfloor,\qquad
 B_p=\lceil hma/p\rceil .
\]
All three numbers are computed from the input word. In particular
\(H_p=0\) except possibly at the at most \(x^{1/4}\) rational
primes \(p\le x^{1/4}\).

\begin{lemma}[The omitted odd prime powers]
\leavevmode\label{lem:higherpowers668}
At every node and translate,
\[
 0\le \sum_{p\in\mathcal P}(U_p-U_{p,1}-H_p)\log p
     \ll hma\,x^{-1/2}\log(2x).
\]
\end{lemma}
\begin{proof}
The defining ramp is nonnegative and bounded by \(a/p^j\)
per edge and high embedding. Put \(P=Z^{1/3}\). For \(p>P\)
the entire sum of \(p^{-j}\) for \(j\ge3\) is at most \(2p^{-3}\).
Integral comparison of \(\sum_{n>P}(\log n)n^{-3}\) gives
\(O(P^{-2}\log(2P))\). For \(p\le P\) the first omitted odd
power exceeds \(Z\). Its tail has sum at most \(2/Z\), so this
range contributes \(O(P\log(2P)/Z)\). Both bounds are
\(O(x^{-1/2}\log(2x))\). Multiplication by \(hma\) proves the
claim without a hypothesis on primes in a short interval.
\end{proof}

For this fixed count choose any nonempty finite list of its
ordinary signed columns having the word weights just used.
One column suffices to define the quantities below. In a common
normal coefficient field \(J\) of degree \(d\) let
\(\mathfrak a_c\) be the fractional ideal generated by all
the pairings of column \(c\) with \(eS\). The zero ideal is allowed.
Use \(C_c,D,m_G\) from Theorem~\ref{thm:fullsignal452}, including
the stabilizer factors and the denominator of the small projector.
For \(p\in\mathcal P\) define
\[
 \begin{split}
 t_{c,p}&=d^{-1}v_p(N_{J/\QQ}\mathfrak a_c)+v_p(C_cDm_G),\\
 v_p^*&=\min\!\left\{B_p,\,
              \max\!\left(0,\min_c t_{c,p}-g_p\right)\right\},\\
 N(x,u)&=\prod_{p\in\mathcal P}p^{d v_p^*}.
 \end{split}
\]
Here \(v_p(N0)=+\infty\). Thus \(N(x,u)\) is defined even when
an entire projected table vanishes. At a prime where all entries
vanish the exponent used is exactly \(dB_p\).

\begin{theorem}[An integer upper comparison for one translate]
\leavevmode\label{thm:projectedinteger668}
The exponents \(d v_p^*\) are nonnegative integers. For the
actual field and word at the node put
\[
 M(x,u)=\frac{2}{Q\delta}
       \left(A_{\rm sig}(x)-W_{d_0}(u)
                   -2\sum_{p\in\mathcal P}H_p\log p\right).
\]
Uniformly in the translate and the chosen columns,
\[
 \Gg_\psi(x)\ \le\ \frac{4\log N(x,u)}{dQ\delta}
                    -M(x,u)+O_{\psi,\eps}(x^{1/2+2\eps}).
\]
Enlarging \(J\) preserves \(\log N(x,u)/d\).
At any chosen subset \(\mathcal E\subset\mathcal P\) one can
use \(v_p^*=B_p\) without reading a projected coefficient there.
The resulting increase in the right side is at most
\[
 O_{\psi,\eps}\left(x^{1+\eps}
          \sum_{p\in\mathcal E}\frac{\log p}{p}\right)
                     +O_{\psi,\eps,A}(x^{-A})
 \qquad(A>0).
\]
\end{theorem}
\begin{proof}
The norm of a nonzero fractional ideal is positive rational.
For its normalized valuation one has
\(d^{-1}v_p(N\mathfrak a_c)=
\sum_{w\mid p}[J_w:\QQ_p]v_w(\mathfrak a_c)/d\),
where \(v_w(p)=1\). This follows from the usual product of residue
degrees with integral ideal exponents. Apply
Theorem~\ref{thm:fullsignal452} on the image lattice to the
single column \(c\). At every active prime it gives
\(t_{c,p}\ge U_p\). For a zero column this holds with infinite
left side. Hence \(t_{c,p}-g_p\ge U_{p,1}\) at every active
prime. The other bound \(B_p\ge U_{p,1}\) follows by dropping
every threshold at the first power. At an inactive prime the
first-power term is zero. We obtain \(U_{p,1}\le v_p^*\) at all
primes. All finite valuations in the definition have denominator
dividing \(d\), and the upper bound is an integer. This proves
the integrality assertion. Ideal extension multiplies its norm
valuation by the relative degree, since the local degrees add in
a tower. Division by the new \(d\) therefore preserves all
\(t_{c,p}\), including for the zero ideal.

Insert the bound for \(U_{p,1}\) and
Lemma~\ref{lem:higherpowers668} in the exact centered identity
of Theorem~\ref{thm:polycenter460}. The signs give
\(-A_{\rm sig}(x)+W_{d_0}(u)\) there. Multiplication by \(2/(Q\delta)\)
gives the displayed main term. Here \(Q=hmb\), \(a\le2xb\) and
\(\delta\ge x^{-\eps}\), so the normalized higher-power error
is \(O_{\psi,\eps}(x^{1/2+2\eps})\). The remainder
\(o(b/x^A)\) in that identity can use any fixed \(A\).
Finally Corollary~\ref{cor:count111} compares
\(2\Delta_u/(Q\delta)\) with \(\Gg_\psi(x)\) at this individual
translate, with an error \(O_{\psi,\eps}(x^{1/2+\eps/2})\).

Replacing an exponent by its upper bound adds at most
\(4(Q\delta)^{-1}\sum_{\mathcal E}(hma/p+1)\log p\).
In the term from the ceiling, the sum of logarithms is at most
\(L\log(2L)\). Since \(n\) grows polynomially and
\(b=\exp(x^\eps)\), \(L\log(2L)/(Q\delta)=O_{\psi,\eps,A}(x^{-A})\)
for every fixed \(A\). This gives the claimed allowance for
exceptional primes.
\end{proof}

\begin{corollary}[Integral tests at one translate per scale]
\leavevmode\label{cor:intervalintegers668}
Use the nodes and weights of Theorem~\ref{thm:threeaverages667}
on the polynomial-conductor intervals. At each node \(x_j\) choose
one translate and let \(N_j,d_j,M_j,Q_j,\delta_j\) denote the data
of Theorem~\ref{thm:projectedinteger668} there. Put
\(E_k=(\sigma-\theta)^{-k-1}\) and select an integer \(H\ge1\)
with
\[
 H E_k \ \ge\ 1+\sum_j d_j\sum_{p\in\mathcal P_j}B_{j,p}\log p .
\]
For \(l=0,1,2\) form the positive rational integers
\[
 P_l=\prod_j N_j^{\lceil 4H w_{lj}/(d_j Q_j\delta_j)\rceil}.
\]
If for every rational frequency and arbitrarily distant common
intervals at that frequency one has, simultaneously,
\[
      H^{-1}\log P_l-\sum_j w_{lj}M_j\le B_\tau E_k
          \qquad(l=0,1,2),
\]
every open-strip zero of the family has real part at most
\(\theta\), for \(1/2+4\eps\le\theta<1\). The choice of \(H\)
uses only finite input data and the positive error allowed in
the test.
\end{corollary}
\begin{proof}
Since \(0\le\log N_j\le d_j\sum_{p\in\mathcal P_j}B_{j,p}\log p\),
rounding all exponents up changes each normalized sum by a number
between zero and \(E_k\). Multiply the node inequality by the
positive weights. Equivalently, interpolate the resulting node
majorants of \(\Gg_\psi\) by Lemma~\ref{lem:convexcount666}.
The interpolation error is \(O_{\psi,\eps}(x^{1/2+4\eps})\) on a
mesh with fixed constant. Its integral and that of the node error
against the three kernels are \(O_{\psi,\eps}(E_k)\), by the
gamma integral in Theorem~\ref{thm:threeaverages667}.
That theorem's pole argument now gives the conclusion. No
identification between the coefficient fields of different nodes
is needed to multiply their rational integers.
\end{proof}

The finite exponents retain the original signed projections and
the ramified prime in the centering term. The estimate still
needed is an upper bound on the logarithms of these integers
after that centering, uniform along the specified intervals.


\subsection{The projected period at primes between the square root and the scale}
\label{sec:mediumsignal670}

On the polynomial-conductor intervals fix the two auxiliary primes
\(q,d_0\) of Theorem~\ref{thm:conductor232}. Let \(\mathcal P_0\)
be the rational primes prime to \(2Fqd_0\) which split in \(K_q\)
and are inert in \(K_d\). Put
\(\mathcal I_x=\{p\in\mathcal P_0:\sqrt{x}<p\le x\}\).
We give a comparison requiring projected integral coefficients
only at these primes. The translate in those coefficients can
be chosen separately at each prime.

Write \(w_{\mathbf1}=1,\ w_\eta=1/(3o+1)\) for the other characters
of \(A/A_0\), and use the primitive representatives
\(\alpha_{ij\eta}\) of Theorem~\ref{thm:orientation46}. Let
\(c_\psi(n)\) be as in Lemma~\ref{lem:convexcount666}.
For these fixed auxiliary primes define the real constants
\[
 \begin{split}
 D^*&=\sum_{\alpha_{ij\eta}\ne\mathbf1}(-1)^j w_\eta
             \real\frac{L'(1,\alpha_{ij\eta})}{L(1,\alpha_{ij\eta})},\\
 H^*&=\sum_{p\in\mathcal P_0}\log p
           \sum_{\substack{v\ge3\\v\ {\rm odd}}}\frac{c_\psi(p^v)}{p^v}.
 \end{split}
\]
The second series converges absolutely. Put
\[
 M_0(x)=x(\log x-1-\gamma-D^*-4H^*)
        -4x\sum_{\substack{p\in\mathcal P_0\\p\le\sqrt{x}}}
               \frac{c_\psi(p)\log p}{p}.
\]

\begin{proposition}[The centered first-power sum in a single range]
\leavevmode\label{prop:primerange670}
Uniformly on these common intervals,
\[
 \Gg_\psi(x)=4\sum_{p\in\mathcal I_x}
          c_\psi(p)(x/p-1)\log p-M_0(x)
                 +O_{\psi,\eps}(x^{1/2+\eps}).
\]
\end{proposition}
\begin{proof}
Form \(\sum_{i,j,\eta}(-1)^jw_\eta\real G_{\alpha_{ij\eta}}(x)\).
The \(i=j=0\) terms are exactly \(\Gg_\psi(x)\).
There are \(3o\) other terms, each bounded by
\(O(x^{1/2+\eps/2})\) by the auxiliary twist estimate on the
common interval. In the defining ramps, \(q,d_0>x^{401}\)
and so neither prime occurs. Coprimality of their conductors
with \(F\) makes the sum over \(i,j\) at \(p^v\) equal to the
multiplier
\((1+\xi_q(p)^v)(1-\xi_d(p)^v)\). It is four exactly at odd
\(v\) and primes with the indicated two signs and is zero
elsewhere. The primes dividing \(2F\) and the prime three split
in \(K_d\) on sufficiently distant intervals, by its prescribed
congruences. The remaining coefficient is \(c_\psi(p^v)\).
The principal correction is \(-x\log x+(1+\gamma)x\) and the
sum of the other linear corrections is \(xD^*\). This gives
the expression with all odd powers directly from the primitive
character formula.

For \(p\le\sqrt x\) at the first power, replacing \(x/p-1\)
by \(x/p\) costs \(O(\sqrt x\log(2x))\) after summation since
\(c_\psi(n)<2\). The odd powers \(v\ge3\) contribute
\(4xH^*+O(x^{1/3}\log^2(2x))\). Indeed there are at most
\(x^{1/3}\log_2 x\) pairs with \(p^v\le x,\ v\ge3\), so the
discarded unit in each ramp costs at most the stated error.
For \(p^v>x\), the tail of
\(\sum p^{-v}\log p\) is \(O(x^{-2/3}\log(2x))\) by the
two-range estimate in Lemma~\ref{lem:higherpowers668}, now with
cutoff \(Z=x\). Substitution proves the formula. In this formula
the arguments of the ramps are \(x\). The ramified prime which
appears at the much larger word arguments in
Theorem~\ref{thm:polycenter460} exceeds this cutoff.
\end{proof}

For \(p\in\mathcal I_x\) let \(z_p\) denote the \(A\)-coordinate
of its Frobenius. With \(T,T_2\) of
Theorem~\ref{thm:rangenew50}, set \(k_p=h\delta c_\psi(p)\).

\begin{lemma}[A translate-independent integral multiplicity]
\leavevmode\label{lem:primeincidence670}
The integer \(k_p\) counts the high embeddings carried to the
distinguished set by the first power at \(p\), for every
translate \(u\). In particular \(0\le k_p\le h\).
\end{lemma}
\begin{proof}
In the product coordinates of Theorem~\ref{thm:orientation46},
the required condition on \(t\in T\) is \(t+z_p\in u+T_2\).
The sign pair is \((0,1)\), so the coordinate change in that
theorem gives this condition also for odd \(\psi\). Expanding
the indicator on \(A\) gives the count
\[
 \frac1n\sum_{\theta\in\widehat A}\theta(z_p-u)
       \left(\sum_{t\in T}\theta(t)\right)
       \left(\sum_{z\in T_2}\overline{\theta(z)}\right).
\]
If \(\theta\) is nontrivial on \(A_0\), the product of the two
sums vanishes by the subgroup invariance and residue balance in
Theorem~\ref{thm:rangenew50}. Otherwise the second sum is \(h_2\),
\(\theta(u)=1\), and the first sum is \(hw_\theta\) by the
prescribed proportions. On these unramified primes the surviving
Fourier sum is \(c_\psi(p)\), since conjugate terms pair.
Multiplication by \(hh_2/n\) proves the assertion.
\end{proof}

Put \(X'=\exp(x^{\eps/3})\), \(a=\lceil xb\rceil\) and define
two nonnegative integers from the common input at \(p\in\mathcal I_x\) by
\[
 G_p=\left\lfloor m(h-k_p)\left(a/p-2b/(nX')\right)\right\rfloor,
 \qquad B_p=\lceil m k_p a/p\rceil .
\]
Here \(x\) is sufficiently large. Choose any finite list of column
and translate pairs from the actual ordinary words with these
parameters. For a pair \((u,c)\) write \(t_{u,c,p}\) for the
normalized valuation of its projected entry ideal, including
\(C_cDm_G\), in Theorem~\ref{thm:projectedinteger668}.
Use one coefficient field of degree \(d\) for this finite list.
At each \(p\) any subset of the pairs may be read. Put
\[
 a_p^*=\min\big\{B_p,\ \max(0,\min_{(u,c)\ {\rm read\ at}\ p}
                                   t_{u,c,p}-G_p)\big\},\qquad
 N_0(x)=\prod_{p\in\mathcal I_x}p^{d a_p^*},
\]
taking the minimum of an empty list to be \(+\infty\).
If the whole list is empty use \(d=1\).

\begin{theorem}[A prime-by-prime choice of projected entries]
\leavevmode\label{thm:primeprojections670}
The number \(N_0(x)\) is a positive integer, and uniformly over
these choices of actual projected columns,
\[
 \Gg_\psi(x)\le\frac{4\log N_0(x)}{dQ\delta}-M_0(x)
                        +O_{\psi,\eps}(x^{1/2+\eps}).
\]
Replacing any selections by the empty list on a set
\(\mathcal E\subset\mathcal I_x\) raises the right side by at most
\(O_\psi(\sum_{p\in\mathcal E}\log p)+O_{\psi,\eps,A}(x^{-A})\)
for every fixed \(A>0\). The three integral tests of
Corollary~\ref{cor:intervalintegers668} hold as a sufficient
condition with \(N_0,M_0\) in place of \(N,M\), summing the
rounding allowance only over \(\mathcal I_x\) with the \(B_p\)
of this subsection.
\end{theorem}
\begin{proof}
All target embeddings at the first power for the chosen signs
belong to \(\Phi_1\). On its nondistinguished part the word
threshold is at most \(2b/(nX')\) eventually. In fact it differs
from \(b/(nX')\) by at most \(2+2mH_0\), which is polynomial in
\(x\), as in Theorem~\ref{thm:centerwords451}.
At a distinguished place the threshold lies between \(b\) and
\(b+\sqrt b\) by Corollary~\ref{cor:count111}.
Since \(a/p\ge b\) on \(\mathcal I_x\), there are \(h-k_p\)
nondistinguished ramps at every edge, each at least
\(a/p-2b/(nX')>0\). The other \(k_p\) ramps are at least
\((a/p-b-\sqrt b)_+\). Thus the full word sum \(U_p\) for
\emph{each} translate is at least
\[
 G_p+m k_p(a/p-b-\sqrt b)_+ .
\]
This also holds if \(U_p=0\), when both terms on the right
vanish. At active primes the single-column estimate of
Theorem~\ref{thm:fullsignal452} applies to each \(t_{u,c,p}\).
At inactive primes the maximum with zero suffices for the
distinguished lower bound. Indeed if a read translate is inactive,
the displayed lower bound forces both \(G_p=0\) and
\(m k_p(a/p-b-\sqrt b)_+=0\), independently of the translate.
The upper cutoff \(B_p\) also dominates the latter quantity.
This proves
\[
        m k_p(a/p-b-\sqrt b)_+\ \le\ a_p^*
                  \ \le\ m k_p a/p+1 .
\]
The exponents become integers after multiplication by \(d\) by
the ideal norm formula of Theorem~\ref{thm:projectedinteger668}.

Using \(k_p=h\delta c_\psi(p)\), \(Q=hmb\), and \(a\ge xb\),
the lower inequality bounds \(c_\psi(p)(x/p-1)\) above by
\(a_p^*/(Q\delta)+c_\psi(p)/\sqrt b\). The sum of the errors
times \(\log p\) is \(O_{\psi,\eps,A}(x^{-A})\) for any fixed
\(A\), since \(p\le x\). Proposition~\ref{prop:primerange670}
gives the asserted upper bound. At a prime switched to its
upper cutoff the increase is at most \(m k_p(b+\sqrt b)+1\)
before normalization, by the same two inequalities. After
division by \(Q\delta\) and summation this gives the stated cost.
Finally the ceiling and logarithm argument of
Corollary~\ref{cor:intervalintegers668} applies to these integers
with exactly the same positive kernels and interpolation error.
\end{proof}

The geometric term now reads integral coefficients only in
\((\sqrt x,x]\), with prime-dependent choices of translates on the
same field. The outstanding upper bound compares its logarithm
with \(M_0(x)\) to square-root accuracy after averaging. The two
node inequalities above by themselves allow a total excess of
order \(x\log(2x)\).


\subsection{How many primes the integral comparison must detect}
\label{sec:detected-prime-mass-671}

Fix \(\psi,\eps\) and the polynomial-conductor intervals of
Section~\ref{sec:mediumsignal670}. At a node put
\[
 s_p(x)=\frac{B_p-a_p^*}{Q\delta},\qquad
 R(x)=4\sum_{p\in\mathcal I_x}
                 [c_\psi(p)-s_p(x)]_+\log p,\qquad
 H_0(x)=\frac{4\log N_0(x)}{dQ\delta}-M_0(x).
\]
Thus \(s_p(x)=0\) when no projected column improves the elementary
upper bound \(B_p\). We measure how often this can happen if the
proposed upper comparison is to hold.

\begin{lemma}[The excess in the upper comparison]
\leavevmode\label{lem:excessmass671}
Uniformly over all choices of entries, on every node,
\[
             H_0(x)=\Gg_\psi(x)+R(x)+O_{\psi,\eps}(x^{1/2+\eps}).
\]
On the interior \(2Y\le x\le Y^{3/2}/2\) one also has
\[
          4\sum_{p\in\mathcal I_x}c_\psi(p)\log p
                            =(1+o_{\psi,\eps}(1))x .
\]
Both estimates are unconditional.
\end{lemma}
\begin{proof}
Write \(B_p=m k_p a/p+\eta_p\) with \(0\le\eta_p<1\).
Since \(a=xb+O(1)\), \(k_p=h\delta c_\psi(p)\), the normalized
upper cutoff differs from \(c_\psi(p)x/p\) by at most
\(c_\psi(p)/(bp)+1/(Q\delta)\). The lower bound for \(a_p^*\)
in Theorem~\ref{thm:primeprojections670} consequently gives
\[
 0\le s_p(x)\le c_\psi(p)(1+b^{-1/2})+1/(Q\delta).
\]
Replacing \(c_\psi(p)-s_p(x)\) by its positive part costs
\(O_{\psi,\eps,A}(x^{-A})\) after summing with \(\log p\), for
any fixed \(A\). Subtract the prime sum of
Proposition~\ref{prop:primerange670} from the definition of
\(H_0\) and use the cutoff estimate just proved. This gives the
first formula.

For the second, during a given interval keep the auxiliary
primes fixed. Consider the nonnegative ramp sum
\[
 F(y)=4\sum_{p\in\mathcal P_0}\log p
       \sum_{v\ge1\ {\rm odd}}c_\psi(p^v)(y/p^v-1)_+ .
\]
By the character expansion in Proposition~\ref{prop:primerange670},
\[
 F(y)=y\log y-(1+\gamma+D^*)y+o_{\psi,\eps}(y)
                          \qquad(Y\le y\le Y^{3/2}).
\]
Indeed each \(\Gg_\psi(y)/y\) tends to zero for fixed \(\psi\). Every open-strip zero has real part strictly
below one, and dominated convergence applies in the absolutely
summable explicit formula. The auxiliary twists contribute \(O_\psi(y^{1/2+\eps/2})\)
uniformly on the interval. The coefficient \(D^*\) may depend
on the two primes but is constant across this interval.

The function \(F\) is convex. For \(0<\eta<1/4\) its right
derivative at an interior \(x\) is between the backward and
forward difference quotients using steps \(\eta x\).
For \(y\log y\) both such slopes equal \(\log x+1+O(\eta)\).
The linear term is exact and the errors in the slopes are
\(o_{\psi,\eps}(1)/\eta\), uniformly in the interior. Thus
\(xF'_+(x)-F(x)=x(1+O(\eta)+o_{\psi,\eps}(1)/\eta)\).
First let \(Y\) tend to infinity and then \(\eta\) to zero.
A ramp at \(n\le x\) contributes exactly one to
\(x(d/dx)_+(x/n-1)_+-(x/n-1)_+\), also at the endpoint.
It follows that \(xF'_+(x)-F(x)\) is four times the sum of
\(c_\psi(p^v)\log p\) for \(p^v\le x\), \(p\in\mathcal P_0\),
\(v\) odd. The primes below \(\sqrt x\) at the first power cost
\(O(\sqrt x\log(2x))\), and all higher powers cost
\(O(x^{1/3}\log^2(2x))\) by the counting estimate in
Proposition~\ref{prop:primerange670}. This proves the second formula.
\end{proof}

Write \(w_j=\int_U^{3U/2}K_k(u)\lambda_j(e^u)\,du\) for the
unmodulated positive weights of Theorem~\ref{thm:threeaverages667},
and put \(L_k=(\sigma-1)^{-k-1}\).

\begin{theorem}[A necessary bound on the undetected prime mass]
\leavevmode\label{thm:missedmass671}
Fix \(1/2+4\eps\le\theta<1\) and a uniform mesh constant.
There is a number \(\delta_*>0\), depending on the fixed
family and \(\theta,\eps\), with the following property.
On any sequence of common intervals on which
\[
              \sum_j w_jH_0(x_j)\le B(\sigma-\theta)^{-k-1},
\]
one has
\[
 \sum_j w_jR(x_j)\ll_{\psi,\eps,\theta,B}Y^{-\delta_*}L_k,
 \qquad
 4\sum_jw_j\sum_{p\in\mathcal I_{x_j}}c_\psi(p)\log p
                         =(1+o_{\psi,\eps}(1))L_k .
\]
Constants may also depend on the mesh constant. For any fixed
\(0<\lambda\le1\), the fraction of this last weighted mass on
which \(s_p(x_j)\le(1-\lambda)c_\psi(p)\) is
\(O_{\psi,\eps,\theta,B}(\lambda^{-1}Y^{-\delta_*})\).
In particular the integral tests cannot hold along such intervals
with only the elementary upper cutoffs at every node.
\end{theorem}
\begin{proof}
Let \(\Phi_\psi\) be the transform in Lemma~\ref{lem:mellin}.
The closed disk about \(\sigma\) of radius \(\sigma-1\) lies in
\(\real s\ge1\), where there are no poles. The possible poles
are locally finite in \(\real s>0\), and the transform is
holomorphic at one. Compactness therefore allows a strictly
larger pole-free closed disk of radius \(r<\sigma\). Cauchy's
inequality gives \(|A_k(\sigma)|\ll_\psi r^{-k}\), at this
single real center and for every \(k\).

Apply Lemma~\ref{lem:convexcount666} pointwise on each mesh cell
with the interpolation weights, and integrate against \(K_k\).
It bounds the weighted sum of the nodal \(\Gg_\psi(x_j)\)
\emph{from below} by the integral of \(\Gg_\psi(e^u)\) on the
window minus \(O((\sigma-1/2-4\eps)^{-k-1})\).
The tail proof in Theorem~\ref{thm:threeaverages667} identifies
that integral with \(A_k(\sigma)+O_\psi(\sigma^{-k})\).
Also \(\sum_jw_j x_j^{1/2+\eps}\ll(\sigma-\theta)^{-k-1}\)
since adjacent node ratios are at most two on distant windows.
Combining the hypothesis and Lemma~\ref{lem:excessmass671}
gives
\[
 0\le\sum_j w_jR(x_j)
      \ll_{\psi,\eps,\theta,B} r^{-k}+(\sigma-\theta)^{-k-1}.
\]
Both \(r\) and \(\sigma-\theta\) exceed \(\sigma-1\). As
\(k=\lceil5\sigma\log Y/4\rceil\), division by \(L_k\)
gives the asserted positive power saving.

For the denominator of the fraction,
\(\sum_j\lambda_j(x)x_j=x\) exactly. Thus
\(\sum_jw_jx_j=\int_U^{3U/2}K_k(u)e^u\,du
=(1+o(1))L_k\). The omission of nodes outside the interior
in the second assertion of the lemma is harmless here, both
for \(x_j\) and for the prime mass, which is
\(O(x_j\log(2x_j))\) on the whole interval. Indeed an endpoint
node outside the interior, or any such node near either end,
has a hat function supported within \(O(1)\) in log scale of
one of the two endpoints, since adjacent ratios are at most two.
On those strips \(x_j\log(2x_j)\ll e^{2u}\) wherever its hat is
nonzero. The tail estimates in Lemma~\ref{lem:kernel} remain
\(O(\sigma^{-k})\) upon shifting their cutoff by a fixed
distance in \(u\), since this changes the exponential majorants
in that proof only by constant factors. The sum of all the hats is one.
Consequently the interior estimate of Lemma~\ref{lem:excessmass671}
gives the second formula after summation. Each prime in the
stated exceptional set contributes at least
\(\lambda c_\psi(p)\) to the positive part defining \(R\).
Division of the first formula by the second gives the fraction.
For the unmodulated test the three-weight hypothesis at
\(\tau=0,\ l=0\) is twice the hypothesis here, up to the upward
rounding in Corollary~\ref{cor:intervalintegers668}. The rounding
is nonnegative. This proves the last assertion for the integral
tests as well.
\end{proof}

This gives a quantitative obstruction on the very projected
lattices used for the comparison. A proof using these integral
tests must recover almost all the indicated prime mass, to the
displayed precision. The statement supplies a necessary condition
on the valuations and leaves its verification on the period
columns open.


\subsection{The common divisor at every prime in the middle range}
\label{sec:commonsaturation672}

Keep the actual polynomial-conductor fields and the set
\(\mathcal I_x\) of Section~\ref{sec:mediumsignal670}. Here \(k\)
is the initial rank of a theta word. Write \(L_u\subset\Phi_1\)
for the \(h_2\) distinguished embeddings at translate \(u\).
At an edge \(i\) of a word \(w\), write \(\nu_i\) for its
partition tuple and \(s_i\) for its target rank. Define
\(g_{u,w,p}=\sum_{i=1}^m V_p(\nu_i)\), where
\(V_p\) is the maximum over the two orientations of each even
Frobenius cycle in Theorem~\ref{thm:cycleweyl461}.

\begin{theorem}[Saturation before division of the period columns]
\leavevmode\label{thm:rawsaturation672}
Fix \(\psi\) and \(0<\eps<1/52\). At every sufficiently large
node and every \(p\in\mathcal I_x\), for every actual word and
every column on that word allowed in
Theorem~\ref{thm:primeprojections670},
\[
          t_{u,c,p}\ \ge\ g_{u,w,p}\ \ge\
          \frac{m k h_2 b}{2p}\ >\ G_p+B_p .
\]
Consequently the capped exponents in that theorem satisfy
\(a_p^*=B_p\), for every choice of columns or translates at
each prime.
\end{theorem}
\begin{proof}
The word ranks and the number of parts are bounded by powers of
\(x\). On each of the \(m\) edges the first \(k\) parts at
\(\tau\in L_u\) equal \(b+O(x^C)\). On \(\Phi_0\) all parts
except possibly the first have size \(O(bx^C/X'+x^C)\), with
\(X'=\exp(x^{\eps/3})\). The first is bounded by
\(a+O(bx^C/X'+x^C)\). These estimates follow from the partition
recurrence in the proof of Theorem~\ref{thm:factorcomp512}.
In detail, padding by zeros and subtracting the return shift
cancels the base weight on each newly introduced part. On the
original parts the two shifts at successive edges leave only
half the rank difference. Their total and the offsets in the
word letters are polynomial. Inserting the initial parts in
\eqref{eq:letters111} gives the stated bounds. In particular
the \emph{sum} of the parts at every \(\mu\in\Phi_0\) is at
most \(2a\), since the polynomial number of errors is \(o(b)\).

On the indicated signs the permutation \(\pi=\sigma_pc_E\)
interchanges \(\Phi_0\) and \(\Phi_1\). Choose \(\Phi_1\) as
the source half on every cycle. For each \(\tau\in L_u\), at
the destination \(\mu=\pi\tau\) retain just the \(k\) indices
from \(\tau\), of exponent one and row numbers \(1,\ldots,k\).
For each such index the list entry on edge \(i\) is at least
\[
 \max\{0,\lceil(\nu_{i,\tau,l}-s_i(p-1))/p\rceil\}
         \ \ge\ \frac{b-O(x^C)}p-s_i
         \ \ge\ \frac{3b}{4p}.
\]
The last inequality holds uniformly because \(p\le x\) and
\(b=\exp(x^\eps)\). The sum of the largest \(k\) entries is
at least the sum of these \(k\) entries. The penalty for choosing
them is at most the sum of all parts at \(\mu\), hence at most
\(2a\). Since \(pa/(kb)=O(x^2/k)=o(1)\), their net contribution
is at least \(kb/(2p)\). Distinct \(\tau\)'s yield distinct
destinations. All other destinations contribute nonnegatively,
by the empty choice. Summing first over destinations and then
over the edges proves the second inequality.

For the first inequality, the operator at edge \(i\) maps the
integral Weyl lattice into \(p^{V_p(\nu_i)}\) times the next
one, by Theorem~\ref{thm:cycleweyl461}. The integral frames
match at consecutive vertices. Clearing the coefficient
denominators, the finite stabilizers and the ordinary projector
therefore gives a multiple of \(p^{g_{u,w,p}}\) on every path,
also after taking the signed sum and pairing with each integral
test. These are exactly the denominators included in \(t_{u,c,p}\).
Corollary~\ref{cor:inertcorrection461} verifies the same assertion
at every coefficient place over \(p\) in the adapted frames.
In an arbitrary marked product frame the scalar change has
absolute norm one, as proved in Theorem~\ref{thm:fullsignal452}.
Thus the normalized rational valuation satisfies the asserted
inequality in either frame. A zero column has valuation
\(+\infty\).

Finally the definitions of the two cutoffs give
\(G_p+B_p\le mha/p+1\). Since \(h\le n\),
\(h_2\ge nx^{-\eps}\), \(a=xb+O(1)\), and
\(k\ge x^{9/2+10\eps}\), the ratio of \(mkh_2 b/(2p)\) to
\(mha/p+1\) tends uniformly to infinity. This also proves the
assertion for an empty list, by the convention in the definition
of \(a_p^*\).
\end{proof}

\begin{corollary}[Size of the undivided comparison on the common intervals]
\leavevmode\label{cor:rawfailure672}
On \(2Y\le x\le Y^{3/2}/2\) one has \(H_0(x)=(1+o(1))x\).
For the unmodulated weights of Theorem~\ref{thm:missedmass671},
\[
             \sum_j w_j H_0(x_j)=(1+o(1))(\sigma-1)^{-K-1},
\]
where \(K=\lceil5\sigma\log(Y)/4\rceil\) is the order of the
kernel. In particular the proposed upper test for the
undivided columns fails on every sufficiently distant interval
for any fixed \(\theta<1\) in the range of that theorem.
\end{corollary}
\begin{proof}
The preceding theorem gives \(s_p(x)=0\) for every \(p\in\mathcal
I_x\). Both identities of Lemma~\ref{lem:excessmass671} now give
the interior asymptotic since \(\Gg_\psi(x)=o_\psi(x)\).
For the weighted statement its proof and that of
Theorem~\ref{thm:missedmass671} give
\(4\sum_jw_j\sum_{p\in\mathcal I_{x_j}}c_\psi(p)\log p
=(1+o(1))(\sigma-1)^{-K-1}\). The evaluation of this last
sum there uses only the prime mass and the mesh. Also
\(\sum_j w_j x_j=(1+o(1))(\sigma-1)^{-K-1}\), and
\(\sup_{x\ge Y}(|\Gg_\psi(x)|+x^{1/2+\eps})/x\) tends to zero.
Apply the first identity of the lemma at every node.
Finally \(((\sigma-\theta)/(\sigma-1))^{K+1}\) tends to infinity.
\end{proof}

The failed comparison reads the valuations before removal of
the rational divisors in Lemma~\ref{lem:periodinteger463}. The
centered count in Section~\ref{sec:missing} and the valuations
after this division are separate open questions.


\subsection{The residual valuation of a divided period column}
\label{sec:dividedprimecomparison673}

We remove the divisor whose effect was exhibited in
Theorem~\ref{thm:rawsaturation672}. On the same fields and at
the same nodes, choose any finite list of signed projected columns.
Each list member \(c\) is attached to one word \(w=w(c)\) at
one translate \(u\). Scalar combinations within a fixed word
and its physical weights have their prescribed denominator
\(C_c\). Write \(d\) for the common coefficient field degree,
and \(e_{u,c,p}=t_{u,c,p}-g_{u,w,p}\) for the residual valuation,
with value \(+\infty\) for a zero column. By
Lemma~\ref{lem:periodinteger463} the integer
\[
 b^{\rm per}_{u,w}=\prod_{p}p^{g_{u,w,p}-\lceil U_p(u)\rceil}
\]
divides each cleared column in the adapted frames. Here the
product is over the unramified primes \(p\ge5\) inert in \(K_d\),
with the extension convention of that lemma. Division is by a
rational scalar on the entire column. Let \(t^\circ_{u,c,p}\)
be the valuation with the original denominator corrections on
this divided column. In the definition of \(a_p^*\) use
\(t^\circ\) in place of \(t\), calling the resulting exponent
\(a_p^\circ\). Write \(e_p=\min_{(u,c)}e_{u,c,p}\) over the
entries read at \(p\), with the same empty-list convention.
All sums in this subsection run over \(p\in\mathcal I_x\).

\begin{lemma}[The word threshold after centering]
\leavevmode\label{lem:residualbaseline673}
Put \(A_p=Q\delta c_\psi(p)(x/p-1)\). Then \(A_p\ge0\), and
uniformly on the common intervals,
\[
 \frac1{Q\delta}\sum_p\log p\,
     \max_u\big|\lceil U_p(u)\rceil-G_p-A_p\big|
            \ll_{\psi,\eps} x^\eps\log(2x).
\]
\end{lemma}
\begin{proof}
By Lemma~\ref{lem:primeincidence670} exactly \(k_p=h\delta
c_\psi(p)\) high embeddings at each edge have distinguished
first destination. There the threshold lies in \([b,b+\sqrt b]\).
Since the positive part is 1-Lipschitz and \(x/p\ge1\), replacing
each such ramp by \(b(x/p-1)\) costs at most \(\sqrt b+1/p\).
At each of the other \(h-k_p\) destinations the threshold lies
in \([0,2b/(nX')]\), and \(a/p\ge b\). Subtracting the
unrounded expression defining \(G_p\) leaves between zero and
\(2mh b/(nX')\) from these destinations. The two roundings cost
at most two altogether. Every remaining odd power \(j\ge3\)
contributes at most \(mha/p^j\), and hence all of them together
contribute at most \(2mha/p^3\). These bounds apply to each
translate with the same constants. Thus the absolute difference
in the lemma is bounded by
\[
       m k_p(\sqrt b+1/p)+2mh b/(nX')+2+2mha/p^3.
\]
After division by \(Q\delta=hmb\delta\) and summation with
\(\log p\), the first three terms are \(O_{\psi,\eps,A}(x^{-A})\)
for every fixed \(A>0\). This uses \(\delta\ge x^{-\eps}\) and
the superpolynomial growth of \(b\) and \(X'\). For the last
term use \(a/b\le2x\) and
\(\sum_{n>\sqrt x}n^{-3}\log n\ll x^{-1}\log(2x)\), by
integration. This proves the bound.
\end{proof}

\begin{theorem}[The excess remaining after the rational division]
\leavevmode\label{thm:residualvaluation673}
The number \(N^\circ(x)=\prod_p p^{d a_p^\circ}\) is a positive
integer. Uniformly over the chosen words, translates and subsets
of entries read at individual primes,
\[
 \begin{aligned}
 \frac{\log N^\circ(x)}{dQ\delta}
  ={}&\sum_p c_\psi(p)(x/p-1)\log p\\
     &+\sum_p\min\{c_\psi(p),e_p/(Q\delta)\}\log p
                 +O_{\psi,\eps}(x^\eps\log(2x)).
 \end{aligned}
\]
The same formula holds if every word at a given translate is
divided by the common integer \(B_u=\gcd_w b^{\rm per}_{u,w}\).
In particular either normalization gives
\[
 \frac{4\log N^\circ(x)}{dQ\delta}-M_0(x)
   =\Gg_\psi(x)+4\sum_p\min\{c_\psi(p),e_p/(Q\delta)\}\log p
                     +O_{\psi,\eps}(x^{1/2+\eps}),
\]
using the integer for the chosen normalization on the left.
\end{theorem}
\begin{proof}
Removing the word divisor subtracts exactly
\(g_{u,w,p}-\lceil U_p(u)\rceil\) from the normalized ideal
valuation. Thus
\(t^\circ_{u,c,p}=\lceil U_p(u)\rceil+e_{u,c,p}\), where
\(e_{u,c,p}\ge0\). This also holds for a zero column, in the
extended sense. The finite norm valuations have denominator
dividing \(d\), so that capping at \(B_p\) and taking the
maximum with zero gives a nonnegative exponent in
\(d^{-1}\mathbb Z\).

Put \(C_p=Q\delta c_\psi(p)\). The unrounded upper cutoff is
\(mk_pa/p\); since \(a=xb+O(1)\),
\[
          |B_p-(A_p+C_p)|\le 1+mk_p/p .
\]
After normalization and summation the right side is
\(O_{\psi,\eps,A}(x^{-A})\) for any fixed \(A\).
For a nonempty selection with finite minimum, replace
\(\lceil U_p(u)\rceil-G_p\) by \(A_p\) before taking the minimum.
The change is at most the maximum absolute error in the lemma.
Taking the positive part and capping do not enlarge this
error. Replacing the cap by \(A_p+C_p\) contributes at most the
cutoff error above. As \(A_p,e_p\ge0\), the resulting exponent
is \(A_p+\min(C_p,e_p)\). If the list is empty or only zero
columns are read, both prescriptions take their upper cutoffs
and only the cutoff error occurs. Sum over the primes and
apply the lemma.

By Theorem~\ref{thm:wordgcd510} every quotient
\(b^{\rm per}_{u,w}/B_u\) divides the same \(D_x\), for all
\(u,w\), with \(\log D_x\ll x^C X'\). Thus the valuation using
\(B_u\) lies between \(t^\circ_{u,c,p}\) and
\(t^\circ_{u,c,p}+v_p(D_x)\). This bound persists under the
prime-dependent minima and cutoffs, including when the list is
empty. The increase of the normalized logarithm is bounded by
\(\log D_x/(Q\delta)=O_{\psi,\eps,A}(x^{-A})\). This proves the
first formula for both choices. Proposition~\ref{prop:primerange670}
gives the second.
\end{proof}

The nonnegative excess in the last display yields the same
sufficient three upper tests as Corollary~\ref{cor:intervalintegers668}
after replacing its integer and center by \(N^\circ,M_0\),
and using its rounding procedure with the cutoffs \(B_p\) on
\(\mathcal I_x\). Indeed the displayed error integrates to at
most the error allowed there. Even an exact evaluation \(e_p=0\)
at every prime would recover the centered prime sum to this
precision. An upper bound on that sum, or independently on the
centered logarithm of the divided columns, is still required.

\section{Gram operators in the first theta lifts}\label{nav:theta-gram}

\subsection{The Petersson pairing as a global Bessel period}
\label{sec:globalbessel707}

We identify the linear functionals measured by
Theorem~\ref{thm:translates706}. Write \(S=U\rtimes H\) for
the same Bessel group, \([S]=S(F)\backslash S(\mathbb A_F)\),
\(J=K(M)\), \(J_S=J\cap S(\mathbb A_{F,f})\), and
\(j_S=\operatorname{vol}_{S_f}(J_S)\). Give \(S_\infty\)
its previous Haar measure, with the compact \(H_\infty\)
of mass one. Use the product of this and the finite Haar
measure on \([S]\), with counting measure on \(S(F)\).
Write \(g_0\in G_\infty\) for the tuple of stationary real
points. At each real place put \(c_{\lambda_v}(h)=
\langle w_v,\pi_v(h)w_v\rangle\), where here \(w_v\) denotes
the unit coefficient vector. Put
\[
 c_\infty(h)=\prod_v c_{\lambda_v}(h_v),\qquad
 C(h)=d_\pi c_\infty(h),\qquad
 \lambda_*=\prod_v(d_v d_{\pi_v}/m_v)^{1/2}.
\]
Here \(m_v>0\) is the full real integral in
Theorem~\ref{thm:fullrealperiod698}, and \(d_\pi=\prod_v d_{\pi_v}\).
Let \(v_*\) be the tensor product of the real value vectors
used in Proposition~\ref{prop:averagedschur684}, and write
\(\sigma=\sigma_\infty\) for the unitary Bessel action on
their value space.

\begin{lemma}[The compact Bessel quotient at principal level]\leavevmode
\label{lem:besselquot707}
The map \((s,c)\mapsto(s,h_c)\) induces a disjoint union
\[
              [S]/J_S=\coprod_{c\in\mathcal C}
                                \Gamma_c\backslash S_\infty .
\]
Its integration rule is \(j_S\) times the sum of the real
quotient integrals. In particular it is a compact quotient.
For \(x=(x_c)_c\in\mathcal H_X\) the prescription
\[
          W_x(s,h_c)=\langle x_c,\sigma(s)v_*\rangle
\]
defines a continuous function on this union, conjugate
linear in \(x\).
\end{lemma}
\begin{proof}
Given a finite point \(u h\in U_fH_f\), write
\(h=\eta h_c k_H\) by the definition of \(\mathcal C\).
Density of \(U(F)\) in \(U_f\), proved in
Lemma~\ref{lem:residuelifts701}, permits a rational unipotent
\(u_0\) with \(h^{-1}u_0^{-1}uh\in J\). Thus \(u h\) is in
\(u_0\eta h_cJ_S\). If two representatives are equivalent,
their rational middle component equates their \(H\)-classes,
so the classes coincide. Its stabilizer at that class is
trivial: it preserves a positive definite integral lattice
and is principal at an odd split mark, as in
Proposition~\ref{prop:rowmean690}. The remaining rational
stabilizer is exactly \(\Gamma_c=U(F)\cap h_cJh_c^{-1}\).
Take a measurable real fundamental domain for each such
stabilizer. Its product with \(h_cJ_S\), over the classes,
is a fundamental domain upstairs, up to null boundaries.
Product Haar gives the asserted integration rule. The real
unipotent quotient is compact: the upper triangular central
filtration has successive additive lattices of full rank,
as in the covolume calculation of Lemma~\ref{lem:unipotentvolume703}.
One can choose compact parallelepipeds in these quotients
and lift and multiply them in the order of the filtration.
Since \(H_\infty\) is compact and the class set finite, the
whole quotient is compact. Finally \(\sigma\) is trivial
on \(\Gamma_c\) by the conductor identity used in
Theorem~\ref{thm:besselmatrix701}. The displayed matrix
coefficient consequently descends and is continuous there.
\end{proof}

\begin{theorem}[The period represented by the full series]\leavevmode
\label{thm:besselperiod707}
Let \(Y\) be a bounded continuous \(J\)-invariant function in
\(L^2([G])\) reproduced by the archimedean convolution
\(Y(g)=\int_{G_\infty}C(h)Y(g(h,1))\,dh\). For all \(x\)
in the source of the projected Poincar\'e map,
\[
 \langle\mathcal P x,Y\rangle
   =\kappa_*\int_{[S]} W_x(r)\,Y(r(g_0,1))\,dr,\qquad
 \kappa_*=\sqrt{j_M/v_U}\,\frac{\lambda_*}{j_Sd_\pi}>0 .
\]
The same \(g_0,\kappa_*\) work for all finite translates of
the test \(Y\) and for \(T(e_0e_1)Y\).
\end{theorem}
\begin{proof}
Use the convention that the first variable of an inner
product is conjugate linear. Expanding the real average of
Proposition~\ref{prop:averagedschur684}, changing \(s\) to
\(s^{-1}\), and taking tensor products gives
\[
  \langle x_c,F(y)\rangle
    =\lambda_*\int_{S_\infty}
           c_\infty(g_0^{-1}s^{-1}y)W_x(s,h_c)\,ds .
\]
Inversion preserves Haar on \(S_\infty\), since its unipotent
radical is unimodular and its compact factor acts with
modulus one. The coefficient averages converge absolutely
place by place by that proposition. This also verifies the
order of the factors in \(g_0^{-1}s^{-1}y\).

Unfold one seed in the inner product of
Theorem~\ref{thm:poincaregram704}. Its absolute periodization
is in \(L^2([G])\), hence its product with \(|Y|\) is
integrable. Its finite support contributes \(j_M\). Summing
the coordinates leaves
\[
 \langle\mathcal P x,Y\rangle
   =\sqrt{j_M/v_U}\sum_c\int_{\Gamma_c\backslash G_\infty}
              \langle x_c,F(y)\rangle Y(y,h_c)\,dy .
\]
Indeed \(Y(y,h_c)\) is left \(\Gamma_c\)-invariant by
rational invariance and the principal level. The class
weight in the seed is a partition for precisely this
quotient.

Insert the preceding integral over \(s\). In the double
integral take the quotient by diagonal left translation
of \(\Gamma_c\) on \(S_\infty\times G_\infty\). The product
\[
  W_x(s,h_c)c_\infty(g_0^{-1}s^{-1}y)Y(y,h_c)
\]
is invariant under this action. Either a fundamental set
in the \(y\) coordinate or one in the \(s\) coordinate gives
a fundamental set in the product. The equality of the
absolute integrals in the two choices follows by partition
into countably many translates and Tonelli. With the
fundamental set in the \(s\) coordinate the absolute integral
is at most
\[
   \|Y\|_\infty\|c_\infty\|_1
                 \int_{\Gamma_c\backslash S_\infty}|W_x(s,h_c)|\,ds
        <\infty .
\]
The coefficient is integrable by
Proposition~\ref{prop:averagedschur684}, and the last
integral is finite by Lemma~\ref{lem:besselquot707}. We
may now make the same exchange for the complex integral.
After substituting \(y=sg_0h\), the inner integral over
\(G_\infty\) equals \(d_\pi^{-1}Y(sg_0,h_c)\) by
reproduction. Apply the quotient integration rule to the
sum over \(c\) and divide by \(j_S\). This gives the formula.
Right translations at the marked primes commute with real
convolution and preserve level since \(J\) is normal in
\(K\); their finite linear combinations preserve the
stated class of tests.
\end{proof}

\begin{corollary}[The dimension of the projected period functionals]\leavevmode
\label{cor:periodrank707}
On the closed archimedean type space in \(L^2([G])^J\),
the functionals initially given on bounded continuous tests by
\[
 Y\ \longmapsto\ \kappa_*\int_{[S]} W_x(r)
                         (T(e_0e_1)Y)(r(g_0,1))\,dr
\]
extend continuously. Their span has dimension at least
\((1-4b^{-8})D_X\). Allowing also finite translates gives
dimension at least \((1-2b^{-4})^2 L D_X\), with
\(L\) as in Theorem~\ref{thm:translates706}.
\end{corollary}
\begin{proof}
Convolution by \(C\) is an orthogonal projection in the
right regular action. Indeed \(C(h^{-1})=\overline{C(h)}\);
Schur orthogonality applied to the same unit coefficient
twice gives \(C*C=C\) with the factor \(d_\pi\) used above.
Since \(C\in L^1\), both identities act on \(L^2\).
The range is the type space in question. Continuous,
compactly supported, \(J\)-invariant functions on \([G]\)
are dense in \(L^2([G])^J\), by density on the locally
compact quotient followed by compact averaging over \(J\).
Convolving them by \(C\) gives a dense set in its range.
These convolved tests are bounded and continuous by
dominated convergence and reproduce pointwise by Fubini.

On that dense set the theorem and self-adjointness identify
the functionals with \(Y\mapsto\langle\mathcal P_e x,Y\rangle\).
The representing vectors lie in the type space by
Theorem~\ref{thm:poincaregram704}, and yield the asserted
extensions there. For a finite translate of the functional
apply it to \(T(s)^*Y\); its representing vector is
\(T(s)\mathcal P_e x\). The two bounds now follow from
Theorems~\ref{thm:projectedrank705} and~\ref{thm:translates706}.
All constants and periods use the fields and nodes of those
theorems. Estimating the integral lattice of the composite
theta lift requires control of its arithmetic normalization
and of its values on these functionals.
\end{proof}


\subsection{The first theta lift of the projected Poincar\'e series}
\label{sec:firstlift708}

Keep \(G=U(Y_0)\), and put \(G_1=U(Y_1)\), where the dimensions
of the two Hermitian spaces are \(r\) and \(r+t\), \(t=5r\).
Use a tame twist of the shorter construction. Its additional
finite level places are disjoint from the two marked primes.
Let \(K_\Phi(g,h)\), for \(g\in G(\mathbb A_F)\) and
\(h\in G_1(\mathbb A_F)\), denote the theta series of the Weil
translate \(\omega(g,h)\Phi\). We use the bilinear convention
for the lift, integrating the input form times \(K_\Phi\).
Its datum on the first edge is the specified joint harmonic
times the Gaussian at infinity. It is the point datum of the
inclusion modulo the marked and twist places times the vacuum
at finite places. In particular in a rational Schr\"odinger
polarization it is Schwartz--Bruhat. The estimates below
also apply to arbitrary Schwartz--Bruhat data of this pair.

\begin{lemma}[Absolute growth of the first theta kernel]\leavevmode
\label{lem:thetagrowth708}
Put \(d=f r(r+t)\). In sign coordinates at the real places,
write \(H(y)=\prod_{v\mid\infty}\max(1,\|y_v\|_{\rm op})\).
For \(g_f\) and \(h\) in fixed compact sets and \(y\in G_\infty\),
the sum of the absolute values in \(K_\Phi((y,g_f),h)\) is
at most \(C_0 H(y)^{2d+2}\). Here \(C_0\) may depend on the
node, the datum and the compact sets. On the same family, at
every sufficiently large node,
\[
               \int_{G_\infty}|C(y)|H(y)^{2d+2}\,dy<\infty .
\]
\end{lemma}
\begin{proof}
The underlying symplectic space is the restriction of scalars
of the Hermitian Hom space with its alternating trace pairing.
A rational Lagrangian has real dimension \(d\) after restriction
to \(\mathbb Q\). Choose rational coordinates on it. The
rational indices selected by the compact support of the
finite data are contained in a fractional lattice in
\(\mathbb R^d\). This can be done uniformly on the sets in the
lemma: a finite Schwartz--Bruhat vector in the local Weil
representation is fixed by an open subgroup, hence has only
finitely many translates under a compact set. The same holds
at the finite adeles by the restricted tensor construction.
Finitely many compact finite supports lie in a common
multiple of the standard compact lattice after clearing
denominators. Their sup norms are bounded as well.

At infinity let \(Q_i\) and \(P_i\) be coordinate multiplication
and \((2\pi i)^{-1}\) times coordinate differentiation on
\(\mathbb R^d\), respectively. Under conjugation by the
unitary Weil operator these \(2d\) infinitesimal Heisenberg
operators transform by the underlying symplectic matrix or
its inverse. This follows by differentiating the covariance
of translations and modulations defining the Weil
representation. In the unitary dual pair the symplectic
action on the Hom space is \(A\mapsto hAy^{-1}\). In sign
coordinates \(\|y_v^{-1}\|_{\rm op}=\|y_v\|_{\rm op}\);
the matrices of \(h_\infty\) and their inverses are bounded
on the chosen compact set. Fixed changes to real symplectic
coordinates change the norm by a constant. Thus every
coefficient in the transformation is at most a constant
times \(H(y)\). The determinant splitting characters have
absolute value one and do not change conjugation.

For any word of length \(j\le 2d+2\) in \(Q_i,P_i\),
its \(L^2\)-norm on the transformed datum is at most
\(C_1 H(y)^j\). Indeed move the unitary operator to the
left, expand each factor into a linear combination, and
bound the finitely many fixed word norms on the original
Schwartz vector. Apply the Euclidean Sobolev inequality
with \(d+1\) derivatives to the transformed function and
to each of its products by a coordinate to power \(d+1\).
Every derivative involved is a word covered by this estimate.
It follows that the value at \(u\in\mathbb R^d\) is bounded by
\(C_2 H(y)^{2d+2}(1+\|u\|)^{-d-1}\). The sum of these
last radial weights over a fixed fractional lattice converges,
as its ball of radius \(R\ge1\) has \(O(R^d)\) points.
Finite sums of elementary data follow by addition.

Finally \(f\ll_{q_0,d_K}x^{11}\log(2x)\), where \(d_K\)
denotes the fixed quadratic-field parameter, and
\(x^5\le r\le3x^5\). Hence \(2d+2+9r^2<r^8\) eventually.
Lemma~\ref{lem:coefficientdecay697}, with the unitary bound
on \(\|y_v\|\le2\), bounds
\(|c_{\lambda_v}(y_v)|\max(1,\|y_v\|)^{2d+2}\)
by a node-dependent constant times
\((2/\|y_v\|)^{9r^2}\). This last majorant is integrable
by Lemma~\ref{lem:cosetpacking699}. Take products over the
real places and multiply by the finite scalar \(d_\pi\).
\end{proof}

Let \(Q_J\) denote right averaging over \(J\) with probability
Haar. For a function of the first variable put
\[
   (\mathcal C B)(g)=\int_{G_\infty} C(y) B(g(y,1))\,dy,
   \qquad B_h(g)=\overline{K_\Phi(g,h)} .
\]
The averaging \(Q_J\) permits twists ramified away from
\(M\), where the unipotent series still has its original
level. Let \(I_0\) be the finite index of the principal level
at the marks and the twist places in the standard finite
compact of \(Y_0\); it is the normalization of the first
edge. Write
\(\Theta_\Phi u(h)= I_0\int_{[G]}u(g)K_\Phi(g,h)\,dg\)
whenever this integral converges.

\begin{theorem}[An integral formula for the first projected lift]\leavevmode
\label{thm:firstlift708}
For the prescribed first-edge data, every \(h\) and every
\(x\in\mathcal H_X\), the integral defining
\(\Theta_\Phi(\mathcal P_e x)(h)\) converges absolutely and
\[
  \overline{\Theta_\Phi(\mathcal P_e x)(h)}
       = I_0\kappa_* \int_{[S]} W_x(r_0)
           \big(\mathcal C\,T(e_0e_1)Q_J B_h\big)
                           (r_0(g_0,1))\,dr_0 .
\]
The integral on the right, including the inner real
convolution, is absolutely convergent locally uniformly in
\(h\). The formula holds also for the unprojected map with the
finite projection omitted.
\end{theorem}
\begin{proof}
First use the unprojected series. We verify absolute
unfolding before taking the finite average. By the seed
formula of Theorem~\ref{thm:poincaregram704}, the integral
of its absolute periodization against \(|B_h|\) unfolds by
Tonelli. Finite integration is over \(h_cJ\). The average
of \(|B_h|\) over this coset, as a function of \(y\), is
left \(\Gamma_c\)-invariant, since \(B_h\) is automorphic
in its first variable and \(h_c^{-1}\Gamma_ch_c\subset J\)
at finite places. Majorize
\(\|F(y)\|\) using the absolute value of the integral for
\(F\) in Theorem~\ref{thm:besselperiod707}. Moving the
fundamental domain from \(y\) to \(s\) as in that theorem
now bounds the absolute integral by a finite sum of
constants times
\[
   \int_{\Gamma_c\backslash S_\infty}\int_{G_\infty}
      |c_\infty(y)|\, (Q_J|B_h|)(sg_0y,h_c)\,dy\,ds .
\]
The majorant here can be integrated on a compact set of
representatives for the outer quotient. For those \(s\),
\(H(sg_0y)\ll H(y)\); all finite arguments remain in the
compact coset just used. Lemma~\ref{lem:thetagrowth708}
gives integrability with a uniform bound for \(h\) in compact
sets. The same reasoning holds after any of the finitely many
right translations at \(M\).

For the complex integral, right \(J\)-invariance of the
series replaces \(B_h\) by \(Q_J B_h\). To move the finite
projection from the first to the second factor in the
Petersson pairing, change variables separately for each
finite translate and use self-adjointness of the group
algebra element. The preceding estimates justify these
changes even when \(B_h\) itself is outside \(L^2([G])\).
Repeat the complex unfolding of
Theorem~\ref{thm:besselperiod707} through the substitution
\(y'=sg_0y\). The resulting inner integral is
\(d_\pi^{-1}(\mathcal C T(e_0e_1)Q_J B_h)(sg_0,h_c)\)
by the definition of \(\mathcal C\). The quotient rule and
normalizations in that theorem give the displayed identity.
All estimates are uniform in \(h\) on compact sets as
asserted.
\end{proof}

\begin{corollary}[An integer acting on the first theta values]\leavevmode
\label{cor:firstprojection708}
Write \(|\mathcal G|e_0e_1=\sum_{s\in\mathcal G}n_s s\).
At the marked split components on \(Y_1\) choose
\(k_{1,s}\) reducing to \(\operatorname{diag}(s,1_t)\);
use the identity outside the marks. For the first-edge
datum and the same twists and nodes one has
\[
 |\mathcal G|\,\Theta_\Phi(\mathcal P_e x)(h)
     =\sum_s n_s\,\Theta_\Phi(\mathcal P x)(h k_{1,s}),
 \qquad n_s\in\mathbb Z,\qquad
               \sum_s|n_s|\le|\mathcal G|^{3/2}.
\]
Consequently any additive coefficient lattice containing the
listed unprojected translates contains the scaled projected
value. The cost \(\log|\mathcal G|\) of this integer is
\(\ll_{q_0,d_K} x^{21}\log(2x)\log b\) at the common nodes.
\end{corollary}
\begin{proof}
Let \(k_{0,s}\) lift \(s\) on \(Y_0\). At each marked
split component the point datum is the indicator, in
the finite Schr\"odinger matrix coordinates, of reduction
to the inclusion \(\iota:y\mapsto (y,0)\). Acting by
\((k_{0,s},k_{1,s})\) makes the substitution
\(A\mapsto k_{1,s}^{-1}Ak_{0,s}\), which preserves this
condition because \(\operatorname{diag}(s,1_t)\iota=\iota s\).
For completeness, this substitution on a split polarization
is the action of the linear change of variables preserving
the position subspace. In the Schr\"odinger model its
unitary normalization is the square root of the absolute
determinant of the substitution, equal to one on units.
The dual component acts by the inverse transpose. The tame
determinant character of the pair on the integral elements
is the character of \(\det(k_{1,s})/\det(k_{0,s})\) on
reduction: the twist uses the same tame character on both
groups, with inverse on the source. The rank-dependent
unramified factors are trivial on units. The two residue
determinants agree, so this scalar is one. At the twist
places the pair is the identity. The full datum is therefore
fixed. The resulting kernel identity reads
\[
 K_\Phi(g,h k_{1,s})=K_\Phi(g k_{0,s}^{-1},h).
\]
After the change \(g\mapsto g k_{0,s}\) in the lift integral
it gives \(\Theta_\Phi(T(s)\mathcal P x)(h)
=\Theta_\Phi(\mathcal P x)(h k_{1,s})\).
Absolute convergence holds by the theorem and its proof.
The integer coefficients follow from
Theorem~\ref{thm:definitelevels678}, which proves the
group-algebra assertion independently of the space on
which the idempotent acts. The sum of the squares of the
coefficients of \(e=e_0e_1\) equals its identity coefficient,
at most one: take that coefficient in \(e^*e=e\) and use the
normalized regular trace. Cauchy--Schwarz gives the asserted
sum of absolute values after multiplication by \(|\mathcal G|\).
The size estimate for \(|\mathcal G|\) is in the same theorem.
This controls the denominator introduced by this finite
projection. The arithmetic size of the unprojected values
and their behavior under the remaining lifts are additional
parts of the integral comparison.
\end{proof}


\subsection{The finite Gram operators in the first and auxiliary lifts}
\label{sec:finitegram709}

For the first theta lift and for each of the \(m\) auxiliary
single lifts in the shorter construction, the source has
rank \(r\) and the target has rank \(n'=r+t=6r\). We compute
the finite-place operator formed with the matrix coefficient
of their theta datum. This operator occurs in the doubling
integral. At each finite place \(w\) of \(F\) give the
stabilizer \(K_w\) of the standard source lattice mass one.
Give the integral matrix lattice in a split Schr\"odinger
model mass one as well.

\begin{lemma}[The coefficient at the finite level places]\leavevmode
\label{lem:sharpplaces709}
At a marked or twist place write \(q\) for the residue
cardinality and \(J_w=\ker(K_w\to\mathrm{GL}_r(\mathbb F_q))\).
For the inclusion point datum of the first edge, the source
matrix coefficient is \(q^{-n'r}{\bf1}_{J_w}\). For an
auxiliary lift it is the same at the marked places and is
\(v_r(q){\bf1}_{K_w}\) at a place of its tame modulus,
where \(v_r(q)=\prod_{i=1}^r(1-q^{-i})\).
\end{lemma}
\begin{proof}
These are split places. In the chosen component write \(A\)
for an \(n'\)-by-\(r\) matrix and \(A_0\) for its \(r\) rows
in the specified inclusion. All data in question are zero
outside integral matrices. The point datum imposes
\(A\equiv\iota\); the block datum at a twist place has
value \(\chi(\det\overline{A_0})\) if \(\overline{A_0}\)
is invertible and zero otherwise. Here \(\chi\) is the local
tame character with the source normalization of the lift.
Thus if both \(A\) and \(Ag\) are in the support, both
\(A_0\) and \(A_0g\) are integral with unit determinant.
Their quotient gives \(g\in K_w\). For the point datum it
also gives \(g\equiv1\), and every such \(g\) fixes the
function. Its squared norm is the measure of one residue
class, \(q^{-n'r}\). For the block datum, substituting
\(Ag\) multiplies the function by \(\chi(\det\bar g)\).
The inverse determinant character on the source in the
Weil action cancels that scalar. The remaining squared
norm is the fraction of invertible top blocks, namely
\(|\mathrm{GL}_r(\mathbb F_q)|/q^{r^2}=v_r(q)\).
This proves the assertions for the actual matrix
coefficients, with zero at all other group elements.
\end{proof}

Away from these places let \(L_w=(\mathcal O_E\otimes_{\mathcal
O_F}\mathcal O_{F_w})^r\) in the finite lattice frame.
Use the local vacuum of squared norm one. Write \(a_w(g)\)
for its source matrix coefficient, including the unitary
splitting character. In the standard unramified normalization
it equals one on \(K_w\) and is bi-invariant under \(K_w\).

\begin{lemma}[The mass away from the integral stabilizer]\leavevmode
\label{lem:vacuumtail709}
For every such place, of residue cardinality \(q\) in \(F\),
one has
\[
   |a_w(g)|\le [L_w:L_w\cap gL_w]^{-n'/2},\qquad
   \sum_{gK_w\ne K_w}|a_w(g)|\le128q^{-r-2}
\]
for \(r\ge2\).
\end{lemma}
\begin{proof}
Recall the lattice normalization of the vacuum. The Hom
lattice is self-dual for its Heisenberg commutator pairing,
with the additive character scaled by the different; the
target and source lattice frames are standard. On each
component its trace dual is the inverse different times the
Hermitian dual, so the scaling cancels that different.
The split or unramified quadratic lattice model suffices
at all these places. Indeed \(E\) contains both \(\mathbb Q(i)\)
and \(\mathbb Q(\sqrt{-d_K})\) with relatively prime
discriminants: \(d_K\equiv7\bmod8\). For any rational
prime one of the two fields is unramified there, and inertia
fixes that field. Complex conjugation acts nontrivially on
both, so cannot belong to inertia in the abelian field
\(E\). Thus \(E/E^+\) is unramified at finite places.

Here is the coefficient estimate in the self-dual lattice
model. The unit vacuum spans the eigenspace for the
integral Heisenberg lattice with its vacuum character. To
see that this space is a line, use the realization by
functions equivariant on the left by that lattice. Right
action by a lattice element on the coset of \(z\) introduces
its commutator with \(z\). Agreement with the character for
every lattice element forces \(z\) to lie in the dual lattice,
which is the lattice itself. Equivariance then determines the
function from its one value on the identity coset. The
standard compact splitting fixes this unit vector: integral
isometries preserve the lattice and act by substitution in
this model; the determinant twists away from the marked and
tame modulus places are unramified. The lattice character can
also be described at dyadic unramified places by using an
integral sesquilinear form whose antisymmetrization is the
given alternating form. Locally choose an integral element
of relative trace one, possible by surjectivity of the trace
in an unramified quadratic extension (lift a unit trace
from the residue field and divide by that unit).
Multiplying a purely imaginary trace coefficient by this
element gives the required sesquilinear coefficient. In this
model its bicharacter on integral vectors is trivial, and
the substitution by an integral isometry preserves it.

Project \(\omega(g)\) of the vacuum back onto the vacuum line
by averaging the lattice action divided by its character,
with probability Haar. The squared norm of this projection
is the integral of the corresponding lattice coefficients
of the translated vacuum. Such a coefficient vanishes outside
the intersection of the original and translated Hom lattices:
commuting the translating vector past an element of the
translated lattice multiplies it by a nontrivial phase if
that vector is outside its self-dual lattice. Inside the
intersection the absolute value is at most one. The integral
is at most the reciprocal intersection index. On Hom the
source acts by precomposition; in matrix coordinates there
are \(n'\) copies of the dual source lattice. The index on a
dual lattice exchanges the two indices of an intersection.
They are equal for \(g\) unitary, since its determinant has
absolute norm one on the whole local quadratic algebra.
Taking a square root gives the first inequality.

Over a discrete valuation ring of residue size \(Q\), the
number \(A_r(j)\) of sublattices of index \(Q^j\) in a free
rank \(r\) lattice satisfies
\[
 A_r(j)=\!\sum_{j_1+\cdots+j_r=j} Q^{\sum_{i=1}^r(i-1)j_i}
       \le 4Q^{(r-1)j}.
\]
Triangular normal form over the valuation ring gives the
formula, with the diagonal powers in the order for which
there are \(i-1\) free entries modulo the \(i\)-th power.
For the bound divide by \(Q^{(r-1)j}\) and sum each of the
first \(r-1\) indices freely. The resulting factor is at most
\(\prod_{l\ge1}(1-2^{-l})^{-1}<4\): keep the first three
factors and use \(\prod_{l\ge4}(1-2^{-l})\ge1-\sum_{l\ge4}2^{-l}\).
Duality gives the same count for overlattices of prescribed
index in the ambient vector space.

At a nonsplit place the coset is determined by \(gL_w\).
If the intersection with \(L_w\) has index \(Q^j\), it
has that same index in \(gL_w\). First choose the intersection,
then the overlattice. There are at most \(16Q^{2(r-1)j}\)
choices. The first inequality gives total mass at most
\(\sum_{j\ge1}16Q^{-(r+2)j}\), as \(n'=6r\), and \(Q\ge q\).
At a split place the coset is determined by a single
rank \(r\) lattice in the chosen component. Write \(q^u,q^v\)
for the two indices of its intersection with the standard
lattice. The other component is the inverse dual, so the
index on the full quadratic algebra is \(q^{u+v}\).
There are at most \(16q^{(r-1)(u+v)}\) possibilities for
this pair of indices. The sum away from the integral coset
is at most
\(\sum_{j\ge1}16(j+1)q^{-(2r+1)j}\).
Both geometric sums are at most \(128q^{-r-2}\) for \(r\ge2\).
\end{proof}

\begin{theorem}[The finite factor in the first lift and the auxiliary lifts]\leavevmode
\label{thm:finitegram709}
Let \(S_f\) here denote the set of marked and twist places.
For the relevant source representation by unitary right
translations \(T\), on vectors fixed by all \(K_w\) outside
\(S_f\), put
\[
  \mathcal R=\sum_{g\in G(\mathbb A_{F,f}^{S_f})/K^{S_f}}
           \left(\prod_{w\notin S_f}a_w(g_w)\right)T(g),
      \qquad K^{S_f}=\prod_{w\notin S_f}K_w .
\]
This series converges in operator norm on the indicated
space. It is self-adjoint and
\[
              \|\mathcal R-I\|_{\rm op}\le512f\,2^{-r}
                         =:\varepsilon_r
\]
at all sufficiently large nodes, for every choice of the
twist. Including the local integrals at the level places
multiplies it on level-fixed inputs by the following scalars.
For an auxiliary lift with tame set \(\mathcal Q\) the scalar is
\[
       c_{\rm aux}=\frac{M^{-fn'r}}{|\mathcal G|}
                         \prod_{w\in\mathcal Q}v_r(q_w).
\]
For the first edge it is
\[
       c_{\rm first}=
       \frac{M^{-fn'r}}{|\mathcal G|}
           \prod_{w\in\mathcal Q}
                \frac{q_w^{-n'r}}{|\mathrm{GL}_r(\mathbb F_{q_w})|}.
\]
These are the finite Gram factors prior to the scalar
normalizing the lift. A stable-level input on \(Y_0\) fixed
by the maximal compact at a twist place is in particular
fixed by its principal subgroup.
\end{theorem}
\begin{proof}
View each summand as an operator from the fixed space into
the ambient Hilbert space, of norm at most the absolute
coefficient. Its value is independent of the right coset
representative. Bi-invariance of the coefficient and left
reindexing by \(K^{S_f}\) show that the sum maps back into
the fixed space. Lemma~\ref{lem:vacuumtail709} bounds the
sum of the absolute coefficients other than the identity by
\[
       \prod_{w\notin S_f}(1+128q_w^{-r-2})-1
       \le \exp\!\left(128 f\sum_{n=2}^{\infty}n^{-r-2}\right)-1 .
\]
Indeed at most \(f\) places lie over any rational prime \(p\),
with \(q_w\ge p\). The first term plus an integral bounds the
last series by \(2^{-r-2}(1+2/(r+1))\). Now \(f2^{-r}\to0\).
Using \(e^z-1\le2z\) for sufficiently small \(z\ge0\)
gives the stated bound and absolute convergence. Equivalently
the sum is convolution by the bi-invariant integrable
coefficient, since every right coset has volume one.
Unitarity gives \(a_w(g^{-1})=\overline{a_w(g)}\), so inversion
in this unimodular group gives self-adjointness.

At each level place integrate the coefficient of
Lemma~\ref{lem:sharpplaces709}. On the indicated fixed vectors
the integral over \(J_w\) acts by its volume
\(|\mathrm{GL}_r(\mathbb F_{q_w})|^{-1}\); on \(K_w\)-fixed
vectors integration over \(K_w\) acts by one. There are
\(f\) marked places of each residue size \(p_i\). Multiplying
these factors proves both formulas.
\end{proof}

\begin{corollary}[Determinants of the normalized finite factors]\leavevmode
\label{cor:finitedet709}
On any nonzero finite-dimensional subspace of the indicated
fixed space, the compression of \(\mathcal R\) has positive
determinant, and the absolute logarithm of that determinant
divided by the subspace dimension is at most
\(2\varepsilon_r\). Summing these normalized logarithmic
errors for all \(m\) auxiliary lifts costs at most
\(1024mf\,2^{-r}\), uniformly in all their tame sets.
\end{corollary}
\begin{proof}
For large nodes \(\varepsilon_r<1/2\). Every Rayleigh quotient
of \(\mathcal R\) on a unit vector belongs to
\([1-\varepsilon_r,1+\varepsilon_r]\) by the theorem. This
also holds after orthogonal compression, so the absolute
logarithm of each eigenvalue is at most \(2\varepsilon_r\).
Sum over the eigenvalues and then over the lifts. This
estimate controls the finite convolutions after division by
their displayed scalars. The real local factor and the
integral lattice in the global lift enter separately.
\end{proof}


\subsection{The real Gram factor and the source of the bilinear lift}
\label{sec:realgram710}

The first lift pairs a form bilinearly with its theta
kernel. We identify the two types in this pairing. Inner
products are linear in the second variable. Write
\(c_\lambda(g)=\langle w,\pi_\lambda(g)w\rangle\) for the
coefficient of the unit vector on the prescribed highest line.

\begin{lemma}[The two conjugate spaces of series]\leavevmode
\label{lem:bilinearsource710}
For \(e=e_0e_1\) define a linear map on the conjugate
complex Hilbert space by
\[
      P^+:\overline{\mathcal H_X}\longrightarrow L^2([G])^J,
        \qquad P^+(\bar x)=\overline{T(e)\mathcal P x}.
\]
Its singular values are those of \(T(e)\mathcal P\).
Its values are cuspidal and fixed by convolution with
\(\prod_{v\mid\infty}d_{\pi_v}\overline{c_{\lambda_v}}\).
On their span and its span of finite translates,
\[
     \langle Z,T(g)Y\rangle
         =\left(\prod_{v\mid\infty}c_{\lambda_v}(g_v)\right)
                           \langle Z,Y\rangle \qquad(g\in G_\infty).
\]
\end{lemma}
\begin{proof}
Conjugation is antiunitary on the automorphic quotient and
commutes with pullback by right translations. It commutes with
the finite projection because \(e\) has rational coefficients
by Corollary~\ref{cor:firstprojection708}. It also preserves
vanishing of unipotent constant terms. Conjugate the
reproducing identity of Theorem~\ref{thm:poincaregram704} and
use Theorem~\ref{thm:projectedrank705}. The two conjugations
preserve the eigenvalues of the Gram operator with multiplicity.

At one place put \(b_\lambda=d_\pi\bar c_\lambda\). Then
\(b_\lambda*\delta_g*b_\lambda=c_\lambda(g)b_\lambda\).
Indeed the convolution at \(h\) equals
\[
 d_\pi^2\int\overline{c_\lambda(y)}
    \langle\pi_\lambda(h)w,\pi_\lambda(y)\pi_\lambda(g)w\rangle\,dy.
\]
The formula follows by Schur orthogonality. Each integral
converges by Cauchy--Schwarz. At these parameters \(b_\lambda\)
is integrable by Lemmas~\ref{lem:coefficientdecay697}
and~\ref{lem:cosetpacking699}. Its convolution acts by a
self-adjoint projection, by the formula at \(g=1\) and the
identity \(b_\lambda(h^{-1})=\overline{b_\lambda(h)}\).
Insert the projection on both sides of the matrix entry.
Take products over the places. Finite translations commute
with this real projection.
\end{proof}

At a real source \(U(p,q)\), put \(r=p+q\), \(n'=6r\),
\(l=\min(p,q)\), and \(K=U(p)\times U(q)\). Write \(z_{ai}\)
for the Fock coordinates, where \(a\) indexes the target
and \(i\) the source. The polynomial inner product is
\(\langle z^\nu,z^\xi\rangle={\bf1}_{\nu=\xi}\nu!\).
The Gaussian is included when passing to Weil vectors. We
use the following property of a joint harmonic on the
prescribed lines:
\[
       \partial_{ai}\partial_{aj}P=0
       \quad\text{if \(i,j\) have opposite source signs, for every \(a\)}.
       \tag{*}
\]
For products of flag minors, it suffices that all factors
using a given target index use source indices of one sign.
Each monomial then satisfies \((*)\). In the bilinear
convention the source parameter of the joint harmonic is
\(\eta=-\lambda\): the conjugate of the automorphic highest
line is highest in the opposite compact order and has all
torus weights negated. Thus \(c_\eta=\bar c_\lambda\).
We prove the following estimates for every nonzero \(P\)
with \((*)\) in that compact highest line. The estimates
retain its Fock norm, so they apply at any choice of scale.

\begin{lemma}[The oscillator coefficient on a positive Cartan subgroup]\leavevmode
\label{lem:oscillatorcartan710}
Let \(a(u)\) be the product of the \(l\) disjoint hyperbolic
boosts on pairs of opposite source signs, with parameters
\(u_1\ge\cdots\ge u_l\ge0\). Put
\[
 H_i=n'+\sum_{a=1}^{n'}(z_{ai_+}\partial_{ai_+}
                         +z_{ai_-}\partial_{ai_-}).
\]
If \(P_1,P_2\) satisfy \((*)\), then
\[
 \langle P_1,\omega(a(u))P_2\rangle
       =\left\langle P_1,\prod_{i=1}^l
                        (\cosh u_i)^{-H_i}P_2\right\rangle .
\]
In particular \(A_P(g)=\langle P,\omega(g)P\rangle\) is in
\(L^1(U(p,q))\). If \(z\in U(p,q)\) is positive Hermitian and
\(\log\|z\|_{\rm op}\le u\), then
\[
   \|P\|^2 e^{-(n'l+\deg P)u}\le A_P(z)\le\|P\|^2 .
\]
\end{lemma}
\begin{proof}
A compact source change of basis mixes coordinates only
within each sign in each target row, by a unitary linear
substitution or its conjugate. It preserves \((*)\) and
degree. One paired source boost acts infinitesimally by
\(E_i-F_i\), where
\[
 E_i=\sum_a \xi_a z_{ai_+}z_{ai_-},\qquad F_i=E_i^*,
                           \qquad |\xi_a|=1 .
\]
Indeed the source boost or its inverse mixes the two entries
of a row with coefficients \(\cosh u,\pm\sinh u\).
Conjugating the negative coordinate to obtain the Fock
polarization makes the mixed term interchange multiplication
and its adjoint, differentiation. The commutators of the
displayed quadratic operator with multiplication and
differentiation give these infinitesimal changes, with
phases according to the coordinate bases. Its scalar term
is zero: the symmetric ordering of the quadratic
multiplication and differentiation terms has no diagonal
commutator in two distinct coordinates. The determinant
splitting character is one along the boost.

Direct differentiation gives
\([F_i,E_i]=H_i,\ [H_i,E_i]=2E_i,\ [H_i,F_i]=-2F_i\).
For \(V\) killed by \(F_i\), the boost acts by
\[
             \exp((\tanh u)E_i)(\cosh u)^{-H_i}V .
\]
To verify it, differentiate this vector in \(u\) and use
\(F_i e^{tE_i}=e^{tE_i}(F_i+tH_i+t^2E_i)\). Its derivative
is \((E_i-F_i)\) times the vector and its initial value is
\(V\). All differentiations hold in Fock norm on compact
\(u\)-intervals: for two variables the squared norm series
of \(e^{t\xi zw}\) times a monomial has terms bounded by
\(|t|^{2j}\) times a polynomial in \(j\), and \(|t|<1\).
The same bound allows derivatives and finite tensor
products. Uniqueness of the differentiable orbit under the
unitary group proves the formula. Distinct boost pairs
commute. Property \((*)\) gives \(F_iP_1=F_iP_2=0\).
Move each nonconstant creation term onto \(P_1\) by its
adjoint to obtain the stated pairing.

For \(z=k a(u)k^{-1}\), apply the formula to \(k^{-1}P\).
The \(H_i\) are diagonal, and their sum on a monomial of
this vector is at most \(n'l+\deg P\). Since
\(\cosh u_i\le e^u\), this proves the two inequalities.
For \(k_1a(u)k_2\) the same identity with two compact
translates bounds \(|A_P|\) by
\(\|P\|^2\prod_i(\cosh u_i)^{-n'}\). The Cartan integration
formula for \(U(p,q)\) has radial density proportional to
\[
 \prod_{i<j}\bigl(\sinh(u_i-u_j)\sinh(u_i+u_j)\bigr)^2
          \prod_i(\sinh u_i)^{2|p-q|}\sinh(2u_i).
\]
The restricted roots \(e_i\pm e_j,e_i,2e_i\) have respective
real multiplicities \(2,2|p-q|,1\) from the complex
off-diagonal entries, unpaired complex coordinates and
imaginary diagonal entry. The density on the chamber is at
most a constant times \(\prod_i e^{2(r-1)u_i}\). Since
\(n'>2(r-1)\), the preceding bound is integrable.
If \(l=0\) the group is compact and the products are empty.
\end{proof}

\begin{theorem}[A quantitative real factor for the bilinear lift]\leavevmode
\label{thm:realgram710}
Let \(P\) be as above, and put \(\mu=\eta+\rho-2\rho_c\),
\(\alpha=\eta+\rho\) in the positive-root order of \(\eta\).
Let \(V_1\) be the Haar volume of \(\log\|g\|_{\rm op}\le1\)
in \(U(p,q)\), and put
\[
 \delta=(1+n'l+\deg P+r(\max\alpha-\min\alpha))^{-1}.
\]
The real local factor in the bilinear pairing satisfies
\[
 \mathfrak a(P,\eta):=\int_{U(p,q)}\overline{c_\eta(g)}
                   A_P(g)\,dg
 \ \ge\ \frac{e^{-1}\|P\|^2V_1}{\dim V_\mu}
                     (\delta/2)^{2pq}>0 .
\]
For fixed real data it is the same for every finite tame twist.
\end{theorem}
\begin{proof}
Give \(K\) probability Haar. In polar coordinates \(g=kz\),
write \(d\nu\) for the quotient measure in \(z>0\). It is
invariant under conjugation by \(K\). With \(M_\mu=1/\dim
V_\mu\), the formula of Proposition~\ref{prop:polarcoefficient697}
gives
\[
      \mathfrak a(P,\eta)=M_\mu\int_{z>0}
             \frac{\langle P,\omega(z)P\rangle}
                  {m_{-\alpha}(z)}\,d\nu(z).
\]
To see the identity, truncate at \(\log\|z\|\le R\) and use
the polar formula in the decreasing order of \(\eta\).
For its compact variable \(u\), orthogonality on the
irreducible compact copy generated by \(P\) gives
\[
 \int_K\overline{m_\mu(ku)}
       \langle P,\omega(k)X\rangle\,dk
                   =M_\mu\langle\omega(u)P,X\rangle .
\]
The component of \(X\) perpendicular to that compact copy
contributes zero. Take \(X=\omega(z)P\) and change \(z\)
to \(uzu^{-1}\). The denominator becomes \(m_{-\alpha}(z)\).
The remaining compact integral uses
\(\int_K m_\mu(u)\omega(u^{-1})P\,du=M_\mu P\), again by
orthogonality. This proves the truncated formula. Its left
side converges absolutely by Lemma~\ref{lem:oscillatorcartan710}
and \(|c_\eta|\le1\). Its right integrand is positive by
the same lemma and positivity of the principal minors of
\(z\). Monotone convergence gives the identity.

On \(\log\|z\|\le\delta\), one has \(\det z=1\), and every
principal minor of size \(h\) is at most \(e^{h\delta}\).
The weight-minor formula in Proposition~\ref{prop:polarcoefficient697}
bounds \(m_{-\alpha}(z)\) above by
\(\exp(r(\max\alpha-\min\alpha)\delta)\). Together with
Lemma~\ref{lem:oscillatorcartan710} this bounds the ratio in
the integral below by \(e^{-1}\|P\|^2\). The radial volume
of this neighborhood is at least \(V_1(\delta/2)^{2pq}\).
Indeed scale the Cartan coordinates in the unit chamber by
\(\delta\). The root values \(x\) there belong to \([0,2]\);
hence \(\sinh(\delta x)\ge\delta x\ge(\delta/2)\sinh x\).
The sum of the root multiplicities plus the \(l\) radial
coordinates is \(2pq\). The density in the lemma gives the
stated volume ratio, including the compact case. All terms
in the defining integral are real-place data; the finite
tame twist has real weight zero.
\end{proof}


\subsection{Petersson norms of the first and auxiliary theta lifts}
\label{sec:firstgram711}

Let \(S,T\) be one of the pairs of ranks \(r,6r\) in
Theorem~\ref{thm:finitegram709}. In this subsection \(S\)
denotes the source Hermitian space. Write \(G=U(S)\),
\(H=U(T)\) and \(K_\Phi(g,h)=\theta(\omega(g,h)\Phi)\).
The splitting characters and the global additive character
are those of the pair. Put
\[
 a(u)=\langle\Phi,\omega(u,1)\Phi\rangle,\qquad
 \Theta Y(h)=I\int_{[G]}Y(g)K_\Phi(g,h)\,dg .
\]
Here \(I\) is the scalar multiplying this lift, including
its level index when present. We retain the Haar scalar in
the doubling formula. This is needed when comparing lifts
on different Hermitian spaces.

\begin{proposition}[The doubled integral at the chosen split point]\leavevmode
\label{prop:doubledkernel711}
For the data with a full-column-rank support at a marked
split component, and the real polynomials of
Lemma~\ref{lem:oscillatorcartan710}, one has
\(a\in L^1(G(\mathbb A_F))\) and
\[
 \int_{[H]}\overline{K_\Phi(g',h)}K_\Phi(g,h)\,dh
       =c_{S,T}\sum_{\gamma\in G(F)}a(g'^{-1}\gamma g).
\]
The left integral and the series are absolutely convergent.
The positive constant \(c_{S,T}\) depends on measures and
the normalization of the oscillator inner product. For the
same pair and measures it is the same for all tame twists.
\end{proposition}
\begin{proof}
The oscillator representation for \(S\oplus S^-\), where
the second form is negated, is the tensor product of the
one for \(S\) and its complex conjugate. In position
coordinates this says that the alternating character on
the negative summand is the conjugate character. The two
global theta sums multiply by summing the two rational
coordinates independently. A rational change of polarization
preserves that identity by adelic Poisson summation with
self-dual measure. Thus the integrand is the doubled theta
series at the restriction of \((g,g')\) to the two summands.

We use the convergent unitary Siegel--Weil formula, in
\href{https://msp.org/pjm/2013/264-1/pjm-v264-n1-p11-s.pdf}
{S.~Yamana, \emph{The Siegel--Weil formula for unitary
groups}, Pacific J.~Math.~264 (2013), Theorem~2}.
The doubled source is split of Witt index \(r\), the target
has dimension \(6r\) and its rational Witt index is at
most \(3r\). Thus dimension minus Witt index exceeds \(r\).
The theta integral converges absolutely, and its defining
Siegel Eisenstein series is in its half-plane of absolute
convergence: the parameter is \((6r-r)/2>r/2\). The finite
data are Schwartz--Bruhat, and the data at the identity
at infinity are polynomials times Gaussians. The formula
applies as an identity of functions on the doubled group
and hence at \((g,g')\). In the split doubled space the
Siegel cosets index its rational maximal isotropic subspaces
\(L\), each once, by extending a basis of such a subspace
to a hyperbolic frame. Its summand is evaluation at zero
in the position model belonging to \(L\).

Here is the role of the marked point. If the projection
of \(L\) onto \(S^-\) has a kernel, choose \(0\ne v\in S(F)\)
with \((v,0)\in L\), where rational vectors are taken over
\(E\). At a marked split place take its nonzero component
in the position coordinates. The functional for \(L\) is
invariant under the Heisenberg operators in its momentum
Lagrangian: in its own model these multiply by characters
with value one at zero. In the split matrix realization on
the two factors this includes the multipliers
\[
                   (A,A')\longmapsto
                         \psi_w\bigl(\xi(Av)\bigr)
\]
for every covector \(\xi\) of the target component. Indeed
the corresponding momentum tensors are the products of the
chosen vector in \(L\) with covectors in the dual target
component. The evaluation of such a tensor on \((A,A')\)
is \(\xi(Av)\), with no term in \(A'\). A nonzero fixed
scalar in the trace pairing can be absorbed into \(\xi\).
Consequently this local functional is supported on
\(Av=0\). To see the support assertion directly, for every
point with \(Av\ne0\) choose \(\xi\) for which the multiplier
is different from one. It is constant on a neighborhood.
Invariance annihilates each test supported there; partition
a compact test support into finitely many such open pieces.
For the datum of Lemma~\ref{lem:sharpplaces709} the marked
matrix is injective throughout its support. This remains
true after precomposition by the invertible source element
\(g_w\). The local term and hence this Siegel summand vanish.

Every remaining \(L\) projects isomorphically onto \(S^-\).
Writing it as a graph, isotropy makes its other projection
an isometry because the two forms are opposites. Thus these
graphs are indexed by \(\gamma\in G(F)\). For the identity
graph, the change to diagonal and difference coordinates
followed by Fourier transformation in the diagonal
coordinate takes the tensor \(f(x)\overline{f'(y)}\), at
the new zero, to \(\int f(x)\overline{f'(x)}\,dx\). Indeed
the difference is set to zero, and the zero Fourier
coefficient integrates the remaining variable. The same
calculation with a rational isometry on one factor gives
\(\langle\omega(g',1)\Phi,\omega(\gamma g,1)\Phi\rangle\),
up to the common measure normalization. Rational determinants
have adelic absolute value one; the global splitting and
the twisting characters are trivial on rational elements.
The possible inverse in the choice of graph index just
reindexes the sum. With self-dual position measures the
identity-graph pairing is the positive \(L^2\) pairing.
Changing its normalization or the quotient Haar multiplies
the formula by a fixed positive scalar. The target twisting
character cancels between the two kernels, and on the source
any splitting change contributes its unitary character to
the coefficient \(a\). This proves the assertion about
\(c_{S,T}\).

Finally \(a\) is a product of the local coefficients.
Lemma~\ref{lem:oscillatorcartan710} gives absolute
integrability at infinity, Lemma~\ref{lem:sharpplaces709}
at the level places, and the product bound in
Theorem~\ref{thm:finitegram709} at the other places, including
the identity coset with mass one. These estimates prove
the stated \(L^1\) property.
\end{proof}

\begin{theorem}[A Petersson estimate for the actual theta lift]\leavevmode
\label{thm:globalgram711}
The lift initially defined on compactly supported \(L^2\)
functions on \([G]\) extends continuously to \(L^2([G])\).
For the real line of parameter \(\lambda\) and the fixed
spaces at the finite places in Theorem~\ref{thm:finitegram709},
put
\[
 B_\Phi=|I|^2c_{S,T}\,c_{\rm lev}
       \prod_{v\mid\infty}\mathfrak a(P_v,-\lambda_v),\qquad
  c_{\rm lev}=
   \begin{cases}c_{\rm first}&\text{on the first edge},\\
                 c_{\rm aux}&\text{on an auxiliary edge}.
   \end{cases}
\]
Use the nonzero matched harmonic \(P_v\) on each edge.
With \(\varepsilon_r=512f2^{-r}\), for \(Y,Z\) in these
source spaces the lift satisfies
\[
 \langle\Theta Z,\Theta Y\rangle
      =B_\Phi\langle Z,\mathcal RY\rangle,\qquad
 (1-\varepsilon_r) B_\Phi\|Y\|^2
     \le\|\Theta Y\|^2\le(1+\varepsilon_r)B_\Phi\|Y\|^2 .
\]
For \(Y=P^+\bar x\) and its indicated finite translates
this extension agrees almost everywhere with the original
pointwise bilinear theta integral.
\end{theorem}
\begin{proof}
On compact subsets of the source quotient the functions
\(g\mapsto K_\Phi(g,\cdot)\) have bounded \(L^2([H])\)
norm. Indeed its norm squared is the diagonal of
Proposition~\ref{prop:doubledkernel711}, locally bounded
by local uniform convergence of the Siegel Eisenstein
series in its open half-plane of absolute convergence.
One can lift each compact subset to a compact set of
representatives since the rational group is discrete.
The sections are measurable as Hilbert vectors, by joint
measurability and separability of \(L^2([H])\). Thus the
lift of a compactly supported \(L^2\) function exists
as a Hilbert-space integral, since such a function is
integrable on its support. For two such inputs take
their inner product and insert the proposition.

To unfold the resulting series absolutely, the integral
of the sum of absolute graph terms times \(|Z(g')Y(g)|\)
over the two source quotients equals
\[
 \int_{G(\mathbb A_F)} |a(u)|
      \int_{[G]} |Z(g')Y(g'u)|\,dg'\,du
       \le\|a\|_1\|Z\|_2\|Y\|_2 .
\]
This uses rational counting measure and right invariance
of quotient Haar. It holds by Tonelli before the bound.
Unfolding the complex integral now gives
\[
 \langle\Theta Z,\Theta Y\rangle
   =|I|^2c_{S,T}\int_{G(\mathbb A_F)}a(u)
          \langle Z,T(u)Y\rangle\,du .
\]
This proves boundedness and the formula on all of \(L^2\)
by density. For the conjugated projected series \(P^+\bar x\)
the proof of Theorem~\ref{thm:firstlift708} bounds the
integral of the absolute value against the kernel for each
\(h\). Its estimate uses absolute seeds and is unchanged
by conjugating the input. It applies to the required finite
translates by shifting the compact set of finite arguments.
Cut off on increasing compact sets exhausting the quotient.
Dominated convergence gives the original theta integral for
each fixed \(h\); \(L^2\) convergence and a subsequence
identify it almost everywhere with the bounded extension.

The real type projection on both inputs compresses
\(T(u_\infty)\) to \(\prod_v c_{\lambda_v}(u_v)\).
This is the Schur calculation in Lemma~\ref{lem:bilinearsource710};
equivalently it is the highest-line matrix entry on each
copy of the source discrete series. Finite translations
preserve this real type. Hence the real integral at \(v\)
is exactly \(\mathfrak a(P_v,-\lambda_v)>0\) by
Theorem~\ref{thm:realgram710}. Fubini applies by \(a\in L^1\).
Integrating over the remaining level and unmarked places
gives \(c_{\rm lev}\mathcal R\) by
Theorem~\ref{thm:finitegram709}. The operator norm estimate
there gives both inequalities.
\end{proof}

\begin{corollary}[First-lift rank and comparison across twists]\leavevmode
\label{cor:firstliftrank711}
At the common nodes
\[
   \operatorname{rank}(\Theta P^+)
          =\operatorname{rank}(P^+)\ge(1-4b^{-8})D_X .
\]
Its full span of lifted marked translates has dimension
at least \((1-2b^{-4})^2 L D_X\), with \(L\) as in
Theorem~\ref{thm:translates706}. For two tame twists
on the same edge with common real datum, choose any fixed
nonzero \(d\)-dimensional subspace \(U\) on which both
finite-level assertions apply. Write \(D_i(U)\) for
the determinant of the Gram matrix of lift \(i\) on an
orthonormal basis of \(U\). Then
\[
 \left|\frac1d\log\frac{D_1(U)}{D_2(U)}
       -\log\frac{|I_1|^2c_{{\rm lev},1}}
                      {|I_2|^2c_{{\rm lev},2}}\right|
 \le\log\frac{1+\varepsilon_r}{1-\varepsilon_r}.
\]
\end{corollary}
\begin{proof}
All factors of \(B_\Phi\) are positive for the nonzero
normalization \(I\) of these lifts. At large nodes
\(\varepsilon_r<1/2\), so the lift is injective on the
specified source, including its marked translates. Apply
Theorems~\ref{thm:projectedrank705} and~\ref{thm:translates706}
after the conjugation of Lemma~\ref{lem:bilinearsource710}.
For the last estimate every eigenvalue of the compressed
Gram matrix for lift \(i\) is between \(B_{\Phi_i}
(1-\varepsilon_r)\) and \(B_{\Phi_i}(1+\varepsilon_r)\).
Multiply over eigenvalues. The target Haar constant and
every real integral are identical for the two twists and
cancel. This is a comparison of Petersson determinants
on a common source. To use it in the integral comparison
of Section~\ref{sec:missing} one still needs arithmetic
control of the images in the target coefficient lattices.
\end{proof}

\section{Marked determinants and signed coefficient matrices}\label{nav:chapter-marked-period-determinants}

\subsection{One marked period determinant at all complex embeddings}
\label{sec:allplacesminor477}

Fix one of the fields \(E_x\), one translate and \(s\ge1\) signed
columns of common small type in Theorem~\ref{thm:commoncompactgram447}.
Use its projected tests and its mass metric, including the same
normalization of the columns on both sides. In particular one may
use the divided columns \(\widetilde Y_j/b_j^{\rm per}\) of
Theorem~\ref{thm:primitivegram463}. Suppose the projected columns
are independent. Then their Gram determinants are positive at
every complex embedding: the rational quadrature formula expresses
them as positive sums of squared absolute values of algebraic
minors, and a nonzero minor stays nonzero at every embedding.

Write \(Z(t_1,\ldots,t_s)\) for the determinant obtained by
pairing the columns with the projections of the indicated small
tests. Let \(\mathcal V\) be the finite list of row tests
\(t_{a,g}\) in the quadrature formula, with positive rational
weights \(q_t\). Put
\[
 q_*=\sum_{t\in\mathcal V}q_t
     =r_\mu\sum_a|\Gamma_a|^{-1},\qquad
 L_0=g_1g_2T,\qquad M=\max(1,\max_a|\Gamma_a|).
\]
Fix one coefficient place \(w_p\) above each of \(p=g_1,g_2,T\).
Denote by \(b_p\) the least valuation at that place on the image
of the integral test lattice under the projector, at exterior
degree \(s\). The \(b_p\) here are valuations of period ideals;
they may include the chosen column denominators. The row functional
of each \(t\in\mathcal V\) has norm at most \(\sqrt M\).

\begin{lemma}[Perturbations attaining all three finite minima]\leavevmode
\label{lem:markedperturb477}
There are tests \(u_1,\ldots,u_s\) of functional norm at most
\(3L_0\sqrt M\) at every complex embedding such that
\[
          v_{w_p} Z(u_1+L_0v_1,\ldots,u_s+L_0v_s)=b_p
                   \qquad(p=g_1,g_2,T)
\]
for any \(v_i\) integral in the small test lattices at all three
chosen places. In particular the assertion holds for integral
linear combinations of tests of \(\mathcal V\).
\end{lemma}
\begin{proof}
For each \(p\) choose \(s\) frame rows attaining \(b_p\), possible
by Theorem~\ref{thm:frameperiods446}. For the first row, the second
row, and so on, combine the three lifts by the Chinese remainder
coefficients between zero and \(L_0-1\) equal to one at their
own rational prime and zero at the others. These are exactly the
tests of Corollary~\ref{cor:threeminor447}, giving the displayed
norm bound and \(v_{w_p}Z(u_1,\ldots,u_s)=b_p\). On perturbing
the rows, each additional determinant term contains at least one
factor \(L_0\). Its remaining rows are all permissible at \(w_p\),
so before that factor its valuation is at least \(b_p\).
It therefore cannot change the least valuation of the original
term. This argument uses only \(v_{w_p}(L_0)>0\), including for
ramified coefficient places.
\end{proof}

\begin{theorem}[One algebraic period determinant at all embeddings]\leavevmode
\label{thm:allplacesminor477}
Let \(J\) be a number field of definition for the periods in these
tests and put \(d=[J:\QQ]\). There is one determinant \(z\in J\)
of the given projected columns such that \(v_{w_p}(z)=b_p\) at
the three chosen places and, for every \(\sigma:J\hookrightarrow\CC\),
\[
 \frac{s!}{(4q_*)^s}\det G_\sigma
   \ \le\ \frac{|\sigma z|^2}{L_0^{2s}}
   \ \le\ \big((3+2d^2)^2 M\big)^s\det G_\sigma .
\]
All row coefficients used in the construction are rational
integers applied to the preceding semilocally integral tests.
\end{theorem}
\begin{proof}
At each embedding select a tuple
\((t_1^\sigma,\ldots,t_s^\sigma)\) of quadrature tests maximizing
the absolute determinant. Its value \(B_\sigma>0\) satisfies
\[
             B_\sigma^2\ge \frac{s!}{q_*^s}\det G_\sigma .
\]
Indeed the Gram formula sums over unordered subsets of size
\(s\). The sum of the products of their weights is at most
\(q_*^s/s!\), by expanding the \(s\)-th power of the sum and
retaining only tuples with distinct indices. This bound needs
no lower estimate for an individual quadrature weight.

Choose common rows \(r_i\) inductively. Before step \(i\),
for each \(\sigma\) use the first \(i-1\) selected rows and
complete them with \(t_i^\sigma,\ldots,t_s^\sigma\). Their
determinant at \(\sigma\) is nonzero, initially by definition.
With the other rows fixed call the determinant as a functional
of row \(i\) by \(l_\sigma\), in this embedding. We choose
\[
      r_i=u_i+L_0\sum_{\tau:J\hookrightarrow\CC}
                           c_\tau t_i^\tau,\qquad
                   c_\tau\in\{0,1,\ldots,2d\}.
\]
Choose the \(c_\tau\) for this step independently and uniformly
for the counting argument. Fix \(\sigma\) and the choices other
than \(c_\sigma\). Consecutive values of \(l_\sigma(r_i)\) differ
by the nonzero number \(L_0 l_\sigma(t_i^\sigma)\). At most one
choice lies in the open complex disk of radius
\(L_0|l_\sigma(t_i^\sigma)|/2\) about zero, by the triangle
inequality. The fraction failing for any embedding is at most
\(d/(2d+1)<1\). Hence a single choice keeps the new determinant
at least \(L_0/2\) times the previous one in absolute value at
each embedding. Iteration gives
\(|\sigma Z(r_1,\ldots,r_s)|\ge (L_0/2)^sB_\sigma\), proving
the lower estimate.

The choices are integral at the marks, so the lemma gives all
three valuations. Each selected row has functional norm at most
\((3+2d^2)L_0\sqrt M\) by the triangle inequality. Apply the
wedge of the row functionals to the wedge of the projected
columns at \(\sigma\). Its norm is at most the product of the
functional norms, and the squared norm of the column wedge is
\(\det G_\sigma\). This proves the upper estimate for the very
same \(z=Z(r_1,\ldots,r_s)\).
\end{proof}

\begin{corollary}[A uniform comparison for divided signed columns]\leavevmode
\label{cor:markedheight477}
There is a finite family \(\mathcal Z\subset J\) of actual period
determinants all attaining the three \(b_p\), such that, with
\(m_\sigma=L_0^{-s}\max_{z\in\mathcal Z}|\sigma z|\),
\[
            \frac{s!}{q_*^s}\det G_\sigma
              \ \le\ m_\sigma^2
              \ \le\ (16M)^s\det G_\sigma .
\]
For the base small level on the polynomial-conductor fields in
Theorem~\ref{thm:primitivegram463}, take \(s=q_c\) and the divided
columns. With its \(K_u\), define
\[
       H_{\mathcal Z}=
         \frac{2}{s[J:\QQ]}\sum_\sigma\log m_\sigma+K_u .
\]
Then \(H_{\mathcal Z}=H(G^\circ)+O_{\psi,\varepsilon}(x^C)\)
for a fixed \(C\), uniformly on that interval. The error is
independent of the coefficient-field degree and is
\(o_{\psi,\varepsilon,A}(b/x^A)\) for every fixed \(A>0\).
\end{corollary}
\begin{proof}
Form the list of determinants with rows \(u_i+L_0\epsilon_i t_i\),
where \(t_i\in\mathcal V\) and \(\epsilon_i=\pm1\) vary independently.
The lemma supplies the finite assertions and every row has norm
at most \(4L_0\sqrt M\). At a fixed embedding choose the maximizing
tuple in the preceding proof. Average these determinants multiplied
by \(\prod_i\epsilon_i\) over all choices of the signs. Multilinearity
kills every term containing an unperturbed row and leaves precisely
\(L_0^s Z(t_1,\ldots,t_s)\). The maximum of the absolute values
is at least the absolute value of that average. The bound for
\(B_\sigma\) and the exterior norm upper bound give both inequalities.

After taking logarithms the upper error in \(H_{\mathcal Z}-H(G^\circ)\)
is at most \(\log(16M)\); the lower error is at worst
\(-\log^+ q_*\), where \(\log^+q=\max(0,\log q)\).
We verify their sizes. Corollary~\ref{cor:threeminor447} bounds
\(\log M\) by a fixed power of \(x\). The Weyl dimension formula
for \(r_\mu\) has at most \([E_x:\QQ]n^2\) factors, each at most
the largest width plus \(n\), where \(n\) is the small rank.
The widths are \(O(xb+x^{C_1})\) and the degrees and ranks are
polynomially bounded, as in Theorem~\ref{thm:commoncompactgram447}.
Thus \(\log r_\mu\ll x^{C_2}\).

The class mass is bounded above by the number of small classes.
Put \(N=[E_x:\QQ]n\) and take a cyclotomic conductor \(N_E\)
for \(E_x\) from Corollary~\ref{cor:shortcenter460}.
The full-level bound in the proof of Lemma~\ref{lem:mass32}
is \(\log h_{\max}\ll N^4\log(2NN_E)\), obtained by recording
the trace-lattice Gram matrix and the matrix of a primitive
field generator in a bounded integral basis. Passing to the
base residue markings adds at most \(L_0^{N^2}\) classes per
full-level class. Indeed at each of the marked rational primes
the kernel of full residue reduction on the underlying integral
lattice lies in its small level, and a residue matrix has at
most \(p^{N^2}\) possibilities. The small objects and their
row frames are those in Section~\ref{sec:markedline350}.
The logarithms of \(N_E,L_0\) have polynomial bounds on the
interval, by Corollaries~\ref{cor:shortcenter460}
and~\ref{cor:threeminor447}. Hence \(\log^+q_*\ll x^{C_3}\).
If additional row levels are imposed, the same exact inequalities
hold, with the logarithm of their further index added to this
class-count bound. Finally \(b=\exp(x^\varepsilon)\) dominates
every fixed power of \(x\). This proves the claimed error.
\end{proof}

Consequently \(H_{\mathcal Z}\) can replace \(H(G^\circ)\) in
Theorem~\ref{thm:primitivegram463} at its stated precision on the
base level. This compares the archimedean and marked finite tests
on the given signed periods. The absolute upper estimate for
their values in Section~\ref{sec:missing} and the required
independence for the chosen input columns remain to be established.


\subsection{Translation averages of the signed theta columns}
\label{sec:signedcovariance478}

Fix \(E_x\), a translate, and ordinary words in the directed tower of
Theorem~\ref{thm:common111} with a common small test space.
Work at one complex embedding. Write \(V\) for the dual of this
projected test space, with the Hermitian norm of
Theorem~\ref{thm:commoncompactgram447}. An initial table for a word
gives a vector in \(V\) by the signed path formula of
Theorem~\ref{thm:covectorpaths352}. We compute a family of functions
whose values at the origin are precisely these vectors. The other
values will locate the cancellation which is possible at the origin.

On edge \(j\) choose one reference theta vacuum \(\Theta_j\) on its
principally polarized trace lattice \(\Lambda_j\).
Use orthonormal matrix coordinates for its Hermitian metric at the
CM embeddings, so the metric weight on the edge is
\(\exp(-c_j|z|^2)\), \(c_j>0\). On a conjugated edge use the
conjugate complex space and coordinates throughout. The finite
expansion of Proposition~\ref{thm:classden312} expresses each
class section, in the common matrix coordinates, as a finite sum
of \(U_{a_j}\Theta_j\) times scalars. Here \(a_j\) is a torsion
class modulo \(\Lambda_j\), and \(U\) is the Hermitian translation
of Lemma~\ref{lem:mean74}. Indeed the lattice intertwiner is a
sum over the quotient by the intersection of two lattices, and a
finite position condition is a finite character sum. Products
of the resulting translations are translations times roots of
unity. Choose one lift for each residue \(a_j\), absorbing the
change of its scalar in the coefficient. Finitely many such
sections have the same multiplier on a sublattice
\(\Gamma_j=L_j\Lambda_j\). One may take \(L_j\) divisible by
all the translation orders, since the ratio of two period
multipliers on \(\Lambda_j\) is their commutator character.

\begin{lemma}[Inner products of translated derivatives]\leavevmode
\label{lem:translatedjets478}
On one edge let
\[
 f_{\alpha,a}(w)=\partial_z^\alpha
                   (U_w U_a\Theta)(0).
\]
With probability Haar on \(\mathbb C^d/\Gamma\) the functions,
interpreted with their common unitary multiplier in the variable
\(w\), satisfy
\[
 \int f_{\alpha,a}(w)\overline{f_{\beta,b}(w)}\,dw
  ={\bf1}_{\alpha=\beta}{\bf1}_{a=b}
              c^{|\alpha|}\alpha!\,\|\Theta\|^2 .
\]
The norm on the right uses probability Haar and the theta metric
on \(\mathbb C^d/\Lambda\).
\end{lemma}
\begin{proof}
For a holomorphic section \(H\), the derivative at zero is
\(e^{-c|w|^2/2}\prod_l(\partial_{w_l}-c\bar w_l)^{\alpha_l}H(w)\).
In the metric on the torus the adjoint of
\(\partial_{w_l}-c\bar w_l\) is \(-\partial_{\bar w_l}\).
The opposite-face boundary terms cancel by the multiplier law.
Commuting this lowering operator across the raising operators in
the second factor gives \(c\) for a matching coordinate.
If a lowering operator remains it kills the holomorphic section.
If only raising operators remain, move one in the other direction
and use holomorphy of the first section. The surviving factor is
therefore \({\bf1}_{\alpha=\beta}c^{|\alpha|}\alpha!\) times the
inner product of the two undifferentiated sections.

For \(a\ne b\) modulo \(\Lambda\), unimodularity of the polarization
gives a period \(\lambda\in\Lambda\) whose commutators with \(U_a\)
and \(U_b\) differ. Otherwise \(a-b\) would belong to the lattice
dual to \(\Lambda\) for the alternating form, which is \(\Lambda\).
The period translation acts unitarily on the common cover and has
the two sections as eigenvectors with distinct unitary eigenvalues.
Their inner product is zero. A translate has the same metric norm
as \(\Theta\) by translation invariance. Passing to a finite cover
preserves the probability integral of its squared magnitude.
\end{proof}

For column \(i\) introduce independent matrix variables \(D_j\),
one matrix at each CM embedding on each edge. Expand its prescribed
path formula with these rules. Keep the actual value of the
initial table at the composite covector of each path, including
the condition that its components be nonzero. At each edge make
the polynomial substitution \(X=D_j^{\mathsf t}Y_{[r]}\) after
the specified determinant twist, as in \eqref{eq:sub81}. Keep the
monomials in \(D_j\) unevaluated. Apply the terminal coefficient
extraction, the small test and its projection, and include the
class masses and all finite scalars. On expanding the theta
sections, group together paths with the same tuple \(a=(a_j)_j\)
of torsion residues. The coefficient on \(D^\alpha\) in that
group is a vector \(v_{i,\alpha,a}\in V\), since it is linear
in the small test. The coupling index on the last two maps is
summed once. These vectors require finite scalar and polynomial
operations. In particular, the signed path sum takes place
before any absolute values. Any column divisor or normalization
can be included in its coefficients.

Put \(Y_i(\boldsymbol w)=\sum_{\alpha,a}v_{i,\alpha,a}
\prod_j f_{\alpha_j,a_j}(w_j)\), summing the finite expansion.
Then \(Y_i(0)\) is the original projected column. Indeed replacing
each monomial in \(D_j\) by its undivided derivative is exactly
the class operator. This replacement commutes with polynomial
operations on the other edge variables and with linear terminal
restriction. It is the construction in
\srcloc{Formalism/LocalToGlobal/KernelLattices/ThDistC}{27}.
The variable \(\boldsymbol w\) translates the sections at the
chosen representatives. This definition also makes sense when
such a translation requires a change of adelic level to be
realized as an automorphic test.

\begin{theorem}[Covariance for the signed paths]\leavevmode
\label{thm:signedcovariance478}
Write \(h=\prod_j\|\Theta_j\|^2\) and
\(q_\alpha=\prod_j c_j^{|\alpha_j|}\alpha_j!\). For a pair of
columns in these common coordinates one has the operator identity
\[
  \int Y_i(\boldsymbol w)Y_l(\boldsymbol w)^*\,d\boldsymbol w
       =h\sum_{\alpha,a}q_\alpha\,
                          v_{i,\alpha,a}v_{l,\alpha,a}^* .
\]
The integral is over the product of the common covers. If \(i,l\)
belong to distinct binary words of Theorem~\ref{thm:common111}
at this translate, the integral vanishes. For \(i=l\) its range
is exactly \(\operatorname{span}\{v_{i,\alpha,a}\}\).
\end{theorem}
\begin{proof}
Expand the finite sums and use the lemma on each edge. This proves
the formula. On the diagonal its quadratic form is a positive sum
of squared absolute values. Its kernel is the orthogonal complement
of the indicated span.

To see what happens between words, fix a test at either end of
the operator identity, a torsion tuple \(a\), and every
differential monomial on the edges after the first. The remaining
polynomial in the first matrix is a sum of matrix coefficients of
one compact column type. Indeed in rank \(k\) the first substitution
evaluates an input flag polynomial of the first partition \(\nu\)
at \(D_1^{\mathsf t}Y_{[k]}\). Right multiplication \(D_1\mapsto
D_1g\) replaces its argument by \(g^{\mathsf t}D_1^{\mathsf t}Y_{[k]}\).
The translates of flag polynomials of column weight \(\nu\) under
this action have the Schur type indexed by \(\nu\), by the
highest-weight realization in Section~\ref{sec:chain81} (equivalently
by the Cauchy decomposition of matrix polynomials). The transpose
convention gives the same indexing at every word. Changing the
two matrix frames on this edge acts by left and right invertible
matrices and preserves its isotypic component. Contractions at
later vertices and the finite signed sums are linear in these
first-edge polynomials.

At a fixed CM embedding the first partition is obtained by
subtracting the half-sum of positive roots from the initial
parameters and adding the common input determinant shift. This
is the weight calculation in Lemma~\ref{lem:wordfirstrow450}.
The shifts on this edge and the translate are independent of the
bits. Formula \eqref{eq:letters111} shows that at a differing bit
the corresponding entries of the partitions differ by \(\Delta\).
They therefore index inequivalent representations of \(U(k)\).
The Fischer inner product on the first matrix is invariant under
unitary changes of its columns, by its Gaussian integral formula.
Its pairing on inequivalent types vanishes by Schur's lemma.
The factors \(c_1^{|\alpha_1|}\) merely rescale each homogeneous
degree; different degrees already pair to zero. Summing the
pairings over the fixed data gives the claimed zero operator.
\end{proof}

\begin{corollary}[Independence under separate torsion translations]\leavevmode
\label{cor:signedspan478}
For \(s\) prescribed columns let \(V_i=\operatorname{span}
\{v_{i,\alpha,a}\}\). There are torsion points \(\boldsymbol w_i\)
of the product of common covers for which
\(Y_1(\boldsymbol w_1),\ldots,Y_s(\boldsymbol w_s)\) are
independent if and only if
\[
      \dim\sum_{i\in I}V_i\ \ge |I|
        \qquad(I\subseteq\{1,\ldots,s\}).
\]
The orders in these covers may be powers of any one fixed prime.
\end{corollary}
\begin{proof}
The span of all values of \(Y_i\) is \(V_i\). Containment in
\(V_i\) follows from the expansion, and a functional annihilating
all the values lies in the kernel of the diagonal covariance.
The same reasoning and continuity apply to values on the dense
set of prime-power grids. A change of lift in the cover changes
a value by a common unit scalar and preserves its line.

Independent choices from the \(V_i\) require the displayed
inequalities. For sufficiency use induction on \(s\), the case
\(s=1\) being immediate. If equality holds for a nonempty proper
subset \(I\), choose its vectors by induction and divide the
ambient space by their span, which is \(\sum_{i\in I}V_i\).
The dimension inequalities for \(I\cup J\) give the required
inequalities for the remaining indices in the quotient. Lift
their independent choices. If every proper nonempty subset has
strict inequality, choose one nonzero vector in the last space
and quotient the ambient space by the line it spans. Each of the
other nonempty sums loses at most one dimension and induction
applies. Finally express each chosen vector as a finite linear
combination of torsion values. The nonzero exterior product
expands into exterior products of those values, one of which
must be nonzero.
\end{proof}

At the prescribed first and later tests the point in this family
is \(\boldsymbol w=0\) for every word. Orthogonality in the
integral above allows vanishing at that point. It determines the
signed coefficient spaces against which the fixed evaluation
must be tested.


\subsection{An upper bound at the prescribed theta tests}
\label{sec:fixedexterior479}

Retain the actual signed coefficient vectors of
Theorem~\ref{thm:signedcovariance478}. The arithmetic test evaluates
every translated column at the origin. To bound that evaluation,
group the terms by their torsion tuple before sampling. Each
group has a multiplier on the reference trace lattice itself.
Thus the length of the shortest dual vector needed in the estimate
comes from the original Hermitian lattice, even when the common
cover used to compare different multipliers is large.

Let \(d_j\) be the complex dimension of the matrix coordinates at
edge \(j\), including all CM embeddings, and put \(D=2\sum_jd_j\).
Write \(K\) for the largest total degree \(\sum_j|\alpha_j|\)
among monomials of the prescribed columns, taking zero for an
empty list of monomials. Write \(f_0=[E_x:K_d]\). Define
\[
 N=\max\left\{1,\left\lceil
     8\sqrt{\frac{\pi\,\operatorname{rd}(E_x)}{2f_0}
                                      (K+2D+2)}\right\rceil\right\},
 \qquad \beta=\tfrac32 N^D .
\]
For each column \(i\), let \(\mathcal A_i\) be the finite set
of tuples \(a\) in its expansion. Write \(Y_{i,a}\) for the
part of \(Y_i\) indexed by \(a\).

\begin{lemma}[Evaluation for a fixed tuple of torsion shifts]\leavevmode
\label{lem:singleshift479}
Let \(B:V\to W\) be a linear map to a finite-dimensional
Hermitian space. For each \(i,a\),
\[
 \|BY_{i,a}(0)\|^2\le
 \beta\int\|BY_{i,a}(\boldsymbol w)\|^2\,d\boldsymbol w
   =\beta h\sum_\alpha q_\alpha\|Bv_{i,\alpha,a}\|^2 .
\]
Here the integral uses probability Haar on the product of the
reference tori \(\mathbb C^{d_j}/\Lambda_j\).
\end{lemma}
\begin{proof}
In one coordinate of \(W\), remove the unitary metric factor
\(\exp(-\sum_j c_j|w_j|^2/2)\) from the expression for
\(BY_{i,a}(\boldsymbol w)\). The result is a sum of raising
derivatives of the fixed holomorphic product section
\(\bigotimes_j U_{a_j}\Theta_j\). All terms have this section's
multiplier. The commutator in Lemma~\ref{lem:translatedjets478}
implies that \(K+1\) antiholomorphic derivatives in any common
direction annihilate the sum. Indeed each such derivative either
removes one raising operator or acts on the holomorphic section
and gives zero.

Scale the edge coordinates by \(\sqrt{c_j}\) to make the common
metric parameter one. For the scaled trace lattice on each edge
the alternating form is
\(\pi^{-1}\operatorname{Im}\langle z,w\rangle\).
Unimodularity identifies its Euclidean dual with its multiple
by \(i/\pi\), up to a sign. The lower bound of
Theorem~\ref{thm:lattice55} for the squared length of a nonzero
trace vector in the scaled lattice is
\(2\pi f_0/\operatorname{rd}(E_x)\). Hence the real dual minimum
of the product has length at least
\(\sqrt{2f_0/(\pi\operatorname{rd}(E_x))}\). These are the
normalizations in the proof of Theorem~\ref{thm:level85}.
Changing a unitary multiplier by a character leaves this period
lattice and the adjoint and commutator identities intact.

Apply Lemma~\ref{lem:grid85} with this minimum, dimension \(D\),
and degree \(K\). The grid average of the squared magnitude is
at most \(3/2\) times its integral. The origin is one of the
\(N^D\) grid points, and every summand is nonnegative. This proves
the inequality in each coordinate of an orthonormal basis.
Summation gives the assertion for \(W\). Orthogonality of the
derivative indices at fixed \(a\) gives the stated integral.
\end{proof}

At an embedding \(\sigma\) use its own small-space metric,
finite coefficients \(v_{i,\alpha,a}^{(\sigma)}\), vacuum
norm product \(h_\sigma\), and differential factors
\(q_{\alpha,\sigma}\). To form these coefficients one transports
the original scalar coefficients together with the phases of the
translated germs. Theorem~\ref{thm:trans57} supplies a phase
common to every derivative order of each such germ. This also
allows the fixed marked product normalization of
Proposition~\ref{thm:classden312}. For
\(\boldsymbol a\in\prod_{i=1}^s\mathcal A_i\), put
\[
 E_{\boldsymbol a,\sigma}
 = h_\sigma^s\sum_{\alpha_1,\ldots,\alpha_s}
   \left(\prod_{i=1}^s q_{\alpha_i,\sigma}\right)
   \left\|v_{1,\alpha_1,a_1}^{(\sigma)}\wedge\cdots\wedge
           v_{s,\alpha_s,a_s}^{(\sigma)}\right\|^2,\qquad
 B_\sigma=\sum_{\boldsymbol a}\sqrt{E_{\boldsymbol a,\sigma}} .
\]
The sum over each derivative index can be taken over its finite
support. The factors inside the exterior product include the
finite masses, coupling signs, terminal restriction and finite
projection of the original data.

\begin{theorem}[A bound for the original exterior period]\leavevmode
\label{thm:fixedexterior479}
For \(s\ge1\), the Gram matrix \(G_\sigma(0)\) of the actual
columns at the prescribed tests satisfies
\[
                    \det G_\sigma(0)\ \le\ \beta^s B_\sigma^2 .
\]
The formula applies in every rank. The right side uses undifferentiated
reference vacuum norms and finite coefficient operations.
\end{theorem}
\begin{proof}
Fix \(\boldsymbol a\) and consider the exterior product of the
\(s\) corresponding column parts, with a separate variable
\(\boldsymbol w_i\) in each part. With all the other factors
fixed, exterior multiplication is a linear map of the remaining
vector into \(\bigwedge^s V\). Applying the lemma in that
variable bounds the squared norm at zero by \(\beta\) times
its integral. Repeating this argument and using Fubini bounds
the squared norm with every variable zero by \(\beta^s\) times
the product-torus integral. Expanding the exterior product by
multilinearity, two distinct lists of multi-indices are
orthogonal in at least one variable by
Lemma~\ref{lem:translatedjets478}. The diagonal terms give
precisely \(E_{\boldsymbol a,\sigma}\).

At the origin, multilinearity also expresses the actual exterior
product as the sum of these products over all \(\boldsymbol a\).
The triangle inequality bounds its norm by
\(\beta^{s/2}B_\sigma\). The squared exterior norm is the Gram
determinant, including when all maximal minors vanish.
\end{proof}

\begin{corollary}[The finite coefficients in the centered estimate]\leavevmode
\label{cor:fixedheight479}
Use the divided period columns of Theorem~\ref{thm:primitivegram463}
and include each division in \(v_{i,\alpha,a}^{(\sigma)}\).
On the polynomial-conductor interval one has
\(\log\beta=o_{\psi,\varepsilon,A}(b/x^A)\) for every fixed \(A>0\),
uniformly in the translate and word list. If these columns are
independent at the prescribed tests, their centered count satisfies
\[
 \Delta_u\le K_u+\frac{2}{s[J:\QQ]}\sum_\sigma\log B_\sigma
             -A_{\rm sig}(x)+W_{d_0}(u)-\rho_x
             +o_{\psi,\varepsilon,A}(b/x^A).
\]
\end{corollary}
\begin{proof}
The number of embeddings of \(E_x\), the ranks and the number
of edges have polynomial bounds in \(x\). On edge \(j\) the total
degree is the sum of the sizes of its partitions. Their largest
entries are \(O(xb+x^{C})\), by the construction of the word
weights and the rank shifts in Lemma~\ref{lem:wordfirstrow450}.
Thus \(D\) and \(\log(2+K)\) are polynomially bounded. The
conductor bound of Corollary~\ref{cor:shortcenter460} gives the
same for \(\log\operatorname{rd}(E_x)\), using the
conductor--discriminant formula. These bounds are uniform at the
conjugate trace tori and independent of \(s\). They give
\(\log\beta\le C'x^{C''}\), which is \(o(b/x^A)\) since
\(b=\exp(x^\varepsilon)\). Independence makes all \(B_\sigma\)
positive by the theorem. Substitute the logarithmic inequality
in the definition of \(H(G^\circ)\) and apply
Theorem~\ref{thm:primitivegram463}.
\end{proof}

This upper bound evaluates the prescribed tests themselves. It
leaves the sizes of the indicated signed finite coefficient sums
and vacuum norms to be compared with \(A_{\rm sig}(x)\) at the
precision required in Section~\ref{sec:missing}.


\subsection{The fixed projected period and its coefficient matrices}
\label{sec:projectedmatrices489}

We evaluate the archimedean contribution at the original small
tests on \(E_x\) in the product coordinates of
Lemma~\ref{lem:basis59}. The coefficient of the total derivative
degree will be explicit. The remaining expression involves positive
matrices built from signed polynomials in a finite-dimensional
space.

Fix an embedding \(\sigma\) of a common coefficient field \(J\),
and \(s\) projected columns in the same Hermitian space \(V\).
At edge \(j\) of column \(i\) let \(n_{ij}\) be the complex dimension
of its product coordinates, \(k_{ij}\) its differential degree,
and \(c_{ij,\sigma}=\pi/A_{ij,\sigma}\) the metric parameter of
its marked elliptic curve \(C_{ij}^\sigma\). Use conjugate coordinates
on conjugate edges. Put \(\alpha!=\prod_{j,h}\alpha_{jh}!\).
The finite expansion of the fixed column is
\[
 Y_i=\sum_{a\in\mathcal A_i}\ \sum_{\substack{\alpha\\
              |\alpha_j|=k_{ij}}}v_{i,a,\alpha}
       \prod_{j,h} J_{\alpha_{jh}}(U_{a_{jh}}\Theta_{t_{ij}}),
       \qquad J_m(F)=F^{(m)}(0).
\]
Indeed, on each edge expand each class section by the product formula
preceding Theorem~\ref{thm:classden312}. The linear change of matrix
coordinates in each term replaces the edge polynomial by a homogeneous
polynomial of the same degree \(k_{ij}\), the size of its partitions.
Apply the signed path formula of Theorem~\ref{thm:signedcovariance478}
with the same terminal coefficient and the same projection. Collect now
by the full tuple \(a\) of elliptic translations and by the derivative
indices \(\alpha\) in product coordinates. This gives the vectors
\(v_{i,a,\alpha}\) by polynomial operations, including the chain-rule
coefficients and all normalizing scalars. In the last two maps there is
one sum over their common intermediate class. The phases of the transported
sections are included before collecting coefficients. Thus the displayed
identity uses the germ of each prescribed section at its given point.
We suppress \(\sigma\) on the vectors and take any required rational
column division in them. Empty expansions give zero columns.

For every nonempty \(\mathcal A_i\) choose positive reals
\(\mu_i(a)\) with sum one. Form the positive semidefinite operator
\[
 B_i(\mu_i)=\sum_{a,\alpha}{\alpha!\over \mu_i(a)}
                    v_{i,a,\alpha}v_{i,a,\alpha}^*
 \quad\text{on }V .
\]
Here \(vv^*\) takes \(h\) to \(v\langle v,h\rangle\), with the inner
product linear in its second argument. For any positive semidefinite
operators \(B_1,\ldots,B_s\) define
\[
 D_s(B_1,\ldots,B_s)
     =[t_1\cdots t_s]\det\nolimits_V
                         \left(I+\sum_{i=1}^s t_iB_i\right).
\]
The square brackets take the coefficient of exactly the specified
monomial.

\begin{theorem}[An upper bound at the fixed projected tests]\leavevmode
\label{thm:projectedmatrices489}
Write \(G_\sigma=(\langle Y_i,Y_j\rangle)_{i,j}\). For each column set
\[
 T_{i,\sigma}=\prod_j
 \left[
   \left({2\exp(L_{C_{ij}^\sigma})\over
                   g_{ij,\sigma}(t_{ij}^\sigma)^2}\right)^{n_{ij}}
       c_{ij,\sigma}^{\,k_{ij}}H_{n_{ij}}(k_{ij})\right],
\]
where \(L_C\) is from Theorem~\ref{thm:greenenergy488} and \(H_n(k)\)
is from Lemma~\ref{lem:sphere60}. Then
\[
    Y_iY_i^*\ \le\ T_{i,\sigma}B_i(\mu_i),\qquad
    \det G_\sigma\ \le\
        \left(\prod_i T_{i,\sigma}\right)D_s(B_1,\ldots,B_s).
\]
The inequalities are valid for every choice of the probabilities,
including at a rank-deficient projected table.
\end{theorem}
\begin{proof}
For a single elliptic coordinate \(z\), Hermitian translation and
Lemma~\ref{lem:green486} give
\[
 |(U_w\Theta_t)(z)|^2
 \le \exp(c|z|^2)\,\sup_C g^2/g(t)^2
 \le \exp(c|z|^2)\,2\exp(L_C)/g(t)^2
\]
with every lift \(w\) of the prescribed shift. Take the product over
the \(n\) coordinates of an edge and denote the holomorphic product
by \(F\). Its homogeneous Taylor part of order \(k\) has squared
Fischer norm \(\sum_{|\alpha|=k}|\partial^\alpha F(0)|^2/\alpha!\).
The sphere bound in Lemma~\ref{lem:sphere60} bounds this by
\((2e^{L_C}/g(t)^2)^n c^k H_n(k)\). Multiplying over the independent
edge variables gives a uniform upper bound \(T_{i,\sigma}\) for the
sum of the squared jet products divided by \(\alpha!\), for each fixed
tuple \(a\). For \(h\in V\), apply Cauchy--Schwarz to
\(\langle h,Y_i\rangle\) with factors
\(\sqrt{\alpha!/\mu_i(a)}\) and their reciprocals.
The first squared sum is \(\langle h,B_i h\rangle\). The second is
at most \(\sum_a\mu_i(a)T_{i,\sigma}=T_{i,\sigma}\).
This proves the operator inequality.

To see the effect on the determinant, write each \(B_i\) as a sum
of rank-one operators \(u u^*\). Cauchy--Binet, applied to the
matrix with columns \(\sqrt{t_i}\,u\), shows that \(D_s\) is the
sum of \(\|u_1\wedge\cdots\wedge u_s\|^2\) over one column from
each list. In particular it is nonnegative, additive in each
operator, and increasing in each for the positive semidefinite order.
For the list \(Y_iY_i^*\) its value is precisely \(\det G_\sigma\).
The second assertion follows from the first. If an original list is
empty take its \(B_i=0\).
\end{proof}

For a positive operator \(B\) on \(V\) let \(e_s(B)\) be the sum of
all products of \(s\) distinct eigenvalues, counted with multiplicity.
Put \(e_s(B)=0\) if \(s>\dim V\). Define
\[
 \mathcal C(B_1,\ldots,B_s)=
   \inf_{\substack{\lambda_i>0\\\prod_i\lambda_i=1}}
                                  e_s\!\left(\sum_i\lambda_i B_i\right).
\]
This expression can be bounded from above by choosing any positive
scalings of product one. The optimization loses at most a fixed
constant per column in the logarithmic estimate.

\begin{lemma}[Capacity of the coefficient matrices]\leavevmode
\label{lem:capacity489}
For all \(s\ge1\),
\[
 {s!\over s^s}\mathcal C(B_1,\ldots,B_s)
 \ \le\ D_s(B_1,\ldots,B_s)\ \le\ \mathcal C(B_1,\ldots,B_s).
\]
If \(\mathcal C>0\), its logarithm is the infimum of a convex
function of the real variables \(\log\lambda_i\) with zero sum.
\end{lemma}
\begin{proof}
Write \(n=\dim V\) and assume \(s\le n\), the other case being immediate.
The polynomial \(P(t)=e_s(\sum_i t_iB_i)\) is homogeneous of degree
\(s\). Its coefficients are nonnegative by the same rank-one
expansion, and \([t_1\cdots t_s]P=D_s\). Every positive evaluation
therefore bounds \(D_s\prod t_i\) from above. On writing
\(t_i=e^{y_i}\), the logarithm is the log of a sum of exponentials
of affine functions. Its Hessian is the covariance matrix of their
exponents with their normalized positive weights, hence is positive
semidefinite whenever \(P\ne0\).

For the reverse bound we give Gurvits's capacity argument
\cite{XXVI-XX-gurvits489}. The polynomial \(\det(zI+\sum_i t_iB_i)\) is nonzero
when all its variables have positive imaginary part, since the
imaginary part of its quadratic form is positive definite.
A polynomial with this nonvanishing property is called stable.
Its derivatives and real specializations are stable or zero. For
a derivative in one variable of degree \(m>0\), its leading
coefficient, as a polynomial in that variable, is the locally
uniform limit upon setting that variable to \(iR\) and dividing by
\((iR)^m\) as \(R\to\infty\). It is nonzero as a polynomial and
hence has no zeros in the product of upper half-planes by Hurwitz's
theorem. The univariate degree thus stays \(m\) there, and the
logarithmic derivative has strictly negative imaginary part by
factoring the polynomial into linear terms. This proves the derivative
assertion. Real specialization follows by taking a limit from the
upper half-plane. The version of Hurwitz used on a product follows
by restricting to complex lines through a putative zero inside a
small polydisk and using the univariate result. Since \(P\) is the
coefficient of \(z^{n-s}\), differentiation in \(z\) and setting
\(z=0\) give its stability up to the zero case.

Let \(P\) now denote any homogeneous stable polynomial of degree
\(q\) in \(q\) variables with nonnegative real coefficients. Its
capacity is \(\inf_{t>0}P(t)/\prod t_i\), equivalent by homogeneity
to the infimum at \(\prod t_i=1\). Assume its capacity is positive.
Freeze the first \(q-1\) variables at positive real values \(t'\).
The resulting polynomial \(R(x)\) has only nonpositive real roots
by stability, real coefficients, and positivity on \(x>0\).
Its degree \(m\) is positive since a positive capacity excludes
a constant in \(x\). A root at zero of multiplicity at least two
would likewise force capacity zero. If zero is a simple root,
\(\inf_{x>0} R(x)/x=R'(0)\). If \(R(0)>0\), write
\(R(x)=R(0)\prod_{\nu=1}^m(1+x/b_\nu)\) with \(b_\nu>0\).
The arithmetic--geometric mean bounds it above by
\(R(0)(1+xR'(0)/(mR(0)))^m\). Minimizing this upper bound divided
by \(x\) gives
\[
 R'(0)\ \ge\ g_m\inf_{x>0}R(x)/x,\qquad
 g_1=1,\qquad g_m=((m-1)/m)^{m-1}\quad(m\ge2).
\]
The same inequality holds for the simple-root case. The numbers
\(g_m\) decrease with \(m\), so \(g_m\ge g_q\).
Divide by the product of the fixed \(t'\) and take their infimum.
The derivative at zero retains capacity at least \(g_q\) times the
old capacity. It is again stable and homogeneous with nonnegative
coefficients. Repeat in all variables. The final scalar is the
coefficient of their product and \(\prod_{q=1}^s g_q=s!/s^s\).
This proves the lower bound, also valid when the capacity is zero.
\end{proof}

\begin{corollary}[The finite matrix term on the period interval]\leavevmode
\label{cor:commoncapacity489}
On the polynomial-conductor interval of
Theorem~\ref{thm:primitivegram463} take its \(s=q_c\) divided
columns. Form \(B_{i,\sigma}\) from the actual projected coefficient
vectors above and put
\(\mathcal C_\sigma=\inf_{\mu_i}\mathcal C(B_{1,\sigma}(\mu_1),
\ldots,B_{s,\sigma}(\mu_s))\). Set
\[
 L_{d_0}=\log(\sqrt{d_0}/(2\pi))-\gamma_{\rm E}+D_{\xi_{d_0}},
 \qquad \bar k={1\over s}\sum_{i,j} k_{ij}.
\]
If the columns are independent at the prescribed tests, then,
uniformly on these fields and translates, for every fixed \(A>0\),
\[
 \begin{aligned}
 \Delta_u\ \le{}& K_u+\bar k L_{d_0}
             +{1\over s[J:\mathbb Q]}\sum_\sigma\log\mathcal C_\sigma
             -A_{\rm sig}(x)+W_{d_0}(u)-\rho_x\\
             &+o_{\psi,\varepsilon,A}(b/x^A).
 \end{aligned}
\]
\end{corollary}
\begin{proof}
The two preceding results bound \(\det G^\circ_\sigma\) by
\(\mathcal C_\sigma\prod_iT_{i,\sigma}\), also after the infimum
since its prefactor is independent of all probabilities. A positive
determinant gives \(\mathcal C_\sigma>0\). For each edge,
Theorem~\ref{thm:mean56} gives mean \(\log c_{ij,\sigma}=L_{d_0}\)
over \(J\). The elliptic curves and nonzero marks in the product
presentation are the transported copies of the marked CM model;
complex conjugation permutes its embeddings. By
Lemma~\ref{lem:cmorbit487} and Theorem~\ref{thm:greenenergy488}, the
mean log of the first factor inside the square brackets defining
\(T_{i,\sigma}\), including its exponent, is
\[
 n_{ij}\big(\log2+1+\log(\sqrt{d_0}/(2\pi))
                      +D_{\xi_{d_0}}-2 b_K(t_{ij})\big).
\]
Here \(b_K(t_{ij})\ge0\). The upper bound after dropping this
nonpositive term is \(O_{\psi,\varepsilon}(n_{ij}x^{\varepsilon/4})\)
by Corollary~\ref{cor:interval56} and its bounds on \(d_0,D_{\xi_{d_0}}\).
Also \(\log H_{n_{ij}}(k_{ij})\le
n_{ij}(1+\log(1+k_{ij}/n_{ij}))\) by the sphere lemma.
The sum of the coordinate dimensions along a word and the log of
its largest degree plus two are bounded by powers of \(x\), uniformly
in the word, by the rank and weight bounds in
Corollary~\ref{cor:fixedheight479}. Thus, after division by \(s\),
all these upper errors together are \(O_{\psi,\varepsilon}(x^C)\)
for some \(C\) depending only on the fixed parameters. Since
\(b=\exp(x^\varepsilon)\) this is \(o(b/x^A)\). Insert the height
bound into Theorem~\ref{thm:primitivegram463}.
\end{proof}

Consequently the archimedean term at these tests has an upper bound
in terms of the signed finite coefficient matrices and the explicit
degree contribution \(\bar k L_{d_0}\). Section~\ref{sec:missing}
still requires a sufficiently small bound relative to
\(A_{\rm sig}(x)\), a nonzero determinant at the fixed tests, and
the exceptional strict-length estimate on this same interval.


\subsection{A metric bound for the fixed signed coefficients}
\label{sec:dualmetric491}

In Corollary~\ref{cor:commoncapacity489} the signed coefficient
vectors come from the original projected tests. We express their
contribution using a single positive Hermitian operator at each
embedding. This also gives an upper bound from any chosen
orthogonal decomposition of the actual test space.

\begin{lemma}[Selecting a degree in a positive determinant]\leavevmode
\label{lem:selectdegree491}
Let \(B\) be a positive semidefinite Hermitian operator in finite
dimension. For every integer \(s\ge1\),
\[
 e_s(B)\ \le\ \inf_{t>0}{\det(I+tB)\over t^s}
       \ \le\ {e^s s!\over s^s}\, e_s(B).
\]
The constant is independent of the dimension.
\end{lemma}
\begin{proof}
Write \(P(t)=\det(I+tB)=\sum_{j=0}^r d_jt^j\), where \(r\)
is the rank and \(d_j=e_j(B)>0\) for \(0\le j\le r\). All
roots of \(P\) are negative real numbers. Every positive
evaluation bounds its degree-\(s\) term from above. For \(r<s\)
let \(t\to\infty\) to obtain zero. Suppose \(s\le r\).
The ratios \(b_j=jd_j/d_{j-1}\) decrease weakly with \(j\).
Indeed, by Rolle's theorem each derivative of \(P\) also has
all its roots negative, with multiplicities allowed. The log
of a positive polynomial with such a factorization has
nonpositive second derivative on the positive real axis and at
zero. Applied to \(P^{(j-1)}\) at zero this gives
\((j!d_j)^2\ge (j-1)!d_{j-1}(j+1)!d_{j+1}\), as required.
Take \(t=s/b_s\).
For \(j>s\), each ratio of successive \(t^j d_j\) is at most
\(s/j\), and for \(j\le s\) the ratio is at least \(s/j\).
Multiplying the inequalities starting from \(t^s d_s\) in
either direction gives
\[
        {t^j d_j\over t^s d_s}
                  \le {s!\over s^s}{s^j\over j!}.
\]
Sum over \(j=0,\ldots,r\) and use the exponential series.
\end{proof}

Fix the embedding and abbreviate its test space to \(V\), of
dimension \(n\). Retain the complete vectors \(v_{i,a,\alpha}\)
of Section~\ref{sec:projectedmatrices489}, including the division
of each column when using \(G^\circ_\sigma\). For a positive
definite Hermitian operator \(H\) define
\[
 \begin{aligned}
 d_{ia}(H)&=\sum_\alpha\alpha!\,
                     \langle v_{i,a,\alpha},Hv_{i,a,\alpha}\rangle,\\
 \Psi_\sigma(H)&=\operatorname{tr}H-\log\det H-n+s
             +2\sum_i\log\!\left(\sum_{a\in\mathcal A_i}
                                             \sqrt{d_{ia}(H)}\right).
 \end{aligned}
\]

\begin{theorem}[One Hermitian metric for the signed coefficient matrices]\leavevmode
\label{thm:dualmetric491}
Let \(D_{\rm mix}=\inf_\mu D_s(B_1(\mu_1),\ldots,B_s(\mu_s))\),
with one probability distribution on the original tuples for
each word as in Theorem~\ref{thm:projectedmatrices489}. Then
\[
 \begin{gathered}
 D_{\rm mix}\ \le\ \mathcal C_\sigma
   \ \le\ \exp\!\left(\inf_{0<H\le I}\Psi_\sigma(H)\right)
   \ \le\ {e^s s!\over s^s}\mathcal C_\sigma\ \le\ e^sD_{\rm mix} .
 \end{gathered}
\]
The infimum of \(\Psi_\sigma\) over all positive definite
operators has the same value. A zero coefficient column gives
zero in all three expressions in exponential scale.
\end{theorem}
\begin{proof}
Assume each word has some nonzero coefficient. Define
\[
 Q=\inf_{\mu,\,\lambda_i>0}
       {\det(I+\sum_i\lambda_i B_i(\mu_i))\over\prod_i\lambda_i}.
\]
Write \(\lambda_i=t\lambda_i'\), with \(\prod_i\lambda_i'=1\).
For fixed \(\lambda'\) and \(\mu\) apply
Lemma~\ref{lem:selectdegree491} to \(\sum_i\lambda_i'B_i(\mu_i)\).
Taking the remaining infima gives
\(\mathcal C_\sigma\le Q\le (e^s s!/s^s)\mathcal C_\sigma\).
Lemma~\ref{lem:capacity489} gives the two inequalities involving
\(D_{\rm mix}\).

For any positive semidefinite \(W\),
\[
  \log\det(I+W)=\min_{H>0}
       \{\operatorname{tr}(H(I+W))-\log\det H-n\}.
\]
Indeed the difference between the right-hand expression and the
left-hand side is the sum of \(u-1-\log u\) over the positive
eigenvalues \(u\) of \((I+W)^{1/2}H(I+W)^{1/2}\).
It is nonnegative and equality holds for \(H=(I+W)^{-1}\le I\).
In \(\log Q\) insert this identity and take the infima in the
opposite order, all being infima over independent sets. Once \(H\)
and \(\mu\) are fixed, minimizing in the \(i\)-th scalar gives
\[
 \inf_{\lambda_i>0}
       \{\lambda_i\operatorname{tr}(HB_i(\mu_i))-\log\lambda_i\}
                          =1+\log\operatorname{tr}(HB_i(\mu_i)).
\]
Next use \(\operatorname{tr}(HB_i(\mu_i))=\sum_a d_{ia}(H)/\mu_i(a)\).
Cauchy--Schwarz gives the minimum, in the sense of an infimum,
as \((\sum_a\sqrt{d_{ia}(H)})^2\). For positive \(d_{ia}\)
use weights proportional to their square roots; zero terms
can be assigned arbitrarily small positive total mass.
Substitution gives \(\log Q=\inf_H\Psi_\sigma(H)\), including
the value \(-\infty\). The same proof starting with \(0<H\le I\)
gives the restricted formula since that set contains every
minimizer \((I+W)^{-1}\). For a zero column the mixed coefficient
vanishes and in \(Q\) its scalar can be taken arbitrarily large.
This gives the asserted limiting interpretation.
\end{proof}

\begin{corollary}[An estimate using any chosen test metric]\leavevmode
\label{cor:orthometric491}
In the setting of Corollary~\ref{cor:commoncapacity489} with
independence at the fixed tests, take any \(0<H_\sigma\le I\)
on each actual projected test space. Then
\[
 \begin{aligned}
 \Delta_u\ \le{}&K_u+\bar k L_{d_0}
      +{1\over s[J:\mathbb Q]}\sum_\sigma\Psi_\sigma(H_\sigma)
      -A_{\rm sig}(x)+W_{d_0}(u)-\rho_x\\
      &+o_{\psi,\varepsilon,A}(b/x^A).
 \end{aligned}
\]
In particular let \(\Pi_b\) be the orthogonal projections of any
decomposition of \(V\), of ranks \(m_b\), and let \(0<h_b\le1\).
For \(H=\sum_b h_b\Pi_b\) the entire metric term is
\[
 \begin{aligned}
 E_{ia,b}&=\sum_\alpha\alpha!\,\|\Pi_bv_{i,a,\alpha}\|^2,\\
 \Psi_\sigma(H)&=s+\sum_b m_b(h_b-1-\log h_b)
          +2\sum_i\log\sum_a \left(\sum_b h_b E_{ia,b}\right)^{1/2}.
 \end{aligned}
\]
\end{corollary}
\begin{proof}
Insert \(\log\mathcal C_\sigma\le\Psi_\sigma(H_\sigma)\) in
Corollary~\ref{cor:commoncapacity489}. The spectral decomposition
of the chosen \(H\) gives the displayed expression, keeping the
projection and all path sums in each vector before taking its
squared norm. For comparison with the exact mixed coefficient,
the previous theorem shows that optimization over all metrics
costs at most one in the normalized log per column. This is
\(o(b/x^A)\) on the indicated intervals for every fixed \(A\).
\end{proof}

The finite metric problem is convex along the positive definite
matrix paths \(H(t)=H_0^{1/2}(H_0^{-1/2}H_1H_0^{-1/2})^tH_0^{1/2}\).
Indeed diagonalizing the middle matrix expresses each \(d_{ia}(H(t))\)
as a sum of exponentials of linear functions of \(t\) with nonnegative
coefficients. Positive such sums are log-convex by H\"older's inequality.
Their square roots and the sums of these square roots remain log-convex
by the same inequality. The trace of \(H(t)\) is convex by the
same diagonalization, and \(\log\det H(t)\) is affine.
This gives the asserted convexity for nonzero columns. Estimating
these signed coefficient energies at the required size and proving
independence at the original tests remain arithmetic questions.

\section{Terminal Fourier components in the period height}\label{nav:chapter-terminal-fourier-heights}

\subsection{The connected half of the terminal Fourier index}
\label{sec:terminalhalf496}

Use the common terminal prime \(T\) and the unrefined integral
Weyl lattices of Theorem~\ref{thm:rankonecost455}. Thus \(T\ge5\)
is good and splits completely in the fields of the terminal test.
On an edge with source matrix rank \(r\), the local condition at
each chosen terminal half is \(\xi'\bar X=\xi\), where the two
covectors at its vertices are nonzero. The factor \(T^{-r}\)
in Lemma~\ref{lem:terminalfourier454} is given by the actual
normalized average over \(u\in\mathbb F_T^r\). We calculate its
contribution before summing over the intermediate covectors.

\begin{lemma}[Summing a connected frequency block]\leavevmode
\label{lem:terminalblock496}
Fix a coefficient place \(w\mid T\), one such block, and its
output covector \(\xi'\). Suppose the terminal frequencies in
that block belong to the connected CM half at \(w\).
In a rational orthonormal product basis adapted to their
subspace, the complete multiplier with the denominator \(T^r\)
on the product vacuum has integral undivided derivatives in all
the constrained coordinates. This remains true term by term
after the translations at the other finite primes. The choices
of bases can be imposed independently on all blocks over \(T\).
\end{lemma}
\begin{proof}
At the residue position lattice, \(u\mapsto
(\bar X\mapsto\xi'\bar X u)\) injects \(\mathbb F_T^r\) into
its dual matrix space: some entry of \(\xi'\) is a unit, and
that row of \(\bar X\) can vary freely. Lift \(\xi'\) to a
primitive integral covector and take a basis of the domain
of \(u\). The resulting \(r\) rank-one functionals extend
to an integral basis of the dual lattice. Indeed projection to
the row with that unit entry gives a left inverse after division
by the unit. Take the dual basis on the position side. Pairing
corresponding vectors in these two bases gives an orthonormal
basis in the split Hermitian model, as in the proof of
Theorem~\ref{thm:rankonecost455}. Do this at every required
block and approximate the combined local basis to precision
\(T^2\) by a rational orthonormal basis. The Cayley
approximation of Theorem~\ref{thm:adapted314} permits this
open congruence at both halves. Its vacuum intertwiner has
index prime to \(T\).

In these product coordinates the \(r\) modulation directions
are in distinct elliptic coordinates modulo the integral
periods, each generating one pure CM \(T\)-line. We verify the
phases needed when their terms are summed. At each coordinate
choose the primitive period \(\omega_1\) of
Lemma~\ref{lem:cyclicneighbor495} to agree with the chosen
local generator after multiplication by \(T\), to precision
\(T^2\); explicitly prescribe \(2\omega_1\) modulo \(T^2\)
times the period lattice. The primitive vector lifting argument
in that lemma permits precisely these two split congruences.
The difference from the original local modulation, after
division by \(T\), lies in \(T\) times the local period lattice.
On the orbit of the \(T\)-local vacuum under all these modulations
this difference acts trivially. It fixes the vacuum and its
pairings with all points of denominator \(T\) are integral.
The local Heisenberg cocycle is trivial for these pairs as
well: the split form pairs integral coordinates with the dual
coordinates, and the two factors of \(T\) cancel. This uses
the product coordinates generating the trace lattice at \(T\).

Here is the same check for the rational lift at the other
primes. In a scalar product coordinate write the period
\(\omega_1=\delta l\), with \(l\in\mathcal O_K\) in the
standard elliptic lattice \(\delta\mathcal O_K\).
The rational lift of the generator is \(2l/T\).
Its Hermitian norm \(4l\bar l/T^2\) is an even integer, since
the opposite component of \(l\) is divisible by \(T^2\).
Thus the corrected translation \(W(2jl/T)\) in the proof of
Theorem~\ref{thm:one61} agrees with \(U_{2j\omega_1/T}\).
At primes other than \(T\) the lift is integral, including
its factor two at the dyadic prime. Its period action on the
reference vacuum is trivial there in the corrected convention;
the correction is trivial on pairs of integral periods.
The product-coordinate directions have zero mutual pairing.
Equivalently, on this orbit the physical modulation for a
generator is a scalar times the indicated cyclic translation,
and any change of scalar compatible with the group law is a
character. Such a character only shifts the Fourier label.
This verification accounts for the possible phases before
any arithmetic triangle inequality is used.

The terminal multiplier on the reference product is therefore
the product of \(r\) complete cyclic projections \(S_{k_j}\)
of Lemma~\ref{lem:cyclicneighbor495}, with the additive
labels supplied by \(\xi\). The remaining factors are the
original scalar vacua. Operators supported at other finite
primes commute with the entire terminal multiplier, by the
factorization of the finite Heisenberg action. A prime-to-\(T\)
translation applied after a cyclic projection can be represented
on its common cover by a point of prime-to-\(T\) order.
Indeed the map from the common cover to the original curve
has kernel of order \(T\), so gives an isomorphism on
prime-to-\(T\) torsion. At \(T\) choose a lift that is integral
in the common period lattice. On reducing a rational lift
modulo that lattice, the period acts by the character of the
specified section, a root of unity. This gives exactly the
\(U_{b_0}S_{k_j}\) of Theorem~\ref{thm:cyclicunit495}, up
to a unit. Thus all undivided derivatives are integral,
including those at a zero label. The other-prime vacuum
intertwiners are finite sums of such translations with
coefficients integral at \(T\). The unit comparison in
Lemma~\ref{lem:unitframe455} permits returning to a common
locally integral frame.
\end{proof}

\begin{theorem}[The terminal cost in the prescribed Weyl lattices]\leavevmode
\label{thm:terminalhalf496}
For an edge \(j\) on column \(i\) and a coefficient place \(w\mid T\)
put \(B_w=\{\tau:P_{\tau,w}\in S_T\}\), using the conjugate
coefficient place on a conjugated edge. Remove the separately
specified nonterminal scalar factors, and include the Fourier
denominators at \(T\). Each matrix entry of the terminally
restricted class operator between the integral Weyl lattices
has valuation at least
\[
                   -\sum_{\tau\in B_w}
                                  (r_j+\nu_{ij,\tau,1}) .
\]
The bound holds for the complete indicator on the original
matrix lattice with its specified covectors.
\end{theorem}
\begin{proof}
By the split position and frequency identification in
Lemma~\ref{lem:terminalfourier454}, the primary shift at \(v\)
lies over the conjugate of its selected position prime.
Consequently it is in the \'etale CM half at the coefficient
place when \(P_{\tau,w}\in S_T\), and in the connected half
otherwise. This is the conjugate annihilator computation of
Theorem~\ref{thm:terminaldegrees454}. The map
\(\tau\mapsto v=P_{\tau,w}\cap F\) is bijective.

Keep the Fourier expansion in the blocks indexed by \(B_w\).
It costs \(T^{-r_j}\) at each such block. Fix one tuple of its
frequencies. Use the bases of Lemma~\ref{lem:terminalblock496}
in all connected blocks. In each \'etale block with nonzero
frequency choose the first dual coordinate equal to its
primitive rank-one functional as in
Theorem~\ref{thm:rankonecost455}. These choices are simultaneous
since they are in distinct local blocks. The connected blocks
with their denominators now have integral derivatives of all
orders. In an \'etale block only its first coordinate can cost
one in valuation per derivative, and none do if its frequency
is zero.

In a physical block \(\tau\) each derivative matrix is an
integral combination of regular operators and, for a matching
nonzero \'etale frequency, the extra term \(U_\tau\partial_a\)
in its first coordinate, where \(U_\tau\) has rank one.
The regular operators are ordinary derivatives on the integral
coordinates and \(T\partial_a\) on the first coordinates
with a possible loss. This decomposition follows by reduction
of the adapted basis at \(T\). Derivatives entering such a
first coordinate from a different block have coefficients
in \(T\mathcal O_w\). Each flag minor is affine along the
rank-one direction \(U_\tau\), even after the source matrix
substitution. A Weyl test has \(\nu_{ij,\tau,1}\) flag minors
in each product. Expanding in the separate rank-one variables
therefore leaves at most that many extra derivatives in block
\(\tau\). This is precisely the differential argument of
Theorem~\ref{thm:rankonecost455}, now on the product with
the connected blocks already summed. It applies to each
integral source generator and dual test. Restoring the
denominator on the \'etale blocks and summing their finitely
many terms proves the claim.
\end{proof}

\begin{corollary}[A smaller terminal term on the common intervals]\leavevmode
\label{cor:centeredterm496}
In Theorem~\ref{thm:primitivegram463} the terminal summand of
\(K_u\) can be replaced by
\[
        \frac{\log T}{q_c}\sum_{i,j}
                   \left(fr_j+\sum_\tau\nu_{ij,\tau,1}\right).
\]
Here \(f=[E_x^+:\mathbb Q]\) and \(q_c\) is the number of columns
in the specified independent list. The saving in the upper
bound for \(\Delta_u\) is
\((f\log T/q_c)\sum_{i,j}r_j\).
\end{corollary}
\begin{proof}
For a fixed \(\tau\), the two conjugate halves at \(T\) are
exchanged by conjugation of the coefficient places; equivalently
the pullbacks \(P_{\tau,w}\) cover the primes of \(E_x\) over
\(T\) with their degree weights. Exactly half of this weight
has \(P_{\tau,w}\in S_T\). The theorem thus contributes at
most \(\tfrac12(fr_j+\sum_\tau\nu_{ij,\tau,1})\log T\)
per edge in the averaged denominator estimate. The two factors
in the Gram calculation give the asserted coefficient. Each
bound was uniform in the covectors and in the coupled
intermediate class. The nonarchimedean triangle inequality
therefore carries it through the prescribed class sums and
compositions. The same projector denominator, scalar division
by \(b_i^{\rm per}\), and archimedean matrix in
Theorem~\ref{thm:primitivegram463} then give the claim on each
of its original fields and intervals.
\end{proof}
The corollary reduces the terminal contribution for any
independent list admitted in that theorem. Establishing the
required projected rank and uniform upper bound for the
resulting Gram height remains necessary in Section~\ref{nav:I-28}.


\subsection{Cancellation of the explicit terminal index in the height}
\label{sec:terminalheight497}

Fix the nonzero covectors at the vertices of a path as in
Lemma~\ref{lem:terminalfourier454}. This time perform the full
cyclic Fourier sums of Lemma~\ref{lem:terminalblock496} at
every terminal position. For a given class, choose its rational
orthonormal basis adapted to the \(r_j\) frequency directions
at each of the \(f\) positions on edge \(j\). The construction
of that lemma applies to both CM halves. It identifies the
terminal section on each of the \(fr_j\) constrained coordinates
with one \(S_k\); the good-reduction orientation was used only
in evaluating its finite valuation. Expand the other local
operators on the product vacuum as finite sums of translations
at primes distinct from \(T\). Return the scalar for each
choice of basis to the original normalization according to
Theorem~\ref{thm:framechange314}. It has absolute value one
at all complex embeddings.

\begin{lemma}[Norm and pointwise bound for the complete projection]\leavevmode
\label{lem:packetnorm497}
Let \(h=\|\Theta_t\|_{\rm av}^2\) on the original normalized
elliptic curve, with the metric density \(\exp(-c|z|^2)\),
\(c=\pi/A\), and probability Haar measure. For every \(k\), the
average squared norm of \(S_k\) on its elliptic period lattice
is \(h/T\). The same holds for \(U_b S_k\) for any torsion lift
\(b\). At every complex embedding these normalized factors,
formed with the original finite Fourier multipliers, satisfy
\[
        |(U_b S_k)(z)|^2
          \le e^{c|z|^2}\,\frac{B_K h}{T},\qquad
        B_K=e\max\{1,2\sqrt{|D_K|}/\pi\},
\]
using the area and reference norm at that embedding. The
corresponding bound for an unprojected translated elliptic
factor is \(e^{c|z|^2}B_K h\).
\end{lemma}
\begin{proof}
On the common cover by \(\Gamma\) of
Lemma~\ref{lem:cyclicneighbor495}, all \(T\) summands have
squared norm \(h\) by translation invariance and invariance of
Haar averages under finite covers. They are mutually orthogonal:
the unitary periods from the original lattice act on them by
distinct characters. Consequently the normalized sum with
coefficient magnitudes \(1/T\) has squared norm \(h/T\).
Its metric density descends to \(\mathbb C/\Lambda'\).
Both lattices have area \(A\); translating the function again
preserves its average norm.

For transport under embeddings we give the scalar calculation.
Take \(N=2M\) as in Theorem~\ref{thm:divisionphase493}, with a
common \(M\) divisible by \(T\) killing every shift under
consideration, enlarging it to kill a lift from the common
cover if necessary. Under an embedding the rational quotient
in that theorem transports exactly when \([w/(2N)]\) is
transported. For \(w=2\omega_1/T\) this point is
\([\omega_1/(TN)]\). Its conjugate is represented by a primitive
period \(\omega_1'\) divided by \(TN\), since its exact order
is \(TN\). To lift a primitive residue pair modulo \(TN\),
complete it to a unimodular matrix modulo each prime power by
elimination at a unit entry, and combine the elementary matrices
by the Chinese remainder theorem. Lifting them gives an
integral unimodular matrix. Semilinearity for CM preserves
the condition of belonging to the prescribed pure line at
level \(T^2\), with the prime transported. Use the multiples
\(jw'\) of this same representative \(w'=2\omega_1'/T\) in
all the terms, and a representative \(b'\) transporting
\([b/(2N)]\). The composition phase is transported as well:
\[
       \exp(\pi i E(b,w))
           =e_{N^2}([b/(2N)],[w/(2N)])^2 .
\]
The defining multiplier for \(U_bU_{jw}\) uses this phase or
its inverse to power \(j\). Functoriality of the Weil pairing
and additivity of division points thus give a common translate
of a full cyclic sum on the conjugate curve. The character
\(\zeta_T^{-k}\) conjugates to another \(T\)-th root of unity.
This proves the norm assertion there including its relative
phases.

For the pointwise estimate write either elliptic period lattice
as \(\gamma I\), where \(I\) is a fractional \(\mathcal O_K\)-ideal.
Every nonzero \(\alpha\in I\) satisfies
\(|\alpha|^2=|N_{K/\mathbb Q}\alpha|\ge N(I)\), since
\((\alpha)I^{-1}\) is integral. The area of \(I\) is
\(\sqrt{|D_K|}N(I)/2\). Thus the shortest period has square
length at least \(2A/\sqrt{|D_K|}\), uniformly for the old
ideal and for the new ideal
\(\mathfrak p^{-1}\bar{\mathfrak p}I\).
Choose a disk of squared radius
\(r^2=\min(A/\pi,A/(2\sqrt{|D_K|}))\), which embeds up to its
boundary in the elliptic torus. At any chosen center apply the
mean-value inequality to the squared modulus of the holomorphic
Hermitian translate. Its integral over the disk is at most
\(e^{cr^2}A\) times its squared average norm. Division by
\(\pi r^2\) bounds the metric value there by \(B_K\) times that
norm. Undoing the metric gives the formulas.
\end{proof}

For the exact expansion at an embedding \(\sigma\), let \(a\)
index the remaining data of a summand of column \(i\): the
covectors, the class path with its coupled intermediate vertex,
and the translations outside \(T\). Allow separate indices
for different choices of product coordinates. After expanding
the original homogeneous differential and the prescribed
projection, write \(v'_{i,a,\alpha}\) for the resulting coefficient
vector, including the frame scalar, class masses and signs. Use
the divided columns if the Gram matrix is \(G^\circ\).
The scalar factors differentiated in that summand are the
complete \(U_bS_k\) and the unprojected \(U_b\Theta_t\) just
constructed. The Fourier coefficient \(1/T\) is already inside
each \(S_k\). For positive \(\mu_i(a)\) of sum one put
\[
  B_i'=\sum_{a,\alpha}\frac{\alpha!}{\mu_i(a)}
                 v'_{i,a,\alpha}(v'_{i,a,\alpha})^* .
\]
An empty sum gives zero. Let \(n_{ij},k_{ij}\) be the number
of elliptic coordinates and the prescribed total derivative
degree of edge \(j\). Set \(h_{ij,\sigma}\) equal to the squared
average norm of its original marked elliptic factor, \(c_{ij,\sigma}
=\pi/A_{ij,\sigma}\), and
\[
  Q_{i,\sigma}=\prod_j
       (B_K h_{ij,\sigma})^{n_{ij}}c_{ij,\sigma}^{\,k_{ij}}
                          H_{n_{ij}}(k_{ij}),\qquad
  d_i^T=f\sum_j r_j .
\]
Here \(H_n(k)\) is the explicit sphere factor of
Lemma~\ref{lem:sphere60}. Use the corresponding factor norms
in the product if the elliptic marks differ by coordinate.

\begin{theorem}[Archimedean estimate after the terminal sums]\leavevmode
\label{thm:packetgram497}
For the prescribed Gram matrix at any embedding and a list
of \(s\) columns, the coefficient matrices just defined satisfy
\[
     \det G_\sigma\ \le\
         T^{-\sum_i d_i^T}
           \left(\prod_i Q_{i,\sigma}\right)
                              D_s(B_1',\ldots,B_s').
\]
\end{theorem}
\begin{proof}
Exactly \(fr_j\) scalar factors on edge \(j\) are projected,
since the fixed nonzero \(\xi'\) gives \(r_j\) independent
functionals at each position. Their product, including the
unconstrained coordinates, is bounded in squared magnitude by
\[
       T^{-fr_j}(B_Kh_{ij,\sigma})^{n_{ij}}
               \exp(c_{ij,\sigma}\|z\|^2)
\]
by Lemma~\ref{lem:packetnorm497}. Lemma~\ref{lem:sphere60}
bounds the sum of the squared derivatives of degree \(k_{ij}\)
divided by their multi-index factorials by this constant
times \(c_{ij,\sigma}^{k_{ij}}H_{n_{ij}}(k_{ij})\).
Taking the product over the separate edge variables bounds
the corresponding sum for each fixed \(a\) by
\(T^{-d_i^T}Q_{i,\sigma}\). Vector Cauchy--Schwarz with weights
\(\alpha!/\mu_i(a)\), as in Theorem~\ref{thm:projectedmatrices489},
gives \(Y_iY_i^*\le T^{-d_i^T}Q_{i,\sigma}B_i'\).
Monotonicity and multilinearity of the mixed determinant
\(D_s\) give the claimed bound.
\end{proof}

\begin{corollary}[The explicit terminal Fourier index cancels]\leavevmode
\label{cor:terminalheight497}
Write \(K_u^{(0)}\) for \(K_u\) in Theorem~\ref{thm:primitivegram463}
with its displayed terminal summand omitted. For an independent
list of \(q_c\) divided projected columns, the upper bound in
that theorem on the same fields and intervals holds with its
\(H(G^\circ)\) replaced by
\[
 \begin{aligned}
 &K_u^{(0)}
     +\frac{\log T}{q_c}\sum_{i,j,\tau}\nu_{ij,\tau,1}\\
 &\quad+\frac{1}{q_c[J:\mathbb Q]}
   \sum_{\sigma:J\hookrightarrow\mathbb C}
     \log\left(\left(\prod_i Q_{i,\sigma}\right)
                         D_{q_c}(B_1',\ldots,B_{q_c}')\right).
 \end{aligned}
\]
\end{corollary}
\begin{proof}
Use Corollary~\ref{cor:centeredterm496} at the finite places.
The terms involving \(fr_j\) there sum to
\(q_c^{-1}\sum_i d_i^T\log T\). Theorem~\ref{thm:packetgram497}
subtracts exactly this number in the averaged logarithm of
the Gram determinant. The stated first-row term and the
displayed coefficient bound give the result.
\end{proof}
The matrices \(B_i'\) still include all allowed intermediate
covectors and depend on \(T\). Bounding their contribution and
exhibiting the required rank for the original signed projection
are separate arithmetic questions. The cancellation here concerns
the specified terminal normalization in the height on \(E_x\).


\subsection{Rank of the terminal derivative table on nonzero covectors}
\label{sec:puncturedrank498}

Let \(p=T\) and use a connected terminal frequency block in
Lemma~\ref{lem:terminalblock496}. Keep the nonzero output
covector \(\xi'\) fixed. After choosing the basis of its
\(r\)-dimensional frequency subspace, the input covector is a
linear functional on the Fourier parameter \(u\). Write
\(k=(k_1,\ldots,k_r)\in\mathbb F_p^r\) for its coordinates, so
that its original coefficient is
\(\exp(-2\pi i\sum_h k_hu_h/p)\). Use the compatible cyclic
translations in the proof of that lemma. The input condition at
the vertex is \(k\ne0\). First consider a single product term
from the translations at the primes other than \(T\).

\begin{lemma}[The moment of a connected Fourier projection]\leavevmode
\label{lem:momentres498}
For each scalar coordinate \(h\) there is a nonzero residue
\(\beta_h\), depending on the differential and the frequency
generator, such that for every prime-to-\(p\) translate \(U_b\)
from the common cover and all \(n\ge0,\ k_h\in\mathbb F_p\),
\[
 \overline{J_n(U_b S_{k_h})}
     =\overline{J_0(U_b S_{k_h})}\,(\beta_h k_h)^n .
\]
We use \(0^0=1\). The factor of order zero is a unit for
\(k_h\ne0\). It is a unit also at \(k_h=0\) if the image of
\(b\) on \(\mathbb C/\Lambda'\) differs from the mark \(t'\).
\end{lemma}
\begin{proof}
Write \(z'=\alpha z\) and \(S_0=c\Theta'_{t'}\) as in the
proof of Theorem~\ref{thm:cyclicunit495}. In the connected
orientation \(v(\alpha)=1\) and \(c\) is a unit. The points
\([\lambda_k]\) in Lemma~\ref{lem:cyclicneighbor495} are in
the \'etale \(p\)-line on the new curve and
\([\lambda_k]=k[\lambda_1]\). Indeed the period pairing with
the chosen frequency identifies \(\Lambda/\Gamma\) with the
character group, and \(\Gamma\subset\Lambda'\).
Use this identification after scaling the curve by \(\alpha\).
The prime-to-\(p\) point \(b\) contributes zero to the additive
quantity \(L_1\) in Theorem~\ref{thm:residuejet492}, where the
annihilator can be chosen with \(p\)-part exactly \(p\).
Multiplication of its residue formula by
\(\overline{\alpha/p}^{\,n}\) gives the assertion with
\[
          \beta_h=\overline{\alpha/p}\, L_1(\alpha\lambda_1).
\]
This is nonzero by full \'etale order. The scalar relating the
two translated sections multiplies both the derivative and
the value by the same unit, so the equality retains the
original phase. At label zero the point on the new curve has
prime-to-\(p\) order. Injectivity of prime-to-\(p\) reduction
and the order-zero divisor calculation in
Theorem~\ref{thm:residuejet492} give the last assertion.
\end{proof}

Fix prime-to-\(p\) displacements \(b_h\). For a finite list
\(\mathcal A\) of multi-indices use the actual rectangular table
\[
   M_{\alpha,k}=\prod_{h=1}^r J_{\alpha_h}(U_{b_h}S_{k_h}),
           \qquad \alpha\in\mathcal A,\quad k\ne0 .
\]
Define \(R(0)=0\) and \(R(m)=1+((m-1)\bmod(p-1))\) for
\(m>0\), and apply \(R\) coordinatewise. Put
\(\mathcal B=R(\mathcal A)\subseteq\{0,\ldots,p-1\}^r\).

\begin{theorem}[Rank on the full set of nonzero terminal covectors]\leavevmode
\label{thm:puncturedrank498}
If each \(b_h\) misses its mark at label zero, then the rank of
the reduction of \(M\) is
\[
        |\mathcal B|-
          {\bf1}_{\{0,p-1\}^{r}\,\subseteq\,\mathcal B}.
\]
The unshifted product vacuum satisfies the hypothesis. For
arbitrary \(b_h\), restrict to the labels
\((\mathbb F_p^\times)^r\). On these columns the residue rank
is the number of distinct componentwise classes of \(\alpha\)
modulo \(p-1\). In either case some minor of the stated size
is a unit at this coefficient place when the rank is positive.
\end{theorem}
\begin{proof}
By Lemma~\ref{lem:momentres498}, scaling a row by
\(\prod_h\beta_h^{-\alpha_h}\) and a column by the inverse
of its product of order-zero values leaves the matrix of the
functions \(k\mapsto k^\alpha\). These scalings are invertible
on the indicated labels. On all of \(\mathbb F_p\) one has
\(X^m=X^{R(m)}\) as functions. The monomials with each exponent
in \(\{0,\ldots,p-1\}\) form a basis of all functions on
\(\mathbb F_p^r\) over the residue field. For \(r=1\) this is
polynomial interpolation at \(p\) distinct elements. Tensor
products give the result for every \(r\).

On deleting the origin the only functions in the kernel are
multiples of its indicator, whose unique expansion in this
basis is \(\prod_h(1-X_h^{p-1})\).
It lies in the span of the selected monomials precisely when
all its \(2^r\) monomials are selected, proving the first
formula. At the unshifted point each new mark is nonzero
two-torsion. On the multiplicative torus the same interpolation
uses the exponents \(0,\ldots,p-2\) in each coordinate and
positive or zero powers reduce modulo \(p-1\). The ordinary
rank criterion over the residue field gives a nonzero minor;
the original matrix is integral, so this minor lifts to a unit.
\end{proof}

\begin{corollary}[All homogeneous orders at one terminal block]\leavevmode
\label{cor:homogeneousrank498}
Assume the order-zero unit hypotheses of
Theorem~\ref{thm:puncturedrank498}. If the rows run over all
monomials of one degree \(d>0\), the residue rank on \(k\ne0\)
is the number of \(b\in\{0,\ldots,p-1\}^r\) for which
\[
           0<|b|\le d,\qquad |b|\equiv d\pmod{p-1}.
\]
For \(d<p\) it is \(\binom{d+r-1}{r-1}\). For \(d\ge r(p-1)\)
it is \((p^r-1)/(p-1)\). Degree zero has rank one.
More generally, for any prescribed linear combinations of
these monomial rows with integral coefficients, their residue
rank is the dimension of the span of the corresponding reduced
polynomials evaluated on \(k\ne0\), after replacing each
variable by \(\beta_h X_h\).
\end{corollary}
\begin{proof}
For \(d>0\) every reduced monomial has at least one positive
exponent. Conversely one can increase any positive exponent
of \(b\) by \(d-|b|\) when this is a nonnegative multiple of
\(p-1\). This proves the counting formula and shows that its
span excludes the delta function at zero. For \(d<p\) no
exponents change. For \(d\ge r(p-1)\), every nonconstant
reduced monomial of the given scalar weight occurs. On the
punctured vector space all functions satisfying
\(F(tk)=t^d F(k)\) have dimension \((p^r-1)/(p-1)\), one
arbitrary value on each projective line. Extend such functions
by zero at the origin. Their unique reduced polynomials have
zero constant term, and only monomials of weight \(d\bmod(p-1)\)
occur by uniqueness under scalar substitution. This proves
the large-degree formula. In the last assertion the same unit
factor from the order-zero values multiplies every row in a
given column, so the polynomial computation also applies to
the stated linear combinations.
\end{proof}
These are the covectors in the terminal indicator
\(\xi'\bar X=\xi\). When two adjacent class operators share a
covector, the global period contracts their two tables in that
index. The rank calculation supplies integral minors before that
contraction; the rank of the contracted, signed projection in
Section~\ref{nav:I-28} requires its actual coefficients.


\subsection{The signed covector sum on a connected terminal block}
\label{sec:weylmoments500}

Fix a field, word, edge and terminal prime \(p=T\) of
Theorem~\ref{thm:common111}. Work at a coefficient place where
the chosen terminal frequency block is connected. On a conjugated
edge make this choice on its conjugate. Keep the output covector
\(\xi'\ne0\) and the unrefined frames fixed. The input at this
block ranges over \(k\in\mathbb F_p^r\setminus\{0\}\), where
\(r\) is the incoming rank. Let \(\kappa\) be the residue field of a
coefficient valuation ring containing the finitely many scalars in
the calculation. Bars denote reduction to this field.

Use the complete cyclic projections and their Fourier labels in
Lemma~\ref{lem:momentres498}. In a product term \(\gamma\) of the
outgoing section let \(b_{\gamma,h}\) be the prime-to-\(p\)
displacement in its \(h\)-th constrained coordinate. Set
\[
 W_\gamma(k)=\prod_{h=1}^{r}
                  \overline{J_0(U_{b_{\gamma,h}}S_{k_h})}.
\]
Any phase multiplying the whole translated product and depending
on \(k\) is included in \(W_\gamma(k)\) by its residue. This
convention keeps the same phase in every differential order.
The formulas below also allow \(W_\gamma(k)=0\).

\begin{lemma}[The Weyl symbol in the input covector]\leavevmode
\label{lem:weylresidue500}
In the adapted basis one can take the following reduction
simultaneously for all input covectors. Evaluate an integral
Weyl symbol on the translated product at this edge. Leave the
unconstrained product derivatives as formal commuting variables
\(D^0\). Its residue after evaluation in the \(r\) constrained
coordinates is
\[
                   W_\gamma(k)\,Q(D^0,k),
        \qquad \deg_k Q\le \nu_{\tau,1},
\]
where \(\tau\) is the matching physical block and \(\nu_{\tau,1}\)
is the first row length of its partition. The polynomial \(Q\)
is obtained by integral polynomial substitution and reduction.
One may then apply any integral linear functional to the
coefficients in the complementary derivatives.
\end{lemma}
\begin{proof}
Extend a primitive lift of \(\xi'\) to a basis of the dual target
space over \(\mathbb Z_p\). Tensor this basis with a basis of the
source. It gives a basis of the dual position matrix lattice,
whose first \(r\) vectors are the functionals \(X\mapsto\xi'Xu_h\).
The dual position basis consists of the dual simple tensors.
In either basis the first \(r\) matrices have one common tensor
factor. Consequently every linear combination of these matrices
has rank at most one. This also holds after transposition. Use
these choices in the product basis construction of
Lemma~\ref{lem:terminalblock496}, with the same approximation
modulo \(p^2\). The rational basis is fixed as \(k\) varies.
In the physical chain rule the matrices of the first \(r\)
directions reduce to \(U_1,\ldots,U_r\) with this common
factor in the matching block. They reduce to zero in every
other block. Indeed reduction selects the position or the dual
position component of the split basis at the given embedding.
Distinct local blocks are independent summands of the split
lattice. Taking the tensor basis on both sides ensures the rank
property for either choice of component.

For a monomial in the product derivatives,
Lemma~\ref{lem:momentres498} replaces each constrained derivative
variable in its residue by \(\beta_h k_h\), with the common factor
\(W_\gamma(k)\) taken only once. Linearity gives this rule for
the original symbol, including at a zero label. Thus in
the matching physical matrix the substitution is
\[
        D_\tau^0\longmapsto D_\tau^0+
                                \sum_{h=1}^r U_h\beta_h k_h .
\]
The correction is an outer product with one factor constant
over the polynomial ring, so it has rank at most one. By
Lemma~\ref{lem:jetweyl453} a source generator and a dual test
use \(P_\tau([D_\tau A_\tau^{\mathsf t},0])\), followed by
coefficient extraction in the elementary parameters of
\(A_\tau\). Right multiplication and zero padding preserve the
common outer factor. Each flag minor is affine in \(k\) because
two columns from the outer product give zero in its determinant
expansion by columns.
Expansion by rows gives the same conclusion when the common
factor is on the other side. Every term in the flag polynomial
has \(\nu_{\tau,1}\) minors. The symbols at the other embeddings
are constant in \(k\). The product, parameter extraction and
integral linear combinations give the asserted degree bound.
All reductions take place in the polynomial coefficients before
evaluating complementary derivatives, which proves the last
assertion.
\end{proof}

To apply the formula to a contraction, choose integral
coordinates for the incoming table at the shared vertex.
In expanding the sum, the index \(\gamma\) can record a
coordinate of that table, a product term of the outgoing
section, and any fixed coupling at the other marked primes.
Write \(Y_{\gamma,i}(k)\) for the incoming values for column
\(i\), including any assigned signed scalars. The adjacent map
is allowed to depend arbitrarily on \(k\).
The expansion uses the unrefined stabilizer masses of
Theorem~\ref{thm:covectorpaths352}. For reduction, take a common
scalar for the finite list of incoming values and another for
all evaluations of the complementary derivative monomials and
remaining coefficients. Choose them so the two lists are
integral. Such scalars exist by finiteness. They give the same
overall normalizing factor in every summand. The assertion
concerns the resulting integral matrix and its specified
normalization. In particular it supplies no bound for the cost
of those two scalars.

\begin{theorem}[Contraction on the lines of the covector space]\leavevmode
\label{thm:weylmoments500}
For each line \(L\) in \(\mathbb F_p^r\) choose \(v_L\in
L\setminus\{0\}\). For \(0\le e\le D=\nu_{\tau,1}\) form the
signed moments
\[
 B_{\gamma,i,e}(L)=
     \sum_{a\in\mathbb F_p^\times}
              a^e W_\gamma(av_L)\,\overline{Y_{\gamma,i}(av_L)}.
\]
For a fixed outgoing Weyl row \(j\), expand the polynomial of
Lemma~\ref{lem:weylresidue500} after the complementary
evaluations as \(\sum_{e=0}^{D}Q_{j,\gamma,e}(X)\) in its
homogeneous parts. The residue of the matrix after summing over
the entire shared covector in this block has entry
\[
             \sum_{\gamma,L}\sum_{e=0}^{D}
                  Q_{j,\gamma,e}(v_L)\, B_{\gamma,i,e}(L).
\]
For several edge shapes the same formula uses their individual
degree bounds. It applies with any remaining indices fixed, and
subsequently to their original finite linear combinations with
integral coefficients. A nonzero minor of this matrix gives a
minor with unit determinant in the normalized integral matrix.
\end{theorem}
\begin{proof}
The indicator of the edge is \(\xi'\bar X=k\) in the fixed
frames. By its complete Fourier expansion every outgoing
product term has precisely the cyclic scalar labels used in the
lemma. All dependence on the same covector from the previous
edge is in the values \(Y_{\gamma,i}(k)\). Holding the common
coupling fixed until the two incident factors have been
multiplied retains the signs of that summand. Apply the lemma
after the indicated denominator clearing. For \(k=a v_L\) the
homogeneous part of degree \(e\) takes the value
\(a^e Q_{j,\gamma,e}(v_L)\). Summing over \(a\), then over
the lines, uses each admissible \(k\) exactly once. Replacing
\(v_L\) by \(t v_L\) multiplies the two factors by \(t^e\)
and \(t^{-e}\). This proves the identity independently of the
representatives. The minor assertion follows by reduction of
the determinant.
\end{proof}

\begin{corollary}[Visible moments at the prescribed word weights]\leavevmode
\label{cor:orbitkernel500}
On the original common intervals \(D+1<p-1\) for every
sufficiently large \(x\), uniformly in the chosen word, edge
and place. On a fixed line let \(s_L\) be the dimension over
\(\kappa\) of the span of the functions
\(a\mapsto W_\gamma(a v_L)\) for a list of outgoing product
terms. When a common scalar incoming table on this line varies,
all normalized Weyl rows in this list together have residue
rank at most \(\min(p-1,s_L(D+1))\), with \(D\) the maximum
first row length on the list. Here any additional coefficients
varying with \(a\) must be included in the functions defining
\(s_L\).

For one unshifted product and \(D+1<p-1\), the common kernel
of the \(D+1\) moments on the line has dimension \(p-D-2\).
Its intersection with the ordinary augmentation condition
\(\sum_{a}Y(av_L)=0\) has dimension at least \(p-D-3\).
Every member of this intersection has an integral lift satisfying
the same augmentation condition before reduction.
\end{corollary}
\begin{proof}
By Lemma~\ref{lem:wordfirstrow450} and the weight formulas
following \eqref{eq:letters111}, \(D\ll xb+x^C\) for a fixed
\(C\), where \(b=\exp(x^\varepsilon)\). Theorem~\ref{thm:prime210}
gives \(\log p=(17+o(1))m\log b\) and \(m\to\infty\).
Hence \((D+1)/(p-1)\to0\). On a single line the row functions
in the theorem lie in the sum, over \(0\le e\le D\), of
\(a^e\) times the indicated \(s_L\)-dimensional space.
The dimension of that sum proves the upper bound.

For the unshifted product every \(W_\gamma(a v_L)\) is nonzero
by Theorem~\ref{thm:puncturedrank498}. Multiplication by this
weight is an invertible diagonal operator. The functions
\(1,a,\ldots,a^D\) are independent on \(\mathbb F_p^\times\),
since a polynomial of degree at most \(D<p-1\) vanishing there
is zero. Subtracting the rank \(D+1\) of the moment map from
\(p-1\) gives the kernel dimension. One more linear equation removes at
most one dimension. Lift all but one value arbitrarily and take
the last to be the negative of their sum. This lifts the given
residue vector with exact augmentation.
\end{proof}
The moments concern the actual shared index on one outgoing
block. On that block even the unshifted integral test can lose
many incoming residue directions after contraction. To obtain
the projected rank in Section~\ref{nav:I-28} one must establish
survival for the prescribed incoming columns with their actual
weights and the remaining class sums.


\subsection{Scalar orbits through the signed theta chain}
\label{sec:scalarchain501}

Fix \(p=T\), the word and the terminal generator of
Theorem~\ref{thm:common111}. Use the original initial tables
whose values sum to zero on each nonzero scalar orbit at each
\(v\mid T\), with the other coordinates held fixed. Extend these
tables by zero at any zero covector. The local character
decomposition in Corollary~\ref{cor:fullperiodmoments400}
is integral at \(T\). We compute a gain in the finite estimate
each time such an input passes through a terminal incidence
condition in the chain.

At a fixed coefficient place \(w\mid p\) let \(\omega\) denote
the Teichm\"uller lift of \(\mathbb F_p^\times\), and write
\(d=p-1\). Use the convention \(H(a k)=\omega(a)^t H(k)\) for
a scalar character of an incoming table. We take \(1\le t<d\).
For a fixed nontrivial additive character \(\psi\) put
\[
 g(t)=\frac{d-t}{d},\qquad
 G_t(c)=\sum_{a\in\mathbb F_p^\times}\omega(a)^t\psi(-ac).
\]
One has \(G_t(0)=0\) and \(v_w(G_t(c))=g(t)\) for \(c\ne0\),
in the normalization \(v_w(p)=1\). Here is a direct check of
the sign convention in Theorem~\ref{thm:fouriersmith338}.
Changing \(c\ne0\) changes the sum by a root of unity, so use
the additive terms \(\zeta^a\) for a primitive \(p\)-th root
\(\zeta\). Let \(\lambda=\zeta-1\), with \(v_w(\lambda)=1/d\)
by \(\prod_{a=1}^{p-1}(1-\zeta^a)=p\). The coefficient of
\(\lambda^n\) after expanding \((1+\lambda)^a\) is
\(\sum_{a=1}^{p-1}\omega(a)^t\binom{a}{n}\in\mathbb Z_p\).
For \(0\le n<d-t\) its reduction vanishes by summing powers
of exponents from \(t\) to \(t+n<d\) on \(\mathbb F_p^\times\).
For \(n=d-t\) the reduction is \(-1/n!\). Thus the term of
index \(d-t\) has uniquely smallest valuation. The assertion
at zero follows by summing the nontrivial scalar character.

\begin{proposition}[The orbit sum at each edge of a fixed character]\leavevmode
\label{prop:orbitsum501}
At each terminal place prescribe a nontrivial scalar character
\(\chi_v\) of the initial table, in the convention just given.
At every vertex choose representatives for the lines of nonzero
covectors. At the terminal vertex choose the actual marked
covector for its line. On a local edge with consecutive
representatives \(\eta,\eta'\) and source dimension \(r\), the
full signed period in this character may be computed by summing
over the intermediate lines and using the local multiplier
\[
 \begin{aligned}
 {\cal K}_{\eta,\eta',\chi_v}(X)
  &=\sum_{a\in\mathbb F_p^\times}\chi_v(a)\,
                        {\bf1}_{\,\eta'\bar X=a\eta}\\
  &=p^{-r}\sum_{u\in\mathbb F_p^r}
      \psi(\eta'\bar X u)\sum_{a\ne0}\chi_v(a)\psi(-a\eta u)
 \end{aligned}
\]
on the integral matrix lattice. The terms with \(\eta u=0\)
vanish. For one product vacuum in the adapted coordinates of
Lemma~\ref{lem:packetnorm497}, the complete local section for
this multiplier has average squared norm
\[
           (p-1)p^{-r}\prod_{h=1}^r h_h
\]
on the constrained coordinates, where \(h_h\) is the
corresponding original elliptic average squared norm.
This holds at every complex embedding and after any common
translation on these coordinates.
\end{proposition}
\begin{proof}
Write the covector at the \(j\)-th vertex as \(a_j\eta_j\).
Its incidence with the next vertex is
\(\eta_{j+1}\bar X_j=(a_j/a_{j+1})\eta_j\).
Take \(a_m=1\) at the marked final covector. Substitution of
\(s_j=a_j/a_{j+1}\) is a bijection on the \(m\) scalar
variables, and \(\chi_v(a_0)=\prod_{j=0}^{m-1}\chi_v(s_j)\).
Thus all scalar sums separate with exactly the first multiplier
displayed above. Perform the substitution independently at
each terminal place. The path formula of
Theorem~\ref{thm:covectorpaths352} has already used the
unrefined stabilizer masses. The matrix cosets at other
places, their signed coefficients, and the coupling coset
common to the last two edges are fixed during this substitution.
The theta section is linear in each local test. Summing over
the other data now gives the stated formula for the signed
period. The finite Fourier expansion of the indicator gives
the second equality. Scalar character orthogonality removes
the terms specified.

For the complex norm the section is
\(\sum_{a\ne0}\chi_v(a)\prod_h S_{a\eta_h}\). Work on the
product of the common covers for the cyclic projections.
In each coordinate the \(p\) frequency translates of the
vacuum are orthogonal with equal squared norm \(h_h\), by
their distinct period characters. Their normalized finite
Fourier transform therefore gives \(p\) orthogonal sections
of squared norm \(h_h/p\). Distinct \(a\) give distinct
label vectors since \(\eta\ne0\). Taking tensor products and
then summing squared norms proves the formula. Translations
are unitary on the common covers. At a complex embedding the
finite Fourier characters and their phases transport as in
Lemma~\ref{lem:packetnorm497}. The scalar characters continue
to have absolute value one.
\end{proof}

\begin{theorem}[Divisibility from the incoming scalar orbits]\leavevmode
\label{thm:augterminal501}
Use the unrefined integral Weyl lattices of
Theorem~\ref{thm:terminalhalf496}. For column \(i\) and edge
\(j\) let \(B_w\) be the embeddings with \'etale frequency
at the effective coefficient place on that edge, and
\(v(\tau)=P_{\tau,w}\cap E_x^+\). Suppose the incoming table
has valuations at least \(h\) in the common lattice and scalar
characters \(\chi_v(a)=\omega(a)^{t_v}\). After summing all
input covectors, the terminal part of this edge has entries
of valuation at least
\[
       h-\sum_{\tau\in B_w}
                   \bigl(r_j+\nu_{ij,\tau,1}-g(t_{v(\tau)})\bigr).
\]
The separately specified nonterminal scalars and class masses
are excluded here. For every separately zero-sum input the
same bound holds with \(1/(p-1)\) at each occurrence of
\(g(t_v)\). These gains add through all the consecutive
edges of the original signed period.
\end{theorem}
\begin{proof}
Keep the output covectors and the data away from \(T\) fixed.
In each \'etale frequency block use the second formula of
Proposition~\ref{prop:orbitsum501}. For each input line the
incoming value at its representative is integral after the
common input bound is removed. At a nonzero frequency pairing
the scalar sum contributes precisely \(g(t_v)\) in valuation.
The reciprocal \(p^{r_j}\) contributes \(-r_j\). Its remaining
modulation \(X\mapsto\psi(\eta'\bar X u)\) has rank one,
as a frequency matrix, and is independent of the input scalar.
Keep the complete cyclic projections in the connected blocks.
There each term with a fixed input scalar is integral in every
constrained differential order by Lemma~\ref{lem:terminalblock496}.
The sum with the character coefficients remains integral.

The adapted bases of Theorem~\ref{thm:terminalhalf496}
apply to the product of these choices. At an \'etale block
only a single rank-one direction can cost a valuation of one per
derivative, up to \(\nu_{ij,\tau,1}\) derivatives. The scalar
Gauss factors multiply the entire differential symbol. Their
valuation gain survives that estimate and the return to the
common frame. Summing over the remaining frequencies and
input lines by the valuation inequality proves the bound.

For a fixed matrix coset the coefficient map reads the
zero-extended table at \(\xi'\bar X\). Scaling \(\xi'\) by \(a\)
scales this pullback by \(a\), so scalar types and separate zero
sums propagate, including at a fixed coupling coset. Conjugation
preserves nontriviality of the character. Induction adds the gains
along the chain. The masses and denominators are those of
Theorem~\ref{thm:terminaldegrees454}.

The scalar projectors at \(w\) have unit denominators, powers of
\(p-1\). They preserve the input bound and the refined lattice
on each line, since a scalar preserves the covector kernel. The
zero-sum condition kills the trivial character at each place.
Every remaining summand gains \(g(t_v)\ge1/(p-1)\), and addition
preserves this bound.
\end{proof}

\begin{corollary}[A gain for the divided projected columns]\leavevmode
\label{cor:augheight501}
For a list of \(q_c\) independent projected columns with
separately zero-sum initial tables, one may subtract
\[
            \frac{f\log T}{q_c(T-1)}\sum_{i,j}1
\]
from the upper bound for their contribution to \(\Delta_u\)
in Corollary~\ref{cor:terminalheight497}. If all initial
tables have the quadratic scalar character at every terminal
place, one may subtract \((f\log T/(2q_c))\sum_{i,j}1\)
instead.
\end{corollary}
\begin{proof}
Scaling by \(z_i^{-1}\) as in Theorem~\ref{thm:terminaldegrees454}
puts the initial values in the unrefined lattice with the same
scalar type. Apply the final lowering and column division of
Theorem~\ref{thm:primitivegram463}. For each \(\tau\) the places
with \'etale frequency have total normalized degree \(1/2\)
by Corollary~\ref{cor:centeredterm496}. Summing their gains for
both period ideals and dividing by \(q_c\) gives the saving.
The quadratic exponent is \((T-1)/2\) at every completion.
Theorem~\ref{thm:packetgram497} supplies the archimedean bound.
\end{proof}
The determinant estimate assumes independence. Section~\ref{nav:I-28}
still requires that independence and the full height bound on
the prescribed projections.


\subsection{Mean Gauss valuations in the projected period height}
\label{sec:meangauss504}

Keep the terminal prime \(p=T\), the word, the marked generators,
and the integral frames of Theorem~\ref{thm:augterminal501}.
Put \(F=E_x^+\) and \(f=[F:\QQ]\). For a given column fix a place
\(v\) of \(F\) above \(p\). Its residue field is \(\mathbb F_p\).
The scalar action on nonzero input covectors is by
\(\mathbb F_p^\times\). Use a finite Galois coefficient field \(J\)
containing \(E_x\), the test data and all scalar character values.
The probability of a coefficient place \(w\mid p\) is
\([J_w:\QQ_p]/[J:\QQ]\). This is the weighting in the global
height, and enlarging \(J\) gives the same average by the tower
formula for local degrees.

At the effective coefficient place of an edge let
\(\mathcal E_v\) be the event that the block belonging to \(v\)
has \'etale frequency. For a nontrivial scalar character \(\chi\)
let \(g_w(\chi)=(p-1-t)/(p-1)\), where \(\chi=\omega_w^t\),
\(1\le t<p-1\), in the convention of
Proposition~\ref{prop:orbitsum501}. The valuation of \(p\) is one.
On a conjugated edge use the effective place and character of
that edge in both the event and \(g_w\).

\begin{lemma}[The smallest mean on a set of coefficient places of weight one half]\leavevmode
\label{lem:halfgauss504}
One has \(\Pr(\mathcal E_v)=1/2\). If \(\chi\) has order \(m>1\),
then
\[
       \mathbb E_w[{\bf1}_{\mathcal E_v}g_w(\chi)]\ge A(m).
\]
Here \(A(2)=1/4\), and for \(m>2\),
\[
 A(m)=\frac{1}{m\varphi(m)}
              \sum_{\substack{1\le a<m/2\\(a,m)=1}}a
                         \ \ge\ \frac1{12}.
\]
These are the exact lower bounds when only the marginal
distribution of \(g_w(\chi)\) and the weight of the event are known.
\end{lemma}
\begin{proof}
For fixed \(w,v\) there is a unique embedding \(\tau\) in the
CM type with \(\tau^{-1}w\cap F=v\), by complete splitting and
the bijection in Theorem~\ref{thm:terminalhalf496}. Conjugation
on \(J\) fixes the totally real field \(\tau(F)\) pointwise.
Thus the same \(\tau\) occurs for the conjugate place. On
\(\tau(E_x)\) it is CM conjugation, so the two pullbacks to
\(E_x\) are the partners above \(v\). Exactly one belongs
to \(S_T\). Conjugation preserves the local degree. The same
pairing applies when the effective place at an edge is conjugated.
This gives the half weight for each individual \(v\).

Fix a generator of \(\mathbb F_p^\times\). Its value under \(\chi\)
is a primitive \(m\)-th root. Since \(m\mid p-1\), \(p\) splits
completely in \(\QQ(\zeta_m)\). Reduction at all its places takes
that value through the primitive \(m\)-th roots in \(\mathbb F_p\),
each once. The total local degree over every such place in \(J\)
is the same. Thus \(g_w(\chi)\) is uniformly distributed on
\(\{a/m:1\le a<m,\ (a,m)=1\}\). For \(m>2\) precisely half
of these values lie below \(1/2\), by pairing \(a\) with \(m-a\).
Among events of weight \(1/2\) the integral is smallest when
it puts all its mass on those values. For \(m=2\) the function
is identically \(1/2\).

For the numerical bound let \(\mu\) be the M\"obius function,
\(\operatorname{rad}(l)=\prod_{q\mid l}q\), and let \(\omega(l)\)
count the distinct prime divisors of \(l\). For \(m>2\) the
endpoint \(m/2\) makes no contribution and inversion gives
\[
 m\varphi(m) A(m)
       =\sum_{d\mid m}\mu(d)d\sum_{1\le b\le m/(2d)} b .
\]
Writing \(u=m/d\), the inner sum equals \((u^2+2u)/8\)
for even \(u\) and \((u^2-1)/8\) for odd \(u\).
The quadratic terms together give \(m\varphi(m)/8\).
For odd \(m\), the remaining terms sum to
\(-\prod_{q\mid m}(1-q)/8\). For \(4\mid m\), each squarefree
\(d\mid m\) leaves even \(u\), and the remaining terms vanish
by \(\sum_{d\mid m}\mu(d)=0\). For \(m=2l\) with \(l>1\)
odd, the even-\(u\) linear terms vanish by the same cancellation
over the divisors of \(l\). The odd-\(u\) corrections then give
\(2\prod_{q\mid l}(1-q)/8\). Consequently
\[
 8A(m)=
 \begin{cases}
 1-(-1)^{\omega(m)}\operatorname{rad}(m)/m^2,&m\text{ odd},\\
 1,&4\mid m,\\
 1+(-1)^{\omega(l)}\operatorname{rad}(l)/l^2,
                                         &m=2l,\ l>1\text{ odd}.
 \end{cases}
\]
In both odd arguments \(l\) or \(m\) is at least three and
\(\operatorname{rad}(l)\le l\). This proves the lower bound.
It is attained by the displayed marginal formula when \(m=6\).
\end{proof}

\begin{theorem}[A terminal saving independent of the order of a pure scalar character]\leavevmode
\label{thm:meangauss504}
Let \(q_c\) divided projected columns be independent as in
Corollary~\ref{cor:terminalheight497}. Suppose each initial
table belongs to a nontrivial scalar character at each \(v\mid T\),
of order \(m_{iv}\) on column \(i\).
From the upper bound in that corollary one may subtract
\[
   \frac{2\log T}{q_c}\sum_{i,j}\sum_{v\mid T} A(m_{iv})
           \ \ge\ \frac{f\log T}{6q_c}\sum_{i,j}1 .
\]
Here \(j\) runs over the edges of the word for column \(i\).
If the columns can be paired with matching edges, terminal halves
and effective places, and inverse scalar characters at every \(v\),
the saving is at
least \((f\log T/(2q_c))\sum_{i,j}1\).
\end{theorem}
\begin{proof}
At a coefficient place \(w\), the improvement at \(v\) from
Theorem~\ref{thm:augterminal501} is
\({\bf1}_{\mathcal E_v}g_w(\chi_{iv})\) in valuation. That
theorem applies to the complete edge after the incoming scalar
sum, with the common class at two consecutive edges held fixed
until their contributions are multiplied. The finite bounds
therefore add over edges. Averaging over \(w\), multiplying by
\(\log T\), and including both factors in the Gram bound gives
the prefactor \(2\log T/q_c\), as in
Corollary~\ref{cor:augheight501}. Lemma~\ref{lem:halfgauss504}
bounds each remaining average. There are exactly \(f\) places
\(v\) because \(T\) splits completely in \(F\).

For a matching pair one has \(g_w(\chi)+g_w(\chi^{-1})=1\)
at the same effective place, including for quadratic \(\chi\).
The event is the same at the paired edges, hence its mean
contribution summed over the two columns is \(1/2\).
The saving per column averaged over that pair follows after
the same factor of two in the Gram bound.
\end{proof}

\begin{corollary}[Bounds in terms of the scalar support of the prescribed input]\leavevmode
\label{cor:packetsave504}
For an independent list with separately zero-sum initial tables,
let \(s_{iv}\) be the number of distinct scalar characters
occurring at \(v\) in column \(i\), over a splitting field.
In the first bound of Theorem~\ref{thm:meangauss504} each
\(A(m_{iv})\) may be replaced, for a general input, by
\[
       \max\left\{\frac{1}{2(T-1)},\
                           \frac{1}{3136s_{iv}^{\,2}}\right\}.
\]
For a single supported character the larger bound of the
theorem applies directly.
\end{corollary}
\begin{proof}
The projectors onto scalar characters have denominators dividing
a power of \(T-1\), which is a unit at \(T\). They preserve the
integral lower bound of the input. Each component propagates
through the incidence condition with its scalar type as in
Theorem~\ref{thm:augterminal501}. For their sum use the minimum
of the gains over the \(s=s_{iv}\) supported characters at this
place. A nonzero column has \(s\ge1\) and augmentation excludes the
trivial character. Every \(g_w\) is at least \(1/(T-1)\), which
proves the first option after averaging over the half event.

For \(0<y<1\) and any character order \(m\ge2\), its marginal
distribution gives
\[
 \Pr\{g_w(\chi)\le y\}
  =\frac{\#\{1\le a\le my:(a,m)=1\}}{\varphi(m)}
  \le y+\frac{2^{\omega(m)}}{\varphi(m)} .
\]
Indeed M\"obius inversion replaces each integer part by its
argument with an error of at most one per squarefree divisor.
The probability is zero if \(my<1\). For all \(m\),
\(\sqrt{m}\,2^{\omega(m)}/\varphi(m)\le6\). At a prime power
\(q^a\) the factor is at most \(2\sqrt q/(q-1)\), which is
less than one for \(q\ge7\). The product of the three possible
factors at \(2,3,5\) is \(\sqrt{30}<6\).
Thus in both ranges \(\Pr\{g_w(\chi)\le y\}\le7\sqrt y\).
Set \(y=(28s)^{-2}\). The union of the exceptional events
for all supported characters has weight at most \(1/4\).
The \'etale event has weight \(1/2\), leaving weight at least
\(1/4\) where all the gains exceed \(y\). Its contribution to
the mean is at least \(y/4=1/(3136s^2)\). Linearity of the
average and the valuation inequality for the sum prove the claim.
\end{proof}

These bounds retain the original scalar supports and the same
integral projections. The archimedean coefficient term and the
first-row terms in Corollary~\ref{cor:terminalheight497}, as
well as independence for the columns needed in
Section~\ref{sec:missing}, remain in the comparison.


\subsection{The CM half in the terminal character average}
\label{sec:cmaverage508}

Use the actual CM field \(K=\QQ(\sqrt{-D})\) of the common period
construction, with \(D\) the auxiliary prime congruent to seven
modulo sixteen. Its discriminant is \(-D\). Write \(\lambda\)
for its primitive quadratic Dirichlet character of conductor \(D\).
The CM embeddings of \(E_x\) extend one chosen embedding of \(K\)
by Theorem~\ref{thm:pol49}. The terminal prime is \(T=\pi\bar\pi\)
with \(\pi\in\mathcal O_K\), by Theorem~\ref{thm:common111}.
In particular the condition \(P_{\tau,w}\mid\pi\) in
Theorem~\ref{thm:terminalhalf496} depends only on the restriction
of the effective coefficient place \(w\) to \(K\). It specifies
one of the two primes of \(K\) above \(T\), the same one for
all \(\tau\) on this edge. Denote this event by \(\mathcal E\).

\begin{lemma}[The joint distribution of the scalar exponents and the CM half]\leavevmode
\label{lem:jointcm508}
At one place \(v\mid T\) let the nonempty support of a separately
zero-sum input consist of characters \(\rho^e\) for \(e\) in a
set \(I\subseteq\{1,\ldots,m-1\}\), where \(\rho\) generates
the group spanned by the support and has order \(m\mid T-1\).
Write \([k]_m\) for the least nonnegative residue modulo \(m\)
and put
\[
      H_I(a)=\min_{e\in I}\frac{[ea]_m}{m}
                      \qquad (a\in(\ZZ/m\ZZ)^\times).
\]
Then the mean \(\mathbb E_w[{\bf1}_{\mathcal E}
\min_{e\in I}g_w(\rho^e)]\) in the terminal estimate is
\[
 \begin{cases}
 \displaystyle\frac{1}{2\varphi(m)}
              \sum_{(a,m)=1}H_I(a), &D\nmid m,\\[6pt]
 \displaystyle\frac{1}{\varphi(m)}
              \sum_{\substack{(a,m)=1\\\lambda(a)=\epsilon}}H_I(a),
                                                    &D\mid m .
 \end{cases}
\]
The sums use one complete residue system. In the second line the
sign is determined at any single coefficient place \(w_0\)
on the chosen half by \(\epsilon=\lambda(mg_{w_0}(\rho))\).
If \(I=-I\bmod m\), either line equals the first expression.
\end{lemma}
\begin{proof}
Let \(a(w)=m g_w(\rho)\), a unit residue as in
Lemma~\ref{lem:halfgauss504}. The definition of the Gauss
valuation uses the negative of the scalar character exponent
modulo \(T-1\). Multiplying that exponent by \(e\) and taking
its positive fractional representative gives
\(g_w(\rho^e)=[ea(w)]_m/m\). These numerators are nonzero
since the characters in the support are nontrivial.

The prime \(T\) splits completely in both \(K\) and
\(\QQ(\zeta_m)\). Its primes in the compositum of the two
fields are thus indexed simply transitively by the Galois
group. Each receives the same total local-degree weight from
the coefficient field. The restriction to \(K\) fixes
\(\mathcal E\), and the restriction to \(\QQ(\zeta_m)\)
fixes \(a(w)\). If \(D\nmid m\), then \(D\) is unramified in
\(\QQ(\zeta_m)\), whereas it ramifies in \(K\). Hence the
intersection is \(\QQ\), and these two uniform choices are
independent. If \(D\mid m\), then \(K\) is the quadratic
subfield of \(\QQ(\zeta_D)\) cut out by \(\lambda\). For
example, the quadratic Gauss sum at \(D\) has square \(-D\)
and transforms by \(\lambda\). A Galois element sending the
prime at \(w_0\) to the prime at \(w\) fixes \(K\) exactly
when its exponent on \(\zeta_m\) is in \(\ker\lambda\).
The inverse of that exponent multiplies \(a(w_0)\) to give
\(a(w)\). This proves the coset formula and the sign rule.
For inverse-stable support \(H_I(-a)=H_I(a)\). The two cosets
are interchanged by \(-1\), since \(K\) is imaginary, proving
the last assertion. The argument uses the effective place
throughout, so it applies equally to a conjugated edge.
\end{proof}

\begin{theorem}[The exact mean for a pure scalar character at the terminal prime]\leavevmode
\label{thm:classgauss508}
For a character of order \(m>1\) the mean
\(G=\mathbb E_w[{\bf1}_{\mathcal E}g_w(\chi)]\) is \(1/4\)
when \(D\nmid m\). For \(D\mid m\) it is
\[
 G=\frac14-\frac{\epsilon h(K)}
                     {w_K\varphi(m)}
                   \prod_{\substack{q\mid m\\q\ {\rm prime},\,q\ne D}}
                              (1-\lambda(q)),
\]
where \(h(K)\) is the class number, \(w_K=|\mu(K)|\), and
\(\epsilon\) is as in the lemma for \(\rho=\chi\). In particular
every character of even order has \(G=1/4\). For every order
one has, uniformly in \(v\), the word and the edge,
\[
       |G-1/4|\le \frac{h(K)}{w_K(D-1)}
                        \ll\frac{\log(2D)}{\sqrt D}.
\]
\end{theorem}
\begin{proof}
Use \(I=\{1\}\) in the lemma. The mean of \(a/m\) over all
unit residues is \(1/2\) by symmetry, including \(m=2\).
For \(D\mid m\), replace the indicator of the coset by
\((1+\epsilon\lambda(a))/2\). Its correction to \(1/4\)
is \(\epsilon/(2m\varphi(m))\) times
\(\sum_{1\le a\le m,\ (a,m)=1} a\lambda(a)\).

Put \(B_\lambda=D^{-1}\sum_{a=1}^D a\lambda(a)\).
For \(R\) a multiple of \(D\), periodicity and
\(\sum_{a=1}^D\lambda(a)=0\) give
\(\sum_{b=1}^R b\lambda(b)=R B_\lambda\).
Inclusion-exclusion for the distinct prime divisors of \(m\)
other than \(D\) now gives
\[
   \frac1m\sum_{\substack{1\le a\le m\\(a,m)=1}}a\lambda(a)
      = B_\lambda\prod_{\substack{q\mid m\\q\ {\rm prime},\,q\ne D}}
                          (1-\lambda(q)).
\]
Indeed a squarefree divisor \(u\) of those primes contributes
\(\mu(u)u\lambda(u)\sum_{b\le m/u} b\lambda(b)/m
=\mu(u)\lambda(u)B_\lambda\).
The identity \(B_\lambda=-2h(K)/w_K\) is the imaginary
quadratic class number formula at zero. In detail,
\(L(0,\lambda)=-B_\lambda\) follows by splitting the Dirichlet
series into Hurwitz zeta functions and using
\(\zeta(0,a/D)=1/2-a/D\). The factorization
\(\zeta_K=\zeta L(\,\cdot\,,\lambda)\), \(\zeta(0)=-1/2\),
and the analytic class number formula
\(\zeta_K(0)=-h(K)/w_K\) give the asserted identity.

The prime two splits in \(K\), as \(-D\equiv1\bmod8\).
Therefore \(2\mid m\) kills the correction whenever \(D\mid m\).
For an odd \(m\) divisible by \(D\), each factor at \(q\ne D\)
has magnitude at most two, while \(\varphi(q^a)\ge2\) for
every odd prime power. The remaining denominator
\(\varphi(D^a)\) is at least \(D-1\). This gives the first
uniform bound. For the second, Minkowski gives an integral
ideal of norm at most \(2\sqrt D/\pi\) in each class of \(K\).
The number of quadratic ideals of norm \(b\) is at most the
divisor function of \(b\), by comparison of the split, inert
and ramified Euler factors with those of \(\zeta(s)^2\).
Summing up to \(R\) gives at most \(R(1+\log R)\) ideals,
by counting pairs of positive integers with product \(\le R\).
This bounds \(h(K)\) as required.
\end{proof}

\begin{corollary}[The terminal contribution on the original common intervals]\leavevmode
\label{cor:cmcomp508}
For an independent list of \(q_c\) divided projected columns
whose input is in one nontrivial scalar character at every
terminal place, one may subtract at least
\[
    \frac{f\log T}{q_c}
       \left(\frac12-\frac{2h(K)}{w_K(D-1)}\right)
                                    \sum_{i,j}1
\]
from the upper bound of Corollary~\ref{cor:terminalheight497}.
If every order is even or is prime to \(D\), the coefficient
in parentheses may be taken to be \(1/2\). For the family
of Theorem~\ref{thm:conductor232} it is at least \(1/2-o(1)\)
uniformly as the intervals tend to infinity. For a general
separately zero-sum input the allowable subtraction on any
block is obtained by multiplying the finite sum in
Lemma~\ref{lem:jointcm508} by \(2\log T/q_c\).
\end{corollary}
\begin{proof}
Apply Theorem~\ref{thm:augterminal501} at each effective
coefficient place and average the gains as in
Theorem~\ref{thm:meangauss504}. For the pure inputs use the
exact means just calculated. There are \(f=[E_x^+:\QQ]\)
blocks per edge, and the two period factors in the determinant
give the stated multiplier. For a sum of characters use the
integral scalar idempotents at \(T\) and the smallest of the
supported gains on that block. The joint-distribution lemma
evaluates its average. In the selection in
Theorem~\ref{thm:conductor232}, \(D\equiv-1\) modulo all odd
primes up to a parameter tending to infinity, so the prime
\(D\) tends to infinity as well. The preceding uniform
class number estimate proves the limiting statement on these
same fields.
\end{proof}

This determines the averaged scalar-character factor of
the terminal Gauss sums. The first-row derivative costs, the
archimedean coefficient matrices and the required projected rank
still enter the bound for the full period determinant.


\subsection{Compact frame pairings in the projected period}
\label{sec:framepair509}

At a fixed archimedean embedding, the upper bound of
Theorem~\ref{thm:packetgram497} uses the Fischer inner products of
the polynomials in the differential variables. We compute these
inner products with the initial forms and the terminal tests kept
in the formula. This extends the scalar calculation at the identity
frame in Theorem~\ref{thm:factor83}. It also gives cross terms
between paths with different intermediate coefficient maps.

For one embedding of the CM type let \(r\le s\) be consecutive
ranks. Give the space \(H^-\) of input flag polynomials the norm
\[
                  \|Q\|^2=\int_{U(r)}|Q(u)|^2\,du ,
\]
with Haar probability, and use the corresponding norm in rank
\(s\). Restriction to the unitary group determines a polynomial,
by Zariski density. A determinant twist on this group preserves
the norm. If the input has column weight \(\lambda\), let \(n\)
be the twist giving the partition \(\nu=\lambda+(n,\ldots,n)\)
used in the differential symbol. The function \(Q\det^n\)
extends to a polynomial of this column weight on the full
matrix space. In \eqref{eq:sub81}, view the differential matrix
as an independent \(s\)-by-\(r\) matrix \(Z\). Its polynomial
operator is
\[
   (C(Z)Q)(Y)=(Q\det^n)(Z^{\mathsf t}Y_{[r]}).
\]
Write \(H^+\) for the compact coefficient space of column weight
\((\nu,0,\ldots,0)\) in rank \(s\), and \(d_+=\dim H^+\).
For several CM embeddings use the tensor products of these spaces
and maps, with a separate \(Z_\tau\) at each embedding. Set
\[
   c=\prod_{\tau}\prod_{h=1}^r(r-h+1)_{\nu_{\tau,h}},
   \qquad (a)_k=\prod_{l=0}^{k-1}(a+l),\qquad (a)_0=1 .
\]
Throughout this subsection \(A^*\) is the Hermitian adjoint in
the stated metrics.

\begin{lemma}[The covariance of one polynomial substitution]\leavevmode
\label{lem:framecov509}
Give the entries of the matrices \(Z_\tau\) independent standard
complex Gaussian distributions of density \(\pi^{-1}e^{-|z|^2}\).
For every endomorphism \(A\) of \(H^-\), one has
\[
           \mathbb E\, C(Z)A C(Z)^*
                     ={c\over d_+}\operatorname{tr}(A)\,I_{H^+}.
\]
\end{lemma}
\begin{proof}
First use one embedding and \(A=QQ^*\). For \(Y\in U(s)\)
the columns of \(Y_{[r]}\) are orthonormal. Each row of
\(Z^{\mathsf t}Y_{[r]}\) is a standard Gaussian row of length \(r\)
by unitary invariance, and different rows are independent. Write
this square Gaussian matrix as \(uB\) by Gram--Schmidt, with
upper triangular \(B\) of positive diagonal. Conditional on
\(B\), left unitary invariance makes \(u\) Haar. Write \(t_h=B_{hh}\).
Conditional on the first \(h-1\) Gaussian columns, \(t_h^2\)
is a sum of \(r-h+1\) independent exponential variables of mean
one, by projecting the next column to the orthogonal complement.
Its moment of order \(\nu_h\) is \((r-h+1)_{\nu_h}\), as follows
from the generating function \((1-z)^{-(r-h+1)}\).
The flag identities under upper triangular substitution give
\[
          |(Q\det^n)(uB)|^2
                    =|Q(u)|^2\prod_{h=1}^r t_h^{2\nu_h}.
\]
Successive conditioning from the last column gives
\(\mathbb E |(C(Z)Q)(Y)|^2=c\|Q\|^2\), independent of \(Y\).
Integrating over \(Y\) gives the trace of
\(\mathbb E(C(Z)QQ^*C(Z)^*)\).

This operator commutes with left translation by \(U(s)\), since
changing \(Y\) to \(gY\) changes \(Z\) to \(g^{\mathsf t}Z\),
which has the same law. The column-weight space \(H^+\) is
irreducible under this left action. Indeed the Cauchy
decomposition of matrix polynomials, used in
Section~\ref{sec:chain81}, pairs one row module with the matching
column module. Selecting the right upper-unipotent invariants of
the given torus weight selects its unique highest line. Thus
Schur's lemma makes the operator scalar, and its trace fixes the
scalar. Polarization gives the assertion for \(Q_1Q_2^*\);
these maps span all endomorphisms. For several embeddings,
independence gives the tensor product identity on simple tensors
of endomorphisms, which span the full endomorphism space.
\end{proof}

Fix one of the words and the embedding \(\sigma\). Expand the
actual column as in Theorem~\ref{thm:packetgram497}, keeping the
two last edges' shared class or coupling coset as a single choice.
For a summand \(\gamma\), write \(x_\gamma\) for its initial
coefficient vector. At edge \(j\) denote the two compact spaces
by \(H_j^-,H_j^+\), the latter dimension by \(d_j\), and the
constant in the lemma by \(c_j\). Write
\(M_{\gamma j}:H_j^+\to H_{j+1}^-\) for the intervening
coefficient transport, and \(L_\gamma:H_m^+\to V\) for the
final linear map to the actual projected test space. Include the
scalar of this summand in \(L_\gamma\), including its phase,
class measure and any column division. One may use an identity
map at an intermediate vertex. Determinant characters on frames
are included in the maps or in the weights of the adjacent
spaces. Use conjugate spaces and conjugate variables on the
conjugated factors. More explicitly, conjugation of a product
conjugates every preceding coefficient and variable in that
product. Moving the bars to these factors gives linear maps
between the corresponding conjugate spaces and preserves the
Gaussian law and the adjoint formulas.

Group summands only when their ordered scalar elliptic derivative
functions and the unitary changes from product variables to
physical matrix variables agree on every edge, at this embedding.
Write \(a\) for such a group. Finer groups, including individual
summands, are allowed. The definition permits covectors and
adapted bases to vary throughout the original sum. In one group
the coefficient polynomial of path \(\gamma\) has the form
\[
 P_\gamma(Z)=L_\gamma C_m(Z_m)M_{\gamma,m-1}
               C_{m-1}(Z_{m-1})\cdots
               M_{\gamma1}C_1(Z_1)x_\gamma .
\]
In fact keeping the derivatives formal in the class formula
gives exactly the determinant multiplication and substitution
in \eqref{eq:sub81}. Expanding translations in the finite local
tests changes the scalar elliptic functions to which one applies
these derivatives. The remaining coefficient operations are the
initial evaluation, the class transports, and the terminal
lowering and projection of Theorem~\ref{thm:covectorpaths352}.
All are independent of the formal differential variables.
An orthonormal product basis from Lemma~\ref{lem:basis59} changes
the variables by a unitary matrix on the full edge space. One
common such change in a group preserves all its Fischer inner
products, even when it mixes matrix coordinates.

\begin{theorem}[The signed coefficient matrix for a group of paths]\leavevmode
\label{thm:framepair509}
Return to the product variables and write the sum of \(P_\gamma\)
in one group \(a\) as \(\sum_\alpha v_{a,\alpha}Z^\alpha\).
Its positive semidefinite coefficient matrix is
\[
 \begin{aligned}
 S_a&=\sum_{\alpha}\alpha!\,v_{a,\alpha}v_{a,\alpha}^*\\
 &=\frac{\prod_{j=1}^m c_j}{d_m}
   \sum_{\gamma,\delta\in a}
       \operatorname{tr}(x_\gamma x_\delta^*)
       \left(\prod_{j=1}^{m-1}
          \frac{\operatorname{tr}(M_{\gamma j}M_{\delta j}^*)}{d_j}
       \right)L_\gamma L_\delta^* .
 \end{aligned}
\]
In particular the sum retains the cross terms after the signed
projection. At a vertex where the two transports are the actions
of elements \(g_\gamma,g_\delta\) in the same unitary
representation of dimension \(d_j\), the trace factor there is
its character at \(g_\gamma g_\delta^{-1}\) divided by \(d_j\).
\end{theorem}
\begin{proof}
For a standard complex Gaussian entry,
\(\mathbb E[Z^b\bar Z^e]={\bf1}_{b=e}b!\) by polar integration.
Independence therefore identifies the first line with the
expectation of the outer product of the sum of the path
polynomials with itself. For a pair of summands start with
\(x_\gamma x_\delta^*\). Integrate \(Z_1\) by
Lemma~\ref{lem:framecov509}. Its value is a scalar multiple
of the identity of \(H_1^+\). Transport on the two sides
changes that identity to \(M_{\gamma1}M_{\delta1}^*\). Its
trace enters upon integrating \(Z_2\). Each subsequent
integration performs the same operation, since the edge
variables are independent. The last transports take the last
identity to \(L_\gamma L_\delta^*\). Multiplying the constants
and summing over the ordered pairs proves the formula. For a
unitary action its adjoint is the action of the inverse, giving
the character assertion.
\end{proof}

\begin{corollary}[A bound for the prescribed determinant without derivative indices]\leavevmode
\label{cor:frameheight509}
For each column \(i\) attach its index to the preceding data.
Use \(Q_{i,\sigma}\) and \(d_i^T\) from
Theorem~\ref{thm:packetgram497}. Put
\[
  A_{i,\sigma}=\sum_\gamma
    \|x_\gamma\|\,
    \frac{\|L_\gamma\|_{\rm HS}}{\sqrt{d_{im}}}
    \prod_{j<m}\frac{\|M_{\gamma j}\|_{\rm HS}}{\sqrt{d_{ij}}},
   \qquad F_i=\prod_{j=1}^m c_{ij}.
\]
The endpoint norm uses the given metric on \(V\). For any
positive probabilities \(\mu_i(a)\) on the groups of each
column, the bound of Theorem~\ref{thm:packetgram497} holds with
\(B'_i=\sum_a S_{ia}/\mu_i(a)\), computed by
Theorem~\ref{thm:framepair509}. In particular the actual Gram
matrix at the prescribed tests satisfies
\[
       \det G^\circ_\sigma
         \ \le\ \prod_i
             \bigl(T^{-d_i^T}Q_{i,\sigma} F_i A_{i,\sigma}^2\bigr).
\]
Both assertions apply on the original period fields for every
translate and every word. Furthermore, for every fixed \(B>0\),
uniformly on the intervals of Theorem~\ref{thm:primitivegram463},
\[
  \log F_i
   =\sum_{j,\tau,h}\nu_{ij,\tau,h}
                  (\log\nu_{ij,\tau,h}-1)
                    +o_{\psi,\varepsilon,B}(b/x^B),
\]
where a zero part contributes zero.
\end{corollary}
\begin{proof}
Within one group the scalar elliptic germs are common. The
Cauchy--Schwarz estimate for their derivatives in the proof of
Theorem~\ref{thm:packetgram497} therefore applies with the
summed coefficient polynomial. Giving each group its own
probability yields the first assertion by the same operator
inequality and monotonicity of the mixed determinant.

For a single path take the trace of the theorem's identity.
Its square root is \(\sqrt{F_i}\) times precisely this path's
summand in \(A_{i,\sigma}\). The jet bound on that path multiplies
this square root by at most \((T^{-d_i^T}Q_{i,\sigma})^{1/2}\).
Sum over the paths by the triangle inequality for the norm in
\(V\). Hadamard's inequality for their resulting column norms
proves the displayed bound, including dependent lists.

For \(a\ge1\) and \(k\ge0\),
\((a)_k=k!\binom{a+k-1}{a-1}\). The binomial factor lies
between one and \((a+k)^{a-1}\). Comparison of the increasing
function \(\log t\) with its integrals over successive unit
intervals gives \(\log(k!)=k\log k-k+O(\log(2+k))\) for \(k\ge1\).
Thus each log factor differs from the asserted main term by
\(O(a\log(2+a+k))\), including \(k=0\). The number of
embeddings and edges and their ranks are bounded by powers of
\(x\). All parts satisfy \(k\ll xb+x^C\) on the given words,
by Lemma~\ref{lem:wordfirstrow450} and the partition ordering.
Since \(b=\exp(x^\varepsilon)\), the total error is bounded
by a fixed power of \(x\), which gives the stated uniform error.
\end{proof}

The factorials in this estimate are explicit for the signed
columns at the original tests. The endpoint maps, initial forms
and class sums still affect their height. Section~\ref{sec:missing}
also requires a lower bound for the rank of the evaluated periods
and the uniform cohomology estimate on the same intervals.

\section{Common divisors of the projected columns}\label{nav:chapter-projected-period-divisors}

\subsection{A common divisor for the signed theta columns}
\label{sec:commondiv510}

Use the polynomial-conductor intervals of
Theorem~\ref{thm:primitivegram463}. Recall \(b=\exp(x^\varepsilon)\),
\(X'=\exp(x^{\varepsilon/3})\) and \(w''=b/(nX')\) from the
construction of \eqref{eq:letters111}. At a fixed \(x\) and
translate \(u\), let \(w\) run through the chosen binary words.
Write \(\nu_{w,j,\tau,l}\) for part \(l\) of the partition on
edge \(j\) at the embedding \(\tau\). The two ranks on that
edge are \(r_j\le s_j\), and \(r_*=\max_j r_j\).
Let \(\mathcal P_x\) consist of the unramified primes \(p\ge5\)
in \(E_x\) which are inert in the quadratic field \(K_d\).
Put, as in Lemma~\ref{lem:periodinteger463},
\[
  g_{w,p}=\sum_j V_p(\nu_{w,j}),\qquad
  b^{\rm per}_{u,w}=\prod_{p\in\mathcal P_x}
                  p^{\,g_{w,p}-\lceil U_p\rceil}.
\]
Here \(U_p\) depends on \(u\) and is zero in the nonsignal case
specified in Corollary~\ref{cor:inertcorrection461}. The exponents
are nonnegative. We determine the logarithm of their common
divisor up to an error smaller than \(b/x^A\) for every fixed
\(A>0\).

\begin{lemma}[A uniform cutoff for all row gains]\leavevmode
\label{lem:rowcut510}
For all sufficiently distant intervals, throughout each interval
and for all translates and words,
\[
     g_{w,p}=0\qquad
        \bigl(p\in\mathcal P_x,\quad p>P_x=12r_*nxX'\bigr).
\]
For \(p\le P_x\) and any two words at the same translate one has
\(\lvert g_{w,p}-g_{w',p}\rvert\ll_{\psi,\varepsilon}x^C\)
with an exponent independent of \(p,u,w,w'\).
\end{lemma}
\begin{proof}
The first-row formula in Lemma~\ref{lem:wordfirstrow450} is
\(\nu_{w,j,\tau,1}=u_1(\tau)+e_{j,\tau}-(r_j+s_j-1)/2\).
The least base value in \eqref{eq:letters111} and its following
weight formulas is at least \(w''\), the rank shifts are bounded
by a power of \(x\), and \(u_1(\tau)\ge0\). The largest \(u_1\)
is at most \(\lceil xb\rceil+O(k)\), and \(e_{j,\tau}\le2b\)
eventually. Since \(w''\) grows faster than any fixed power of
\(x\), these facts give
\[
   \min_{w,j,\tau}\nu_{w,j,\tau,1}\ge b/(2nX'),
           \qquad \max_{w,j,\tau,l}\nu_{w,j,\tau,l}\le 3xb .
\]
For two words at this translate every part at the same position
differs by at most \(\Delta\). Indeed the only varying input
entries are the interior \(u_i\), each changing by \(\Delta\)
when its bit changes. Subsequent steps add scalars independent
of the word to the old entries and pad by entries independent
of the word. This is the recurrence for the ordered weights
in the proof of Lemma~\ref{lem:wordfirstrow450}.

Fix a single edge of ranks \(r,s\). For one of the allowed
choices \(H\) of alternating halves of the even cycles of
\(\pi=\sigma_p c_E\), let \(\mu\) lie in the opposite halves.
For each positive odd \(a\) write \(h=\pi^{-a}\mu\). The
list for \(\mu\) in the definition of \(V_p\) has entries
\(D_{l,a}=d_{p^a}(\nu_{h,l},s)\), \(1\le l\le r\).
For integers \(M\ge2,\ s\ge1,\ t\ge0\),
\[
       0\le d_M(t,s)=\max(0,\lceil(t-s(M-1))/M\rceil)
                           \le\lfloor t/M\rfloor .
\]
The upper bound follows by writing \(t=Mq+l'\) with
\(0\le l'\le M-1\le s(M-1)\). Thus the sum of the whole
list is at most \(3rxb\,p/(p^2-1)\le 6rxb/p\).
Every nonempty choice of list entries in the maximum must pay
at least \(\nu_{\mu,1}\) from the opposite partition. For
\(p>P_x\) this already exceeds the sum of the list. All
maxima are zero, proving the cutoff.

For the variation bound write each sorted-list maximum as a
maximum over subsets \(I\) of
\(\sum_{(l,a)\in I}D_{l,a}-\sum_{l\le |I|}\nu_{\mu,l}\),
where parts beyond \(r\) are zero. Use the same finite list
of indices \(p^a\le3xb\) for both words, padding by zeros.
The change of \(D_{l,a}\) has absolute value at most
\(\lceil\Delta/p^a\rceil\), by the ceiling formula.
The change of the subtracted sum is at most \(r\Delta\).
Consequently every subset value changes by at most
\[
   r\Delta+
   r\!\sum_{\substack{a\ {\rm odd}\\p^a\le3xb}}
                  \lceil\Delta/p^a\rceil
     \ \le\ r\bigl(\Delta+2\Delta/p+1+\log(3xb)/\log p\bigr).
\]
Taking maxima preserves a uniform absolute difference bound.
Sum over \(\mu\), take the maximum over \(H\), and sum over the
edges. The numbers of embeddings and edges, the ranks,
\(\Delta\) and \(\log(3xb)\) are bounded by powers of \(x\).
\end{proof}

There is also a formula without ceilings for the main term.
For one cycle of even length \(e\) and a given opposite
embedding \(\mu\), use the infinite list
\[
             t_{l,a}=\nu_{\pi^{-a}\mu,l}/p^a
                \quad(1\le l\le r,\quad a\ge1\ \text{odd}).
\]
Write \(T_k\) for the sum of its \(k\) largest entries,
\(T_0=0\), and \(T_\infty\) for the sum of all its entries.
Set \(A_k=\sum_{l\le k}\nu_{\mu,l}\) for \(0\le k\le r\).
Define \(\overline V_p(\nu)\) by taking the same sum over
opposite embeddings and maximum over \(H\) as for \(V_p\),
using at each \(\mu\) the rational number
\[
       \max\left\{\max_{0\le k<r}(T_k-A_k),\
                                    T_\infty-A_r\right\}.
\]
For an empty collection of even cycles take both gains to be zero.

\begin{proposition}[A finite evaluation of the main term]\leavevmode
\label{prop:continuousrows510}
In this formula one has
\[
    T_\infty=\frac{\displaystyle
       \sum_{\substack{1\le a<e\\ a\ {\rm odd}}}
                      p^{-a}|\nu_{\pi^{-a}\mu}|}
                      {1-p^{-e}} .
\]
For \(1\le k<r\), \(T_k\) is obtained from the entries with
\(a<ke\). For the actual words,
\(\overline V_p(\nu)=0\) for \(p>P_x\), and, for all \(p\),
\[
       0\le \overline V_p(\nu)-V_p(\nu)
          \le f r\bigl(s(1+\log(3xb)/\log p)+2\bigr),
\]
where \(f=[E_x^+:\mathbb Q]\).
\end{proposition}
\begin{proof}
Group the terms by \(a\bmod e\). Each such residue and row gives
a geometric progression of ratio \(p^{-e}\), proving the sum
formula. A positive entry beyond its first \(k\) terms has \(k\)
larger predecessors in the same progression. Zero progressions
can be ignored or padded by zeros. This proves the finite rule
for \(T_k\). Among subsets of size at least \(r\), the penalty
is always \(A_r\), and their sums increase to the sum of the
entire list. This proves that the displayed maximum is the
supremum over all finite subsets. The cutoff proof just given
also bounds the sum of this larger list by \(6rxb/p\).

Increasing a nonnegative summable list by a total amount \(a_0\)
increases the supremum by at most \(a_0\). For \(p^a\le3xb\),
each \(t_{l,a}\) exceeds \(D_{l,a}\) by at most \(s\), by the
ceiling formula. For \(p^a>3xb\) the old entries are zero. If
\(J=\lfloor\log(3xb)/\log p\rfloor\), the sum over \(a>J\)
for any one row is at most \(3xb\,p^{-J-1}/(1-p^{-1})<2\).
Adding these bounds for at most \(f\) lists of \(r\) rows
proves the inequality for each orientation and hence for
their maxima.
\end{proof}

\begin{theorem}[One integer for every theta word at a translate]\leavevmode
\label{thm:wordgcd510}
At each \(x,u\) put \(B_u=\gcd_w b^{\rm per}_{u,w}\).
There is a positive integer \(D_x\), valid for every translate,
such that
\[
      B_u\mid b^{\rm per}_{u,w}\mid B_uD_x,\qquad
                  \log D_x\ll_{\psi,\varepsilon}x^C X'.
\]
For any reference word \(w_0\) in the family the common divisor
satisfies the uniform estimate
\[
  0\ \le\
     \sum_{\substack{p\in\mathcal P_x\\p\le P_x}}
         \left(\sum_j\overline V_p(\nu_{w_0,j})-U_p\right)\log p
         -\log B_u
               \ \ll_{\psi,\varepsilon} x^C X'.
\]
On enlarging \(D_x\) within this bound, the explicit compact
factorials of Corollary~\ref{cor:frameheight509} also have a
common factor \(F_u^-\) with \(F_u^-\mid F_w\mid F_u^-D_x\).
\end{theorem}
\begin{proof}
For \(p\in\mathcal P_x\), set the exponent of \(p\) in \(D_x\)
equal to the maximum over translates of
\(\max_w g_{w,p}-\min_w g_{w,p}\). Give other primes exponent
zero. By Lemma~\ref{lem:rowcut510} each exponent is polynomially
bounded in \(x\) and the support is contained in \(p\le P_x\).
The elementary bound \(\sum_{p\le P_x}\log p\le P_x\log(2P_x)\)
gives the required logarithmic estimate, with an enlarged \(C\).
Subtracting the same \(\lceil U_p\rceil\) from every exponent
at a translate proves the two divisibilities.

The exponent of \(p\) in \(B_u\) is
\(\min_w g_{w,p}-\lceil U_p\rceil\). The difference from the
coefficient of \(\log p\) in the displayed main term is the sum
of three nonnegative quantities,
\[
  \sum_j(\overline V_p(\nu_{w_0,j})-V_p(\nu_{w_0,j})),\qquad
   g_{w_0,p}-\min_w g_{w,p},\qquad \lceil U_p\rceil-U_p .
\]
The first is bounded by a power of \(x\) by the proposition,
the second by the lemma, and the third by one. At primes above
the cutoff both gains and \(U_p\) vanish, since
\(0\le\lceil U_p\rceil\le g_{w,p}\) there. Summation proves
the absolute estimate.

For the compact factor take componentwise minima and maxima of
the parts over words and form \(F_u^-,F_u^+\) from the same
rising factorial products. Each \(F_w/F_u^-\) divides
\(F_u^+/F_u^-\). At any position this last ratio introduces at
most \(\Delta\) positive integer factors, all at most
\(r_*+3xb\). There are polynomially many positions and
translates, by Proposition~\ref{prop:field12} and the rank
bounds in Corollary~\ref{cor:fixedheight479}. Multiplying the
previous \(D_x\) by these ratios for all translates increases
its logarithm by at most a polynomial in \(x\).
\end{proof}

\begin{corollary}[The common division on each projected minor]\leavevmode
\label{cor:commonminors510}
Let \(Y^\circ\) have columns \(\widetilde Y_i/b^{\rm per}_{u,w_i}\)
and let \(Y^B\) have columns \(\widetilde Y_i/B_u\), with the
same integral period evaluations and the same projection. In
any one system of test coordinates at any finite place, write
\(I_h\) for the fractional ideal generated by the size-\(h\)
minors. For every \(h\ge1\),
\[
   D_x^h I_h(Y^\circ)\ \subseteq\ I_h(Y^B)\
                                 \subseteq\ I_h(Y^\circ).
\]
For every independent list of \(q_c\) columns in the original
positive test metric and at every archimedean embedding,
\[
    0\le {1\over q_c}
       \log\frac{\det G^B_\sigma}{\det G^\circ_\sigma}
          \le 2\log D_x=o_{\psi,\varepsilon,A}(b/x^A)
                 \qquad (A>0\text{ fixed}).
\]
\end{corollary}
\begin{proof}
The change multiplies column \(i\) by the positive integer
\(b^{\rm per}_{u,w_i}/B_u\) dividing \(D_x\). Each minor is
therefore multiplied by a divisor of \(D_x^h\), proving both
ideal inclusions even when the ideal is zero. The Gram
determinant is multiplied by the square of the product of the
column factors. Take logarithms and use
\(x^{A+C}X'/b\to0\). The row-gain construction and the common
integer use the fields and weights of these very columns.
\end{proof}

This supplies a division compatible with taking linear
combinations across words. The finite expression in the theorem
computes its size to the precision stated. The comparison of
the remaining signed period height with the prime count, its
required rank, and the exceptional strict lengths are still
needed in Section~\ref{sec:missing}.


\subsection{Large primes in the common period divisor}
\label{sec:divprimes511}

Continue with the common integer \(B_u\) of
Theorem~\ref{thm:wordgcd510}. We evaluate the contribution of the
primes above a fixed power of \(x\) using the Frobenius group of
the same period field. Put \(G=\operatorname{Gal}(E_x/\mathbb Q)\),
\(d_G=|G|\), and let \(S_x\) contain the ramified rational primes
and \(2,3\), with no other primes. Write \(N_x\) for the conductor
of the field. On these intervals \(N_x,d_G\) are bounded by
powers of \(x\), by Theorem~\ref{thm:polycenter460} and the
field construction. For \(g\in G\) put
\[
 A_g(t)=\sum_{\substack{p\le t,\ p\notin S_x\\
                         \operatorname{Frob}_p=g}}\frac{\log p}{p},
                \qquad y_0=\exp((\log(2x))^4).
\]

\begin{lemma}[Reciprocal prime sums in the period field]\leavevmode
\label{lem:fieldmertens511}
Uniformly for \(t\ge y_0\) and \(g\in G\),
\[
       A_g(t)=\frac{\log t}{d_G}+C_g+
                         O_{\psi,\varepsilon}(e^{-c\sqrt{\log t}})
\]
with some \(c>0\) independent of \(x,u,g\). If \(\chi^*\) induces
the character \(\chi\) of \(G\), define \(D_{\bf1}=-\gamma\)
and \(D_\chi=-L'(1,\chi^*)/L(1,\chi^*)\) for \(\chi\ne{\bf1}\)
in this lemma. The constant is
\[
 C_g=\frac{1}{d_G}\sum_{\chi\in\widehat G}\overline{\chi(g)}
 \left(D_\chi-\sum_{p\in S_x}\frac{\chi^*(p)\log p}{p}
       -\sum_p\frac{\chi^*(p)^2\log p}{p(p-\chi^*(p))}\right).
\]
\end{lemma}
\begin{proof}
The real primitive characters with an auxiliary factor have no
exceptional zero near one by Theorem~\ref{thm:long232}. The other
real primitive characters have conductor dividing the fixed \(8F\).
Thus the zero-free bound, the zero count and the explicit formula
in the proof of Lemma~\ref{lem:far14} apply uniformly on these
fields. To check the lower endpoint here, truncate the zero sum
in \eqref{eq:uncond12} at \(\exp(\sqrt{\log t})\).
The low-zero contribution after division by \(t\) is at most
\[
  O_{\psi,\varepsilon}\!\left(
     \log^3(2N_x)
      \exp\!\left(-c_1\min\!\left\{
            \sqrt{\log t},\,\frac{\log t}{\log(2N_x)}
                      \right\}\right)\right).
\]
The tail is \(O((\log(2N_x)+\sqrt{\log t})e^{-\sqrt{\log t}})\)
by the inverse-square ordinate bound of that lemma. Its
elementary terms after division by \(t\) are
\(O_{\psi,\varepsilon}((1+\log t+\log^2(2N_x))/t)\), including
the error from the shifted line. Since \(\log(2N_x)\ll\log(2x)\)
and \(\log t\ge(\log(2x))^4\), these bounds give
\(\lvert G_{\chi^*}(t)\rvert\ll t e^{-c_2\sqrt{\log t}}\)
for some \(c_2>0\).

For \(h=t e^{-c_2\sqrt{\log t}/3}\), take the forward difference
quotient of \(R_{\chi^*}\) from \eqref{eq:Gchi}. It differs from
\(\sum_{n\le t}\Lambda(n)\chi^*(n)/n\) by at most
\((h+1)\log(2t)/t\), by bounding the additional terms with
\(t<n\le t+h\). The two \(G_{\chi^*}\)-values in the quotient
contribute \(O(e^{-c_3\sqrt{\log t}})\) for a smaller \(c_3>0\).
The difference quotient of the main term at the trivial
character is \(\log t-\gamma+O(h/t)\), and for the nontrivial
characters it is exactly \(D_\chi\). To pass to unramified primes,
remove all powers of exponent at least two and the first powers
in \(S_x\). Every prime in \(S_x\) is already below \(y_0\) for
large \(x\). The sum of all higher powers of one prime gives
the last summand of the formula for \(C_g\). The part with
\(p^a>t,\ a\ge2\), in absolute value, is \(O(\log(2t)/\sqrt t)\).
For \(p\le\sqrt t\) the geometric tail is at most \(2\log p/t\)
per prime. For larger \(p\) sum \(2\log p/p^2\) by comparison
with the integral over all integers. Finally character
orthogonality selects \(\operatorname{Frob}_p=g\). Averaging
over \(d_G\) characters preserves the absolute uniform error.
\end{proof}

Choose any one reference word at the fixed translate. Set
\(S_{j,\tau,k}=\sum_{l=1}^k\nu_{w_0,j,\tau,l}\),
\(S_{j,\tau,0}=0\), and \(S_{j,\tau}=S_{j,\tau,r_j}\).
For \(g\) inert in \(K_d\) write \(\pi_g=g c_E\) acting on
the embedding labels as in Theorem~\ref{thm:cycleweyl461}.
For each even cycle \(C\) of this permutation, let \(C_0,C_1\)
be its two alternating halves. Put
\[
 \begin{aligned}
 f_{j,g,\mu}(z)
   &=\max_{0\le k\le r_j}
             \{S_{j,\pi_g^{-1}\mu,k}z-S_{j,\mu,k}\},\\
 H_g(z)
   &=\sum_{j,C}\max_{b'=0,1}\sum_{\mu\in C_{b'}}
                                       f_{j,g,\mu}(z)\\
   &\quad-{\bf1}_{g|K_q=1}
        \sum_{j,t\in T_{\rm hi}}[a z-e_{j,\pi_g t}]_+ .
 \end{aligned}
\]
Here \(a=\lceil xb\rceil\), and \(j\) runs through all edges in
the last sum as well. Set \(H_g=0\) when \(g\) splits in \(K_d\).
Let \(L_u\) denote the distinguished small-place set of
\eqref{eq:letters111}, and let \(O_u\) consist of the embedding
labels outside both \(T_{\rm hi}\) and \(L_u\). For \(g\) inert
in \(K_d\) put
\[
 J_g=\sum_{j,C}\max_{b'=0,1}
           \sum_{\mu\in C_{b'}\cap O_u} S_{j,\pi_g^{-1}\mu}
       -{\bf1}_{g|K_q=1}\,a
           \sum_j\#\{t\in T_{\rm hi}:\pi_g t\notin L_u\} .
\]
Give \(J_g\) value zero for the other \(g\). All these weights use
the original selected field, high set, distinguished set and
translate.

\begin{proposition}[The three prime ranges of the divisor]\leavevmode
\label{prop:middleweights511}
There is an exponent \(C\), depending only on the fixed
parameters, for which \(\lvert J_g\rvert\ll bx^C\). For
\(z\ge0\), \(H_g\) is continuous and piecewise affine with
\(H_g(0)=0\) and Lipschitz constant at most \(O(bx^C)\).
For any sufficiently large fixed \(M\), uniformly for all \(g\)
and \(x^M<t\le y_0\),
\[
              H_g(1/t)=J_g/t+O_{\psi,\varepsilon}(bx^C/X').
\]
At primes \(p>x^M\) in \(\mathcal P_x\), the difference
between \(\sum_j\overline V_p(\nu_{w_0,j})-U_p\) and
\(H_{\operatorname{Frob}_p}(1/p)\) is \(O(bx^C/p^3)\).
\end{proposition}
\begin{proof}
Keep only \(a'=1\) in the lists \(t_{l,a'}\) of
Proposition~\ref{prop:continuousrows510}. The largest \(k\)
entries of the resulting list come from the first \(k\)
parts at \(\pi_g^{-1}\mu\). This gives \(f_{j,g,\mu}(1/p)\),
including the choice \(k=r_j\). The omitted entries have total
at most \(3r_jxb\,p^{-3}/(1-p^{-2})\) per opposite embedding.
The same bound with \(a p^{-3}/(1-p^{-2})\) applies to the
omitted higher odd powers of each pair in \(U_p\), because
all \(e_{j,\tau}\) are positive. The halves can be chosen
independently on different cycles. The supremum over subsets
and maximum over halves change by no more than the uniform
error. All multiplicities are polynomial in \(x\), giving the
assertion at primes.

The constants in the affine functions have nonpositive sign
before taking their maxima, and the choice \(k=0\) is available.
Each maximum is Lipschitz with constant bounded by its largest
slope. All slopes, summed over their indices, are \(O(bx^C)\)
by Lemma~\ref{lem:rowcut510}. The same bound applies to the
subtracted ramps, proving the assertions about \(H_g\) and \(J_g\).

At a high place the first part is at least \(b/2\) for large
\(x\), since \(u_1\) contains \(\lceil xb\rceil\). At a place
in \(L_u\) this lower bound follows from \(B_\tau\ge\lceil b\rceil\)
in the formula for \(e_{j,\tau}\). At the remaining places
the full partition size is \(O(bx^C/X')\), since
\(B_\tau\le\lceil b/(nX')\rceil+1\) there and the other
contributions to the first part are polynomial. These statements
use the rank shifts in Lemma~\ref{lem:wordfirstrow450}.
For \(t>x^M\) with \(M\) sufficiently large, each incoming
partition sum divided by \(t\) is less than \(b/2\).
Thus at the first two kinds of destination every \(f\) is zero.
At \(\mu\in O_u\), choosing \(k=r_j\) and dropping the penalty
give lower and upper bounds within \(S_{j,\mu}\) of
\(S_{j,\pi_g^{-1}\mu}/t\).

In the subtracted ramps the weights \(e_{j,\tau}\) on \(L_u\)
are at least \(b\). At the remaining places, including the
high places, \(e_{j,\tau}=O(bx^C/X')\) by the explicit base
values. Therefore the ramp vanishes when its destination is
in \(L_u\); at the other destinations it differs from \(a/t\)
by \(O(bx^C/X')\). Taking the orientations and the sums proves
the middle-range formula, with the polynomial factors absorbed
in the exponent \(C\).
\end{proof}

\begin{theorem}[The absolute size from polynomially many primes and field constants]\leavevmode
\label{thm:divtail511}
For every fixed \(A>0\) one can take a fixed \(M=M(\psi,\varepsilon,A)\)
and \(z_0=x^M\) such that, uniformly in the translate and the
reference word,
\[
 \begin{aligned}
 \log B_u
   ={}&\sum_{\substack{p\in\mathcal P_x\\p\le z_0}}
             (g_{w_0,p}-\lceil U_p\rceil)\log p\\
      &+\sum_{g\in G}J_g
                \left(\frac{\log y_0}{d_G}+C_g-A_g(z_0)\right)\\
      &+\frac1{d_G}\sum_{g\in G}
                           \int_{y_0}^{P_x} H_g(1/t)\,dt
                         +o_{\psi,\varepsilon,A}(b/x^A).
 \end{aligned}
\]
The integrals use only finitely many affine pieces and ordinary
logarithms. At the same precision the prime series in all the
\(C_g\) may be truncated at \(p\le z_0\), retaining the
logarithmic derivatives at one in Lemma~\ref{lem:fieldmertens511}.
\end{theorem}
\begin{proof}
Eventually \(z_0<y_0<P_x\). By Theorem~\ref{thm:wordgcd510}
the change from \(B_u\) to the divisor of the reference word
costs \(O(x^C X')\) in the logarithm. Keep its exact exponents
at \(p\le z_0\). At the other primes the replacement by
\(\sum_j\overline V_p-U_p\) costs \(O(x^C X')\) by the
nonnegative error estimates in the same theorem. On
\(z_0<p\le P_x\) the first-power approximation in the
proposition costs in total
\[
    O\!\left(bx^C\sum_{n>z_0}\frac{\log n}{n^3}\right)
                  = O\!\left(bx^C(1+\log z_0)/z_0^2\right).
\]
Take \(M\) large enough that this is \(o(b/x^A)\).

Up to \(y_0\), replacement by \(J_g/p\) costs at most
\(O(bx^C y_0\log(2y_0)/X')=o(b/x^A)\). This follows from
\(\log X'=x^{\varepsilon/3}\), which dominates every fixed
power of \(\log x\). The remaining middle sum is exactly
\(\sum_g J_g(A_g(y_0)-A_g(z_0))\). Substitute the lemma at
\(y_0\). Its error, even after summation over \(G\), is at
most \(bx^{C'}e^{-c(\log(2x))^2}=o(b/x^A)\).

On \([y_0,P_x]\) put \(W_g(t)=tH_g(1/t)\). The Lipschitz
estimate and \(H_g(0)=0\) give \(|W_g|\ll bx^C\) and
\(|W_g'(t)|\ll bx^C/t\) on each smooth piece. The sum over
primes in this last range is the Stieltjes integral of
\(W_g\) against \(A_g\) on \((y_0,P_x]\), summed over \(g\).
Partial summation with Lemma~\ref{lem:fieldmertens511} replaces
this by the integral in the statement. The constants \(C_g\)
cancel between endpoints and the integral of the derivative.
The error is bounded by
\[
       O\!\left(d_G bx^C(1+\log P_x)
                               e^{-c(\log(2x))^2}\right).
\]
The pieces are continuous at every junction, so partial
summation adds no interior boundary errors. On each piece of
\(H_g(1/t)\) the integrand is \(a_1+a_2/t\). For each
\(\mu\) the candidates are \(r_j+1\) lines in \(1/t\).
Their pairwise crossings give a finite subdivision. On each
resulting interval the two half-cycle sums cross at most once
unless they coincide. Taking the union of these subdivisions
uses a polynomial number of breakpoints in the numbers of
edges, embeddings and ranks.

Finally the absolute tail of the prime series in each \(C_g\)
after \(z_0\) is \(O((1+\log z_0)/z_0)\), by comparison
with the sum of \(\log n/n^2\). All primes of \(S_x\)
are below \(z_0\) upon enlarging \(M\). Another enlargement
absorbs this last error after multiplication by \(\sum_g|J_g|\).
This gives the claimed accuracy for the original common
divisor on every field of the interval.
\end{proof}

The arithmetic information for this divisor now occurs in the
primes of size at most \(x^M\), the field constants at one and
the explicit integrals. Its value alone gives neither the rank
of the signed periods nor the upper bound for their remaining
height in Section~\ref{sec:missing}.


\subsection{The terminal restriction and the two factorial norms}
\label{sec:twofactorials512}

At one translate use the incoming analytic row and the definite
word \(w\) of \eqref{eq:letters111}. Both use the same field,
embeddings and \(m\) edges. Two metric normalizations enter at the
last definite edge. We first identify their conversion in the
coefficient matrix of Theorem~\ref{thm:framepair509}.

\begin{lemma}[The endpoint in highest-vector normalization]\leavevmode
\label{lem:frameend512}
Let \(V_v\) be a compact flag module of dimension \(d_v\). Use
the Hermitian Fischer form in which its highest polynomial \(p_v\)
has norm one. The squared Haar frame norm of Section~\ref{sec:framepair509}
is \(1/d_v\) times this squared Fischer norm. This holds also after
a determinant twist and for products over the CM embeddings.

In the integral frames let \(L:V_l\to V_m\) be the actual terminal
restriction with its normal-coordinate scalar removed. Here the
indices \(l,m\) denote the two weights. Write \(d_l=\dim V_l\).
For endpoint tests \(R_\gamma,R_\delta:V_m\to V\) in the same
frames and a fixed metric on \(V\), one has
\[
  \frac{(R_\gamma L)(R_\delta L)^{*,{\rm Haar}}}{d_l}
      =s_w\,R_\gamma R_\delta^{*,{\rm F}},
  \qquad s_w=\prod_\tau c_\tau^2\kappa_\tau .
\]
The adjoints on the right use the normalized Fischer metric on
\(V_m\). The identity holds if both composites on the left are
precomposed with the same unitary frame change.
\end{lemma}
\begin{proof}
Evaluation at the identity is the inner product with \(p_v\)
by Lemma~\ref{lem:weyldual351}. The coefficients of \(p_v\)
are real, so the identity gives the Hermitian version with
the form linear in its first entry. Unitary row translation
preserves the inner product on the linear variables and hence
the Fischer metric on their symmetric tensors. Thus
evaluation on unitary frames pairs the vector with the orbit
of \(p_v\). The Haar average of the rank-one orthogonal
projections onto these unit translates commutes with the
unitary action. The flag space is irreducible, as in
Lemma~\ref{lem:framecov509}. The average has trace one, so it
equals \(I/d_v\). Pairing with an arbitrary vector gives the
first assertion. A determinant has absolute value one on the
unitary group. Fubini and polarization give the product
assertion, with \(d_v\) the product dimension.

For a map from \(V_l\) to \(V\), changing the domain metric
from Fischer to Haar multiplies its adjoint by \(d_l\).
Hence the quotient in the display equals
\(R_\gamma L L^{*,{\rm F}}R_\delta^{*,{\rm F}}\), with Fischer
metrics on both coefficient spaces. Theorem~\ref{thm:lowerlattice347}
identifies \(LL^{*,{\rm F}}\) with \(s_w I\). A common unitary
frame change cancels against its adjoint.
\end{proof}

For the endpoint factor in Corollary~\ref{cor:frameheight509},
taking the trace gives
\[
  \frac{\|R_\gamma L\|_{{\rm HS},{\rm Haar}}}{\sqrt{d_l}}
        =\sqrt{s_w}\,\|R_\gamma\|_{{\rm HS},{\rm F}} .
\]
The tests \(R_\gamma\) can contain the contraction at the
small class, its finite projection and the path scalar from
Theorem~\ref{thm:covectorpaths352}. A change of terminal frame
can be kept in the polynomial path or included in this
calculation when it is common to the two factors. In the
unnormalized column \(L_z=Z_w^{\rm norm}L\), so the squared
endpoint has the additional factor \(|\sigma Z_w^{\rm norm}|^2\).
In \(\widetilde Y_w=(z_w Z_w^{\rm norm})^{-1}Y_w\) it cancels.
The factors \(z_w^{-1}\) and \(1/b^{\rm per}_{u,w}\) still enter
the divided column once each. By Corollary~\ref{cor:singular349}
the absolute logarithm of \(s_w\), and the sum of its absolute
finite log valuations, are \(o(b/x^A)\) for every fixed \(A>0\),
even with suprema over the words and translates.

Let \(D(t)=t\log t-t\) for \(t>0\), with \(D(0)=0\).
Let \(F_w\) be the product of rising factorials in
Corollary~\ref{cor:frameheight509}. Define \(\mathcal H_{\rm in}\)
to be the product, over all real places and edges of the incoming
row, of \(H(\alpha)H(\beta)H(\gamma)\) in
Proposition~\ref{prop:fock187}. Thus each incoming hook product
is taken in Gaussian coordinates of dilation \(1/\pi\).
Write \(h=|T_{\rm hi}|\), \(h_2=|L_u|\),
\(a=\lceil xb\rceil\) and \(Y=X'=\exp(x^{\varepsilon/3})\).

\begin{theorem}[The factorials on the given row and word]\leavevmode
\label{thm:factorcomp512}
Set \(N_w=\sum_{j,\tau,l}\nu_{w,j,\tau,l}\), the total definite
degree. On the intervals of Theorem~\ref{thm:primitivegram463},
uniformly in \(u,w\), there is a fixed \(C\) such that
\[
  N_w=m(ha+kh_2b)+O_{\psi,\varepsilon}(bx^C/Y+x^C).
\]
For every fixed \(A>0\) one has
\[
 \begin{aligned}
  \log F_w&=m\{hD(a)+kh_2D(b)\}+o_{\psi,\varepsilon,A}(b/x^A),\\
  \log\mathcal H_{\rm in}
      &=m(k-1)h_2D(b)+o_{\psi,\varepsilon,A}(b/x^A),\\
  \log\frac{F_w}{(b^{\rm per}_{u,w})^2\mathcal H_{\rm in}}
      &=m\{hD(a)+h_2D(b)\}-2\log B_u
                               +o_{\psi,\varepsilon,A}(b/x^A).
 \end{aligned}
\]
Here \(B_u\) is the common integer in Theorem~\ref{thm:wordgcd510}.
\end{theorem}
\begin{proof}
At the first edge the old highest weights are \(u_l\) minus
the corresponding half-sum entries, of absolute difference
\(O(k)\) from \(u_l\). The shift at edge \(j\) adds
\(e_{j,\tau}-s_j/2\) to form its partition. For the next input
one pads by zeros and subtracts \(e_{j,\tau}-r_j/2\).
Each old entry consequently decreases by \((s_j-r_j)/2\).
A newly introduced entry begins at \(-e_{j,\tau}+r_j/2\).
This is the recurrence of Lemma~\ref{lem:wordfirstrow450}.
All rank shifts and \(jH_0\) are polynomial in \(x\), and
\(e_{j,\tau}=B_\tau+(2j-1)H_0\). It follows that at every edge
the original \(k\) parts are \(u_l(\tau)+B_\tau+O(x^C)\).
Every other part is \(O(x^C)\), because the two occurrences
of \(B_\tau\) cancel.

At a high place the top part equals \(a+O(bx^C/Y+x^C)\).
The remaining parts have size \(O(bx^C/Y+x^C)\), including
the part with initial entry \(-\lceil xb/Y\rceil-(k-1)/2\).
At a distinguished place the original \(k\) parts all equal
\(b+O(x^C)\). Outside these two sets all parts are
\(O(bx^C/Y+x^C)\). These bounds follow from the three explicit
base values after \eqref{eq:letters111} and its interior bit
entries. Summing the bounds proves the formula for \(N_w\)
after absorbing the polynomial number of indices into \(C\).

Apply the last assertion of Corollary~\ref{cor:frameheight509}.
Replacing a part listed above by \(a\), \(b\), or zero as
appropriate changes its \(D\)-value by \(O(bx^{C'}/Y+x^{C'})\).
Indeed the integral of \(|\log t|\) on \((0,1)\) is one,
and on \([1,3xb+2]\) the derivative is bounded by
\(\log(3xb+2)\). The number of terms is polynomial and
\(Y\) grows faster than every fixed power of \(x\). This proves
the first logarithmic estimate.

For the incoming hook products the proof of
Proposition~\ref{prop:fock187} gives exactly \(k-1\) parts
\(B_\tau+O(x^C)\) per edge and place. The other parts are
polynomially bounded. Its hook formula contributes polynomial
total error to \(m(k-1)\sum_\tau D(B_\tau)\). At the
distinguished places \(B_\tau=b+O(1)\). At all other places
\(B_\tau=O(bx^C/Y+x^C)\), whose contribution falls within
the stated error. Subtract the resulting two estimates
and use the uniform bound
\(|\log b^{\rm per}_{u,w}-\log B_u|=o(b/x^A)\) in
Theorem~\ref{thm:wordgcd510}. This proves the last formula.
\end{proof}

\begin{corollary}[An upper bound for the integral divisor in degree \(N_w\)]\leavevmode
\label{cor:degreebound512}
For all words and translates on these intervals,
\[
  0\le\log b^{\rm per}_{u,w}\le N_w(\log P_x+O(1)),
  \qquad
  \frac{2\log b^{\rm per}_{u,w}}{\log F_w}
                  =O_{\psi,\varepsilon}(x^{-2\varepsilon/3}),
\]
where \(P_x=12r_*nxY\).
\end{corollary}
\begin{proof}
In the subset formula for \(V_p\) drop all penalties and keep
all entries. At odd exponent \(j'\) each source place occurs
at most once, and each entry is at most
\(\nu_{w,j,\tau,l}/p^{j'}\) by Lemma~\ref{lem:rowcut510}.
Summing gives \(g_{w,p}\le N_w p/(p^2-1)\). The signal
subtraction is nonnegative, and the cutoff is \(P_x\).
For \(y\ge2\), the bound \(\sum_{p\le y}\log p/p\le\log y+O(1)\)
follows from \(\sum_{p\le n}\lfloor n/p\rfloor\log p\le\log(n!)\)
and \(\sum_{p\le n}\log p=O(n)\), with \(n=\lfloor y\rfloor\).
The difference \(p/(p^2-1)-1/p\) contributes a convergent
prime sum after multiplication by \(\log p\). This gives the
first estimate.

Let \(N\) be the polynomially bounded number of part positions.
Since one first part is at least \(b/(2nY)\) by
Lemma~\ref{lem:rowcut510}, we have
\(\log(N_w/N)\ge\log b-\log Y-O(\log x)\).
Each \((r-l+1)_v\) is at least \(v!\), whose logarithm is at
least \(D(v)\) by integral comparison (also for \(v=0\)).
Convexity now gives
\(\log F_w\ge N_w(\log(N_w/N)-1)\). Finally
\(\log b=x^\varepsilon\) and
\(\log P_x=x^{\varepsilon/3}+O(\log x)\). The second estimate
follows.
\end{proof}

The comparison above evaluates the polynomial factors at both
sets of prescribed highest weights. A relation between the full
period norms must also account for the class sums and for the
normalizations accompanying these polynomial factors.


\subsection{The leading term of the integral period divisor}
\label{sec:divconstant513}

On the same fields and intervals, put \(Y=X'\), let
\(o=\operatorname{ord}(\psi)\), and put
\[
             M_x=m k h_2 b,\qquad
             c_\psi=\frac12-\frac{1}{8\cdot 2^{v_2(o)}} .
\]
We compute the leading term of \(\log B_u\) by counting the
destinations of the distinguished places under Frobenius. This
also gives the sign and size of the comparison with the incoming
hook products in Theorem~\ref{thm:factorcomp512}.

Use \(G=\operatorname{Gal}(E_x/\mathbb Q)\).
For \(g\) inert in \(K_d\) write \(\pi_g=gc_E\) acting on the
CM type \(\Phi=\Phi_0\sqcup\Phi_1\), where
\(T_{\rm hi}\subset\Phi_0\) and \(L_u\subset\Phi_1\).
Write \(O_u=\Phi\setminus(T_{\rm hi}\cup L_u)\).
For each even cycle \(C\) of \(\pi_g\) let \(C_0,C_1\) be its
alternating halves. Define the integer
\[
 e_u(g)=\sum_{C\ {\rm even}}\ \max_{i=0,1}
    \#\{t\in L_u\cap C_i:\ \pi_g t\in O_u\}.
\]
Set \(e_u(g)=0\) when \(g\) splits in \(K_d\). In particular
\(0\le e_u(g)\le h_2\).

\begin{lemma}[The prime sum at the scale of the distinguished places]\leavevmode
\label{lem:divmean513}
Uniformly in \(u\) and the reference word on these intervals,
\[
       \frac{\log B_u}{M_x\log Y}
       =\frac1{|G|h_2}\sum_{g\in G} e_u(g)
                                      +o_{\psi,\varepsilon}(1).
\]
\end{lemma}
\begin{proof}
Write \(t_0=\exp((\log(2x))^4)\) and
\(t_1=Y\exp(-(\log(2x))^2)\). Eventually \(t_0<t_1<P_x\).
By Theorem~\ref{thm:rangenew50}, \(h\le n\),
\(h_2\ge n x^{-\varepsilon}\), and \(k\ge x^{9/2}\).
Thus \(mha/M_x=O(x^{1+\varepsilon}/k)=o(1)\).
The degree formula of Theorem~\ref{thm:factorcomp512}
gives \(N_w=(1+o(1))M_x\). By Theorem~\ref{thm:wordgcd510},
we may replace \(B_u\) by \(b^{\rm per}_{u,w}\), with
logarithmic error \(o(b/x)\).

At each contributing prime the exponent is between zero
and \(N_w p/(p^2-1)\) by the proof of
Corollary~\ref{cor:degreebound512}. At \(p\le t_0\) the
contribution is therefore \(O(M_x(1+\log t_0))\). At
\(t_1<p\le P_x\) it is \(O(M_x(1+\log(P_x/t_1)))\).
Indeed partial summation of the Chebyshev bound
\(\sum_{p\le t}\log p=O(t)\) gives
\(\sum_{y<p\le z}\log p/p\ll 1+\log(z/y)\).
Both bounds are \(o(M_x\log Y)\).

It remains to evaluate the middle range. Denote the
first-power weight of Proposition~\ref{prop:middleweights511}
by \(H_g\), and the full size of a partition on edge \(j\) at
\(\mu\) by \(S_{j,\mu}\). For \(t_0<p\le t_1\), all incoming
sizes divided by \(p\) are smaller than \(b/2\), since all
sizes are \(O(bx^C)\) for a fixed \(C\).
At any destination in \(T_{\rm hi}\cup L_u\) the first part
is at least \(b/2\) eventually by the weight calculation of
Theorem~\ref{thm:factorcomp512}. Its first-power maximum is
therefore zero. At a destination in \(O_u\), taking the full
list or dropping the penalty bounds its maximum between
\(S_{j,\pi_g^{-1}\mu}/p-S_{j,\mu}\) and
\(S_{j,\pi_g^{-1}\mu}/p\).

Each source in \(L_u\) has total size \(k b+O(x^C)\) at every
edge, by the same weight calculation. At a high source the
total is \(a+O(bx^C/Y+x^C)\). At every ordinary source or
destination it is \(O(bx^C/Y+x^C)\). For a fixed choice of
halves we may therefore keep just the distinguished sources
with ordinary destinations, giving each weight \(kb/p\).
The total error from all the high sources is at most
\(mha/p+O(bx^{C'}/Y+x^{C'})\), including the polynomial
number of positions. The other substitution errors and the
dropped destination penalties are \(O(bx^{C'}/Y+x^{C'})\).
The subtracted first-power signal is at most \(mha/p\).
These bounds hold for any choice of halves. Maximization on
each cycle at each edge now selects exactly \(e_u(g)\) such
sources per edge, since \(\pi_g\) switches the two halves.
We obtain
\[
 H_g(1/p)=\frac{m k b}{p}\,e_u(g)
       +O\left(\frac{mha}{p}+\frac{b x^{C'}}{Y}+x^{C'}\right)
       =\frac{m k b}{p}\,e_u(g)+o(M_x/p).
\]
The last error is uniform, since \(p/Y\le\exp(-(\log(2x))^2)\).

By Proposition~\ref{prop:middleweights511} the higher odd
powers add \(O(bx^C/p^3)=o(M_x/p)\) uniformly in this range.
Replacing the integer row gains and the signal ceiling by their
continuous versions costs \(O(x^C Y)\) in the total prime sum
by Theorem~\ref{thm:wordgcd510}: the two rounding errors there
are separately nonnegative before taking their difference, so
the bound also holds on this subinterval. Consequently the
middle sum differs from \(mkb\sum_{t_0<p\le t_1}
e_u(\operatorname{Frob}_p)\log p/p\) by \(o(M_x\log Y)\).
Here only the unramified primes outside \(S_x\) are used, as
in Lemma~\ref{lem:fieldmertens511}.

Apply that lemma to \(A_g(t_1)-A_g(t_0)\) for each \(g\).
The constant cancels. After multiplication by \(e_u(g)/h_2\)
and summation, the error is
\(O(|G|e^{-c(\log(2x))^2})=o(1)\), since \(|G|\) is polynomial
in \(x\). Finally \(\log(t_1/t_0)/\log Y\to1\), proving
the assertion.
\end{proof}

\begin{theorem}[The leading constant for the common integral divisor]\leavevmode
\label{thm:divconstant513}
On the intervals of Theorem~\ref{thm:primitivegram463} one has
\[
    \log B_u=(c_\psi+o_{\psi,\varepsilon}(1))M_x\log X',
    \qquad
    \log b^{\rm per}_{u,w}
             =(c_\psi+o_{\psi,\varepsilon}(1))M_x\log X',
\]
uniformly for every translate and every word. In particular
\(c_\psi=3/8\) for every character of odd order.
\end{theorem}
\begin{proof}
Write \(A=\operatorname{Gal}(K_\psi C_{j_x}/\mathbb Q)\) and
\(A_0=\operatorname{Gal}(K_\psi C_{j_x}/K_\psi)\). The two
quadratic coordinates are independent of \(A\) and of each
other by the ramification at \(q,d_0\) in
Theorem~\ref{thm:polycenter460}. Each sign pair occurs with
weight \(1/4\), and at fixed signs the translation in \(A\)
is uniform. Indeed multiplying by \(c_E\) adds a fixed element
to these coordinates, including a possible component in \(A\).
Identify each of \(\Phi_0,\Phi_1\) with \(A\). Then
\(|A|=n\), \(A_0\) is cyclic of two-power order, and \(A/A_0\)
is cyclic of order \(o\).

First take \(g\) inert in \(K_d\) and split in \(K_q\).
Then \(\pi_g\) switches \(\Phi_0,\Phi_1\).
Every cycle is even and its two halves are its intersections
with those sets. All distinguished sources can be selected,
except those sent into \(T_{\rm hi}\). Each such source is
sent uniformly through \(\Phi_0\) by the group average.
The conditional average of \(e_u(g)/h_2\) is therefore
\(1-h/n=1-o(1)\), by Theorem~\ref{thm:rangenew50}.

Now take \(g\) inert in both quadratic fields. Then \(\pi_g\)
preserves each of \(\Phi_0,\Phi_1\) and translates \(\Phi_1\)
by a uniform \(v\in A\). Only its action there matters.
Temporarily count all distinguished sources in the better
half of each even cycle, including those sent to \(L_u\).
Call the resulting count \(e^0(v)\). Removing those destinations
changes each maximum by at most the number removed, so
\[
      0\le e^0(v)-e_u(g)\le\#\{t\in L_u:t+v\in L_u\}.
\]
The average of the right side over \(v\) is \(h_2^2/n\), by
counting all ordered pairs of distinguished places. From the
construction in Theorem~\ref{thm:rangenew50}, the original
distinguished set has integer representatives in
\([0,|A_0|/(8L))\), with \(L\to\infty\). Hence
\(h_2/n=O(1/L+1/n)=o(1)\), also after translation.

Translate within \(\Phi_1\) to view \(L_u\) in \(A_0\).
That same construction makes it exactly balanced in the
\(D\) residue classes modulo \(D A_0\), with \(D\) a power
of two tending to infinity. This means that every such class
contains the same number of its elements. Fix the coset of
\(v\) in \(A/A_0\), and let \(r\) be its order there. Visits of
a \(v\)-cycle to \(A_0\) are spaced by \(r\) steps. If \(r\)
is even, every such cycle is even and all its visits to
\(L_u\) have the same parity. Thus \(e^0(v)=h_2\).

For odd \(r\), multiplication by \(r\) permutes \(A_0\), so
\(rv\) is uniform there. Outside a fraction \(1/D\), its
subgroup is \(2^j A_0\) with \(2^{j+1}\mid D\).
A \(v\)-cycle meeting \(A_0\) then meets it in one coset of
\(2^j A_0\) and is even. Successive visits differ by its
generator \(rv\) and are \(r\) steps apart. Their parities
therefore distinguish the two classes modulo \(2^{j+1}A_0\).
Balance modulo \(D A_0\) gives equal counts in these classes,
so \(e^0(v)=h_2/2\). On the exceptional fraction use
\(0\le e^0(v)\le h_2\).

The fraction of elements of odd order in the cyclic quotient
of order \(o\) is \(2^{-v_2(o)}\). Averaging the last two
cases and restoring the excluded destinations shows that the
conditional mean of \(e_u(g)/h_2\) when both signs are inert is
\[
     1-\frac{1}{2\cdot 2^{v_2(o)}}+O(1/D+h_2/n).
\]
The other two sign pairs contribute zero. The full mean is
\(c_\psi+o(1)\) uniformly in \(u\). Apply
Lemma~\ref{lem:divmean513} and Theorem~\ref{thm:wordgcd510}.
\end{proof}

\begin{corollary}[The integral divisor in the two norm factors]\leavevmode
\label{cor:signedfactor513}
For the original row and every prescribed word,
\[
   \log\frac{F_w}{(b^{\rm per}_{u,w})^2\mathcal H_{\rm in}}
       =-(2c_\psi+o_{\psi,\varepsilon}(1))\,M_x\log X' .
\]
For an embedding \(\sigma\) write \(A^0_{i,\sigma}\) for the
sum of path norm products in Corollary~\ref{cor:frameheight509}
on \(\widetilde Y_i\), before division by \(b^{\rm per}_{u,w_i}\).
For \(q_c\) words the actual Gram determinant of the divided
columns satisfies
\[
 \begin{aligned}
 \det G^\circ_\sigma\ \le{}&
 \exp\!\left(-(2c_\psi+o_{\psi,\varepsilon}(1))\,
                         q_c M_x\log X'\right)\\
 &{}\times\prod_{i=1}^{q_c}
       T^{-d_i^T}Q_{i,\sigma}\,\mathcal H_{\rm in}
                             (A^0_{i,\sigma})^2 .
 \end{aligned}
\]
The bound includes lists of deficient rank.
\end{corollary}
\begin{proof}
Since \(\log a\ll_\varepsilon x^\varepsilon\), the ratio of
\(m(hD(a)+h_2D(b))\) to \(M_x\log X'\) is
\(O(x^{2\varepsilon/3}(x h/h_2+1)/k)=o(1)\) by the size
inequalities at the start of Lemma~\ref{lem:divmean513}.
Theorems~\ref{thm:factorcomp512} and~\ref{thm:divconstant513}
now give the first formula. Each path norm scales by the column
divisor, so \(A_{i,\sigma}=A^0_{i,\sigma}/b^{\rm per}_{u,w_i}\).
Substitute into Corollary~\ref{cor:frameheight509} and sum
the uniform logarithmic error over the \(q_c\) factors.
\end{proof}

Section~\ref{sec:missing} still requires control of the hook,
elliptic and class factors together, sufficient rank at the
prescribed tests, and the strict local length bound.





\subsection{The size of the divisor from five Frobenius powers}
\label{sec:middlerangedivisor674}

At a node and a prime \(p\in\mathcal I_x\), decompose
\(\pi=\sigma_pc_E\) into cycles \(C\). Each cycle has even
length. Put \(T_C=C\cap T_{\rm hi}\), \(L_C=C\cap L_u\).
For either all positive odd \(j\), or the truncated set
\(J_9=\{1,3,5,7,9\}\), write \(J\) for the set of exponents
and define
\[
\begin{split}
 F_C(J)={}&\sum_{t\in T_C}
 \left(\sum_{\substack{j\in J\\ \pi^jt\notin L_C}}\frac{a}{p^j}
       +{\bf1}_{\pi t\in L_C}(a/p-b)_+\right),\\
 R_C(J)={}&\sum_{\mu\in C\cap\Phi_0}
  \left(kb\sum_{\substack{j\in J\\ \pi^{-j}\mu\in L_C}}p^{-j}
                                  -a{\bf1}_{\mu\in T_C}\right)_+ .
\end{split}
\]
We suppress \(J\) when all positive odd powers are included.
The first number uses the sources on \(\Phi_0\), the second
the sources on \(\Phi_1\). Both are explicit in the Frobenius
permutation and the two prescribed sets of embeddings.

\begin{theorem}[The maximizing orientation in the middle prime range]
\leavevmode\label{thm:cycleorientation674}
For all sufficiently large nodes and every \(p\in\mathcal I_x\),
every actual maximizing orientation for \(V_p(\nu_i)\) uses
\(\Phi_1\) as source on each cycle meeting \(L_u\). For every
fixed \(A>0\), uniformly in \(p,u,w\), the divisor of the word
in Lemma~\ref{lem:periodinteger463} satisfies
\[
 v_p(b^{\rm per}_{u,w})=
 m\sum_{C:\,L_C\ne\varnothing}(R_C-F_C)
                      +o_{\psi,\eps,A}(b/x^A).
\]
\end{theorem}
\begin{proof}
On one edge replace the partition at \(t\in T_{\rm hi}\) by
\((a,0,\ldots)\), the partition at \(\tau\in L_u\) by the
rectangle with \(k\) parts \(b\), and all other partitions by
zero. These model rows are allowed to have real entries.
The estimates in the proof of Theorem~\ref{thm:factorcomp512}
give a total absolute row error
\(O_{\psi,\eps}(b x^{C_0}/X'+x^{C_0})\), also when summed over
all the edges and embeddings. Here \(C_0\) is a fixed constant.

For either orientation first pass to the list with entries
\(\nu_{\pi^{-j}\mu,l}/p^j\). Proposition~\ref{prop:continuousrows510}
bounds the error by a polynomial on each edge, also on each
individual cycle with its orientation fixed, as its proof
compares the unsorted lists entry by entry. The contribution at
a destination is the supremum over finite subsets of the sum of
their entries less the sum of that many leading destination
parts, padded with zeros. Perturbing the incoming lists changes
this supremum by at most their total absolute change. The cost
changes by at most the sum of absolute changes to the destination
parts. Each source row appears with total incoming weight at most
\(\sum_{j\ {\rm odd}}p^{-j}\le1\). These facts bound the total
change due to the model rows by twice their total absolute error.
The resulting errors, including the polynomial loss, are
\(o_{\psi,\eps,A}(b/x^A)\) for any fixed \(A\).

In the orientation with source on \(\Phi_0\), a destination
outside \(L_C\) can retain its entire list at zero cost. At a
destination in \(L_C\), a selection of fewer than \(k\) entries
pays \(b\) for each. At most one first-power entry is nonzero
and all entries with \(j\ge3\) are smaller than \(b\), since
\(p>\sqrt x\). A selection of at least \(k\) entries costs
\(kb\), more than the entire incoming sum
\(a p/(p^2-1)\). Thus the value is exactly \(F_C\).
For source on \(\Phi_1\), any nonempty selection pays the same
cost, \(a\) at destinations in \(T_C\) and zero elsewhere.
Comparing the whole incoming list with the empty choice gives
\(R_C\). Taking a supremum permits infinite lists by monotone
convergence of their entry sums.

If \(L_C\ne\varnothing\), the distinct destinations of its first
powers show \(R_C\ge |L_C|(kb/p-a)\). Also
\(F_C\le |T_C|a p/(p^2-1)\le2na/p\). The ratios
\(ap/(kb)=O(x^2/k)\) and \(na/(kb)=O(nx/k)\) tend uniformly
to zero. Hence \(R_C-F_C\ge kb/(2p)\) on each such cycle for
large \(x\), a gap exceeding the errors above on every edge.
The maxima on different cycles are independent. If \(L_C\) is empty
the model has \(R_C=0\) and \(F_C\ge0\), so its maximum equals
\(F_C\), whether or not the actual model perturbation breaks a tie.

It remains to subtract \(\lceil U_p(u)\rceil\). In the word
definition replace the distinguished thresholds by \(b\) and
the other thresholds by zero. At \(p^j\le3xb\) this costs at
most \(\sqrt b+2b/(nX')\) per term by the same threshold bounds
as in Lemma~\ref{lem:residualbaseline673}. There are polynomially
many such terms. Beyond this cutoff the sum of all the model
ramps and the original ramps is \(O(mh)\), by a geometric tail.
The model first and higher powers sum to exactly \(m\sum_C F_C\),
since \(a/p^3<b\). The total error, including the ceiling, is
\(o_{\psi,\eps,A}(b/x^A)\). Subtract this from the actual
maximized sum over edges to obtain the assertion.
\end{proof}

\begin{corollary}[Five odd powers determine the normalized logarithm]
\leavevmode\label{cor:fivepowers674}
Choose a translate \(u(p)\) at each \(p\in\mathcal I_x\).
For the common integer \(B_u\) of Theorem~\ref{thm:wordgcd510},
\[
 \begin{aligned}
 &\frac{1}{Q\delta}\sum_{p\in\mathcal I_x}v_p(B_{u(p)})\log p\\
 &\qquad =
 \frac m{Q\delta}\sum_{p\in\mathcal I_x}\log p
    \sum_{C:\,L_C\ne\varnothing}\bigl(R_C(J_9)-F_C(J_9)\bigr)
   +O_{\psi,\eps}\!\left(x^{-1/2+(10+1/8)\eps}\log(2x)\right).
 \end{aligned}
\]
The error tends to zero. The same formula holds for any selection
of individual word divisors, with the word allowed to depend on
the prime as well.
\end{corollary}
\begin{proof}
At a fixed translate \(v_p(B_u)\) is the minimum over the
words of \(v_p(b^{\rm per}_{u,w})\). The main term of the theorem
is independent of the word and its error is uniform. Its
contribution to the normalized logarithm is
\(O_{\psi,\eps,A}(x^{-A})\) for any chosen \(A\), by using a
larger exponent in the theorem before summing over \(p\le x\).

Deleting the terms \(j\ge11\) in the reverse sums costs at most
\(2kbh_2/p^{11}\) in total per edge. For each \(j\) a source
has a unique destination, and taking the positive part can
only reduce the change. The forward tail is at most
\(2ha/p^{11}\) by the same counting. Both estimates remain valid
upon restricting the cycles in the corollary. By integration,
\(\sum_{n>\sqrt x} n^{-11}\log n\ll x^{-5}\log(2x)\).
After including \(m\) and normalizing by \(Q\delta\), the two
bounds carry the respective factors \(2kn/h\) and
\(2a/(b\delta)\). From Theorem~\ref{thm:rangenew50} these are
\(O(x^{9/2+(10+1/8)\eps})\) and \(O(x^{1+\eps})\).
This gives the displayed error. Its exponent is negative for
\(0<\eps<1/52\).
\end{proof}

This evaluation determines the divisor removed in
Section~\ref{sec:dividedprimecomparison673} over the indicated
primes. The logarithmic sum of residual valuations in
Theorem~\ref{thm:residualvaluation673} also involves the images
of the actual projected maps.

\section{Signed terminal frequencies and their complex norms}\label{nav:chapter-signed-terminal-frequencies}

\subsection{The connected terminal Fourier projection at arbitrary precision}
\label{sec:connectedprecision516}

Fix a translate and a place above the terminal prime \(p=T\) on
the fields of Section~\ref{sec:missing}. We evaluate the connected
frequency coordinates of Lemma~\ref{lem:terminalblock496} in the
original signed period. Write \(\mathcal O\) for the valuation ring
of a finite local coefficient field containing the data, enlarged
when necessary, and normalize \(v(p)=1\). Congruences below are
modulo \(p^M\mathcal O\), including when the coefficient field
is ramified. For \(k\in\mathbb F_p\) let \([k]\) be its representative
in \(\{0,\ldots,p-1\}\).

In one scalar coordinate put \(F_{b,k}=U_b S_k\), with the
complete cyclic projections and prime-to-\(p\) displacements
of Theorem~\ref{thm:cyclicunit495}. Write \(D=d/dz\) in the original
frame and \(J_i(F)=(D^iF)(0)\). On the quotient elliptic curve
\(C'\) let \(z'=\alpha z\) be the coordinate for its good
differential. The connected orientation gives \(v(\alpha)=1\).
The displacement for label one gives a point \(Q\) on \(C'\)
of exact \(p\)-order whose reduction still has order \(p\).
Label \(k\) gives \([k]Q\), by the period pairing in
Lemma~\ref{lem:cyclicneighbor495}. Use \(Z\) for the elliptic
logarithmic derivative in the good differential, as defined in
Lemma~\ref{lem:addext492}.

\begin{lemma}[A common center for all prime-to-\(p\) translates]
\leavevmode\label{lem:connectedcenter516}
The scalar \(\theta=\alpha Z(Q)\) is a unit. For every \(b,k\)
and every \(d\ge0\),
\[
 B_d(b,k)=\big((D-\theta[k])^dF_{b,k}\big)(0)
                         \ \in\ p^d\mathcal O.
\]
In particular the same center \(\theta[k]\) works for all
displacements at the other primes. The assertion holds when
the section vanishes at the origin.
\end{lemma}
\begin{proof}
In the coordinate \(z'\) the section is a scalar times the
translate of the marked theta on \(C'\). The scalar includes
the phase of the translation and has valuation zero in the
connected orientation by Theorem~\ref{thm:cyclicunit495}.
Let \(P_0\) be the sum of its two-torsion mark and the point
induced by \(b\). Its order is prime to \(p\). The point that
enters the logarithmic derivative at label \(k\) is
\(P_k=P_0+[k]Q\). For \(k\ne0\) its reduction is different from
the origin because its \(p\)-primary part has order \(p\).
Lemma~\ref{lem:addext492} gives \(v(Z(Q))=-1\), which proves the
unit assertion.

We need the first logarithmic derivative to precision \(p\) in
\(\mathcal O\). Use the integral extension of \(C'\) by the
additive group constructed in Lemma~\ref{lem:addext492}. Over
points reducing away from the origin it has an integral coordinate
\(s_0\). Addition of a vertical scalar translates that coordinate,
and the unique torsion lift \(\widehat P\) of a torsion point has
coordinate \(Z(P)\). There is an integral lift \(g_Q\) of \(Q\)
with \(s_0=0\). Thus \(\widehat Q=g_Q+\iota(Z(Q))\), where
\(\iota\) denotes the vertical inclusion. The torsion lift of
\(P_0\) is integral. Indeed, start with an integral lift in any
chart at \(P_0\), multiply it by an order \(m\) of \(P_0\)
prime to \(p\), and subtract from that lift the resulting integral
vertical scalar divided by \(m\). The corrected lift is killed
by \(m\). Uniqueness follows from the absence of torsion in the
characteristic-zero additive kernel. It also shows that adding
torsion lifts lifts the sum. Consequently
\[
    \widehat P_k=\widehat P_0+[k]g_Q+\iota([k]Z(Q)).
\]
For \(k\ne0\) the sum of the first two terms is integral and
belongs to the \(s_0\) chart over \(P_k\). Hence
\(Z(P_k)-[k]Z(Q)\in\mathcal O\). Returning to \(z\) gives
\[
  (D\log F_{b,k})(0)-\theta[k]
        =\alpha\bigl(Z(P_k)-[k]Z(Q)\bigr)\in p\mathcal O .
\]
Here \(J_0(F_{b,k})\) is a unit by Theorem~\ref{thm:cyclicunit495}.

In a good Weierstrass model for \(C'\) the second logarithmic
derivative in \(z'\) is \(-u_{\Lambda'}-\wp(P_k)\), and the subsequent
ones are \(-\partial^{j-2}\wp(P_k)\) for \(j\ge3\), with \(\partial\)
the invariant derivation in that differential. These are integral
by the logarithmic derivative calculation of
Theorem~\ref{thm:residuejet492}. Indeed \(P_k\) reduces away
from the pole, \(u_{\Lambda'}\) is integral on the good model,
and its invariant derivation preserves regular functions locally.
In the original coordinate the derivative of order \(j\) acquires
\(\alpha^j\). Thus the logarithmic derivatives at zero of the formal
series \(e^{-\theta[k]z}F_{b,k}(z)\) belong to \(p^j\mathcal O\)
for all \(j\ge1\). Its undivided derivative of order \(d\) is its
constant value times the sum over the partitions of a set of
\(d\) elements, multiplying for each block the logarithmic
derivative of its size. This formula follows by differentiating,
either adding a singleton or enlarging one of the existing blocks.
Each product lies in \(p^d\mathcal O\), proving the assertion for
nonzero labels without dividing by a factorial. At \(k=0\) the
center is zero and Theorem~\ref{thm:cyclicunit495} gives directly
\(v(J_d(F_{b,0}))\ge d\), including at a zero of the section.
\end{proof}

For \(s\) such connected coordinates use variables
\(Z_1,\ldots,Z_s\), units \(\theta_i\), and a tuple of labels
\(\mathbf k\). Write \(c_i=\theta_i[k_i]\).
For a polynomial \(f\) with coefficients in \(\mathcal O\), its
Hasse derivative \(f^{[\mathbf d]}(c)\) is the coefficient of
\(X^{\mathbf d}\) in \(f(c+X)\). It is integral even when ordinary
derivatives divided by factorials would be inconvenient to evaluate.
For displacements \(b_i\) write \(B_{i,d}=B_d(b_i,k_i)\) in the
corresponding frames.

\begin{theorem}[The signed terminal calculation modulo every power of \(p\)]
\leavevmode\label{thm:connectedjets516}
At every integer precision \(M\ge1\) one has
\[
 \left.f(D_1,\ldots,D_s)\prod_i F_{b_i,k_i}(z_i)\right|_{z=0}
 \equiv
 \sum_{|\mathbf d|<M} f^{[\mathbf d]}(c)
                          \prod_{i=1}^s B_{i,d_i}
                                 \pmod {p^M\mathcal O}.
\]
The operator on the left annihilates, modulo \(p^M\), all polynomials in
\[
   I^M,\qquad I=(p,\ Z_1^p-\theta_1^{p-1}Z_1,\ldots,
                         Z_s^p-\theta_s^{p-1}Z_s).
\]
For the full set of labels the finite coefficient ring has the
description
\[
 \mathcal O[Z_1,\ldots,Z_s]/I^M\ \simeq \prod_{\mathbf k\in\mathbb F_p^s}
       \mathcal O[Z_1,\ldots,Z_s]/(p,Z_1-c_1,\ldots,Z_s-c_s)^M ,
\]
where the centers in each factor use its own label. At fixed
labels only the \(\binom{s+M-1}{s}\) Taylor coefficients with
\(|\mathbf d|<M\) enter, and the coefficient with index \(\mathbf d\)
is needed only modulo \(p^{M-|\mathbf d|}\).
\end{theorem}
\begin{proof}
Expand \(f\) as \(\sum_{\mathbf d}f^{[\mathbf d]}(c)(Z-c)^{\mathbf d}\).
The variables act on independent scalar factors, so evaluation
of a centered monomial gives exactly \(\prod_i B_{i,d_i}\). Its
valuation is at least \(|\mathbf d|\) by the lemma. This proves
the congruence, the precision of each coefficient, and vanishing
on the \(M\)-th power of each fixed-label ideal at that label.
Counting monomials of total degree less than \(M\) gives the
stated binomial coefficient.

For different label tuples the ideals
\((p,Z_1-c_1,\ldots,Z_s-c_s)\) are comaximal, since some pair
of centers differs by a unit. Their intersection is \(I\).
Indeed modulo \(p\) the polynomial
\(Z_i^p-\theta_i^{p-1}Z_i\) is the product of the distinct linear
factors \(Z_i-\theta_i[k]\) with unit differences. The Chinese
remainder theorem in each successive variable identifies the
quotient by these polynomials with the product of the evaluation
rings \(\mathcal O/p\) at the tuples. This proof also applies
when \(\mathcal O/p\) has nilpotents. For pairwise comaximal
ideals their product is their intersection, and their powers
are still pairwise comaximal. Taking the \(M\)-th powers thus
gives the asserted product decomposition and the common annihilation.
\end{proof}

For example, in one coordinate and for every \(n\ge1\) the
calculation at precision two is
\[
 J_n(F_{b,k})\equiv c^n J_0(F_{b,k})
    +n c^{n-1}\bigl(J_1(F_{b,k})-cJ_0(F_{b,k})\bigr)
                              \pmod {p^2\mathcal O}.
\]
Here \(c^0=1\) also when \(k=0\). This retains the first correction
to the residue formula at the same terminal character.

\begin{corollary}[Elementary divisors at the prescribed projection]
\leavevmode\label{cor:signedprecision516}
Fix the frames and a finite matrix of the signed projected columns
and tests of Theorem~\ref{thm:covectorpaths352}. Expand each actual
entry over its class paths and their shared covectors. At the
chosen place evaluate the other factors, and choose a common
scalar clearing the coefficients of the remaining polynomials
in the connected derivatives, including the incoming values,
class measures, and any prescribed column divisions. Let \(H\)
be this integral matrix. In each term apply the theorem to its
connected coordinates, and denote the resulting signed matrix
by \(H_M\). At every precision \(M\) the elementary-divisor
valuations of \(H\) and \(H_M\), each truncated at \(M\), agree.
Zero elementary divisors are assigned valuation \(+\infty\).
In particular the number with valuation strictly below \(M\)
is a lower bound on the rank of the original projected columns.
\end{corollary}
\begin{proof}
The integral substitution in \eqref{eq:sub81} is a polynomial
in the commuting scalar derivatives in the exact chosen frames.
After the indicated clearing it remains integral when all other
factors have been evaluated. For conjugated edges use the
conjugated coefficient data at their corresponding places.
Lemma~\ref{lem:terminalblock496} gives precisely the complete
cyclic projections used above, up to the recorded phases. Apply
the congruence to each product term, keeping a shared covector
fixed in both adjacent factors until they have been multiplied.
Summing with the original coefficients gives \(H=H_M\) modulo
\(p^M\) entry by entry. No error is multiplied by the number
of paths in this assertion since \(p^M\mathcal O\) is an ideal.

Their cokernels after reduction modulo \(p^M\) are identical.
If \(\varpi\) is a uniformizer and \(e=v_\varpi(p)\), Smith
reduction of a matrix over \(\mathcal O\) presents this finite
cokernel as a sum of modules of lengths \(\min(v_\varpi(d_i),eM)\),
together with the free target rows reduced modulo \(p^M\).
The matrix dimensions are the same, so uniqueness of elementary
divisors gives precisely the claimed truncated list on dividing
by \(e\). Every divisor of valuation less than \(M\) is nonzero
in characteristic zero. Clearing by a nonzero common scalar
preserves the characteristic-zero rank.
\end{proof}

The calculation applies on the same interval as
Corollary~\ref{cor:smallresidual515}. An estimate for the number
of surviving elementary divisors and the height of the projected
period over that whole interval is still required in
Section~\ref{sec:missing}.


\subsection{Complex norms on the nonzero terminal frequencies}
\label{sec:terminalnorm518}

Fix the terminal prime \(p=T=\varpi\bar\varpi\) of
Theorem~\ref{thm:common111}, an output covector, and one of the
translated product terms in Lemma~\ref{lem:terminalblock496}.
We work at a complex embedding and keep the point of evaluation at
zero. In one constrained elliptic coordinate the lattice is
\(\Lambda=\delta I\), with \(I\) a fractional ideal of the full ring of
\(K_d\). Write \(A\) for its area, \(c=\pi/A\), and
\(\Theta\) for the marked theta section. Its norm uses the Hermitian
metric and Haar probability on \(\mathbb C/\Lambda\). Interchange
\(\varpi,\bar\varpi\) if necessary so that the frequency subgroup
is \(\varpi^{-1}\Lambda/\Lambda\). Choose its generator \(w\) to
precision \(p^2\) as in Lemma~\ref{lem:cyclicneighbor495}. With the
translation operators and phases of that lemma, put
\[
 S_k=p^{-1}\sum_{j=0}^{p-1} e^{-2\pi i k j/p}U_{jw}\Theta,
 \qquad y_k=J_P(U_b S_k),\qquad
 J_P(H)=\sum_{n=0}^K a_n\partial_z^n H(0).
\]
Here \(P(Z)=\sum_{n=0}^K a_n Z^n\) is any complex polynomial and
\(b\) is the fixed prime-to-\(p\) shift of the product term, lifted
to the common cover. Give these polynomials the squared norm
\(\mu(P)=\|\Theta\|^2\sum_{n=0}^K c^n n!|a_n|^2\).

\begin{lemma}[Energy on a split cyclic subgroup]\leavevmode
\label{lem:cyclicsampling518}
Let \(\lambda\) be the shortest nonzero length of the real dual
lattice \(\Lambda^\vee\), and put \(\rho=\lambda/\sqrt c\). Define
\[
 E=50\left(\frac{2\sqrt{K+6}}{\pi\sqrt p\,\rho}\right)^6,
 \qquad N=\max\{1,\lceil8\sqrt{K+6}/\rho\rceil\},
 \qquad B=3N^2/2.
\]
Then
\[
 \left|\sum_{k\in\mathbb F_p}|y_k|^2-\mu(P)\right|
       \le E\mu(P),\qquad
 \max_{k\in\mathbb F_p}|y_k|^2\le (B/p)\mu(P).
\]
Also \(\rho^2\ge 2/(\pi\sqrt{|\operatorname{disc}K_d|})\).
\end{lemma}
\begin{proof}
Finite Fourier orthogonality gives
\(\sum_k|y_k|^2=p^{-1}\sum_j|J_P(U_bU_{jw}\Theta)|^2\).
Commuting the two translations changes each summand inside the
absolute value by a scalar of modulus one, common to all derivative
orders. The last mean therefore samples the periodic function
\(g(z)=|J_P(U_z U_b\Theta)|^2\) on \(\varpi^{-1}\Lambda/\Lambda\).
Its Haar mean is \(\mu(P)\). Indeed the metric representative of
order \(n\) is obtained by applying
\((\partial_z-c\bar z)^n\) to the holomorphic section and inserting
the metric factor. The adjoint is \(-\partial_{\bar z}\).
Commuting a lowering operator past a raising operator contributes
\(c\). Any leftover lowering operator kills a holomorphic section,
so different orders are orthogonal and the squared norm in order
\(n\) is \(c^n n!\|\Theta\|^2\), as in Lemma~\ref{lem:translatedjets478}.

The Fourier estimate proved in Lemma~\ref{lem:grid85} gives
\( |\widehat g(u)|\le\mu(P)(2\sqrt{c(K+6)}/(\pi|u|))^6\) for
\(0\ne u\in\Lambda^\vee\). Its hypothesis holds because each
lowering operator reduces the raising order by one, so \(K+1\)
lowering operators kill the whole section. A unitary character in
the period multiplier cancels in the metric pairing used in the
proof. Averaging retains only the dual of \(\varpi^{-1}\Lambda\),
which is \(\bar\varpi\Lambda^\vee\). Thus the shortest surviving
nonconstant frequency has length \(s=\sqrt p\,\lambda\).
In \(2^h s\le |u|<2^{h+1}s\) there are at most \(25\cdot 2^{2h}\)
such frequencies, by packing disjoint disks of radius \(s/2\).
Summing the sixth-power bound and using
\(25\sum_{h\ge0}2^{-4h}<50\) proves the first inequality.

For the pointwise estimate pass to
\(\Gamma=\bar\varpi\Lambda\). The \(p\) sections \(U_{jw}\Theta\)
on that cover are eigenvectors for a common unitary period
translation with distinct eigenvalues. Indeed the pairing of a
suitable period of \(\Lambda\) with the generator \(w\) is a
primitive \(p\)-th root. These vectors are orthogonal and have
the same norm as \(\Theta\) for probability measure. Hence
\(\|S_k\|^2=\|\Theta\|^2/p\) for every \(k\). By
Lemma~\ref{lem:cyclicneighbor495}, \(S_k\) also has a theta
multiplier on \(\Lambda'=(\bar\varpi/\varpi)\Lambda\), differing
from that of \(S_0\) by a unitary character. Its squared metric
norm on this torus agrees with the one on the cover. The rotation
from \(\Lambda\) to \(\Lambda'\) preserves \(c\) and \(\lambda\).
The same adjoint calculation gives mean \(\mu(P)/p\) for
\(|J_P(U_zU_b S_k)|^2\) on the new torus. Lemma~\ref{lem:grid85}
with the indicated \(N\) bounds its mean on the full \(N\)-torsion
grid by \(3\mu(P)/(2p)\). Its \(N^2\) summands include the
nonnegative value \(|y_k|^2\) at zero, giving the assertion.
The identical grid argument on the original torus gives
\(|J_P(U_b\Theta)|^2\le B\mu(P)\).

Finally the planar dual is \(i\Lambda/A\). If \(0\ne u\in I\),
then \((u)I^{-1}\) is an integral ideal, so \(|u|^2\ge N(I)\).
The area is \(|\delta|^2 N(I)\sqrt{|\operatorname{disc}K_d|}/2\).
Consequently \(\rho^2=\min_{0\ne z\in\Lambda}|z|^2/(\pi A)\)
has the stated lower bound.
\end{proof}

In one scalar coordinate there is also an estimate after subtracting
the mean on \(\mathbb F_p^\times\). Write
\(\bar y=(p-1)^{-1}\sum_{k\ne0}y_k\). Since \(\sum_k S_k=\Theta\),
the squared norm lost in this subtraction is exactly
\[
 (p-1)|\bar y|^2=\frac{|J_P(U_b\Theta)-y_0|^2}{p-1}
 \le \frac{B(1+p^{-1/2})^2}{p-1}\,\mu(P).
\]
Here the triangle inequality uses the two pointwise bounds just
proved. Together with the cost \(|y_0|^2\le B\mu(P)/p\) of deleting
zero, this gives a lower bound for the actual one-coordinate
augmentation, with the original phases of the Fourier sum.

\begin{theorem}[Rank and norms on nonzero terminal covectors]\leavevmode
\label{thm:puncturedisometry518}
For one fixed translated product term let \(r_0\) be the number of
constrained coordinates at the chosen split block. In coordinate
\(h\) use \(S_{h,k_h}, b_h,c_h,K_h,E_h,B_h\) as above. On the
polynomials \(P=\sum_\nu a_\nu Z^\nu\) of degree at most \(K_h\)
in \(Z_h\), put
\[
 \|P\|_{\rm F}^2=\left(\prod_h\|\Theta_h\|^2\right)
      \sum_\nu |a_\nu|^2\prod_h c_h^{\nu_h}\nu_h! .
\]
Let \(T_bP(k)\) be \(P(\partial)\) applied at zero to
\(\prod_h(U_{b_h}S_{h,k_h})(z_h)\). Then
\[
 \prod_h(1-E_h-B_h/p)_+\,\|P\|_{\rm F}^2
 \ \le\sum_{0\ne k\in\mathbb F_p^{r_0}} |T_bP(k)|^2
 \ \le\prod_h(1+E_h)\,\|P\|_{\rm F}^2,
\]
where \(u_+=\max(u,0)\). At the word weights of
Theorem~\ref{thm:common111} these factors lie between \(1-\eta\)
and \(1+\eta\), with \(\eta\le p^{-1+o(1)}\) uniformly as the
interval recedes. In particular the evaluation on the nonzero
covectors is injective on the prescribed polynomial subspace.
\end{theorem}
\begin{proof}
On one coordinate Lemma~\ref{lem:cyclicsampling518} bounds the
squared singular values of the full evaluation map above by
\(1+E_h\) in the indicated polynomial metric. Deleting the label
zero costs at most \((B_h/p)\mu(P)\) for every input, so the map
on \(\mathbb F_p^\times\) has squared singular values at least
\((1-E_h-B_h/p)_+\), including zero eigenvalues if necessary.
On a Cartesian product of label sets the evaluation map is the
tensor product of the one-coordinate maps. This follows on monomials
by differentiation in separate variables, and extends by linearity.
The displayed polynomial metric is the tensor product metric.
Diagonalizing the positive Gram operators gives the products of
the one-coordinate eigenvalues. The nonzero covectors contain
\((\mathbb F_p^\times)^{r_0}\) and are contained in the full grid.
Adding squared absolute values can only increase the quadratic form.
This proves both bounds, including for polynomials with mixed terms.

For the uniform statement take \(K_h\) up to the total degree of
the prescribed polynomial substitution on the edge. At each CM
embedding its degree is the size of its partition and is at most
the row rank times the first row length, by the determinant shift
in Lemma~\ref{lem:wordfirstrow450}. The number of embeddings and
all ranks are bounded by fixed powers of \(x\). The first row
is \(O(x\exp(x^\varepsilon)+x^C)\) by the weight calculation in
Corollary~\ref{cor:orbitkernel500}. Also \(d_0\) is polynomially
bounded on the conductor intervals of Theorem~\ref{thm:conductor232},
so the discriminant estimate in Lemma~\ref{lem:cyclicsampling518}
gives \(\rho_h^{-2}\ll x^{C'}\) at every embedding. It follows
that \(B_h\ll x^{C''}\exp(x^\varepsilon)\) and
\(E_h\ll (B_h/p)^3\). The terminal size estimate in
Theorem~\ref{thm:prime210} gives \(\log p\gg m x^\varepsilon\)
with \(m\to\infty\). Hence \(r_0\max_h(E_h+B_h/p)\le p^{-1+o(1)}\).
The lower product is at least one minus that sum once all factors
are positive. The upper product is at most \(\exp(\sum_h E_h)\),
which gives the stated error. Restricting a positive definite
quadratic form to the prescribed subspace proves injectivity there.
\end{proof}

\begin{corollary}[Gram determinants and pairing with an incoming table]
\leavevmode\label{cor:terminalgram518}
Let \(P_1,\ldots,P_q\) be linearly independent prescribed
polynomials for the same product term, with coefficients independent
of \(k\). Write \(G_{\rm F}\) for their polynomial Gram matrix and
\(G_{\rm term}\) for the Gram matrix of their values on all nonzero
covectors. On the above intervals
\[
 \left|\frac{1}{q}\log\frac{\det G_{\rm term}}{\det G_{\rm F}}\right|
 \le 2\eta\qquad(\eta\le1/2).
\]
On the span of the given inputs let \(T\) be this evaluation
map, and let \(\Pi\) be orthogonal projection onto its image in
the space of tables on nonzero covectors, with counting measure.
For every incoming table \(Y\) one has
\[
 (1-\eta)\|\Pi Y\|^2\le\|T^*Y\|_{\rm F}^2
                       \le(1+\eta)\|\Pi Y\|^2.
\]
\end{corollary}
\begin{proof}
Apply the theorem to every linear combination of the \(P_i\).
Conjugating the resulting Gram inequalities by \(G_{\rm F}^{-1/2}\)
puts each eigenvalue of the relative Gram matrix in
\([1-\eta,1+\eta]\). Taking logarithms and using
\(-\log(1-\eta)\le2\eta\) proves the determinant bound. The
positive eigenvalues of \(TT^*\) are the eigenvalues of \(T^*T\),
and \(TT^*\) is zero on the orthogonal complement of the image.
Its quadratic form on \(Y\) is \(\|T^*Y\|_{\rm F}^2\), proving the
second assertion. For a bilinear incoming contraction take the
complex conjugate of the incoming coefficients in the Hermitian pairing.
\end{proof}

Corollary~\ref{cor:sameprimes517} permits these estimates at the very
primes for which the full strict groups vanish throughout the weight
interval. For the signed period of Section~\ref{nav:I-28}, one must
still bound the projections of the prescribed incoming tables and
the cancellation among terms with different shifts and class paths.


\subsection{The scalar-orbit projection of the terminal derivatives}
\label{sec:scalarprojection519}

In a terminal block the scalar \(a\in\mathbb F_p^\times\) acts on
an entire covector by \(k\mapsto ak\). The incoming tables in
Proposition~\ref{prop:orbitsum501} have mean zero on every such
orbit. We estimate the part of an outgoing theta table that
survives pairing with this space. All translations, phases and
derivatives below use the fixed product of
Theorem~\ref{thm:puncturedisometry518} at the original point.

\begin{lemma}[Translation by a short lift of a frequency]\leavevmode
\label{lem:shortmotion519}
In one scalar coordinate write \(y=(J_P(U_bS_k))_{k\in\mathbb F_p}\)
and \(E_+=E(K+1)\), with \(E\) as in Lemma~\ref{lem:cyclicsampling518}.
There is a unitary operator \(W\) on the space of these tables,
independent of \(P\), such that
\[
 \begin{gathered}
 (Wf)(k)=\gamma e^{2\pi i a k/p}f(k+d),\qquad |\gamma|=1,
 \quad a\ne0,\\
 \|(I-W)y\|\le 4\sqrt{(K+1)/p}\sqrt{1+E_+}\,\|P\|_{\rm F}.
 \end{gathered}
\]
Here \(a,d\in\mathbb F_p\) and the norm on tables uses counting measure.
\end{lemma}
\begin{proof}
Choose a shortest nonzero \(v\in\Lambda\) and set \(h=v/\varpi\).
Its length is smaller than that of any nonzero period, so
\(h=aw+\ell\) with \(a\not\equiv0\pmod p\) and \(\ell\in\Lambda\).
The disk of radius \(|v|/2\) embeds up to its boundary in the
torus of area \(A\). Thus \(c|h|^2=\pi|v|^2/(pA)\le4/p\).
In the finite Fourier sum defining \(S_k\), cyclically shifting the
summation index gives \(U_w S_k=e^{2\pi i k/p} S_k\). A period
translation \(U_\ell\) on the summand indexed by \(j\) acts by its
original period sign times a character of \(j\), by the commutator
calculation in Lemma~\ref{lem:cyclicneighbor495}. Thus it sends
\(S_k\) to that sign times \(S_{k+d}\) for a fixed \(d\). Composing
these two translations and commuting past \(U_b\) introduces a
constant phase independent of \(k\). This proves that
\(J_P(U_hU_b S_k)=(Wy)(k)\) with the indicated unitary \(W\).

To estimate the difference, let \(H_t=U_{th}U_bS_k\) for
\(0\le t\le1\). The translation law in the common coordinate is
\(U_uH(z)=\exp(-c\bar u z-c|u|^2/2)H(z+u)\). Differentiating
along this real segment yields
\(\dot H_t=h\partial_z H_t-c\bar h z H_t\). Hence
\(d J_P(H_t)/dt=J_Q(H_t)\), where
\(Q(Z)=hZP(Z)-c\bar h P'(Z)\). In the factorial metric,
\(\|ZP\|_{\rm F}\le\sqrt{c(K+1)}\|P\|_{\rm F}\) and
\(\|cP'\|_{\rm F}\le\sqrt{cK}\|P\|_{\rm F}\), by shifting each
monomial in the norm formula. The complex sampling lemma applies
along the segment. Its proof for translated sections uses only
unitarity of the multiplier and therefore permits an arbitrary
complex displacement. Also \(U_{th}U_b\) differs from \(U_{th+b}\)
by a common phase. The vector of the derivatives in \(t\) thus
has norm at most \(\sqrt{1+E_+}\|Q\|_{\rm F}\). Integrate from
zero to one and insert the bound on \(|h|\sqrt c\).
\end{proof}

\begin{proposition}[Loss at a single multiplicative character]\leavevmode
\label{prop:mellinloss519}
For every character \(\chi\) of \(\mathbb F_p^\times\), including the
trivial character, put \(u_\chi(0)=0\) and
\(u_\chi(k)=\chi(k)/\sqrt{p-1}\) for \(k\ne0\). If
\(t=\sqrt{(K+1)/p}\le1/2\), then
\[
 |\langle u_\chi,y\rangle|^2
       \le 150(1+E_+)t\,\|P\|_{\rm F}^2 .
\]
\end{proposition}
\begin{proof}
For every integer \(j\) prime to \(p\), the operator \(W^j\) is,
up to a unit scalar, a shift by \(jd\) and a modulation by
\(ja\ne0\). For a nontrivial \(\chi\) the finite-field
translation and modulation bound of Gurevich, Hadani and Sochen
therefore gives \(|\langle u_\chi,W^j u_\chi\rangle|\le2/\sqrt p\).
In their standard split torus the chosen vector is precisely the
normalized character on \(\mathbb F_p^\times\) extended by zero,
with the counting inner product \cite[\S I-B and (I.1)]{XXVI-XX-oscillator519}.
For the trivial character the same inner product, up to a phase,
is the sum of the nonconstant additive character of frequency \(ja\)
over all \(k\) except \(0\) and \(-jd\), divided by \(p-1\).
The complete sum vanishes, so its modulus is at most \(2/(p-1)\).
For \(p\mid j\) the unitary bound one suffices.

Write \(s=1-t\). By unitarity the inverse of \(I-sW^*\) is its
convergent geometric series. Squaring the norm and grouping pairs
of exponents by their difference gives
\[
 \|(I-sW^*)^{-1}u_\chi\|^2
 \le \frac{1}{1-s^2}
     \left(1+\frac{4}{\sqrt p\,t}
              +\frac{2s^p}{1-s^p}\right)\le\frac6t .
\]
Indeed \(\sqrt p\,t\ge1\), \(s^p\le e^{-pt}\le e^{-\sqrt p}\),
\(p\ge5\), and \(1-s^2\ge t\). On the other hand the sampling
lemma gives \(\|y\|\le\sqrt{1+E_+}\|P\|_{\rm F}\). Applying
Lemma~\ref{lem:shortmotion519} in
\(I-sW=tI+s(I-W)\) gives
\(\|(I-sW)y\|\le5t\sqrt{1+E_+}\|P\|_{\rm F}\).
Pairing this last vector with \((I-sW^*)^{-1}u_\chi\) recovers
the original inner product in absolute value. Cauchy--Schwarz
proves the assertion.
\end{proof}

For \(D=\mathbb F_p^{r_0}\setminus\{0\}\) let \(A\) be the
orthogonal projection onto the tables with mean zero on each
line of nonzero multiples. In particular
\[
 (Af)(k)=f(k)-(p-1)^{-1}\sum_{a\in\mathbb F_p^\times}f(ak).
\]
The same \(a\) acts on every coordinate. Use \(T_b\) for the
evaluation map of Theorem~\ref{thm:puncturedisometry518} on \(D\).

\begin{theorem}[Rank after the scalar-orbit projection]\leavevmode
\label{thm:scalaraugmentation519}
On every sufficiently distant common interval, at every complex
embedding and for one fixed translated product of the marked
terminal block,
\[
 (1-\eta')\|P\|_{\rm F}^2\le\|A T_bP\|^2
              \le(1+\eta')\|P\|_{\rm F}^2,
        \qquad \eta'\le p^{-1/2+o(1)} .
\]
This holds for every linear combination of the prescribed
polynomial rows. In particular their rank is preserved by
the evaluation followed by scalar augmentation. More generally,
on removing the same \(J\) specified scalar character types from
every line, the lower bound has loss at most
\(p^{-1+o(1)}+Jp^{-1/2+o(1)}\).
\end{theorem}
\begin{proof}
Let \(D_0=(\mathbb F_p^\times)^{r_0}\). Both \(D_0\) and its
complement in \(D\) are unions of whole scalar orbits. Hence
restricting to \(D_0\) commutes with the projections onto
individual types, and hence also with deleting them. On \(D_0\)
the evaluation is \(T_1\otimes T_R\), where \(T_1\) is scalar
evaluation on the first punctured coordinate and \(T_R\) is the
tensor product for the others. For \(r_0=1\) use the one-dimensional
trivial space as the empty tensor product. Let \(P_\chi\) be
orthogonal projection to the character \(\chi\) in the first
coordinate. It has rank one. Proposition~\ref{prop:mellinloss519}
gives, in the order of Hermitian forms,
\[
  T_1^*P_\chi T_1\le \delta I,\qquad
        \delta=150(1+E_{1,+})\sqrt{(K_1+1)/p}.
\]
Write \(Q_\beta\) for the orthogonal type projections of the
simultaneous scalar action on the remaining coordinates. These
sum to the identity. On the product grid the projection to
total type \(\alpha\) is \(\sum_\chi P_\chi\otimes
Q_{\alpha\chi^{-1}}\), since characters multiply on tensor products.
Pulling this projection back by \(T_1\otimes T_R\) gives an
operator bounded by
\[
  \delta I\otimes\sum_\chi T_R^* Q_{\alpha\chi^{-1}}T_R
         =\delta I\otimes T_R^*T_R .
\]
This uses positivity of each second factor and therefore applies
also to polynomials coupling different coordinates. By the scalar
sampling lemma \(\|T_R\|^2\le\prod_{h>1}(1+E_h)\). The lower
bound on the unprojected product grid in the proof of
Theorem~\ref{thm:puncturedisometry518} is
\(\prod_h(1-E_h-B_h/p)_+\). Subtract at most
\(J\delta\prod_{h>1}(1+E_h)\) for \(J\) removed types, then add
the nonnegative squared norm on \(D\setminus D_0\). The upper
bound follows by contractivity of an orthogonal projection and
the upper bound on \(D\) in the same theorem. In the word range
\(K_h+1=p^{o(1)}\), \(E_{h,+}\le p^{-3+o(1)}\), and the lower
product is at least \(1-p^{-1+o(1)}\). These estimates are uniform
in the word, the shift and the embedding. Taking \(J=1\) and
removing the trivial character gives the result for \(A\).
\end{proof}

\begin{corollary}[Pairing with the full space of separately augmented tables]
\leavevmode\label{cor:multipleaugmentations519}
Take a fixed product term at all split terminal blocks and impose
the zero-mean condition independently at each block, holding the
other covectors fixed. Evaluation followed by orthogonal projection
to this space has all squared singular values in
\([1-p^{-1/2+o(1)},1+p^{-1/2+o(1)}]\) on the prescribed polynomial
subspace with its product factorial metric. The same bounds hold
for the Rayleigh quotients of its Gram matrix relative to the
polynomial Gram matrix on every independent list of rows.
\end{corollary}
\begin{proof}
Each block uses its own scalar group. Averaging in independent
variables commutes, so the projection onto the common zero-mean
space is the tensor product of the single-block projections. The
product evaluation also factors. Diagonalizing the positive Gram
operators multiplies their eigenvalue bounds. The number of blocks
is bounded by a fixed power of \(x\), hence is \(p^{o(1)}\) on
these intervals. Multiplying the bounds in the theorem preserves
the stated error. Restriction to any polynomial subspace preserves
the form inequalities. For an incoming vector with the separate
zero-mean conditions, moving the orthogonal projection across the
pairing leaves that vector fixed. Thus this is the Gram matrix for
pairing with an orthonormal basis of the entire stated space.
\end{proof}

The prime choices in Corollary~\ref{cor:sameprimes517} apply here.
An individual prescribed incoming column can be perpendicular to
the span of these projected rows. The estimate in
Section~\ref{nav:I-28} still requires control of those columns after
the intermediate classes, shifts and their coefficients are summed.


\subsection{Signed sums of torsion translates at the marked prime}
\label{sec:signedtranslates520}

Two translates which have the same displacement on the common
elliptic cover must be added with their original phases. We estimate
the cross terms for distinct displacements on that cover, still at
the origin used to define the period. For one constrained coordinate
use \(\Gamma=\bar\varpi\Lambda\) and \(\Xi=\varpi^{-1}\Lambda\).
Write \(L_bP=(J_P(U_bS_k))_{k\in\mathbb F_p}\). We use the same
marked section and factorial metric for all \(b\). A translation
by a period of \(\Gamma\) acts on every \(S_k\) by the same scalar:
such a period commutes with all \(U_{jw}\) and acts on \(\Theta\)
by its period sign. Thus one can group equal displacements modulo
\(\Gamma\), incorporating the translation phase in their coefficients.

\begin{lemma}[Inner products for distinct torsion displacements]\leavevmode
\label{lem:crossshift520}
Let \(b,b'\) be distinct modulo \(\Gamma\) with
\(M(b-b')\in\Gamma\), and let \(P,Q\) have degree at most \(K\).
Take a basis \(h_1,h_2\) of \(\Xi\) and set
\(H=\max_i|h_i|\). For every integer \(j\ge1\),
\[
 |\langle L_bP,L_{b'}Q\rangle|
 \le (1+E(K+j))\bigl(MH\sqrt{c(K+j)}\bigr)^j
                          \|P\|_{\rm F}\|Q\|_{\rm F} .
\]
One can take \(cH^2\le (4+\pi\sqrt{|\operatorname{disc}K_d|})/p\).
\end{lemma}
\begin{proof}
Write \(e(h,b)\) for the unit scalar in \(U_hU_b=e(h,b)U_bU_h\).
The associated alternating form, up to sign, is
\(\operatorname{Im}(\bar h b)/A\) in the exponential \(e^{2\pi i t}\).
The annihilator of \(\Xi\) under this pairing is \(\Gamma\).
Indeed the area normalization makes \(\Lambda\) self-dual for the
alternating form, and moving the factor \(\varpi^{-1}\) across it
conjugates that factor. As \(b-b'\notin\Gamma\), for some basis
vector \(h\) the number \(\xi=e(h,b-b')\) is a nontrivial
\(M\)-th root of unity. In particular \(M\ge2\) and
\(|\xi-1|\ge2\sin(\pi/M)\ge4/M\).

Form the scalar function of a real variable
\[
 F(t)=\left\langle
   \big(J_P(U_{th}U_bS_k)\big)_k,
   \big(J_Q(U_{th}U_{b'}S_k)\big)_k\right\rangle ,
\]
using the inner product linear in its first argument. Since
\(h=aw+\ell\) with \(\ell\in\Lambda\), \(U_h\) acts on the list
\((S_k)_k\) by one unitary matrix, a modulation and a label
translation up to a common phase, as in Lemma~\ref{lem:shortmotion519}.
Also \(U_{nh}=U_h^n\). Moving this translation past the two shifts
and canceling the common unitary in the pairing gives
\(F(n)=\xi^n F(0)\) for each nonnegative integer \(n\).

In the first argument, differentiation in \(t\) applies the
polynomial operator \(P\mapsto hZP-c\bar h P'\). On degree at
most \(K+j-1\) its norm is at most \(2|h|\sqrt{c(K+j)}\), by
the factorial norm calculation in that lemma. The same holds in
the second argument. Applying the product rule \(j\) times and
the full-grid sampling bound in Lemma~\ref{lem:cyclicsampling518}
at the translated points gives
\[
 \sup_t |F^{(j)}(t)|\le (1+E(K+j))
     \bigl(4|h|\sqrt{c(K+j)}\bigr)^j\|P\|_{\rm F}\|Q\|_{\rm F}.
\]
The \(j\)-th forward difference of \(F\) at zero is
\((\xi-1)^jF(0)\). By iterating the fundamental theorem of calculus
it is also the integral of \(F^{(j)}(t_1+\cdots+t_j)\) on the
unit \(j\)-cube. Division by \(|\xi-1|^j\) proves the claim.

For the basis take a shortest period \(v_1\in\Lambda\), with
length \(s\). It is primitive, so it extends to a lattice basis.
Subtract multiples of \(v_1\) from the second vector to obtain
\(v_2\) with parallel component of magnitude at most \(s/2\).
Its perpendicular component has magnitude \(A/s\). The disk
argument gives \(cs^2\le4\), and the ideal norm bound in
Lemma~\ref{lem:cyclicsampling518} gives \(s^2/A\ge 2/\sqrt{|\operatorname{disc}K_d|}\).
Consequently \(c|v_2|^2\le1+(\pi/2)\sqrt{|\operatorname{disc}K_d|}\).
Take \(h_i=v_i/\varpi\).
\end{proof}

In each coordinate \(h\), take a nonempty finite set \(B_h'\) of
displacements in one coset of the \(M_h\)-torsion on
\(\mathbb C/\Gamma_h\), with one lift per displacement. For the
tuple \(\boldsymbol b\in\prod_h B_h'\) let \(P_{\boldsymbol b}\)
be the polynomial obtained after summing all coefficients belonging
to that tuple. Coefficients may be signed or complex and must be
independent of the shared frequency \(k\). Give the collection
the squared norm \(\|\boldsymbol P\|^2=\sum_{\boldsymbol b}
\|P_{\boldsymbol b}\|_{\rm F}^2\). Let \(S\boldsymbol P\) be the
sum of the corresponding product evaluations \(T_{\boldsymbol b}
P_{\boldsymbol b}\) on the nonzero covectors. We use \(A\) for
subtraction of the scalar-orbit means as in
Theorem~\ref{thm:scalaraugmentation519}.

\begin{theorem}[Norm and rank of a signed sum with torsion displacements]\leavevmode
\label{thm:signedtranslation520}
Take the marked primes, embeddings and derivative degrees on the
common intervals of Corollary~\ref{cor:sameprimes517}. The following
bounds hold uniformly for the cosets, the sets within them and the
grouped coefficient polynomials.
\begin{enumerate}
\item If \(0<\sigma<1/2\) is fixed and all \(M_h\le p^{1/2-\sigma}\),
then \(\big|\|S\boldsymbol P\|^2-\|\boldsymbol P\|^2\big|
\le p^{-2\sigma+o(1)}\|\boldsymbol P\|^2\).
\item If \(0<\sigma<1/4\) is fixed and all \(M_h\le p^{1/4-\sigma}\),
then \(\big|\|AS\boldsymbol P\|^2-\|\boldsymbol P\|^2\big|
\le p^{-2\sigma+o(1)}\|\boldsymbol P\|^2\).
\end{enumerate}
For uniformly \(M_h=p^{o(1)}\) the respective errors are
\(p^{-1+o(1)}\) and \(p^{-1/2+o(1)}\). These estimates also hold
for a product of marked blocks, with each covector nonzero and
the scalar-orbit projections imposed separately. In each case
the evaluated rank of a list of rows equals the rank of its list
of grouped coefficient arrays.
\end{theorem}
\begin{proof}
First work in a single coordinate and write \(n=|B_h'|\le M^2\).
On the direct sum over the shifts, the full-grid map is
\(L:(P_b)_b\mapsto\sum_b L_bP_b\). Lemma~\ref{lem:cyclicsampling518}
bounds each diagonal block of \(L^*L-I\) in operator norm by \(E(K)\).
The off-diagonal blocks have norm at most
\(u_j=(1+E(K+j))(MH\sqrt{c(K+j)})^j\) by
Lemma~\ref{lem:crossshift520}. Using
\(\sum_{b\ne b'}v_bv_{b'}\le(n-1)\sum_b v_b^2\) on the norms
of the inputs gives, for every \(j\ge1\),
\[
   \|L^*L-I\|\le E(K)+(n-1)u_j .
\]
Deleting the label zero costs at most \(nB/p\) in the squared
operator norm of that evaluation functional. This follows by
Cauchy--Schwarz across the \(n\) shifts from the pointwise
bound, with \(B\) of Lemma~\ref{lem:cyclicsampling518}. In the
punctured coordinate the squared norm of the projection onto
any one multiplicative character is bounded, as a form on the
direct sum, by \(150n(1+E(K+1))\sqrt{(K+1)/p}\) times its
input form. This is the same Cauchy--Schwarz argument applied
to Proposition~\ref{prop:mellinloss519}.

On the full or componentwise punctured grid the map on the arrays
is the tensor product of these one-coordinate direct-sum maps.
Indeed on a tensor of lists the product of the sums is the sum
over all shift tuples. Such tensors span the array space, with
exactly the stated direct-sum polynomial metric. The tensor bounds,
restriction to nonzero covectors and scalar-character calculation
in Theorems~\ref{thm:puncturedisometry518} and
\ref{thm:scalaraugmentation519} therefore apply with the three
bounds just obtained. In particular the cost of removing the
trivial total character is at most the one-character loss in one
coordinate times the upper squared norm in the other coordinates.
This follows by summing the orthogonal character projections on
those other coordinates, as in the proof of the latter theorem.

For the first range take a fixed integer \(j\ge 1/\sigma\).
Uniformly \(\sqrt{c(K+j)}H\le p^{-1/2+o(1)}\) by the basis bound
and the growth estimates of Theorem~\ref{thm:puncturedisometry518}.
Then \(nu_j\le p^{1-2\sigma-j\sigma+o(1)}\), and the zero-label
loss is at most \(p^{-2\sigma+o(1)}\). For the second range take
\(j=2\). The off-diagonal error is at most
\(M^4p^{-1+o(1)}\le p^{-4\sigma+o(1)}\), and the character loss
is at most \(M^2p^{-1/2+o(1)}\le p^{-2\sigma+o(1)}\). The other
errors are smaller. When \(M=p^{o(1)}\), taking \(j=2\) gives
the two sharper rates as stated. The number of coordinates and
blocks is bounded by a fixed power of \(x\), hence is \(p^{o(1)}\).
Multiplying the upper and lower factors preserves these rates.
The relative error tends to zero in each range. Applying the
lower inequality to every linear combination of the rows proves
the rank assertion, including when the original list is dependent.
\end{proof}

\begin{corollary}[Displacements at the two auxiliary primes]\leavevmode
\label{cor:flagshifts520}
Fix an integer \(C\ge1\). Within one fixed product model of a
marked block, group the signed terms by the cosets, in the common
covers, of the \((g_1g_2)^C\)-torsion in each constrained coordinate.
Within each group the ranks of the signed polynomial arrays and of
their separately augmented terminal tables agree. The logarithms
of the two Gram determinants on any independent list of \(q\) arrays
differ by at most \(q p^{-1/2+o(1)}\). The other displacements,
the chosen cosets and the lifts of their base points may vary with
the interval.
\end{corollary}
\begin{proof}
Proposition~\ref{prop:ray50} puts each \(g_i\) between \((kb)^8\)
and twice that quantity, with \(b=\exp(x^\varepsilon)\). Thus
\(\log (g_1g_2)^C\ll_C x^\varepsilon+\log x=o(\log p)\) by
Theorem~\ref{thm:prime210}. The order is coprime to \(p\).
Projection of the common cover to the original elliptic curve
has degree \(p\) and so is an isomorphism on torsion of this
order. Explicitly, multiplication by that order is invertible
on its kernel, allowing adjustment of any lift to a unique torsion
lift. Differences of the indicated displacements therefore satisfy
the coset hypothesis in the same cover as the Fourier labels.
The theorem now applies to the signed sums within each group. On
an independent list, conjugating the two Gram inequalities by the
inverse square root of the polynomial Gram matrix bounds each
relative eigenvalue within \(p^{-1/2+o(1)}\) of one. Summing their
logarithms gives the assertion.
\end{proof}
\begin{corollary}[Contraction against the actual refined initial class]
\leavevmode\label{cor:actualsource520}
Fix a word \(c\), an initial coarse class and one unit-basis matrix
coefficient of its infinity type \(\rho\), as in
Proposition~\ref{prop:column110}. At the first edge fix the output
covectors and product coordinates. Collect any set of signed path
contributions whose first-edge displacements meet the hypotheses of
part~2 of Theorem~\ref{thm:signedtranslation520} in these coordinates.
The nonterminal cosets, the remaining edges and the test scalars may
be summed within this collection, keeping each shared index once.
For each output test \(i\) let \(\boldsymbol P_i\) be its grouped
array, including its arithmetic coefficients and class masses. Let
\(G_{\rm F}\) be the Gram matrix of these arrays, and let
\(G_{\rm in}\) be the Gram matrix of the resulting period
functionals, restricted to the separately zero-mean tables on
this actual initial class and matrix coefficient. Use the dual of
the automorphic probability norm on the functionals. Then
\[
 (1-\eta)R_cG_{\rm F}\ \le G_{\rm in}\le\ (1+\eta)R_cG_{\rm F},
 \qquad R_c=M_H(\dim\rho)(p^k-1)^f ,
\]
where \(k\) is the initial unitary rank, \(f=[E_x^+:\QQ]\), and
\(\eta\) is the augmented error in that theorem. On summing over
distinct initial coarse classes or orthogonal matrix coefficients,
the corresponding Gram matrices add with the same bounds.
\end{corollary}
\begin{proof}
By Proposition~\ref{prop:column110} the \(M_H\) initial classes
at the two flag primes have trivial stabilizers. Each has exactly
\(N=(p^k-1)^f\) refined classes indexed by the nonzero covectors,
by Corollary~\ref{cor:rank111}. They have equal Haar mass
\(1/(M_H N)\) in probability measure, since refinement divides
each compact open class into equal-measure cosets. Schur
orthogonality gives archimedean squared norm \(1/\dim\rho\) for
the chosen matrix coefficient. Hence the squared norm of an
arbitrary table \(\phi\) supported on this class and coefficient
is \(R_c^{-1}\sum_k |\phi(k)|^2\), where here the summation index
runs through the tuples of nonzero covectors. All such values can
be prescribed independently. Imposing the zero means restricts
this space precisely to the image of the product projection \(A\).

In the path identity of Theorem~\ref{thm:covectorpaths352} the
initial argument is read once at its input refinement. With the
next covectors fixed, Lemma~\ref{lem:terminalfourier454} expresses
the first incidence at that input by the cyclic projections in
our coordinates. All choices on later edges are independent of
the first input label once the common indices are fixed. Summing
their coefficients with the actual phases therefore gives on this
collection the functional \(\phi\mapsto\sum_k (S\boldsymbol P_i)(k)
\phi(k)\). This formula absorbs the weights of the arithmetic
sum into its coefficient array, so the use of counting measure in
writing it introduces no second class mass. Independent variations
in the unconstrained coordinates may be evaluated first by linearity.

For any table \(y\), the squared dual norm of \(\phi\mapsto
\sum_k y(k)\phi(k)\) on this subspace is \(R_c\|Ay\|^2\). Indeed,
in counting measure its Riesz vector is \(\bar y\), and the
zero-mean projection commutes with conjugation. Scaling the
source metric by \(1/R_c\) scales the dual metric by \(R_c\).
Apply the identity to every linear combination of the tests and
use Theorem~\ref{thm:signedtranslation520}. Disjoint initial
classes have disjoint supports and distinct basis matrix coefficients
are orthogonal by Schur orthogonality. Dual norms on their
orthogonal direct sum therefore add in squares, as claimed.
\end{proof}

\Needspace{8\baselineskip}
Distinct intermediate classes can still contribute to the same
initial class, with different first output covectors, displacements
or product models. Their functionals then add before the dual norm
is taken. The lower bound for this full sum and its comparison with
the incoming analytic row are still required in Section~\ref{nav:I-28}.


\subsection{Inner products of residue pullbacks through different class paths}
\label{sec:pencil521}

In Theorem~\ref{thm:covectorpaths352} the matrices on two paths
through different intermediate classes can be compared in their common
initial and final frames. We calculate their inner product when the
final covector runs over its full table. We then count the markings at
which their initial evaluations can lie on the same line. Throughout,
\(p=T\) is the split terminal prime already chosen on the interval.

At one half let \(U,V\) be vector spaces over \(\mathbb F_p\) of
dimensions \(k,n\). For a nontrivial complex character \(\chi\) of
\(\mathbb F_p^\times\), let \(H_\chi(U^\vee)\) be the space of tables
on \(U^\vee\) satisfying \(h(a\xi)=\chi(a)h(\xi)\) for \(a\ne0\).
All tables vanish at zero and use counting measure. For a linear map
\(C:U\to V\) write \(P_Ch(y)=h(yC)\). This is an operator into
\(H_\chi(V^\vee)\). We use a Hilbert--Schmidt inner product linear in
the second argument.

\begin{theorem}[The Gram matrix of two prescribed residue maps]
\leavevmode\label{thm:twopencils521}
For any two maps \(C,D:U\to V\), one has
\[
 \langle P_C,P_D\rangle_{\rm HS}
   =\frac1{p-1}\sum_{t\in\mathbb F_p^\times}\chi(t)\,
                     p^{n-\operatorname{rank}(D-tC)} .
\]
For \(C\ne0\) its diagonal entry is
\(a_C=p^n(1-p^{-\operatorname{rank}C})/(p-1)\).
Suppose \(C,D\) are nonzero and projectively distinct as matrices.
Let \(h\) be the rank of \(D-ZC\) over \(\mathbb F_p(Z)\), and
\(\delta=\min_{t\ne0}\operatorname{rank}(D-tC)\). Then
\[
 \frac{|\langle P_C,P_D\rangle_{\rm HS}|}{\sqrt{a_Ca_D}}
 \ \le\ \min\left(1,\frac{h p^{-\delta}}{1-p^{-1}}\right),
 \qquad \delta\ge1.
\]
The inner product is zero if the ranks at all nonzero scalars are
equal. For \(D=tC\), \(P_D=\chi(t)P_C\).
\end{theorem}
\begin{proof}
Choose one representative \(u\) of each line in \(U^\vee\). The
functions supported on their respective lines, taking the values
\((p-1)^{-1/2}\chi(a)\) at \(au\), form an orthonormal basis. Compute
the inner product by applying both operators to this basis and summing
over \(y\in V^\vee\). A pair of nonzero evaluations contributes only
when \(yD=t yC\). There is then exactly one such \(t\), and the
contribution is \(\chi(t)/(p-1)\). For each fixed \(t\), the equation
has \(p^{n-\operatorname{rank}(D-tC)}\) solutions. Those with
\(yD=yC=0\) must be removed for each \(t\). Their number is independent
of \(t\) and \(\sum_{t\ne0}\chi(t)=0\), giving the formula. Setting
\(D=C\) gives the diagonal. Homogeneity proves the assertion for
scalar multiples.

A nonzero \(h\)-rowed minor of \(D-ZC\) is a polynomial of degree
at most \(h\). Every specialization of smaller rank is a root, and
every specialization has rank at most \(h\) because the larger
minors vanish identically. Subtract the constant \(p^{n-h}\) in the
character sum. At most \(h\) summands then survive, each of modulus
at most \(p^{n-\delta}\). The two diagonal entries are at least
\(p^n(1-p^{-1})/(p-1)\). Division and the Hilbert--Schmidt
Cauchy--Schwarz inequality give the estimate. A zero specialization
would make the matrices proportional. The assertion about constant
rank follows from the same character sum.
\end{proof}

\begin{corollary}[A bound for signed sums on the full terminal table]
\leavevmode\label{cor:signedpencils521}
Fix common endpoints and frames. For a finite list of paths let
\(C_{i,v}\ne0\) be the composite residue matrices at the selected
halves \(v\mid\pi_T\), with \(f\) such halves. Prescribe a nontrivial
character \(\chi_v\) at each half. Put
\(Q_i=\bigotimes_v P_{C_{i,v}}\), \(a_i=\prod_v a_{C_{i,v}}\).
Let \(A_i\) be their coefficient maps between fixed finite-dimensional
Hilbert spaces, including all finite signs and assigned class masses.
In a fixed coupling on the last two edges the same shared index is used
in both maps. Matrices differing only by local nonzero scalars may
first be collected, multiplying each coefficient by the corresponding
character values. For the resulting list put
\[
 d_{ij}=\#\{v:[C_{i,v}]\ne[C_{j,v}]\},\qquad
 b=\min\{1,k/(p-1)\},\qquad
 e=\max_i\sum_{j\ne i} b^{d_{ij}} .
\]
Then
\[
 (1-e)\sum_i a_i\|A_i\|_{\rm HS}^2
 \ \le\ \left\|\sum_i Q_i\otimes A_i\right\|_{\rm HS}^2
 \ \le\ (1+e)\sum_i a_i\|A_i\|_{\rm HS}^2 .
\]
In particular if every two tuples differ at least at \(d\) halves, a
list of \(M\) terms has \(e\le (M-1)b^d\).
\end{corollary}
\begin{proof}
For the full final table replace the marked covector in the path
recursion of Theorem~\ref{thm:covectorpaths352} by its argument \(y\).
The incidence conditions now give the initial covector \((y_v C_{i,v})_v\).
A zero entry gives a zero value. The coefficient of a fixed path,
including its matrix cosets, is held fixed in this extension. Evaluation
at the prescribed marking recovers the original summand. A path with a
zero composite at a half can be discarded before this calculation.
Across several final classes the calculation applies in each pair of
endpoint frames separately.

Inner products of tensor products multiply. Each unequal matrix line
therefore contributes a normalized inner product of modulus at most
\(b\) by Theorem~\ref{thm:twopencils521}, since \(h\le k\). At the
other halves the modulus is one. Expand the squared norm, writing
\(z_i=\sqrt{a_i}\|A_i\|_{\rm HS}\). The absolute value of the
sum of off-diagonal terms is at most \(\sum_{i\ne j} b^{d_{ij}}z_i z_j\).
The inequality \(2z_i z_j\le z_i^2+z_j^2\) bounds this by
\(e\sum_i z_i^2\). This also verifies the estimate for arbitrary
signed coefficients without an independence assumption on the paths.
\end{proof}

\begin{proposition}[Markings at which two paths can cancel]
\leavevmode\label{prop:markedpencil521}
For prescribed \(C,D\) at one half, choose \(y\) uniformly among
nonzero final covectors. The probability that both \(yC,yD\) are
nonzero and on the same line is
\[
 \frac{1}{p^n-1}\sum_{t\ne0}
       \bigl(p^{n-\operatorname{rank}(D-tC)}
                  -p^{n-\operatorname{rank}[C\ D]}\bigr),
\]
where \([C\ D]:U\oplus U\to V\) has both sets of columns.
For nonproportional \(C,D\) the probability is at most
\(\min\{1,(p-1)p^{-h}+h p^{-\delta}\}\), with \(h,\delta\) as
above. For independent markings at several halves the exact
probabilities multiply. For a list of paths the fraction of markings
at which any two surviving paths read the same tuple of input lines
is at most the sum of these products over unordered pairs.
\end{proposition}
\begin{proof}
A surviving collision determines \(t\) uniquely. The common kernel
of the two pullbacks has exactly \(p^{n-\operatorname{rank}[C\ D]}\)
elements, including zero. Removing it gives the numerator. Among
nonzero \(y\) the probability of annihilating a fixed matrix of rank
\(j\) is \((p^{n-j}-1)/(p^n-1)\le p^{-j}\). Apply this with
\(j=h\) at all regular scalars and with \(j\ge\delta\) at the at most
\(h\) exceptional scalars. Independence here refers solely to uniform
counting on a product of finite sets. Multiplication gives the
product formula, followed by the union bound. At any particular
marking outside this union, different surviving tuples define
orthogonal evaluation functionals on the homogeneous initial table,
since their line supports are disjoint.
\end{proof}

At the particular Bessel marking the formula tests the actual
covectors \(y_v=(\bar\pi\bmod v)\overline{e_N^\vee}\). A counting
bound for the excluded markings leaves their membership to be
checked at this prescribed tuple. Even within one surviving tuple
of input lines the signed arithmetic coefficients still add before
the absolute value is taken.


\subsection{Initial lines at the prescribed marking and their heights}
\label{sec:markedlines522}

Fix \(E=E_x\), its \(f=[E^+:\QQ]\) selected terminal places
\(S\), and the ordinary initial frames of the signed period. Put
\(n=[E:\QQ]=2f\) and \(p=T=\varpi\bar\varpi\), where \(S\) consists
of the places above the chosen generator \(\varpi\) in the imaginary
quadratic field. Thus \(p\) splits completely in \(E\),
\(\varpi\mathcal O_E=\prod_{v\in S}v\), and every complex conjugate
of \(\varpi\) has absolute value \(\sqrt p\). We give bounds at the
specified marking for the same field and prime used in
Corollary~\ref{cor:sameprimes517}.

Write \(\operatorname{red}_S(P)\) for the tuple of projective
reductions of an \(E\)-point \(P=[z]\). At each finite place divide
the coordinates by any local scalar that gives them maximum absolute
value one before reducing. The resulting line is independent of the
choice. Use the absolute projective height
\[
 H_2(P)=\left(\prod_w \|z\|_w^{n_w}\right)^{1/n},
\]
where the finite norms are maxima with \(|p|_w=p^{-1}\) at \(w\mid p\),
the infinite norms are Euclidean, and \(n_w\) is the local degree.
In particular the weight of a complex place is two. The product
formula gives independence from the global representative.

\begin{lemma}[Separation of reductions at the actual terminal halves]
\leavevmode\label{lem:heightsep522}
For distinct projective points \(P,P'\), the number \(c(P,P')\) of
places in \(S\) where their reductions agree satisfies
\[
 c(P,P')\log p\le n\log\big(H_2(P)H_2(P')\big).
\]
In particular reduction on \(S\) is injective among points of height
strictly less than \(p^{1/4}\). The same assertion holds for matrices
regarded as points in projective space. In affine coordinates
\(P=[1:z_2:\cdots:z_k]\) with entries in \(\mathcal O_E\) of
absolute value at most \(B\) at every embedding, two distinct such
points satisfy the stronger bound \(c(P,P')\log p\le n\log(2B)\).
Here \(B\ge1\).
\end{lemma}
\begin{proof}
For representatives \(z,z'\) choose a nonzero minor
\(m=z_i z'_j-z_j z'_i\). At a finite place,
\(|m|_w\le\|z\|_w\|z'\|_w\). At an agreeing place in \(S\), reduction
after the indicated local scalings gives an additional factor
\(p^{-1}\). At each infinite place the determinant inequality on
the two coordinate pairs gives the same bound without that factor.
Multiply over all places using \(\prod_w|m|_w^{n_w}=1\) and \(n_v=1\)
on \(S\). The factor \(p^{-c(P,P')}\) proves the assertion. If all
\(f=n/2\) places agree, the product of the heights is at least
\(p^{1/2}\). For integral affine points choose an index with
\(z_j-z'_j\ne0\). Every agreeing place divides this algebraic integer
and its absolute norm is at most \((2B)^n\). This proves the last
claim.
\end{proof}

\begin{theorem}[Short representatives and many separated initial lines]
\leavevmode\label{thm:terminalreps522}
Let \(N_*\) be the cyclotomic modulus of the interval. In the fixed
residue coordinates of the initial class every tuple of lines in
\(\prod_{v\in S}\mathbb P^{k-1}(\mathbb F_p)\) is represented by an
integral row in \(E^k\) with every entry, in every complex embedding,
of absolute value at most \(nN_*\sqrt p/2\). Its projective height
is at most \(\sqrt{k}\,nN_*\sqrt p/2\).

Fix \(0<\theta<1\) and \(k\ge2\). There is a set of integral affine
rows whose distinct reductions differ at more than \(\theta f\) of
the actual halves. With
\[
 L=\left\lfloor\frac{p^{(1-\theta)/2}}{4nN_*}\right\rfloor
\]
one can take exactly \((L+1)^{n(k-1)}\) rows. On the common
intervals the logarithm of this number, divided by
\(f(k-1)\log p\), is \(1-\theta+o(1)\), uniformly in \(x\).
\end{theorem}
\begin{proof}
Put \(K=\QQ(\zeta_{N_*})\), which contains \(E\) by
Theorem~\ref{thm:conductor232}. The traces to \(E\) of the powers of
\(\zeta_{N_*}\) span \(E\) over \(\QQ\). Indeed the powers span
\(K\), and trace is surjective because its restriction to \(E\) is
multiplication by \([K:E]\). Select a rational basis
\(b_1,\ldots,b_n\) among these traces. They are algebraic integers
with \(|\sigma b_j|\le [K:E]\le N_*\) in every embedding.

The Chinese remainder theorem supplies an integral lift \(z\) of
any tuple in \(\prod_{v\in S}\mathbb F_p\). Expand \(z/\varpi\) in
the basis \(b_j\), with rational coordinates \(c_j\). Subtract
\(\varpi\sum_j m_j b_j\) from \(z\), rounding each \(c_j\) to a nearest
integer \(m_j\). This preserves every prescribed residue and leaves
absolute value at most \(\sqrt p\sum_j|\sigma b_j|/2\) at every
embedding. Choose a nonzero representative of each desired line at
each place and apply this construction to each coordinate. The
resulting integral row is locally primitive at every selected place.
Its finite height factors are at most one, giving the height bound.
This argument only uses a rational basis of integral elements.

For the separated family take first entry one and, independently for
each other entry, any sum \(\sum_j m_j b_j\) with
\(m_j\in\{0,\ldots,L\}\). Linear independence gives the stated
count. Every nonconstant entry is bounded by \(B=nN_*L\), which is
at least one on sufficiently distant intervals. A nonzero difference
of entries has norm at most \((2B)^n\), and \(2B<p^{(1-\theta)/2}\).
Lemma~\ref{lem:heightsep522} gives \(c<(1-\theta)f\). Finally \(n,N_*\)
have polynomial bounds in \(V\), while \(\log p=(17+o(1))m\log b\)
with \(b=\exp(x^\varepsilon)\) by Theorem~\ref{thm:prime209}. Thus
\(\log(4nN_*)=o(\log p)\) uniformly. Substituting in \(L\) proves
the logarithmic count.
\end{proof}

\begin{proposition}[Contraction and the integral image on a prescribed set of lines]
\leavevmode\label{prop:lineimage522}
Fix a coarse initial class \(a\), a tuple of nontrivial scalar
characters with product \(\chi\) on \(D=(\mathbb F_p^\times)^f\),
and a finite set \(\mathcal B\) of distinct tuples of input lines.
Take the coefficient field to contain the character values. For each
\(B\in\mathcal B\) fix any representative \(q_B\) in the actual
residue frames. For the prescribed final marking and projected test
\(F\), set
\[
 W_{B,F}=\sum_{\gamma:\,[\eta_\gamma]=B}
              \chi(t_\gamma) W_{\gamma,F},\qquad
             \eta_\gamma=t_\gamma q_B,\quad t_\gamma\in D.
\]
The sum is over every surviving path starting at \(a\) with the
indicated tuple, including all intermediate classes and the coupled
last edges of Theorem~\ref{thm:covectorpaths352}.
On the homogeneous integral initial tables supported on \(\mathcal B\)
the image of the period as an \(\mathcal O_J\)-module is exactly
\[
                    \sum_{B\in\mathcal B} W_{B,F}(D_{a,B}),
\]
with \(D_{a,B}\) the actual coefficient lattice on the line orbit.
On a fixed orthonormal matrix coefficient let \(w_{B,F}\) denote the
scalar value of the functional there. The squared dual norm on this
complex source in the automorphic probability metric is exactly
\[
     \frac{R_c}{(p-1)^f}\sum_{B\in\mathcal B}|w_{B,F}|^2,
     \qquad R_c=M_H(\dim\rho)(p^k-1)^f .
\]
For multiple tests the identity applies to every linear combination.
It therefore gives a positive-semidefinite lower bound on the Gram
matrix of the full projected period on the ordinary initial source.
\end{proposition}
\begin{proof}
A homogeneous table on the orbit of \(q_B\) has values
\(\phi(t q_B)=\chi(t)z_B\). The initial class has trivial stabilizer,
and its lattice on each orbit is constant as shown before
Theorem~\ref{thm:pointlattice353}. The unit character values make
integrality equivalent to \(z_B\in D_{a,B}\), including at places
dividing \(p-1\). Every such choice is allowed on each orbit
independently. Nontriviality at each half makes every local scalar
sum zero. Substitution in the path identity gives exactly
\(\sum_B W_{B,F}(z_B)\), proving the image formula. For one
orthonormal coefficient, Corollary~\ref{cor:actualsource520} gives the
squared source norm \((p-1)^f R_c^{-1}\sum_B|z_B|^2\).
Taking the dual norm proves the second formula. Restriction to a
subspace can only decrease a dual norm, also for each linear
combination of the tests. This is the asserted matrix inequality.
\end{proof}

\begin{corollary}[Paths with small projective height at the fixed marking]
\leavevmode\label{cor:markedheight522}
Work in common residue coordinates at the initial class. Suppose the
marked rows of a collection of paths are the reductions of nonzero
rows \(P_i\) over \(E\) of height less than \(p^{1/4}\). Distinct
projective \(P_i\) then contribute to disjoint terms of the sum in
Proposition~\ref{prop:lineimage522}. In particular this holds for
compatible integral matrix lifts \(A_{i,j}\) giving the prescribed
nonzero terminal pullbacks, if the reductions of the projective
products are the marked input lines and
\[
                H_2(y_i)\prod_{j=0}^{m-1} H_2(A_{i,j})<p^{1/4}.
\]
Here \(y_i\) lifts the marked terminal row and the matrix height uses
the entries as projective coordinates. The same conclusion holds for
the entire separated family in Theorem~\ref{thm:terminalreps522}.
\end{corollary}
\begin{proof}
Matrix multiplication satisfies the maximum-entry norm inequality
at finite places by the ultrametric inequality. At infinite places
the Euclidean norm for rows and the Frobenius norm for matrices are
submultiplicative, by Cauchy--Schwarz. Taking the weighted products
bounds the height of the nonzero product row by the displayed left
side. The separation lemma gives distinct tuples, and the integral
and norm formulas use disjoint input orbits on these tuples.
\end{proof}

Thus all actual tuples admit bounded integral coordinate lifts, and
many tuples admit mutually separated smaller ones. The signed
coefficient \(W_{B,F}\) at each tuple still includes paths of every
height reducing to it. A lower estimate for those coefficient sums
and their comparison with the lattice indices across the tests are
required for Section~\ref{sec:missing}.

\clearpage
\section{References}\label{nav:bibliography-XXVII}
\begingroup\def\section*#1{}\begin{thebibliography}{99}\small
\bibitem{XXVI-XX-bayer447} C. Bayer and J. Teichmann, \emph{The proof of Tchakaloff's theorem}, arXiv:math/0502473 (2005). \href{https://arxiv.org/abs/math/0502473}{Authors' preprint}.
\bibitem{XXVI-XX-gurvits489} L. Gurvits, \emph{Van der Waerden/Schrijver--Valiant like conjectures and stable (aka hyperbolic) homogeneous polynomials}, Electron. J. Combin. \textbf{15} (2008), R66. \href{https://arxiv.org/abs/0711.3496}{Author's preprint}.
\bibitem{XXVI-XX-oscillator519} S. Gurevich, R. Hadani and N. Sochen, \emph{The finite harmonic oscillator and its applications to sequences, communication, and radar}, IEEE Trans. Inform. Theory \textbf{54} (2008), 4239--4253, equation (I.1) and \S I-B. \href{https://people.math.wisc.edu/~sgurevich/IEEE-Journal.pdf}{Authors' paper}.
\end{thebibliography}\endgroup
\end{document}
